\documentclass[11pt]{article}

\usepackage[a4paper,margin=27mm,headheight=22pt]{geometry}
\usepackage[T1]{fontenc}
\usepackage[utf8]{inputenc}
\usepackage{lmodern}
\usepackage{microtype}
\usepackage{amsmath,amssymb,amsthm,mathtools,mathrsfs}
\DeclareMathOperator{\Log}{Log}
\usepackage{booktabs,longtable,tabularx,array,multirow}
\usepackage{enumitem}
\usepackage[dvipsnames]{xcolor}
\usepackage{fancyhdr}
\usepackage{titlesec}
\usepackage[most]{tcolorbox}
\usepackage{url}
\usepackage{seqsplit}
\usepackage{graphicx}
\usepackage{float}
\usepackage[colorlinks=true,linkcolor=MidnightBlue,citecolor=BrickRed,urlcolor=RoyalBlue]{hyperref}
\usepackage{bookmark}
\usepackage[nameinlink,noabbrev]{cleveref}

\definecolor{keionavy}{HTML}{17365D}
\definecolor{keiobluegray}{HTML}{49677F}
\definecolor{keiopale}{HTML}{F3F7FA}
\definecolor{keiogreen}{HTML}{1B6E43}
\definecolor{keioxblue}{HTML}{315A8A}
\definecolor{keioviolet}{HTML}{694A86}
\definecolor{keioopen}{HTML}{9B2C2C}
\definecolor{keiorule}{HTML}{CBD5DF}
\colorlet{navy}{keionavy}
\colorlet{bluegray}{keiobluegray}
\colorlet{pale}{keiopale}
\colorlet{passgreen}{keiogreen}
\colorlet{xpassblue}{keioxblue}
\colorlet{cpassviolet}{keioviolet}
\colorlet{openred}{keioopen}
\colorlet{rulegray}{keiorule}

\hypersetup{
  pdftitle={The Polyexponential Trace-Tower Program: Moving Levels, Faber and Inverse Monodromy, Positive Spectra, and Arithmetic Barriers},
  pdfauthor={KEIO},
  pdfsubject={Complete Research Monograph, Version 3.1},
  pdfkeywords={polyexponential, polylogarithm, Kaluza sign criterion, Faber polynomial, inverse monodromy, moving level, Euler constant, Padé, continued fraction, Jacobi operator, Hankel determinant}
}

\pagestyle{fancy}
\fancyhf{}
\fancyhead[L]{\small\color{keiobluegray}The Polyexponential Trace-Tower Program}
\fancyhead[R]{\small\color{keiobluegray}Complete Research Monograph v3.1}
\fancyfoot[C]{\small\color{keiobluegray}\thepage}
\renewcommand{\headrulewidth}{0.3pt}

\titleformat{\part}[display]
  {\centering\huge\bfseries\color{keionavy}}
  {Part \thepart}{0.7em}{}
\titleformat{\section}
  {\Large\bfseries\color{keionavy}}{\thesection.}{0.7em}{}
\titleformat{\subsection}
  {\large\bfseries\color{keionavy}}{\thesubsection.}{0.7em}{}
\titlespacing*{\part}{0pt}{1.5ex}{2.5ex}
\titlespacing*{\section}{0pt}{2.4ex plus .5ex minus .2ex}{1.1ex}
\titlespacing*{\subsection}{0pt}{2.0ex plus .4ex minus .2ex}{0.8ex}

\setlength{\parindent}{0pt}
\setlength{\parskip}{0.48em}
\setlength{\emergencystretch}{3em}
\setlist[itemize]{leftmargin=1.55em,itemsep=.22em,topsep=.35em}
\setlist[enumerate]{leftmargin=1.75em,itemsep=.3em,topsep=.35em}
\renewcommand{\arraystretch}{1.14}

\newtheoremstyle{keio}
  {0.9em}{0.7em}{\itshape}{}% above, below, body, indent
  {\bfseries\color{keionavy}}{.}{0.55em}{}
\theoremstyle{keio}
\newtheorem{theorem}{Theorem}[section]
\newtheorem{proposition}[theorem]{Proposition}
\newtheorem{corollary}[theorem]{Corollary}
\newtheorem{lemma}[theorem]{Lemma}
\newtheorem{conjecture}[theorem]{Conjecture}
\theoremstyle{definition}
\newtheorem{definition}[theorem]{Definition}
\newtheorem{problem}[theorem]{Problem}
\theoremstyle{remark}
\newtheorem{remark}[theorem]{Remark}

\newcommand{\Q}{\mathbb Q}
\newcommand{\Qbar}{\overline{\mathbb Q}}
\newcommand{\C}{\mathbb C}
\newcommand{\R}{\mathbb R}
\newcommand{\N}{\mathbb N}
\newcommand{\Z}{\mathbb Z}
\newcommand{\F}{\mathbb F}
\newcommand{\Ein}{\operatorname{Ein}}
\newcommand{\Cin}{\operatorname{Cin}}
\newcommand{\Ei}{\operatorname{Ei}}
\newcommand{\Ci}{\operatorname{Ci}}
\newcommand{\Chi}{\operatorname{Chi}}
\newcommand{\Eexp}{\mathcal E_{\mathrm{exp}}}
\newcommand{\trdeg}{\operatorname{trdeg}}
\newcommand{\Tr}{\operatorname{Tr}}
\newcommand{\Norm}{\operatorname{N}}
\newcommand{\LP}{\mathcal{LP}}
\newcommand{\ord}{\operatorname{ord}}
\newcommand{\Span}{\operatorname{span}}
\newcommand{\mult}{\operatorname{mult}}
\newcommand{\calG}{\mathcal G}
\newcommand{\lcm}{\operatorname{lcm}}
\newcommand{\diag}{\operatorname{diag}}
\newcommand{\dist}{\operatorname{dist}}
\newcommand{\FPTr}{\operatorname{FPTr}}
\newcommand{\Var}{\operatorname{Var}}
\newcommand{\Prob}{\mathbb P}
\newcommand{\Gal}{\operatorname{Gal}}
\newcommand{\stirling}[2]{\genfrac{[}{]}{0pt}{}{#1}{#2}}
\providecommand{\tightlist}{\setlength{\itemsep}{0pt}\setlength{\parskip}{0pt}}

\newcommand{\UPASS}{\textcolor{keiogreen}{\textsf{U-PASS}}}
\newcommand{\XPASS}{\textcolor{keioxblue}{\textsf{X-PASS}}}
\newcommand{\RPASS}{\textcolor{keioxblue}{\textsf{R-PASS}}}
\newcommand{\CPASS}{\textcolor{keioviolet}{\textsf{C-PASS}}}
\newcommand{\OPEN}{\textcolor{keioopen}{\textsf{OPEN}}}

\newtcolorbox{statusbox}{
  colback=keiopale,colframe=keiorule,boxrule=.55pt,arc=1.5mm,
  left=2.8mm,right=2.8mm,top=2mm,bottom=2mm,breakable
}
\newtcolorbox{notebox}{
  colback=white,colframe=keionavy,boxrule=.65pt,arc=1.5mm,
  borderline west={2.2pt}{0pt}{keionavy},
  left=3mm,right=2.8mm,top=2mm,bottom=2mm,breakable
}

\begin{document}

\hypersetup{pageanchor=false}
\begin{titlepage}
\thispagestyle{empty}
\vspace*{18mm}
{\color{keionavy}\rule{\textwidth}{1.2pt}}\\[11mm]
{\Huge\bfseries\color{keionavy}The Polyexponential\\[2mm]Trace-Tower Program\par}
\vspace{6mm}
{\LARGE\bfseries Moving Levels, Faber and Inverse Monodromy,\\
Positive Spectra, and Arithmetic Barriers\par}
\vspace{12mm}
{\large Complete Research Monograph\par}
\vfill
\begin{statusbox}
\textbf{Logical-status convention.}
\UPASS{} denotes an unconditional theorem with a complete proof;
\XPASS{} an unconditional finite theorem supported by exact reproducible
computation; \RPASS{} a theorem relative to the explicitly cited arithmetic
independence input from companion paper [001]; \CPASS{} a proved implication
from a displayed open hypothesis; and \OPEN{} a statement not proved and never
used silently in an unconditional argument.
\end{statusbox}
\vfill
{\Large KEIO\par}
\vspace{2mm}
{\large Version 3.1 --- 1 September 2026\par}
\vspace{7mm}
{\color{keionavy}\rule{\textwidth}{1.2pt}}
\end{titlepage}
\hypersetup{pageanchor=true}
\pagenumbering{roman}

\begin{abstract}
This monograph develops a single trace-tower program around the integer-order
polyexponentials
\[
 E_q(z)=\sum_{n\ge1}\frac{(-1)^{n-1}z^n}{n!n^q}
 \qquad(q\ge1).
\]
At finite level, their inverse-power traces are normalized Faber polynomials.
The trace field collapses to one level coordinate, and two even traces recover
that coordinate exactly; moving through arithmetically independent levels
restores a transversal tower of independent coordinates.  In the opposite
direction, the inverse surfaces have maximal finitary symmetric monodromy and
finite packets of inverse branches and first jets are algebraically free.
This contrast---maximal local branch freedom versus exact global trace
collapse---is the organizing principle of the volume.

The finite Faber polynomials are absolutely indecomposable in every order and
their geometric monodromy reduces uniformly to the alternating/symmetric
dichotomy.  Full symmetric monodromy is proved in even degree, in every prime
degree, in the infinite family \(m=p+2\), at fixed degree and sufficiently
large polyexponential order, and in the exact computer-assisted rectangles
recorded in the reproducibility appendix.  A finite-jet Faber transfer gives
isolated transpositions in wider polylogarithmic--hypergeometric families.
For every integer polylogarithmic order, a strict Kaluza-sign argument proves
all-degree absolute indecomposability and yields symmetric monodromy in every
prime degree and in the family \(m=p+2\); exact fixed-order window audits add
the stated finite odd-degree rectangles.  The logarithmic member further
gives type-III supernecklace polynomials with explicit arithmetic and
collision-geometric consequences.

The same framework supplies positive trace-class carriers, strict total
positivity, moment and Hankel towers, finite Jacobi inverse spectra, and
Gamma-jet spectral transfers.  It retains the exact Ramanujan--Meijer--Gamma
transmutations of the preceding manuscript and the Euler factorial
Padé/continued-fraction results, including full unit-boundary convergence at
\(t=-1\) for \(p=2,3,5,13\) and one rational subsequence converging at the real
place and at every fixed \(p\)-adic place.  Finally, the volume states the
precise approximation, reflexivity, \(p\)-adic height, Stokes-framing, and
differential-separability barriers that remain.  No unconditional theorem
here proves the irrationality or transcendence of Euler's constant, the
Euler--Gompertz constant, or an odd zeta value.
\end{abstract}

\begin{notebox}
\textbf{Companion input and scope.}
Whenever arithmetic independence at algebraic arguments is used, the statement
is marked \RPASS{} and cites the companion paper [001].  The finite and inverse
monodromy theorems are logically independent of that arithmetic input.  The
general \(f_{\sigma,\tau}\) large-\(\tau\) Morse assertion, full odd-composite
geometric \(S_m\), the global derivative-separability conjecture, and the
individual irrationality or transcendence of \(\gamma\) and \(\delta\) remain
open and are not promoted to theorems.
\end{notebox}

\tableofcontents
\newpage
\pagenumbering{arabic}

\part{Universal Traces and Arithmetic Moving Levels}
\section{Introduction and exact scope}

Write
\[
 \Ein(z)=\int_0^z\frac{1-e^{-t}}t\,dt
 =\sum_{r\ge1}\frac{(-1)^{r-1}z^r}{r\,r!},
\]
and let
\[
 \gamma=\lim_{n\to\infty}\left(\sum_{k=1}^n\frac1k-\log n\right),
 \qquad
 \delta=eE_1(1).
\]
Neither \(\gamma\) nor \(\delta\) is known to be irrational; see
Lagarias' survey \cite{Lagarias2013}.  The elementary positive-axis identity
\begin{equation}\label{eq:Ein-connection}
 \Ein(x)=\gamma+\log x+E_1(x)\qquad(x>0)
\end{equation}
places \(\gamma\) at the boundary between algebraic-point \(E\)-values and an
asymptotic connection constant.

The companion paper \cite{KEIO2026I} determined the algebraic relations among
all nonzero algebraic-point \(\Ein\)-values over the pure exponential field,
constructed several finite Euler witnesses, and computed the first two
inverse traces of the divisor of \(\Ein-\gamma\).  The present paper begins
where that finite theory stops.  Its organizing question is:
\begin{quote}
What remains true when the level of the divisor is allowed to move through a
countable algebraically independent family converging to \(\gamma\)?
\end{quote}

The results have four layers.  In the spectral layer, the two even traces
of orders two and four
recover the level exactly, so the entire fixed-level inverse-trace ring is
one-dimensional; moving the level restores one algebraically
independent coordinate per level.  The order-two row has a positive
telescoping structure and hence a genuine operator-theoretic realization.
Algebraic levels of a natural gamma-free inverse map have
transcendental preimages.  All finite linear witnesses are governed by
one normalized-certificate uniqueness lemma.  The dyadic tail is the
unique minimizer in a natural cumulative-positive augmentation cone, while
outside that cone three-point forms are already dense.  The resulting
coefficient-height transition identifies the exact scale at which the first
tail can be cancelled.  In the transmutation layer, the four-term Euler
recurrence from the 2026 Ramanujan Challenge is connected explicitly to the
official Meijer--\(G\) construction, the cubic Pilehrood family, a finite
probability law, a regularized Gamma determinant, and a full rational
zeta-jet tower.  Its denominator and numerator live in a flat dual-number
\(SL_3\) system, equivalently an ordinary \(SL_6\) system.  In the
obstruction layer, the particular residual-prime,
fixed-place, Rodrigues, natural-index P-recursive, and strict-Stokes
constructions studied here yield scoped saturation identities, budget
inequalities, or sharply stated missing inputs.

The principal results are as follows.

\begin{theorem}[Universal trace and transversal independence]
\label{thm:intro-universal}
For \(c\in\C^\times\), let \(Z_c\) be the zero multiset of
\(\Ein(z)-c\).  For every \(m\ge2\),
\[
 T_m(c):=\sum_{\rho\in Z_c}\rho^{-m}=P_m(c^{-1}),
\]
where \(P_m\in\Q[X]\) is monic of degree \(m\).  If
\(\{c_j:j\in J\}\) is algebraically independent over a characteristic-zero
field \(F\), and one integer \(m_j\ge2\) is selected for each \(j\), then
\(\{T_{m_j}(c_j):j\in J\}\) is algebraically independent over \(F\).
Moreover
\[
 c^{-1}=144T_4(c)-144T_2(c)^2-14T_2(c),
\]
so \(F[T_m(c):m\ge2]=F[c^{-1}]\).
The map \(c\mapsto(T_2(c),T_3(c))\) is injective except for the single
unordered pair
\[
 \{6-2\sqrt3,6+2\sqrt3\},
\]
whose common trace pair is \((-1/24,1/96)\).  For each \(m\ge2\), the
nonzero roots of \(P_m\) also give a nonempty finite set of algebraic levels
at which \(T_m\) vanishes and every point of the corresponding divisor is
transcendental.
\end{theorem}

\begin{theorem}[Finite-rank tails and incomplete-Gamma jets]
\label{thm:intro-tail-closures}
Let \(\Gamma\subset\Qbar_{>0}^{\times}\) have multiplicative rank \(r\).
For every finite \(A\subset\Gamma\),
\[
 \trdeg_KK(E_1(\alpha):\alpha\in A)\ge |A|-r-1.
\]
After deleting at most \(r+1\) values, the remaining full \(E_1\)-tail
family is jointly algebraically independent over \(K\).  The same bound
holds for the combined classical family
\(\{E_1(\alpha),\Ei(\alpha),\Ci(\alpha),\operatorname{Si}(\alpha)\}\).

For the regularized lower incomplete-Gamma jets
\[
 R_{\alpha,k}=[s^k]\{s^{-1}-\Gamma_{<}(s,\alpha)\},
\]
every finite \(S\subset\Gamma\times\N_0\) satisfies
\[
 \trdeg_KK(R_{\alpha,k}:(\alpha,k)\in S)\ge |S|-r.
\]
Thus at most \(r\) exceptions need be removed from the complete two-index
lower-jet family.  The corresponding upper jets have defect at most
\(r\) plus the number of distinct jet orders.
\end{theorem}

\begin{theorem}[Dyadic moving zero-trace tower]
\label{thm:intro-dyadic}
For \(N\ge1\), put
\[
 Y_N=(N+1)\Ein(2^N)-N\Ein(2^{N+1}).
\]
Then \(\{Y_N:N\ge1\}\) is countably algebraically independent over \(K\),
\(Y_N\downarrow\gamma\), and
\[
 Y_N-\gamma\sim \frac{N+1}{2^N}e^{-2^N}.
\]
For every choice \(m_N\ge2\), the exact traces
\[
 \tau_{m_N,N}:=\sum_{\Ein(\rho)=Y_N}\rho^{-m_N}
\]
are countably algebraically independent over \(K\).  For fixed \(m\),
\(\tau_{m,N}\to S_m\), the \(m\)-th inverse trace of \(\Ein-\gamma\).
\end{theorem}

\begin{theorem}[Inverse levels and critical coefficients]
\label{thm:intro-inverse}
The function
\[
 \mathscr R(q)=E_1(1)-2E_1(q)+E_1(q^2)
\]
maps \((1,\infty)\) increasingly onto \((0,E_1(1))\).  The inverse image of
every algebraic number in this interval is transcendental.  In particular,
the unique roots \(r_m>1\) of \(\mathscr R(r_m)=1/m\), \(m\ge5\), are all
transcendental.

For each fixed algebraic \(q>1\), the critical three-point cancellation
coefficients \(u_N(q)\) defined in \eqref{eq:uNq} satisfy, for every finite
\(I\),
\[
 \trdeg_KK(u_N(q):N\in I)\ge |I|-1.
\]
\end{theorem}

\begin{theorem}[Positive spectral realizations]
\label{thm:intro-positive}
Set
\[
 \tau_N=\frac1{Y_N^2}-\frac1{2Y_N},\qquad
 d_0=\tau_1,\qquad d_N=\tau_{N+1}-\tau_N\quad(N\ge1).
\]
Then \(d_N>0\), the \(d_N\) are countably algebraically independent over
\(K\), and
\[
 S_2=\sum_{N\ge0}d_N.
\]
The diagonal operator \(\mathsf A e_N=d_Ne_N\) belongs to every Schatten
class and has trace \(S_2\).  Moreover
\[
 \Delta_2(z)=\det\!\left(I-\frac{z^2}{2}\mathsf A\right)
 =\prod_{N\ge0}\left(1-\frac{d_N}{2}z^2\right)
\]
is an order-zero Laguerre--P\'olya entire function.  Its positive zeros are
countably algebraically independent over \(K\), and
\[
 \sum_{\rho\in Z(\Delta_2)}\rho^{-2}=S_2.
\]
The analogous order-four increments define a positive diagonal operator
\(\mathsf A_4\) with trace \(S_4\), and, writing
\(\mathsf A_2=\mathsf A\),
\[
 \gamma=
 \left\{144\Tr\mathsf A_4
 -144(\Tr\mathsf A_2)^2-14\Tr\mathsf A_2\right\}^{-1}.
\]
\end{theorem}

\begin{theorem}[Cross-index witness rigidity]
\label{thm:intro-witness}
Let \(q_*\in(0,1)\) be the real solution of \(\Ein(q_*)=\gamma\).  Let
\(\beta_n\in(0,1)\) be the positive real zero of
\[
 G_n(z)=n\Ein(z)-\Ein(z^n)-(n-1)\gamma,
\]
and, for odd \(n\ge3\), let \(\lambda_n<-1\) be its negative real zero.
Then at most one member of
\[
 \{q_*\}\cup\{\beta_n:n\ge2\}
 \cup\{\lambda_n:n\ge3,\ n\text{ odd}\}
\]
is algebraic.
\end{theorem}

\begin{theorem}[Critical moving-point barrier]
\label{thm:intro-barrier}
Any theorem of the natural form
\[
 |P(F(\xi))|\ge
 \exp\!\bigl(-C(1+|\xi|)^c\Phi(\deg P,\log H(P))\bigr)
\]
intended to hold uniformly for all fixed \(E\)-function systems and positive
integral sample points, or already for all nonzero polynomial values of the
displayed system, must have \(c\ge1\).  The obstruction occurs for the fixed
system
\[
 F=(\Ein(z),\Ein(2z),\Ein(4z))
\]
and the fixed linear form \(P=X_1-2X_2+X_3\).  General sufficiency at
\(c=1\) is not known; at that scale the same example forces at least the
polynomial loss \(\xi^{-1}\).  Within the augmentation-one stratum the
dyadic coefficients uniquely minimize the tail in the cumulative-positive
cone.  Outside that cone three points give dense tail values, while
cancellation of the first active tail requires at least
\(\log H\ge(q-1)q^N+O(1)\).  For consecutive augmentation-one triples this
transition has a matching construction; at natural point cost a normalized
lower bound for that affine template must pay height exponent at least two.

For the dyadic tower, algebraicity of \(S_2\) is
equivalent to the existence of one algebraic coefficient \(s\) for which
\[
 -\log|2sY_N^2+Y_N-2|\sim2^N.
\]
\end{theorem}

\begin{theorem}[Exact Ramanujan transmutation package]
\label{thm:intro-recurrence}
The two distinguished solutions of the Euler recurrence
\eqref{eq:challenge-recurrence} have the finite sums
\eqref{eq:qn-closed}--\eqref{eq:pn-closed}, satisfy the recurrence for every
\(n\ge0\), and obey \(p_n/q_n\to\gamma\).  Their terminating
\({}_2F_2\) state and the official Meijer--\(G\) state are joined by the
rational coboundary \eqref{eq:ram-coboundary}; the normalized conserved
pairing is the finite identity \eqref{eq:ram-bilinear-matrix}.  The
denominator system embeds in the flat \(SL_3\) field
\eqref{eq:Ma}--\eqref{eq:cmf-flatness}, and its zero--first jet lift is the
flat \(SL_6\) system \eqref{eq:SL6-block}.

The Euler--Beta--log transform
\eqref{eq:q-Beta-transform}--\eqref{eq:p-Beta-transform} identifies the
same solutions with the cubic Pilehrood family.  The rational carrier
\(\mathscr K_a\) has denominator and numerator as its zero- and first-jets;
\(\mathscr K_a(s)/\mathscr K_a(0)\) converges to
\(\Gamma(1-s)^3\Gamma(1+s)^2\) and yields the zeta limits
\eqref{eq:zeta-log-jets}.  Its scaled integer numerator is the monic contact
polynomial of \cref{thm:contact-polynomial}.  All recurrence solutions have
the limiting form \eqref{eq:all-initial-limits}.  Moreover, with
\(b=n+4\) and \(D_b=\lcm(1,\ldots,b)\),
\[
 \frac{b!}{D_b}\gcd\!\left(b,\frac{D_b}{b}\right)
 \mid\gcd(p_n,q_n),
\]
and the factorial-six part of the exact Casoratian
\eqref{eq:challenge-casoratian} comes from the determinant of the
denominator-cleared official dual transfer.
\end{theorem}

\begin{theorem}[Arithmetic and Stokes barriers]
\label{thm:intro-new-barriers}
The following statements hold unconditionally.
\begin{enumerate}[label=\textup{(\roman*)}]
 \item The dyadic level \((Y_N)\) and its order-two trace \((\tau_N)\) are
 not P-recursive.
 \item The two-base form of \cref{prop:two-base-obstruction} remains
 exponentially small with
 \[
  -\log\{Y_N(2)-Y_{M_N}(3)\}
  =2^N+N\log2-\log(N+1)+O(1).
 \]
 Thus multiplicative independence of the two bases does not remove this
 moving-point scale.
 \item Every finite multi-level trace field is exactly the
 \(\mathbb G_a\)-invariant field of the gauge shift
 \(G\mapsto G+t,H_i\mapsto H_i-t\).  The completed moving tower breaks this
 formal symmetry through its boundary limit.
 \item Every nondegenerate rank-one real Hilbert quotient obtained from the
 fixed residual-direction constraints is exactly saturated under a candidate
 \(\gamma=a/b\).  For the optimally saddled uncombined Pr\'evost form, the
 lcm-only height-one ledger has deficit
 \(2-\log(3+2\sqrt2)\).
 \item The fixed-prime Rodrigues family \eqref{eq:fixed-p-rodrigues}
 places \(p^{n-\ell}\) in the endpoint coefficient.  Even if this whole
 factor is credited as fixed-place gain, the resulting ledger and the
 archimedean saddle satisfy \eqref{eq:universal-fixed-p-bound}.
 \item The displayed unframed lateral transition of the elementary
 differential operator annihilating \(F=-\Ein\) is gamma-free.  The Euler
 constant occurs in the connection framing and, under
 \(\gamma\in\Qbar\), can be removed by an algebraic shear.  A framed
 Euler--Tate matched-line lifting theorem would imply the
 \(\Qbar\)-linear independence of \(1,\pi,\gamma\), and hence the
 transcendence of \(\gamma\), \(\gamma/\pi\), and \(S_2\).
\end{enumerate}
\end{theorem}

The qualification ``transversal'' in
\cref{thm:intro-universal,thm:intro-dyadic} is essential: one chooses one
trace order at each level.  At a fixed level all trace orders are algebraic
over the same one-dimensional level field.  Likewise, the infinite sum in
\cref{thm:intro-positive} has no automatic arithmetic consequence.  An
infinite sum of algebraically independent transcendental numbers may be
algebraic.

\paragraph{Organization.}
\Cref{sec:input} records the exact input from the companion paper.
\Cref{sec:rank-deficit} proves finite-rank cofinite independence for both
connection coordinates and classical tails, together with the complete
cancellation geometry; \cref{sec:external-realizations} exports the
engine to sine-integral graphs, generalized Dickman transforms, and
regularized incomplete-Gamma jets; the
characteristic-zero algebraic-matroid realization is isolated in
\cref{app:Ein-matroids} because it uses an external linearization theorem.
\Cref{sec:trace-polynomials} develops the universal trace calculus and
identifies its finite multi-level field with the exact shift quotient.
\Cref{sec:dyadic} proves the dyadic and all-order moving towers.
\Cref{sec:spectral} constructs the two positive operators and the
order-two Laguerre--P\'olya determinant.  \Cref{sec:other-towers} gives fixed-coefficient
power-orbit towers, transcendental inverse levels, and the logarithmic
lattice.  \Cref{sec:certificates} proves normalized certificate uniqueness,
cross-index rigidity, and the one-defect cancellation theorem.
\Cref{sec:relative} gives the relative five-class exclusion and the exact
finite-descent criterion.
\Cref{sec:ramanujan-recurrence,sec:cmf} solve the four-term Euler recurrence,
classify all of its limits, and construct the flat \(SL_3/SL_6\) field.
\Cref{sec:meijer-duality,sec:beta-log} give the exact official gauge and the
Euler--Beta--log offset ladder.  \Cref{sec:gamma-jets,sec:ram-probability}
develop the rational Gamma jets, contact polynomial, finite probability
law, and zero geometry; \cref{sec:gamma-determinant} gives the regularized
determinant realization.  \Cref{sec:cross-place-pade} proves the real and
\(p\)-adic Euler--Pad\'e convergence results.
\Cref{sec:barrier} proves the sharp inside/outside augmentation dichotomy and
isolates the critical moving-point obstruction.
\Cref{sec:sondow} replaces an invalid asymmetric-integral heuristic by an
exact evaluation.  \Cref{sec:prevost-grid,sec:fixed-prime-rodrigues}
establish the residual-prime, Ridout, and Rodrigues barriers.
\Cref{sec:stokes-framing,sec:mixed-gevrey,sec:euler-tate} separate unframed
Stokes data from the mixed arithmetic-Gevrey and canonical motivic
framings, prove the conditional dilation and Gamma-jet ranks, and isolate
the strengthened matched-line target.  The final
section records what remains open.

\paragraph{Dependency and status ledger.}
To keep imported input, internal proof, and conditional extrapolation
separate, the logical dependencies are summarized below.
\begin{center}
\small
\begin{tabular}{@{}
 >{\raggedright\arraybackslash}p{0.24\textwidth}
 >{\raggedright\arraybackslash}p{0.32\textwidth}
 >{\raggedright\arraybackslash}p{0.34\textwidth}@{}}
\toprule
Source & Imported statement & Status in this paper\\
\midrule
Companion paper \cite[Thms.~1.4, 5.1 and Lem.~11.2]{KEIO2026I}
& Algebraic independence of algebraic-point resonant \(\Ein_q\)-values
and of the balanced dyadic coordinates
& Standing arithmetic input; the tail, incomplete-Gamma, trace, carrier,
and rigidity deductions are proved here\\
Fischler--Rivoal \cite{FischlerRivoalMixed2024}
& Definitions of mixed functions; conjecture
\(\mathbf E\cap\mathbf D=\Qbar\); conditional linear specialization
& The dilation and all-order Gamma-jet families are constructed here;
every conclusion using the intersection conjecture is explicitly
conditional\\
Official Challenge solution \cite[\S5.2]{RamanujanOfficial2026}
& Official Meijer--\(G\) normalization and Miller-selected limiting direction
& Imported normalization; the rational coboundary, bilinear identity, and
transmutations are proved here\\
Rosen--Sidman--Theran \cite[Thm.~1.10]{RosenSidmanTheran2025}
& Linear representability over an algebraically closed characteristic-zero
field
& Used only in Appendix~\ref{app:Ein-matroids}\\
Powers \cite[\S5]{Powers2025}
& A concluding multisection algebraic-independence assertion
& Used only under an explicit hypothesis in
Remark~\ref{rem:Powers-conditional}; no unconditional result depends on it\\
\bottomrule
\end{tabular}
\end{center}

\section{Input from the companion paper}
\label{sec:input}

Throughout,
\[
 K=\Eexp=\Qbar(e^\alpha:\alpha\in\Qbar),
\]
and \(K^{\mathrm{alg}}\) denotes its relative algebraic closure in \(\C\).
We use the following theorem from \cite[Thm.~5.1]{KEIO2026I}.

\begin{theorem}[Algebraic-point \(\Ein\)-independence]
\label{thm:input-independence}
The countable family
\[
 \{\Ein(\alpha):\alpha\in\Qbar^\times\}
\]
is algebraically independent over \(K\), in the finite-subset sense.
Consequently it remains algebraically independent over
\(K^{\mathrm{alg}}\).
\end{theorem}

For the incomplete-Gamma application below we also use the resonant
extension of this input.  Put
\[
 \Ein_q(z)=\sum_{n\ge1}\frac{(-1)^{n-1}z^n}{n^q n!}
 \qquad(q\ge1).
\]
By \cite[Thm.~1.4]{KEIO2026I}, the family
\begin{equation}\label{eq:resonant-input}
 \{\Ein_q(\alpha):q\ge1,\ \alpha\in\Qbar^\times\}
\end{equation}
is algebraically independent over \(K\), in the finite-subset sense.

The final assertion is standard: algebraic independence is preserved under
algebraic extension of the base field.  We shall repeatedly use the following
linear consequence.

\begin{lemma}[Full-rank linear images]
\label{lem:linear-images}
Let \(x_1,\ldots,x_r\) be algebraically independent over a field \(F\), and
let \(A\in M_{s\times r}(F)\) have row rank \(s\).  Then the \(s\) coordinates
of \(A(x_1,\ldots,x_r)^{\mathsf T}\) are algebraically independent over \(F\).
\end{lemma}

\begin{proof}
Complete the rows of \(A\) to an invertible \(r\)-by-\(r\) matrix over \(F\).
The resulting linear change of variables is an automorphism of the rational
function field \(F(x_1,\ldots,x_r)\).
\end{proof}

For later comparison we also recall the Euler-level traces from
\cite[Thm.~14.5]{KEIO2026I}.  If \(Z_\gamma\) is the zero multiset of
\(\Ein(z)-\gamma\), then
\begin{align}
 S_2&:=\sum_{\rho\in Z_\gamma}\rho^{-2}
 =\frac1{\gamma^2}-\frac1{2\gamma},\label{eq:S2}\\
 S_3&:=\sum_{\rho\in Z_\gamma}\rho^{-3}
 =\frac1{\gamma^3}-\frac3{4\gamma^2}+\frac1{6\gamma},\label{eq:S3}
\end{align}
and
\begin{equation}\label{eq:gamma-recovery}
 \gamma=\frac{24S_2+1}{24S_3+6S_2}.
\end{equation}
In particular \(S_2\) is algebraic if and only if \(\gamma\) is algebraic.

\section{Finite-rank independence and cancellation geometry}
\label{sec:rank-deficit}

The algebraic-point theorem has a useful export form when the arguments lie
in a multiplicative group of finite rank.  We restrict here to positive
algebraic arguments, so that the real logarithm is single-valued and no
additional \(2\pi i\)-coordinate occurs.  Put
\begin{equation}\label{eq:c-alpha}
 c_\alpha=\Ein(\alpha)-\log\alpha=\gamma+E_1(\alpha)
 \qquad(\alpha\in\Qbar_{>0}).
\end{equation}

\begin{theorem}[Finite-rank deficit and cofinite independence]
\label{thm:cofinite-independence}
Let \(\Gamma\subset\Qbar_{>0}^{\times}\) have multiplicative rank \(r\).
For every finite \(A\subset\Gamma\),
\begin{equation}\label{eq:finite-rank-deficit}
 \trdeg_KK(c_\alpha:\alpha\in A)\ge |A|-r.
\end{equation}
Consequently there is a set \(E\subset\Gamma\), \(|E|\le r\), such that
\begin{equation}\label{eq:cofinite-family}
 \{c_\alpha:\alpha\in\Gamma\setminus E\}
 \quad\hbox{is algebraically independent over }K.
\end{equation}

Fix \(\alpha_0\in\Gamma\).  There is likewise a set
\(E_{\alpha_0}\subset\Gamma\setminus\{\alpha_0\}\),
\(|E_{\alpha_0}|\le r\), for which
\begin{equation}\label{eq:cofinite-differences}
 \{E_1(\alpha)-E_1(\alpha_0):
 \alpha\in\Gamma\setminus(E_{\alpha_0}\cup\{\alpha_0\})\}
\end{equation}
is algebraically independent over \(K\).
\end{theorem}

\begin{proof}
Choose \(\lambda_1,\ldots,\lambda_r\) as a basis of the \(\Q\)-vector space
spanned by \(\{\log\alpha:\alpha\in\Gamma\}\); torsion causes no difficulty
because all elements are positive.  For every finite \(A\), all
\(\log\alpha\), \(\alpha\in A\), lie in the \(\Q\)-span of the
\(\lambda_j\).  By \eqref{eq:c-alpha},
\[
 K(\Ein(\alpha):\alpha\in A)
 \subset K(c_\alpha:\alpha\in A,\lambda_1,\ldots,\lambda_r).
\]
The left side has transcendence degree \(|A|\) by
\cref{thm:input-independence}; this proves \eqref{eq:finite-rank-deficit}.

Choose a maximal algebraically independent subfamily of
\(\{c_\alpha:\alpha\in\Gamma\}\).  If its complement contained
\(r+1\) elements, the algebraic dependence of each of them on a finite part
of the maximal family would produce a finite set with deficit at least
\(r+1\), contradicting \eqref{eq:finite-rank-deficit}.  This proves
\eqref{eq:cofinite-family}.

For distinct \(\alpha_1,\ldots,\alpha_m\ne\alpha_0\), the differences
\(\Ein(\alpha_i)-\Ein(\alpha_0)\) are algebraically independent over \(K\):
they are full-rank linear images of the independent coordinates at
\(\alpha_0,\alpha_1,\ldots,\alpha_m\).  Subtracting the logarithmic
differences costs at most the same \(r\) logarithmic coordinates.  The first
argument and the maximality argument therefore apply verbatim.
\end{proof}

The common Euler term costs at most one further transcendence degree.  This
turns the shift bookkeeping into an unconditional theorem about the tails
themselves.

\begin{theorem}[Finite-rank tail defect and cofinite independence]
\label{thm:E1-cofinite}
Let \(\Gamma\subset\Qbar_{>0}^{\times}\) have multiplicative rank \(r\).
For every finite \(A\subset\Gamma\),
\begin{equation}\label{eq:E1-finite-rank-defect}
 \boxed{\displaystyle
 \trdeg_KK(E_1(\alpha):\alpha\in A)\ge |A|-r-1.}
\end{equation}
There is a set \(E\subset\Gamma\), \(|E|\le r+1\), such that
\begin{equation}\label{eq:E1-cofinite}
 \{E_1(\alpha):\alpha\in\Gamma\setminus E\}
\end{equation}
is algebraically independent over \(K\).  In particular at most \(r+1\)
members of the full tail family are algebraic over \(K\), and at most
\(r+1\) members of
\[
 \{\gamma\}\cup\{E_1(\alpha):\alpha\in\Gamma\}
\]
are algebraic over \(K\).  If \(\gamma\in K^{\mathrm{alg}}\), every
occurrence of \(r+1\) in these tail assertions can be replaced by \(r\).
\end{theorem}

\begin{proof}
By \cref{thm:cofinite-independence},
\[
 |A|-r\le\trdeg_KK(c_\alpha:\alpha\in A).
\]
Since \(c_\alpha=\gamma+E_1(\alpha)\), the field on the right is contained
in \(K(\gamma,E_1(\alpha):\alpha\in A)\).  Adjoining \(\gamma\) costs at
most one transcendence degree, which proves
\eqref{eq:E1-finite-rank-defect}.  Choose a maximal algebraically
independent subfamily of the tails.  If its complement contained \(r+2\)
elements, adjoining the finitely many members of the maximal family needed
to algebraize them would contradict \eqref{eq:E1-finite-rank-defect}.
This proves \eqref{eq:E1-cofinite} and the first algebraicity count.

If \(\gamma\) is algebraic over \(K\), adjoining it costs nothing and the
same proof gives defect \(r\).  If it is transcendental, it is not counted
among the algebraic members.  The assertion for the union follows in both
cases.
\end{proof}

\begin{corollary}[The two-prime triangular grid]
\label{cor:two-prime-grid}
For \(n\ge1\), put
\[
 A_n=\{2^k3^\ell:k,\ell\ge0,\ 1\le k+\ell\le n\},
 \qquad |A_n|=\frac{n(n+3)}2.
\]
Among \(\{c_\alpha:\alpha\in A_n\}\), all but at most two can be chosen
jointly algebraically independent over \(K\), and at most two are
algebraic.  Among \(\{E_1(\alpha):\alpha\in A_n\}\), all but at most three
can be chosen jointly algebraically independent, and at most three are
algebraic.  Thus at \(n=20\) the grid has \(230\) entries: at least \(228\)
of the first family, and at least \(227\) of the tail family, belong to
jointly algebraically independent subfamilies.
\end{corollary}

\begin{proof}
The group generated by \(2\) and \(3\) has multiplicative rank two.  Apply
\cref{thm:cofinite-independence,thm:E1-cofinite}; the cardinality is
\(\sum_{j=1}^n(j+1)=n(n+3)/2\).
\end{proof}

\begin{corollary}[Four classical special-function families]
\label{cor:four-classical-cofinite}
Let \(\Gamma\subset\Qbar_{>0}^{\times}\) have multiplicative rank \(r\).
From the combined family
\begin{equation}\label{eq:four-classical-family}
 \{E_1(\alpha),\Ei(\alpha),\Ci(\alpha),
   \operatorname{Si}(\alpha):\alpha\in\Gamma\}
\end{equation}
delete at most \(r+1\) members.  The remaining family is algebraically
independent over \(K\); in particular at most \(r+1\) members of the full
family are algebraic over \(K\).

For \(\alpha=1,\ldots,M\), the \(4M\) displayed values have transcendence
degree at least
\begin{equation}\label{eq:four-classical-integer-count}
 4M-\pi(M)-1
\end{equation}
over \(K\), and at most \(\pi(M)+1\) of them are algebraic over \(K\).
\end{corollary}

\begin{proof}
For each positive \(\alpha\), the four independent coordinates
\[
 \Ein(\alpha),\quad \Ein(-\alpha),\quad
 \frac{\Ein(i\alpha)+\Ein(-i\alpha)}2,\quad
 \frac{\Ein(i\alpha)-\Ein(-i\alpha)}{2i}
\]
are respectively
\[
 \gamma+\log\alpha+E_1(\alpha),\quad
 \gamma+\log\alpha-\Ei(\alpha),\quad
 \gamma+\log\alpha-\Ci(\alpha),\quad
 \operatorname{Si}(\alpha).
\]
Across distinct positive \(\alpha\), the underlying points
\(\{\pm\alpha,\pm i\alpha\}\) are disjoint; hence all these coordinates
are algebraically independent over \(K\) by
\cref{thm:input-independence,lem:linear-images}.  For any finite selection
from \eqref{eq:four-classical-family}, adjoining \(\gamma\) and at most
\(r\) logarithmic generators recovers the same number of independent
coordinates.  The resulting defect is at most \(r+1\), and maximality gives
the cofinite claim.  The integers up to \(M\) generate a multiplicative
group of rank \(\pi(M)\), which proves
\eqref{eq:four-classical-integer-count}.
\end{proof}

The finite-dimensional form identifies exactly where the deficit occurs.
Let \(\alpha_1,\ldots,\alpha_n\) be positive algebraic numbers contained in
a rank-\(r\) group.  Choose logarithmic generators
\(\lambda=(\log g_1,\ldots,\log g_r)^{\mathsf T}\) and an exponent matrix
\(L\in M_{n\times r}(\Q)\) such that
\begin{equation}\label{eq:log-exponent-matrix}
 (\log\alpha_1,\ldots,\log\alpha_n)^{\mathsf T}=L\lambda.
\end{equation}

\begin{theorem}[Linear-image rank deficit]
\label{thm:linear-rank-deficit}
Let \(M\in M_{m\times n}(\Qbar)\) have row rank \(s\), and put
\(t=\operatorname{rank}(ML)\).  With
\(c=(c_{\alpha_1},\ldots,c_{\alpha_n})^{\mathsf T}\), one has
\begin{equation}\label{eq:linear-rank-deficit}
 \boxed{\displaystyle
 \trdeg_KK(Mc)\ge s-t\ge s-r.}
\end{equation}
\end{theorem}

\begin{proof}
Write \(x=(\Ein(\alpha_1),\ldots,\Ein(\alpha_n))^{\mathsf T}\).  Then
\(x=c+L\lambda\).  Choose a row selector
\(R\in M_{s\times m}(\Qbar)\) for which \(RM\) has row rank \(s\).
The \(s\) coordinates of \(RMx\) are algebraically independent by
\cref{lem:linear-images}.  On the other hand,
\[
 RMx=R(Mc)+(RML)\lambda,
 \qquad \operatorname{rank}(RML)\le\operatorname{rank}(ML)=t.
\]
Thus \(\trdeg_KK((RML)\lambda)\le t\), while \(R(Mc)\) is a linear image
of \(Mc\).  Therefore
\[
 s\le \trdeg_KK(Mc)+t,
\]
which is the assertion.
\end{proof}

There is also a sharp statement in the directions where both the Euler
coefficient and all logarithms cancel.

\begin{corollary}[Complete cancellation space]
\label{cor:complete-cancellation-space}
Let
\begin{equation}\label{eq:cancellation-space}
 U=\ker\begin{pmatrix}\mathbf1^{\mathsf T}\\L^{\mathsf T}\end{pmatrix}
 \subset\Qbar^n,
 \qquad
 d=\dim U=n-\operatorname{rank}[\mathbf1,L].
\end{equation}
For any basis \(u_1,\ldots,u_d\) of \(U\), the numbers
\begin{equation}\label{eq:cancellation-values}
 W_j=\sum_{i=1}^n u_{j,i}\Ein(\alpha_i)
     =\sum_{i=1}^n u_{j,i}E_1(\alpha_i)
 \qquad(1\le j\le d)
\end{equation}
are algebraically independent over \(K\).  Thus
\eqref{eq:cancellation-space} is not only the complete space of
gamma-free, log-free coefficient directions: every independent direction
produces a new algebraically independent period.
\end{corollary}

\begin{proof}
The two defining equations for \(U\), together with
\eqref{eq:Ein-connection}, prove the equality in
\eqref{eq:cancellation-values}.  The basis rows have rank \(d\); apply
\cref{lem:linear-images} to the algebraically independent \(\Ein\)-coordinates.
\end{proof}

For a directed graph \(G=(V,E)\), let \(\partial\) be its incidence matrix,
and attach a positive algebraic number \(\alpha_v\) to each vertex.  If
\(L\) is the corresponding exponent matrix, then
\cref{thm:linear-rank-deficit} gives the graph form
\begin{equation}\label{eq:graph-rank-deficit}
 \trdeg_KK\bigl(E_1(\alpha_{t(e)})-E_1(\alpha_{s(e)}):e\in E\bigr)
 \ge |V|-c(G)-\operatorname{rank}(\partial^{\mathsf T}L),
\end{equation}
where \(c(G)\) is the number of connected components.  This is a genuine
incidence-geometric extension of the one-dimensional prime-defect counts
used below.

\section{External realizations of the independence engine}
\label{sec:external-realizations}

We record three consequences that export the algebraic-point input beyond
the exponential-integral language.  They classify relations rather than
asserting new irrationality results for Euler's constant.

\subsection*{Sine-integral graphic matroids}

Let \(G=(V,E)\) be a finite loopless directed multigraph, and choose nonzero algebraic numbers
\(a_v\) such that \(a_u^2\ne a_v^2\) for \(u\ne v\).  Put
\[
 s_v=\operatorname{Si}(a_v)
 =\frac{\Ein(ia_v)-\Ein(-ia_v)}{2i},
 \qquad
 I_e=s_{t(e)}-s_{s(e)}
 =\int_{a_{s(e)}}^{a_{t(e)}}\frac{\sin z}{z}\,dz.
\]

\begin{theorem}[Complete graphic relation ideal]
\label{thm:sine-graphic-ideal}
Let \(\partial_G:K^E\to K^V\) be the incidence map
\(\partial_Ge=e_{t(e)}-e_{s(e)}\).  Then
\begin{equation}\label{eq:sine-graphic-ideal}
 \boxed{\displaystyle
 \ker\!\left(K[X_e:e\in E]\longrightarrow K[I_e:e\in E]\right)
 =
 \left\langle\sum_{e\in E}c_eX_e:
 c\in\ker\partial_G\right\rangle.}
\end{equation}
Consequently
\begin{equation}\label{eq:sine-graphic-trdeg}
 \trdeg_KK(I_e:e\in E)=|V|-c(G),
\end{equation}
and a set of edge integrals is algebraically independent exactly when its
edges form a forest.
\end{theorem}

\begin{proof}
The pairs \(\{ia_v,-ia_v\}\) are disjoint.  The full-rank linear-image
lemma applied to the corresponding algebraically independent \(\Ein\)
values shows that the \(s_v\) are algebraically independent.  The edge
vector is \(I=\partial_G^{\mathsf T}s\).  Hence its coordinate ring is
the symmetric algebra of the incidence row space, whose defining ideal is
generated by the linear cycle relations.  The rank of the incidence matrix
is \(|V|-c(G)\), and a set of incidence columns is independent exactly when
the corresponding edges form a forest in the underlying multigraph.
\end{proof}

This polynomial-relation statement strengthens the disjoint-endpoint
linear-independence results for sine-integral values in
\cite{Delaygue2022} by using the stronger input from the companion paper.

\subsection*{Generalized Dickman log-Laplace values}

For the generalized Dickman distribution of
\cite{GrahovacEtAl2024}, the Laplace transform is
\begin{equation}\label{eq:Dickman-Laplace}
 \psi_{\theta,a}(s)=\exp\{-\theta\Ein(as)\}.
\end{equation}
Take positive algebraic \(\theta_j,a_j,s_j\), set \(x_j=a_js_j\), and
\[
 \Lambda_j=-\log\psi_{\theta_j,a_j}(s_j)=\theta_j\Ein(x_j).
\]

\begin{theorem}[Dickman log-Laplace classification]
\label{thm:Dickman-classification}
For every finite \(J\),
\begin{equation}\label{eq:Dickman-trdeg}
 \boxed{\displaystyle
 \trdeg_KK(\Lambda_j:j\in J)=\#\{x_j:j\in J\}.}
\end{equation}
For each distinct \(x\), choose a representative \(j_x\).  The complete
polynomial relation ideal is
\begin{equation}\label{eq:Dickman-relation-ideal}
 \ker\!\left(K[X_j:j\in J]\to K[\Lambda_j:j\in J]\right)
 =
 \left\langle
 \theta_{j_x}X_j-\theta_jX_{j_x}:
 x_j=x,\ j\ne j_x
 \right\rangle.
\end{equation}
Moreover,
for integers \(m_j\),
\begin{equation}\label{eq:Dickman-multiplicative}
 \boxed{\displaystyle
 \prod_{j\in J}\psi_{\theta_j,a_j}(s_j)^{m_j}=1
 \Longleftrightarrow
 \sum_{j:x_j=x}m_j\theta_j=0
 \quad\text{for every distinct }x.}
\end{equation}
\end{theorem}

\begin{proof}
Choose one representative for each distinct \(x_j\).  The corresponding
\(\Ein(x_j)\) are algebraically independent; every \(\Lambda_j\) is a nonzero
algebraic multiple of its block representative.  This proves the
transcendence degree and the linear defining ideal.  The Laplace values are
positive real numbers, so a multiplicative relation equal to one is
equivalent, after the real logarithm, to a linear relation among the
distinct \(\Ein(x)\); algebraic independence forces the displayed block
conditions.
\end{proof}

The theorem classifies the logarithmic values and the multiplicative
relations equal to one.  It does not claim to classify every polynomial
relation among the Laplace values themselves.

\subsection*{Regularized incomplete-Gamma jets}

For \(\alpha>0\) algebraic, let
\[
 \Gamma_{<}(s,\alpha)=\int_0^\alpha e^{-t}t^{s-1}\,dt,
 \qquad
 \Gamma_{>}(s,\alpha)=\int_\alpha^\infty e^{-t}t^{s-1}\,dt,
\]
initially in their half-planes of convergence and then by continuation.
At \(s=0\), define
\begin{equation}\label{eq:incomplete-gamma-jets}
 R_{\alpha,k}=[s^k]\{s^{-1}-\Gamma_{<}(s,\alpha)\},
 \qquad
 U_{\alpha,k}=[s^k]\Gamma_{>}(s,\alpha)
 \quad(k\ge0).
\end{equation}
This boundary-Laurent and multipoint regime is distinct from the
positive-parameter incomplete-Gamma special values studied, for example,
in \cite{MurtySaha2015}; no priority claim beyond the precise statements
proved below is intended.

\begin{theorem}[Cofinite independence of regularized lower jets]
\label{thm:lower-gamma-jets}
Let \(\Gamma\subset\Qbar_{>0}^{\times}\) have multiplicative rank \(r\),
and let \(S\subset\Gamma\times\N_0\) be finite.  Then
\begin{equation}\label{eq:lower-gamma-jet-defect}
 \boxed{\displaystyle
 \trdeg_KK(R_{\alpha,k}:(\alpha,k)\in S)\ge |S|-r.}
\end{equation}
Consequently, after deleting at most \(r\) members, the entire two-index
family \(\{R_{\alpha,k}:\alpha\in\Gamma,\ k\ge0\}\) is algebraically
independent over \(K\).  In particular at most \(r\) of these values are
algebraic over \(K\).

More explicitly, with \(L_\alpha=\log\alpha\),
\begin{equation}\label{eq:lower-gamma-jet-explicit}
 R_{\alpha,k}=-\frac{L_\alpha^{k+1}}{(k+1)!}
 +\sum_{j=0}^k\frac{(-1)^jL_\alpha^{k-j}}{(k-j)!}
 \Ein_{j+1}(\alpha).
\end{equation}
Thus \(R_{\alpha,0}=c_\alpha\), while
\begin{equation}\label{eq:lower-gamma-at-one}
 R_{1,k}=(-1)^k\Ein_{k+1}(1)
\end{equation}
and the latter all-order family is algebraically independent over \(K\)
without exceptions.
\end{theorem}

\begin{proof}
Splitting \(e^{-t}=1+(e^{-t}-1)\), expanding \(\alpha^s\), and then
expanding \((n+s)^{-1}\) gives \eqref{eq:lower-gamma-jet-explicit}.
For the finite set \(S\), collect the finitely many coordinates
\(\Ein_j(\alpha)\) that occur on the right.  The rows defining the
\(R_{\alpha,k}\) have rank \(|S|\): within each fixed \(\alpha\), the row
of order \(k\) has the new highest coordinate \(\Ein_{k+1}(\alpha)\) with
nonzero coefficient, and different \(\alpha\)-blocks are disjoint.

Let \(X\) be the resulting vector of \(M\) \(\Ein\)-coordinates and let
\(L=K(\lambda_1,\ldots,\lambda_r)\), where the \(\lambda_i\) span the real
logarithms of \(\Gamma\).  Over \(L\), complete the \(|S|\) displayed rows
to an invertible linear coordinate system \((R,Z)\) on \(X\), with
\(|Z|=M-|S|\).  By \eqref{eq:resonant-input},
\[
 M=\trdeg_KK(X)
 \le\trdeg_KL(R,Z)
 \le r+\trdeg_KK(R)+M-|S|.
\]
This proves \eqref{eq:lower-gamma-jet-defect}.  The maximality argument of
\cref{thm:cofinite-independence} then gives the cofinite assertion.
Setting \(L_1=0\) proves \eqref{eq:lower-gamma-at-one}.
\end{proof}

Write \(\Gamma(s)=s^{-1}+\sum_{k\ge0}g_ks^k\).  Since
\(\Gamma=\Gamma_<+\Gamma_>\), one has
\begin{equation}\label{eq:upper-lower-jet-split}
 U_{\alpha,k}=g_k+R_{\alpha,k}.
\end{equation}

\begin{corollary}[Upper-jet defect]
\label{cor:upper-gamma-jets}
For finite \(S\subset\Gamma\times\N_0\), let
\(d(S)=\#\{k:(\alpha,k)\in S\}\).  Then
\begin{equation}\label{eq:upper-gamma-jet-defect}
 \boxed{\displaystyle
 \trdeg_KK(U_{\alpha,k}:(\alpha,k)\in S)
 \ge |S|-r-d(S).}
\end{equation}
For every fixed \(k\), all but at most \(r+1\) members of
\(\{U_{\alpha,k}:\alpha\in\Gamma\}\) form an algebraically independent
family over \(K\).  At \(k=0\), \(U_{\alpha,0}=E_1(\alpha)\), so this
recovers \cref{thm:E1-cofinite}.
\end{corollary}

\begin{proof}
By \eqref{eq:upper-lower-jet-split}, adjoining the \(d(S)\) Gamma
coefficients \(g_k\) to the upper jets generates all lower jets in \(S\).
Combine \eqref{eq:lower-gamma-jet-defect} with the transcendence-degree
inequality.  For fixed \(k\), apply the same maximality argument as before.
\end{proof}

The finite-order shift ledger can now be replaced by a complete invariant
ring.  This is a structural classification, not an irrationality theorem.

\begin{proposition}[Complete finite Gamma-jet quotient]
\label{prop:finite-gamma-jet-quotient}
Fix \(N\ge0\), let
\[
 A_N(s)=1+\sum_{k=0}^Ng_ks^{k+1}\pmod {s^{N+2}},
 \qquad T_k=U_{1,k},\qquad d_k=g_k-T_k,
\]
and put \(\kappa_j=[s^j]\log A_N(s)\).  On the polynomial ring in the
\(g_k,T_k\), define a \(\mathbb G_a\)-action by
\begin{equation}\label{eq:gamma-jet-shift}
 A_N(s)\longmapsto e^{-ts}A_N(s)\pmod {s^{N+2}},
 \qquad d_k\longmapsto d_k.
\end{equation}
Then
\begin{equation}\label{eq:gamma-jet-invariant-ring}
 \boxed{\displaystyle
 \Q[g_0,\ldots,g_N,T_0,\ldots,T_N]^{\mathbb G_a}
 =\Q[\kappa_2,\ldots,\kappa_{N+1},d_0,\ldots,d_N].}
\end{equation}
The analogous equality holds for fraction fields.  At the actual Gamma
split,
\begin{equation}\label{eq:actual-gamma-jet-invariants}
 \kappa_1=-\gamma,qquad
 \kappa_j=\frac{(-1)^j}{j}\zeta(j)\ (j\ge2),qquad
 d_k=(-1)^{k+1}\Ein_{k+1}(1).
\end{equation}
Thus \(\gamma\) is the unique orbit coordinate at every finite order; the
zeta cumulants and resonant \(\Ein\)-coordinates generate all polynomial
and rational invariants.
\end{proposition}

\begin{proof}
The change of variables
\[
 (g_0,\ldots,g_N,T_0,\ldots,T_N)
 \longleftrightarrow
 (\kappa_1,\ldots,\kappa_{N+1},d_0,\ldots,d_N)
\]
is triangular and polynomial in both directions, by the truncated formal
identities \(\log A_N\) and \(A_N=\exp(\log A_N)\).  Under
\eqref{eq:gamma-jet-shift}, \(\kappa_1\mapsto\kappa_1-t\) and every other
displayed coordinate is fixed.  The invariant ring of translation in one
variable is generated by the remaining variables, proving
\eqref{eq:gamma-jet-invariant-ring}.  Finally
\(A_N(s)\) is the truncation of
\(s\Gamma(s)=\Gamma(1+s)\), so
\[
 \log\Gamma(1+s)=-\gamma s+
 \sum_{j\ge2}\frac{(-1)^j}{j}\zeta(j)s^j.
\]
Equation \eqref{eq:upper-lower-jet-split} at \(\alpha=1\), together with
\eqref{eq:lower-gamma-at-one}, gives the formula for \(d_k\).
\end{proof}

\begin{remark}[Correct reading of the all-order ledger]
The gauge orbit has dimension exactly one at every order.  Any additional
unknown excess in a truncated equation count is not another copy of the
Euler shift: the new odd-zeta cumulants are themselves invariant arithmetic
coordinates.  Consequently a growing unknown-minus-equation count is useful
bookkeeping, but is not by itself a no-go theorem.
\end{remark}

The corresponding export to arbitrary characteristic-zero algebraic
matroids uses an external linearization theorem and is therefore separated
from these direct applications; see \cref{app:Ein-matroids}.

\begin{remark}[A conditional density consequence of Powers]
\label{rem:Powers-conditional}
In the notation of Powers' later multisection manuscript
\cite[\S5]{Powers2025}, assume that
\[
 G_1(1),H_1(1),\ldots,G_N(1),H_N(1)
\]
are algebraically independent over \(\Qbar(e)\) for every \(N\).  This
assertion appears in the concluding multisection discussion rather than as
a proved numbered theorem in the cited version.  There
\begin{align*}
 G_n(1)&=\frac{\widetilde\delta^{(n)}-\delta^{(n)}}{2e},\\
 H_n(1)&=\gamma^{(n)}+
 \frac{\widetilde\delta^{(n)}+\delta^{(n)}}{2e}.
\end{align*}
Conditional on it, combine the two generalized Euler--Gompertz rows
\[
 \delta^{(1)},\ldots,\delta^{(N)},\qquad
 \widetilde\delta^{(1)},\ldots,\widetilde\delta^{(N)}.
\]
Their field satisfies
\[
 \trdeg_{\Qbar}\Qbar\bigl(
 \delta^{(1)},\ldots,\delta^{(N)},
 \widetilde\delta^{(1)},\ldots,\widetilde\delta^{(N)}\bigr)
 \ge\left\lfloor\frac{3N}{2}\right\rfloor-1.
\]
Indeed the Gamma-moment identities quoted by Powers give
\[
 \trdeg_{\Qbar}\Qbar(\gamma^{(1)},\ldots,\gamma^{(N)})
 \le N-\left\lfloor\frac N2\right\rfloor+1.
\]
The assumed \(2N\) independent coordinates, together with the
transcendence of \(e\), lie in the field generated by \(e\), the two
Euler--Gompertz rows, and these Gamma moments.  Subtracting transcendence
degrees yields the displayed bound.  Greedy basis selection in the
associated algebraic matroid, ordered as
\(\delta^{(1)},\widetilde\delta^{(1)},\delta^{(2)},\ldots\), then gives a
jointly algebraically independent subfamily of lower density at least
\(3/4\).  This entire paragraph is conditional, and no unconditional
result elsewhere in the paper depends on it.
\end{remark}

\section{Universal level-trace polynomials}
\label{sec:trace-polynomials}

For \(c\in\C^\times\), set
\[
 F_c(z)=\Ein(z)-c,
\]
and denote by \(Z_c\) its zero multiset, with multiplicities.  Since \(F_c\)
has order one, its inverse-power sums converge absolutely from order two
onward.

\begin{theorem}[Universal level-trace polynomial]
\label{thm:universal-trace}
For \(m\ge2\), define
\begin{equation}\label{eq:Pm-def}
 P_m(X):=-m[z^m]\log\bigl(1-X\Ein(z)\bigr).
\end{equation}
Then \(P_m\in\Q[X]\), its constant term is zero, and it is monic of degree
\(m\).  Moreover
\begin{equation}\label{eq:universal-trace}
 \boxed{\displaystyle
 T_m(c):=\sum_{\rho\in Z_c}\rho^{-m}=P_m(c^{-1}).}
\end{equation}
In particular
\begin{align}
 P_2(X)&=X^2-\frac X2,\label{eq:P2}\\
 P_3(X)&=X^3-\frac34X^2+\frac16X,\label{eq:P3}\\
 P_4(X)&=X^4-X^3+\frac{25}{72}X^2-\frac1{24}X.\label{eq:P4}
\end{align}
\end{theorem}

\begin{proof}
A genus-one Hadamard factorization has the form
\[
 F_c(z)=-c\,e^{a_cz}
 \prod_{\rho\in Z_c}
 \left(1-\frac z\rho\right)e^{z/\rho}.
\]
Taking the logarithm at the origin gives, for every \(m\ge2\),
\[
 [z^m]\log\frac{F_c(z)}{F_c(0)}
 =-\frac1m\sum_{\rho\in Z_c}\rho^{-m}.
\]
On the other hand,
\[
 \frac{F_c(z)}{F_c(0)}=1-\frac{\Ein(z)}c.
\]
Comparison proves \eqref{eq:universal-trace}.  Formula
\eqref{eq:Pm-def} is a polynomial in \(X\) of degree at most \(m\).  Since
\(\Ein(z)=z+O(z^2)\), the contribution of
\(X^m\Ein(z)^m/m\) to \eqref{eq:Pm-def} is \(X^m\).  Thus \(P_m\) is monic
of degree \(m\).  The displayed formulas follow by expanding through degree
four.
\end{proof}

\begin{remark}[Multiplicity and genus]
No simplicity assertion is required: all zeros are counted with their
multiplicities.  The first inverse sum is not used.  At order one the
canonical genus-one exponential factor absorbs the non-absolutely convergent
sum, whereas \eqref{eq:universal-trace} is absolute for \(m\ge2\).
\end{remark}

\begin{corollary}[Transversal trace independence]
\label{cor:transversal}
Let \(F\subset\C\) be a field of characteristic zero, and let
\(\{c_j:j\in J\}\subset\C^\times\) be algebraically independent over \(F\).
Choose one integer \(m_j\ge2\) for each \(j\).  Then
\[
 \{T_{m_j}(c_j):j\in J\}
\]
is algebraically independent over \(F\), in the finite-subset sense.
\end{corollary}

\begin{proof}
Fix a finite set \(I\subset J\), and write
\(x_j=c_j^{-1}\) and \(t_j=P_{m_j}(x_j)\).  Since
\[
 P_{m_j}(x_j)-t_j=0,
\]
the field \(F(x_j:j\in I)\) is algebraic over \(F(t_j:j\in I)\).  Hence
\[
 \trdeg_FF(t_j:j\in I)
 =\trdeg_FF(x_j:j\in I)=|I|.
\]
\end{proof}

\begin{theorem}[Exact even-trace recovery and ring closure]
\label{cor:fixed-level-closure}
For every \(c\ne0\),
\begin{equation}\label{eq:even-trace-recovery}
 \boxed{\displaystyle
 c^{-1}=144T_4(c)-144T_2(c)^2-14T_2(c).}
\end{equation}
Consequently, for every characteristic-zero field \(F\subset\C\),
\[
 \boxed{\displaystyle
 F[T_m(c):m\ge2]
 =F[T_2(c),T_4(c)]
 =F[c^{-1}]}
\]
and
\[
 \boxed{F(T_2(c),T_4(c))=F(c).}
\]
\end{theorem}

\begin{proof}
With \(x=c^{-1}\), \eqref{eq:P2} and \eqref{eq:P4} give
\[
 144P_4(x)-144P_2(x)^2-14P_2(x)=x.
\]
This proves \eqref{eq:even-trace-recovery}.  Every trace is a polynomial in
\(x\) by \cref{thm:universal-trace}, while the displayed identity recovers
\(x\) from \(T_2,T_4\).  Taking fraction fields and using \(F(x)=F(c)\)
finishes the proof.
\end{proof}

The preceding recovery identity also gives the exact, rather than merely
heuristic, form of the finite-stage shift obstruction.  It is important to
distinguish this gauge action, which fixes every connection coordinate,
from common translation of the connection coordinates themselves;
augmentation zero pertains to the latter action, not to the former.

\begin{theorem}[Exact finite-stage shift quotient]
\label{thm:exact-shift-quotient}
Let \(B\) be a characteristic-zero field, put
\[
 R=B[G,H_1,\ldots,H_n],\qquad c_i=G+H_i,
\]
and let \(\mathbb G_a\) act by
\begin{equation}\label{eq:gauge-action}
 G\longmapsto G+t,\qquad H_i\longmapsto H_i-t.
\end{equation}
Then
\begin{equation}\label{eq:shift-invariant-ring}
 \boxed{R^{\mathbb G_a}=B[c_1,\ldots,c_n],\qquad
 \operatorname{Frac}(R)^{\mathbb G_a}=B(c_1,\ldots,c_n).}
\end{equation}
If the universal traces are formed at these levels, then
\begin{align}
 B\bigl(T_m(c_i):1\le i\le n,\ m\ge2\bigr)
 &=B(c_1,\ldots,c_n),\label{eq:trace-invariant-field}\\
 B\bigl[T_m(c_i):1\le i\le n,\ m\ge2\bigr]
 &=B[c_1^{-1},\ldots,c_n^{-1}].\label{eq:trace-invariant-ring}
\end{align}
Thus the complete finite multi-level trace field is exactly the rational
quotient by the gauge direction.  It loses one orbit coordinate and no
others.
\end{theorem}

\begin{proof}
The polynomial change of variables \(H_i=c_i-G\) identifies
\(R=B[c_1,\ldots,c_n,G]\), and \eqref{eq:gauge-action} becomes translation
of \(G\) with every \(c_i\) fixed.  In characteristic zero the polynomial
and rational invariants of translation in one variable are respectively
\(B[c_1,\ldots,c_n]\) and \(B(c_1,\ldots,c_n)\), proving
\eqref{eq:shift-invariant-ring}.  Every trace is a polynomial in
\(c_i^{-1}\), while \eqref{eq:even-trace-recovery} recovers \(c_i^{-1}\)
from \(T_2(c_i),T_4(c_i)\).  This proves
\eqref{eq:trace-invariant-field}--\eqref{eq:trace-invariant-ring}.
\end{proof}

\begin{remark}[What the formal quotient does and does not prove]
The elimination ideal of the equations \(c_i=G+H_i\) in \(B[G]\) is zero:
the same finite invariant data admit a model with \(G=0\) and another with
\(G\) transcendental.  Hence invariant algebraic identities alone cannot
decide the arithmetic of the orbit coordinate.  This is a statement about
the formal axiom system, not about the actual complex number \(\gamma\).
Indeed the analytic completion supplies extra boundary data.  If
\(y_N\to g\) in a Hausdorff topological field and a continuous action fixes
every \(y_N\), it fixes \(g\).  Applied to \(Y_N\to\gamma\), and to the two
carrier limits \(\tau_N\to S_2\), \(\sigma_N\to S_4\), this is precisely
where the finite-stage gauge symmetry breaks; \(S_2,S_4\) recover
\(\gamma\) by \eqref{eq:gamma-even-recovery}.
\end{remark}

\begin{corollary}[Arithmetic classification of every Euler trace]
\label{cor:Euler-trace-arithmetic}
For \(m\ge2\), put \(S_m=T_m(\gamma)\).  Then
\[
 S_m\in\Qbar\quad\Longleftrightarrow\quad\gamma\in\Qbar,
 \qquad
 S_m\text{ is transcendental}
 \quad\Longleftrightarrow\quad\gamma\text{ is transcendental}.
\]
Moreover the two even traces recover Euler's constant without exceptions:
\begin{equation}\label{eq:gamma-even-recovery}
 \boxed{\displaystyle
 \gamma=\frac1{144S_4-144S_2^2-14S_2}.}
\end{equation}
In particular
\[
 \gamma\in\Q\quad\Longleftrightarrow\quad S_2,S_4\in\Q.
\]
\end{corollary}

\begin{proof}
If \(S_m=P_m(\gamma^{-1})\) is algebraic, then \(\gamma^{-1}\) is a root
of the nonconstant polynomial \(P_m(X)-S_m\) over \(\Qbar\), and hence is
algebraic.  The converse is immediate.  Formula
\eqref{eq:gamma-even-recovery} is \eqref{eq:even-trace-recovery} at
\(c=\gamma\); it also proves the rationality equivalence.
\end{proof}

\begin{theorem}[The unique two-trace level collision]
\label{thm:T2T3-collision}
For \(c,d\in\C^\times\),
\[
 (T_2(c),T_3(c))=(T_2(d),T_3(d))
\]
holds if and only if either \(c=d\), or
\begin{equation}\label{eq:T2T3-exceptional-pair}
 \boxed{\{c,d\}=\{6-2\sqrt3,6+2\sqrt3\}.}
\end{equation}
In the exceptional case the common pair is
\begin{equation}\label{eq:T2T3-exceptional-value}
 \boxed{(T_2,T_3)=\left(-\frac1{24},\frac1{96}\right).}
\end{equation}
\end{theorem}

\begin{proof}
Put \(x=c^{-1}\) and \(y=d^{-1}\).  From \eqref{eq:P2}, equality of the
second traces gives
\[
 (x-y)\left(x+y-\frac12\right)=0.
\]
If \(x=y\), then \(c=d\).  Otherwise \(x+y=1/2\), and equality of the
third traces, using \eqref{eq:P3}, reduces to
\[
 x^2+xy+y^2-\frac34(x+y)+\frac16=0,
\]
or \(xy=1/24\).  Thus \(x,y\) are the roots of
\[
 X^2-\frac12X+\frac1{24},
\]
and taking reciprocals gives \eqref{eq:T2T3-exceptional-pair}.  Direct
substitution in \eqref{eq:P2} and \eqref{eq:P3} gives
\eqref{eq:T2T3-exceptional-value}.
\end{proof}

\begin{corollary}[Exact vanishing traces with transcendental divisors]
\label{cor:vanishing-trace-levels}
For every \(m\ge2\), define
\[
 \mathcal V_m=
 \{c\in\C^\times:T_m(c)=0\}.
\]
Then
\begin{equation}\label{eq:vanishing-trace-levels}
 \boxed{\mathcal V_m=
 \{x^{-1}:x\ne0,\ P_m(x)=0\}\subset\Qbar^\times.}
\end{equation}
This set is nonempty and finite; its elements arise from exactly \(m-1\)
nonzero roots of \(P_m\), counted with multiplicity.  For every
\(c\in\mathcal V_m\), every zero \(\rho\) of \(\Ein(z)-c\) is
transcendental, although the absolutely convergent inverse trace satisfies
\[
 \sum_{\Ein(\rho)=c}\rho^{-m}=0.
\]
For example, \(\mathcal V_2=\{2\}\).
\end{corollary}

\begin{proof}
The trace identity \eqref{eq:universal-trace} gives
\eqref{eq:vanishing-trace-levels}.  The polynomial \(P_m\) has degree
\(m\), constant term zero, and linear coefficient
\[
 [X]P_m(X)=m[z^m]\Ein(z)=\frac{(-1)^{m-1}}{m!}\ne0.
\]
Hence zero is a simple root and the remaining \(m-1\) roots are nonzero,
counted with multiplicity.  They are algebraic because \(P_m\in\Q[X]\).

Now fix \(c\in\mathcal V_m\).  If a zero \(\rho\) of \(\Ein(z)-c\) were
algebraic, then \(\rho\ne0\) and
\(\Ein(\rho)=c\in\Qbar\), contradicting
\cref{thm:input-independence}.  Thus every point of the divisor is
transcendental.  Finally \(P_2(X)=X(X-1/2)\) gives the example.
\end{proof}

Thus even all inverse traces at one fixed level still contain only one
coordinate.  Within the universal inverse-trace construction, varying the
trace order at a fixed level produces only the one-dimensional field
generated by that level.  A transversal algebraically independent
inverse-trace family therefore requires varying the level; other constructions
may instead vary a coefficient, a power orbit, or an external parameter.

\section{Balanced dyadic levels and all-order moving traces}
\label{sec:dyadic}

For \(N\ge1\), define
\begin{equation}\label{eq:YN}
 Y_N=(N+1)\Ein(2^N)-N\Ein(2^{N+1}).
\end{equation}
The companion paper proved the first assertion below.  The remaining
assertions are the analytic bridge to the Euler level.

\begin{proposition}[Balanced dyadic approach]
\label{prop:dyadic-approach}
The family \(\{Y_N:N\ge1\}\) is countably algebraically independent over
\(K\).  Moreover
\begin{equation}\label{eq:epsilonN}
 Y_N=\gamma+\varepsilon_N,
 \qquad
 \varepsilon_N=(N+1)E_1(2^N)-NE_1(2^{N+1}),
\end{equation}
and
\begin{equation}\label{eq:epsilon-properties}
 \varepsilon_N>0,\qquad
 \varepsilon_N\downarrow0,\qquad
 \varepsilon_N\sim\frac{N+1}{2^N}e^{-2^N}.
\end{equation}
In particular \(Y_N\downarrow\gamma\).
\end{proposition}

\begin{proof}
The algebraic independence is \cite[Lem.~11.2]{KEIO2026I}; it also follows directly from
\cref{thm:input-independence,lem:linear-images}, because the finite row family
\(\bigl((N+1)e_N-Ne_{N+1}\bigr)\) has full row rank.

Substitution of \eqref{eq:Ein-connection} in \eqref{eq:YN} cancels both the
logarithmic and the \(N\log2\) terms and gives \eqref{eq:epsilonN}.  Positivity
follows from
\[
 \varepsilon_N
 =E_1(2^N)+N\{E_1(2^N)-E_1(2^{N+1})\}>0.
\]
For strict decrease put \(g(t)=E_1(2^t)\).  Since
\[
 g'(t)=-(\log2)e^{-2^t},\qquad
 g''(t)=(\log2)^2 2^t e^{-2^t}>0,
\]
strict convexity gives
\[
 \varepsilon_N-\varepsilon_{N+1}
 =(N+1)\{g(N)-2g(N+1)+g(N+2)\}>0.
\]
Finally, the standard positive-axis expansion
\begin{equation}\label{eq:E1-asymptotic}
 E_1(x)=\frac{e^{-x}}x\left(1+O(x^{-1})\right)
 \qquad(x\to+\infty)
\end{equation}
gives the asymptotic in \eqref{eq:epsilon-properties}; see
\cite[\S6.12\(i\)]{DLMF}.
\end{proof}

For \(m\ge2\), let
\begin{equation}\label{eq:tau-mN}
 \tau_{m,N}:=\sum_{\Ein(\rho)=Y_N}\rho^{-m},
\end{equation}
where the solutions are counted with multiplicity.

\begin{theorem}[All-order transversal tower]
\label{thm:all-order-tower}
For every \(m\ge2\) and \(N\ge1\),
\begin{equation}\label{eq:tau-Pm}
 \tau_{m,N}=P_m(Y_N^{-1}).
\end{equation}
For any selection \(m_N\ge2\), the family
\[
 \{\tau_{m_N,N}:N\ge1\}
\]
is countably algebraically independent over \(K\).  For each fixed
\(m\ge2\),
\begin{equation}\label{eq:tau-limit}
 \tau_{m,N}\longrightarrow S_m:=P_m(\gamma^{-1}).
\end{equation}
\end{theorem}

\begin{proof}
Equation \eqref{eq:tau-Pm} is \cref{thm:universal-trace}.  Apply
\cref{cor:transversal} to the independent levels \(Y_N\).  The limit follows
from \cref{prop:dyadic-approach} and continuity of \(P_m(1/x)\) at
\(x=\gamma\).
\end{proof}

\begin{remark}[Why transversal is sharp]
For a fixed \(N\), the even traces already recover the level exactly:
\begin{equation}\label{eq:YN-even-recovery}
 \boxed{\displaystyle
 Y_N^{-1}=144\tau_{4,N}-144\tau_{2,N}^2-14\tau_{2,N}.}
\end{equation}
Thus
\[
 K[\tau_{m,N}:m\ge2]
 =K[\tau_{2,N},\tau_{4,N}]
 =K[Y_N^{-1}].
\]
The full double family is therefore not algebraically independent: the
theorem gives exactly one independent spectral coordinate per moving level.
\end{remark}

The order-two row has additional positivity.  Put
\begin{equation}\label{eq:f-def}
 f(x)=\frac1{x^2}-\frac1{2x}=\frac{2-x}{2x^2},
 \qquad \tau_N:=\tau_{2,N}=f(Y_N).
\end{equation}

\begin{theorem}[Exact positive convergence at order two]
\label{thm:tau2-convergence}
One has
\[
 0<\tau_1<\tau_2<\cdots<S_2,
 \qquad \tau_N\longrightarrow S_2.
\]
More precisely,
\begin{equation}\label{eq:S2-tail-exact}
 \boxed{\displaystyle
 S_2-\tau_N=
 \frac{\varepsilon_N
 \{\gamma(4-\gamma)+(2-\gamma)\varepsilon_N\}}
 {2\gamma^2(\gamma+\varepsilon_N)^2}.}
\end{equation}
Consequently
\begin{equation}\label{eq:S2-tail-asymptotic}
 S_2-\tau_N\sim
 \frac{4-\gamma}{2\gamma^3}
 \frac{N+1}{2^N}e^{-2^N}.
\end{equation}
\end{theorem}

\begin{proof}
Since \(E_1(2)<e^{-2}/2\) and \(\gamma<1\),
\(Y_1=\gamma+2E_1(2)-E_1(4)<\gamma+e^{-2}<2\), so all \(Y_N\) lie in
\((0,2)\).  But
\[
 f'(x)=\frac{x-4}{2x^3}<0\qquad(0<x<2).
\]
Thus \(Y_N\downarrow\gamma\) implies \(\tau_N\uparrow f(\gamma)=S_2\), and
positivity follows from \(f(x)>0\) on \((0,2)\).  Direct subtraction with
\(Y_N=\gamma+\varepsilon_N\) gives \eqref{eq:S2-tail-exact}.  Combine it
with \eqref{eq:epsilon-properties} to obtain
\eqref{eq:S2-tail-asymptotic}.
\end{proof}

\begin{theorem}[Non-P-recursiveness of the dyadic moving sequences]
\label{thm:non-p-recursive}
Neither \((Y_N)_{N\ge1}\) nor \((\tau_N)_{N\ge1}\) is P-recursive over
\(\C\).  Equivalently, neither sequence satisfies a nonzero recurrence
\[
 \sum_{j=0}^r p_j(N)u_{N+j}=0
 \qquad\bigl(p_j\in\C[N]\bigr)
\]
for all sufficiently large \(N\).
\end{theorem}

\begin{proof}
Suppose first that such a recurrence holds for \(Y_N=\gamma+\varepsilon_N\).
Since every polynomial multiple of \(\varepsilon_{N+j}\) tends to zero, we
obtain
\[
 \gamma\sum_{j=0}^rp_j(N)\longrightarrow0.
\]
The sum is a polynomial, hence it vanishes identically.  Let \(j_0\) be the
least index for which \(p_{j_0}\ne0\).  Dividing the remaining relation by
\(p_{j_0}(N)\varepsilon_{N+j_0}\) gives
\[
 1+\sum_{j>j_0}
 \frac{p_j(N)}{p_{j_0}(N)}
 \frac{\varepsilon_{N+j}}{\varepsilon_{N+j_0}}=0.
\]
But \eqref{eq:epsilon-properties} implies, for every fixed \(j>j_0\),
\[
 \frac{\varepsilon_{N+j}}{\varepsilon_{N+j_0}}
 =\exp\{-\bigl(2^j-2^{j_0}\bigr)2^N+O(N)\},
\]
which dominates every rational function of \(N\).  The displayed equality
therefore tends to \(1=0\), a contradiction.

For \(\tau_N=f(Y_N)\), \cref{thm:tau2-convergence} gives
\[
 \tau_N=S_2+f'(\gamma)\varepsilon_N+O(\varepsilon_N^2),
 \qquad f'(\gamma)\ne0.
\]
The same least-index argument applies verbatim.
\end{proof}

\begin{remark}[Scope of the non-P-recursiveness result]
The theorem does not say that no holonomic construction can involve the
finite-stage values.  It says that the natural moving index itself cannot be
packaged as a fixed P-recursive sequence.  Any application of arithmetic
holonomy must therefore replace the moving tower by a genuinely different,
fixed differential or recurrence system rather than merely reindexing
\(Y_N\).
\end{remark}

For orientation, the first values are shown in \cref{tab:dyadic-values}.
The table is illustrative only; no numerical value is used in a proof.

\begin{table}[ht]
\centering
\caption{The dyadic level and its order-two zero trace.}
\label{tab:dyadic-values}
\begin{tabular}{@{}rcc@{}}
\toprule
\(N\)&\(Y_N\)&\(\tau_N\)\\
\midrule
1&0.671237333907806&1.474569389580200\\
2&0.588478390885392&2.037963940792663\\
3&0.577366307471447&2.133831985200994\\
4&0.577215698103968&2.135171685377284\\
5&0.577215664901535&2.135171980841625\\
\midrule
limit&0.577215664901533&2.135171980841646\\
\bottomrule
\end{tabular}
\end{table}

\section{Positive spectral carriers for
\texorpdfstring{\(S_2\) and \(S_4\)}{S2 and S4}}
\label{sec:spectral}

Define
\begin{equation}\label{eq:dN-def}
 d_0=\tau_1,
 \qquad
 d_N=\tau_{N+1}-\tau_N\quad(N\ge1).
\end{equation}

\begin{theorem}[Positive algebraically independent increments]
\label{thm:increments}
The sequence \(\{d_N:N\ge0\}\) is positive and countably algebraically
independent over \(K\).  It satisfies
\begin{equation}\label{eq:S2-sum}
 \boxed{\displaystyle S_2=\sum_{N=0}^\infty d_N}
\end{equation}
and, for \(N\ge1\),
\begin{equation}\label{eq:dN-asymptotic}
 d_N\sim
 \frac{4-\gamma}{2\gamma^3}
 \frac{N+1}{2^N}e^{-2^N}.
\end{equation}
\end{theorem}

\begin{proof}
Positivity follows from \cref{thm:tau2-convergence}.  The finite
transformation
\[
 (\tau_1,\ldots,\tau_{M+1})
 \longleftrightarrow(d_0,\ldots,d_M)
\]
is an invertible triangular \(\Q\)-linear transformation.  The
\(\tau_N\) are algebraically independent by
\cref{thm:all-order-tower}; hence so are the \(d_N\).  The telescoping
identity
\[
 \sum_{N=0}^Md_N=\tau_{M+1}
\]
and \cref{thm:tau2-convergence} give \eqref{eq:S2-sum}.  Finally the ratio
of two successive tails in \eqref{eq:S2-tail-asymptotic} tends to zero, so
\[
 d_N=(S_2-\tau_N)-(S_2-\tau_{N+1})\sim S_2-\tau_N,
\]
which proves \eqref{eq:dN-asymptotic}.
\end{proof}

Let \(\{e_N:N\ge0\}\) be the standard basis of \(\ell^2(\N_0)\).

\begin{corollary}[A positive operator carrier]
\label{cor:operator}
The diagonal operator
\[
 \mathsf A e_N=d_Ne_N
\]
is positive, self-adjoint and trace class, with
\[
 \Tr\mathsf A=S_2.
\]
More strongly,
\[
 \sum_{N\ge0}d_N^p<\infty\qquad(p>0).
\]
Thus \(\mathsf A\in\mathcal S_p\) for every \(p\ge1\), and satisfies the
corresponding quasi-Schatten condition for \(0<p<1\).  Its nonzero
eigenvalues are simple and countably algebraically independent over \(K\).
\end{corollary}

\begin{proof}
All assertions except simplicity follow immediately from
\cref{thm:increments}.  The double-exponential asymptotic
\eqref{eq:dN-asymptotic} gives summability of every positive power.  If
\(d_j=d_k\) for \(j\ne k\), the algebraically independent family would
satisfy a nonzero linear relation; hence the eigenvalues are distinct.
\end{proof}

The Fredholm determinant of the operator gives a scalar entire realization.

\begin{theorem}[Order-zero Laguerre--P\'olya realization]
\label{thm:Delta2}
The product
\begin{equation}\label{eq:Delta2}
 \Delta_2(z)=
 \det\!\left(I-\frac{z^2}{2}\mathsf A\right)
 =\prod_{N=0}^\infty\left(1-\frac{d_N}{2}z^2\right)
\end{equation}
converges locally uniformly and defines an order-zero entire function in the
Laguerre--P\'olya class.  Its zeros are simple and equal to
\[
 \pm r_N,\qquad r_N=\sqrt{\frac2{d_N}}.
\]
The positive half \(\{r_N:N\ge0\}\) is countably algebraically independent
over \(K\), and
\begin{equation}\label{eq:Delta2-trace}
 \boxed{\displaystyle
 \sum_{\rho\in Z(\Delta_2)}\rho^{-2}=S_2.}
\end{equation}
\end{theorem}

\begin{proof}
Since \(\sum d_N<\infty\), the product in \eqref{eq:Delta2} converges
normally on compact sets.  Every finite partial product has only real zeros,
so their locally uniform limit belongs to the Laguerre--P\'olya class; see
\cite[Chapter~VIII]{Levin1996}.  For every \(\eta>0\),
\[
 \sum_{N\ge0}r_N^{-\eta}
 =2^{-\eta/2}\sum_{N\ge0}d_N^{\eta/2}<\infty.
\]
Thus the zero exponent of convergence, and hence the order of the canonical
product, is zero.  Simplicity follows from the distinctness of the \(d_N\).

For a finite index set \(I\), the relations \(d_N=2/r_N^2\) imply
\[
 \trdeg_KK(r_N:N\in I)
 \ge \trdeg_KK(d_N:N\in I)=|I|,
\]
which proves algebraic independence of the positive zeros.  Finally,
absolute convergence and \eqref{eq:S2-sum} give
\[
 \sum_{\rho\in Z(\Delta_2)}\rho^{-2}
 =\sum_{N\ge0}\left(r_N^{-2}+(-r_N)^{-2}\right)
 =\sum_{N\ge0}d_N=S_2.
\]
\end{proof}

\begin{theorem}[A second positive carrier and even-trace recovery]
\label{thm:positive-S4-carrier}
Put
\[
 \sigma_N=T_4(Y_N)=P_4(Y_N^{-1})\quad(N\ge1),
\]
and define
\[
 d_0^{(4)}=\sigma_1,\qquad
 d_N^{(4)}=\sigma_{N+1}-\sigma_N\quad(N\ge1).
\]
Then every \(d_N^{(4)}\) is positive, the family
\(\{d_N^{(4)}:N\ge0\}\) is algebraically independent over \(K\), and
\begin{equation}\label{eq:S4-positive-sum}
 S_4=\sum_{N\ge0}d_N^{(4)}.
\end{equation}
Consequently the diagonal operator
\(\mathsf A_4=\diag(d_N^{(4)}:N\ge0)\) belongs to every Schatten class and
\(\Tr\mathsf A_4=S_4\).  Together with the order-two carrier
\(\mathsf A_2:=\mathsf A\) of \cref{cor:operator}, it gives
\begin{equation}\label{eq:gamma-two-positive-carriers}
 \boxed{\displaystyle
 \gamma=
 \left\{144\Tr\mathsf A_4
 -144(\Tr\mathsf A_2)^2-14\Tr\mathsf A_2\right\}^{-1}.}
\end{equation}
\end{theorem}

\begin{proof}
The earlier bound \(Y_N<2\) and the monotonicity \(Y_N\downarrow\gamma\)
show that \(x_N=Y_N^{-1}\) increases and satisfies \(x_N>1/2\).  From
\eqref{eq:P4},
\[
 P_4(1/2)=\frac1{288},\qquad P_4'(1/2)=\frac1{18},
\]
and
\[
 P_4''(x)=12x^2-6x+\frac{25}{36}>0\qquad(x\ge1/2).
\]
Thus \(P_4\) is positive and strictly increasing on this range, so all
increments are positive.  The \(\sigma_N\) are algebraically
independent by \cref{thm:all-order-tower}; the triangular change from
\((\sigma_1,\ldots,\sigma_{M+1})\) to
\((d_0^{(4)},\ldots,d_M^{(4)})\) preserves this property.  Telescoping and
\(T_4(Y_N)\to T_4(\gamma)=S_4\) prove
\eqref{eq:S4-positive-sum}.  The doubly exponential convergence of \(Y_N\)
gives summability of every positive power of \(d_N^{(4)}\).  Finally apply
\eqref{eq:gamma-even-recovery}.
\end{proof}

\begin{remark}[Three distinct meanings of ``gamma-free'']
Each finite datum \(Y_N,\tau_N,d_N\) is defined by algebraic-point
\(\Ein\)-values without displaying \(\gamma\).  The limit and sum
\[
 S_2=\lim_N\tau_N=\sum_Nd_N
\]
are countable constructions.  They do not prove that \(S_2\) belongs to the
finite field generated by algebraic-point \(\Ein\)-values.  Indeed
\([z^2]\Delta_2(z)=-S_2/2\), so its Taylor coefficients are already beyond
what the finite-stage algebraic-independence theorem controls.  No claim is
made that \(\Delta_2\) is an \(E\)-function, a \(D\)-finite function, or an
arithmetic entire function.
\end{remark}

\begin{remark}[No arithmetic conclusion from the infinite sum]
Algebraic independence controls finite polynomial relations.  It says
nothing by itself about the arithmetic nature of an infinite sum.  Thus
\eqref{eq:S2-sum} and \eqref{eq:Delta2-trace} do not prove that \(S_2\), or
equivalently \(\gamma\), is transcendental.
\end{remark}

\section{Further moving-level towers}
\label{sec:other-towers}

The dyadic tower is not isolated.  This section gives a fixed-coefficient
power-orbit family and a logarithmic lattice.  They have different decay
rates and different independence defects.

\subsection{Fixed-coefficient power orbits}

Fix an integer \(n\ge2\) and a positive algebraic number \(q>1\).  For
\(k\ge0\), put
\begin{equation}\label{eq:Bnk}
 B_{n,k}(q)=
 \frac{n\Ein(q^k)-\Ein(q^{nk})}{n-1}.
\end{equation}

\begin{theorem}[Power-orbit moving tower]
\label{thm:power-orbit}
For fixed \(n\) and \(q\), the family
\(\{B_{n,k}(q):k\ge0\}\) is countably algebraically independent over \(K\).
Moreover
\begin{equation}\label{eq:eta-nk}
 B_{n,k}(q)=\gamma+\eta_{n,k}(q),\qquad
 \eta_{n,k}(q)=
 \frac{nE_1(q^k)-E_1(q^{nk})}{n-1},
\end{equation}
and
\begin{equation}\label{eq:eta-properties}
 \eta_{n,k}>0,\qquad
 \eta_{n,k}\downarrow0,\qquad
 \eta_{n,k}\sim
 \frac n{n-1}\frac{e^{-q^k}}{q^k}.
\end{equation}
For every choice \(m_k\ge2\), the exact traces
\[
 T_{m_k}(B_{n,k}(q))
\]
are countably algebraically independent over \(K\).  For fixed \(m\), they
converge to \(S_m\).  At \(m=2\) they increase to \(S_2\) and yield, by
successive differences, another positive algebraically independent spectral
realization of \(S_2\).
\end{theorem}

\begin{proof}
Let \(X_j=\Ein(q^j)\).  The distinct algebraic points \(q^j\) give an
algebraically independent coordinate family by
\cref{thm:input-independence}.  The row vectors
\[
 v_k=ne_k-e_{nk}\qquad(k\ge0)
\]
are linearly independent: in a finite relation, the coordinate \(nk_0\)
associated with the largest occurring \(k_0\) appears as the terminal
coordinate of that row and as the initial coordinate of no row in the
finite set.  Thus its coefficient vanishes, and descending induction
finishes the proof.  If the finite set consists only of \(k_0=0\), then
\(v_0=(n-1)e_0\ne0\), which is the same conclusion.  Apply
\cref{lem:linear-images}.

The connection formula \eqref{eq:Ein-connection} cancels the logarithmic
terms and gives \eqref{eq:eta-nk}.  For \(x>1\),
\[
 \frac d{dx}\left(
 \frac{nE_1(x)-E_1(x^n)}{n-1}\right)
 =\frac n{n-1}\frac{e^{-x^n}-e^{-x}}x<0.
\]
This proves strict decrease; positivity follows from
\(E_1(x)>E_1(x^n)>0\), and the asymptotic follows from
\eqref{eq:E1-asymptotic}.  The trace assertions follow from
\cref{cor:transversal,thm:universal-trace}.  The order-two construction is
the proof of \cref{thm:increments,thm:Delta2} applied to this decreasing
level sequence: indeed \(B_{n,0}(q)=\Ein(1)<1<2\), since
\((1-e^{-t})/t<1\) on \((0,1]\).
\end{proof}

This tower samples the rank-one witness of the companion paper exactly:
\begin{equation}\label{eq:B-Gn}
 B_{n,k}(q)-\gamma=\frac{G_n(q^k)}{n-1}.
\end{equation}

\begin{remark}[Fixing \(n\) is essential]
The full two-parameter family is not algebraically independent.  For
\(r\ge1\),
\begin{equation}\label{eq:cross-n-relation}
 B_{n^r,k}(q)=\frac{n-1}{n^r-1}
 \sum_{j=0}^{r-1}n^{r-1-j}B_{n,n^jk}(q).
\end{equation}
For example \(B_{4,k}=(2B_{2,k}+B_{2,2k})/3\).
\end{remark}

\subsection{Algebraic-level inverse transcendence}

The first difference between two members of the \(n=2\) power orbit gives
the gamma-free function
\begin{equation}\label{eq:scriptR-def}
 \mathscr R(q)=E_1(1)-2E_1(q)+E_1(q^2)\qquad(q>1).
\end{equation}
It has an unexpectedly rigid inverse on algebraic levels.

\begin{theorem}[Transcendental inverse levels]
\label{thm:inverse-levels}
The map \(\mathscr R:(1,\infty)\to(0,E_1(1))\) is a strictly increasing
bijection.  If \(a\in\Qbar\cap(0,E_1(1))\), then the unique number
\(r_a>1\) satisfying \(\mathscr R(r_a)=a\) is transcendental.

In particular, for every integer \(m\ge5\), the unique root \(r_m>1\) of
\begin{equation}\label{eq:rm-level}
 E_1(1)-2E_1(r_m)+E_1(r_m^2)=\frac1m
\end{equation}
is transcendental, and
\begin{equation}\label{eq:rm-asymptotic}
 r_m=1+\sqrt{\frac em}+\frac e{2m}+O(m^{-3/2}).
\end{equation}
\end{theorem}

\begin{proof}
Direct differentiation gives
\[
 \mathscr R'(q)=\frac2q(e^{-q}-e^{-q^2})>0.
\]
Moreover \(\mathscr R(1)=0\) and
\(\mathscr R(q)\to E_1(1)\) as \(q\to\infty\), proving the bijection.
If \(r_a\) were algebraic, the connection formula would turn
\(\mathscr R(r_a)=a\) into
\[
 \Ein(1)-2\Ein(r_a)+\Ein(r_a^2)=a.
\]
The three arguments are distinct nonzero algebraic numbers, so this relation
contradicts \cref{thm:input-independence}.

Splitting the defining integral at (2) and (3), and using
(1/t>1/2) on (1<t<2) and (1/t>1/3) on (2<t<3), gives
\[
 E_1(1)>\frac1{2e}+\frac{e^{-2}-e^{-3}}3>\frac15.
\]
Thus \(a=1/m\) is admissible for \(m\ge5\).  Finally, writing \(q=1+h\),
direct differentiation at one gives
\[
 \mathscr R(1+h)=e^{-1}(h^2-h^3+O(h^4)).
\]
Series inversion at \(\mathscr R=1/m\) proves
\eqref{eq:rm-asymptotic}.
\end{proof}

\subsection{A logarithmic lattice and density-one transcendence}

For \(r\ge1\), put
\begin{equation}\label{eq:cr}
 c_r=\Cin(r)-\log r=\gamma-\Ci(r).
\end{equation}
These levels approach \(\gamma\) only polynomially, but the loss caused by
the logarithms can be counted by primes.

\begin{lemma}\label{lem:cr-positive}
One has \(c_r>0\) for every integer \(r\ge1\).
\end{lemma}

\begin{proof}
For \(x>0\), set \(c(x)=\gamma-\Ci(x)\).  Since
\(c'(x)=-\cos x/x\), its local minima occur at
\(x_k=\pi/2+2\pi k\).  At the first minimum,
\[
 \Cin(\pi/2)
 \ge \int_0^{\pi/2}\left(\frac t2-\frac{t^3}{24}\right)dt
 =\frac{\pi^2}{16}-\frac{\pi^4}{1536}
 >\log(\pi/2),
\]
using \(1-\cos t\ge t^2/2-t^4/24\) on this interval.  The last strict
inequality follows, for example, from
\(333/106<\pi<355/113\) and
\(\log(1+u)\le u-u^2/2+u^3/3\) for \(0<u<1\).
Thus
\(c(x_0)>0\).  For \(k\ge1\), integration by parts gives
\(\Ci(x_k)\le2/x_k<1/2<\gamma\).  Here we used
\(x_k\ge5\pi/2\) and the elementary bound \(\gamma>1/2\).
Hence all local minima are positive, and so is
\(c(x)\).  In particular \(c_r>0\).
\end{proof}

For each \(r\), let
\begin{equation}\label{eq:Ur}
 U_r=T_2(c_r)=\frac1{c_r^2}-\frac1{2c_r}.
\end{equation}

\begin{theorem}[Prime-defect lattice theorem]
\label{thm:prime-defect}
Let \(\pi(M)\) be the number of primes not exceeding \(M\).  Then
\begin{equation}\label{eq:prime-defect-trdeg}
 \trdeg_KK(c_1,\ldots,c_M)
 =\trdeg_KK(U_1,\ldots,U_M)
 \ge M-\pi(M).
\end{equation}
Consequently at most \(\pi(M)\) of the first \(M\) members of each family
\(\{c_r\}\) and \(\{U_r\}\) are algebraic over \(K\).  Both families are
\(K\)-transcendental on a set of natural density one.
\end{theorem}

\begin{proof}
The companion result \cite[Cor.~5.5]{KEIO2026I} makes
\(\Cin(1),\ldots,\Cin(M)\) algebraically independent over \(K\).
Equivalently, this follows from
\(\Cin(r)=\{\Ein(ir)+\Ein(-ir)\}/2\): the pairs
\(\{ir,-ir\}\) are disjoint, so these averaging rows have full rank.
Since
\[
 \log r=\sum_{p\le M}v_p(r)\log p,
\]
we have
\[
 K(\Cin(1),\ldots,\Cin(M))
 \subset K(c_1,\ldots,c_M,\log p:p\le M).
\]
Therefore
\[
 M\le \trdeg_KK(c_1,\ldots,c_M)+\pi(M),
\]
which gives the lower bound for the \(c_r\).  Equation
\[
 2U_rc_r^2+c_r-2=0
\]
shows that \(K(c_1,\ldots,c_M)\) is algebraic over
\(K(U_1,\ldots,U_M)\), while the reverse field inclusion is rational.
Thus the two transcendence degrees agree.  If \(a\) of the first \(M\)
coordinates were algebraic over \(K\), the transcendence degree would be at
most \(M-a\); hence \(a\le\pi(M)\).  The prime number theorem completes the
density statement.
\end{proof}

The trace \(U_r\) also converges to \(S_2\).  Direct subtraction gives
\begin{equation}\label{eq:Ur-S2-exact}
 U_r-S_2=
 \frac{\Ci(r)\{\gamma(4-\gamma)-(2-\gamma)\Ci(r)\}}
 {2\gamma^2(\gamma-\Ci(r))^2},
\end{equation}
and the standard expansion of \(\Ci\) yields
\begin{align}
 U_r-S_2
 &=\frac{4-\gamma}{2\gamma^3}\frac{\sin r}{r}
 -\frac{4-\gamma}{2\gamma^3}\frac{\cos r}{r^2}\notag\\
 &\quad+
 \frac{6-\gamma}{2\gamma^4}\frac{\sin^2r}{r^2}
 +O(r^{-3}).\label{eq:Ur-S2-asymptotic}
\end{align}
Thus this lattice is generally only \(O(r^{-1})\) from the Euler trace.
The theorem does not decide any specified \(c_r,U_r\) with \(r\ge2\);
\(c_1=\Cin(1)\) is already \(K\)-transcendental.

\section{Normalized certificates and cross-index rigidity}
\label{sec:certificates}

The moving towers concern limits.  We now return to finite witnesses and
show that all normalized linear certificates are unique.

Let
\[
 V=\bigoplus_{\alpha\in\Qbar^\times}K^{\mathrm{alg}}[\alpha]
\]
be the free \(K^{\mathrm{alg}}\)-vector space with finite support on the
nonzero algebraic points, and define
\[
 \Phi:V\longrightarrow\C,
 \qquad \Phi([\alpha])=\Ein(\alpha).
\]

\begin{lemma}[The evaluation intersection]
\label{lem:evaluation-intersection}
The map \(\Phi\) is injective and
\begin{equation}\label{eq:Phi-intersection}
 \Phi(V)\cap K^{\mathrm{alg}}=\{0\}.
\end{equation}
\end{lemma}

\begin{proof}
By \cref{thm:input-independence}, every finite collection of the displayed
\(\Ein\)-values is algebraically independent over \(K^{\mathrm{alg}}\), and
hence linearly independent.  Moreover any nonzero linear form in
algebraically independent variables is transcendental over the base field:
extend it to an invertible linear coordinate system.  This proves both
claims.
\end{proof}

\begin{definition}
A finite \(K^{\mathrm{alg}}\)-certificate for \(\gamma\) is a triple
\((v,c,a)\in V\times(K^{\mathrm{alg}})^\times\times K^{\mathrm{alg}}\)
satisfying
\begin{equation}\label{eq:certificate}
 \Phi(v)=c\gamma+a.
\end{equation}
Its normalized data are \((v/c,a/c)\).
\end{definition}

\begin{theorem}[Normalized certificate uniqueness]
\label{thm:certificate-uniqueness}
If \((v_1,c_1,a_1)\) and \((v_2,c_2,a_2)\) are two certificates, then
\begin{equation}\label{eq:normalized-certificate-unique}
 \boxed{\displaystyle
 \frac{v_1}{c_1}=\frac{v_2}{c_2},\qquad
 \frac{a_1}{c_1}=\frac{a_2}{c_2}.}
\end{equation}
\end{theorem}

\begin{proof}
Subtracting the two normalized identities gives
\[
 \Phi\!\left(\frac{v_1}{c_1}-\frac{v_2}{c_2}\right)
 =\frac{a_1}{c_1}-\frac{a_2}{c_2}\in K^{\mathrm{alg}}.
\]
Apply \eqref{eq:Phi-intersection}.
\end{proof}

This is an exclusion theorem, not an existence theorem.  The certificate
set may be empty.  Its force is that whenever two apparently unrelated
algebraicity events would produce certificates, their normalized supports
must coincide exactly.

We apply it to the real Euler witnesses recalled in the introduction.  The
existence, location and within-index algebraic-zero rigidity of \(q_*\),
\(\beta_n\) and \(\lambda_n\) were established in
\cite[Thm.~14.2, Prop.~14.3, and Thm.~14.4]{KEIO2026I}.

\begin{theorem}[Cross-index witness rigidity]
\label{thm:cross-index}
At most one member of
\begin{equation}\label{eq:witness-set}
 \mathcal W=
 \{q_*\}\cup\{\beta_n:n\ge2\}
 \cup\{\lambda_n:n\ge3,\ n\text{ odd}\}
\end{equation}
is algebraic.  Thus precisely one of the following alternatives holds:
\begin{enumerate}[label=\textup{(\roman*)}]
 \item every member of \(\mathcal W\) is transcendental;
 \item exactly one member is algebraic, and in that event \(\gamma\) and
 \(\delta\) are algebraically independent over \(K\).
\end{enumerate}
\end{theorem}

\begin{proof}
If \(q_*\) is algebraic, its normalized certificate vector is
\[
 v_*=[q_*],\qquad \Phi(v_*)=\gamma.
\]
If an algebraic number \(\alpha\) is a zero of \(G_n\), its normalized
certificate vector is
\begin{equation}\label{eq:G-certificate-vector}
 v_{n,\alpha}
 =\frac n{n-1}[\alpha]-\frac1{n-1}[\alpha^n],
 \qquad \Phi(v_{n,\alpha})=\gamma.
\end{equation}
The real-zero classification places \(\alpha\) in \((0,1)\) or, for odd
\(n\), in \((-\infty,-1)\).  Hence \(\alpha\ne\alpha^n\), and
\(v_{n,\alpha}\) has exactly one positive and one negative support
coefficient.  It cannot equal the one-point vector \(v_*\).

If \(v_{n,\alpha}=v_{m,\beta}\), equality of the positive and negative
support coordinates gives
\[
 \frac n{n-1}=\frac m{m-1},
\]
so \(n=m\), and then \(\alpha=\beta\).  Certificate uniqueness therefore
excludes two distinct algebraic members of \eqref{eq:witness-set}.  The
algebraic independence of \(\gamma,\delta\) in the exceptional case is
\cite[Thm.~14.2 and Thm.~14.4]{KEIO2026I}.
\end{proof}

\subsection{A defect-one family of critical cancellation coefficients}

Fix an algebraic number \(q>1\).  For \(N\ge1\), put
\begin{align}
 D_N(q)&=E_1(q^N)-2E_1(q^{N+1})+E_1(q^{N+2}),\label{eq:DNq}\\
 A_N(q)&=(N+2)E_1(q^{N+1})-(N+1)E_1(q^{N+2}),\label{eq:ANq}\\
 u_N(q)&=\frac{A_N(q)}{D_N(q)}.\label{eq:uNq}
\end{align}
Both numerator and denominator are positive: \(D_N(q)>0\) is the strict
second-difference inequality for the strictly convex function
\(t\mapsto E_1(q^t)\), while monotonicity gives
\(A_N(q)>E_1(q^{N+1})>0\).  The number \(u_N(q)\) is the
unique positive coefficient for which the three-point augmentation-one tail
vanishes:
\begin{equation}\label{eq:uN-tail-zero}
 -u_NE_1(q^N)+(N+2+2u_N)E_1(q^{N+1})
 -(N+1+u_N)E_1(q^{N+2})=0.
\end{equation}

\begin{theorem}[One-defect cancellation family]
\label{thm:uN-defect}
For every finite set \(I\subset\N_{\ge1}\),
\begin{equation}\label{eq:uN-trdeg}
 \boxed{\displaystyle
 \trdeg_K K(u_N(q):N\in I)\ge |I|-1.}
\end{equation}
Consequently at most one member of \(\{u_N(q):N\ge1\}\) is algebraic over
\(K\).  If one member is algebraic over \(K\), all the remaining members are
countably algebraically independent over \(K\), and \(\gamma,\delta\) are
algebraically independent over \(K\).  If \(\gamma\) is algebraic over
\(K\), then the whole family \(\{u_N(q):N\ge1\}\) is countably
algebraically independent over \(K\).

Finally,
\begin{equation}\label{eq:uN-asymptotic}
 u_N(q)\sim\frac{N+2}{q}
 \exp\!\bigl(-(q-1)q^N\bigr).
\end{equation}
\end{theorem}

\begin{proof}
Let \(X_j=\Ein(q^j)\).  The connection formula and
\eqref{eq:uN-tail-zero} give
\begin{equation}\label{eq:uN-certificate}
 -u_NX_N+(N+2+2u_N)X_{N+1}
 -(N+1+u_N)X_{N+2}=\gamma.
\end{equation}
For finite \(I\), let
\(U=\bigcup_{N\in I}\{N,N+1,N+2\}\).  Since \(u_N>0\), the relations
\eqref{eq:uN-certificate}, used in decreasing order of \(N\), solve every
\(X_N\) with \(N\in I\) in terms of \(u_N\), \(\gamma\), and the
coordinates \(X_j\) with \(j\in U\setminus I\).  Hence
\[
 K(X_j:j\in U)
 \subset
 K(u_N:N\in I,\gamma,X_j:j\in U\setminus I).
\]
The left-hand coordinates are algebraically independent over \(K\), so
\[
 |U|\le
 \trdeg_KK(u_N:N\in I)+1+|U|-|I|,
\]
which proves \eqref{eq:uN-trdeg}.  If one \(u_{N_0}\) is algebraic over
\(K\), apply the bound to \(I\cup\{N_0\}\) to obtain full independence of
every finite set with \(N_0\) omitted.

In that exceptional case \eqref{eq:uN-certificate} expresses \(\gamma\) as
a nonconstant \(K^{\rm alg}\)-linear form in three algebraically independent
\(\Ein\)-coordinates.  Together with
\(\delta/e=\Ein(1)-\gamma\), where \(1\) is a fourth distinct algebraic
point, two linearly independent coordinate forms are obtained.  Thus
\(\gamma,\delta\) are algebraically independent over \(K\).  If instead
\(\gamma\in K^{\rm alg}\), the extra ``\(+1\)'' in the preceding
transcendence-degree estimate disappears, proving full independence of the
\(u_N\).

The asymptotic follows from \eqref{eq:E1-asymptotic}: the first terms in
\eqref{eq:DNq} and \eqref{eq:ANq} dominate and
\[
 \frac{E_1(q^{N+1})}{E_1(q^N)}
 \sim q^{-1}\exp\!\bigl(-(q-1)q^N\bigr).
\]
\end{proof}

\section{Relative exclusion and finite trace descent}
\label{sec:relative}

Put
\[
 A=\Ein(1),\qquad B=\Ein(-1),\qquad
 C=\Cin(1)=\frac{\Ein(i)+\Ein(-i)}2.
\]
The numbers (A,B,C) are algebraically independent over \(K\), and
\begin{align}
 E_1(1)&=A-\gamma,&
 \Ei(1)&=\gamma-B,\label{eq:six-identities}\\
 \Ci(1)&=\gamma-C,&
 \Chi(1)&=\gamma-\frac{A+B}{2},&
 \delta&=e(A-\gamma).\notag
\end{align}

\begin{theorem}[Five relative certificate classes]
\label{thm:five-classes}
Among the following five assertions, at most one can hold:
\begin{align*}
 \mathcal C_0&:\ \gamma\in K^{\mathrm{alg}},\\
 \mathcal C_1&:\ E_1(1)\in K^{\mathrm{alg}}
 \quad\Longleftrightarrow\quad \delta\in K^{\mathrm{alg}},\\
 \mathcal C_2&:\ \Ei(1)\in K^{\mathrm{alg}},\\
 \mathcal C_3&:\ \Ci(1)\in K^{\mathrm{alg}},\\
 \mathcal C_4&:\ \Chi(1)\in K^{\mathrm{alg}}.
\end{align*}
Consequently at least four members of
\[
 \gamma,\quad\delta,\quad E_1(1),\quad\Ei(1),\quad\Ci(1),\quad\Chi(1)
\]
are transcendental over \(K\).  If the possible exceptional class is not
\(\mathcal C_1\), then at least five are transcendental over \(K\).
\end{theorem}

\begin{proof}
The equivalence in \(\mathcal C_1\) follows from
\(\delta=eE_1(1)\) and \(e\in K^\times\).  The five assertions produce the
following normalized certificate vectors, respectively:
\[
 0,\qquad [1],\qquad[-1],\qquad
 \frac{[i]+[-i]}2,\qquad
 \frac{[1]+[-1]}2.
\]
For example, if \(E_1(1)\in K^{\mathrm{alg}}\), then
\(\Phi([1])=\gamma+E_1(1)\); if \(\Ci(1)\in K^{\mathrm{alg}}\), then
\(\Phi(([i]+[-i])/2)=\gamma-\Ci(1)\).  These five vectors are distinct, so
\cref{thm:certificate-uniqueness} permits at most one assertion.  Counting
the sizes (1,2,1,1,1) of the five classes gives the final statement.
\end{proof}

\begin{remark}[The necessary correction]
The theorem does \emph{not} say that at most one of the six constants is
algebraic over \(K\).  The pair \(E_1(1),\delta\) forms one relative class:
they are simultaneously algebraic or transcendental over \(K\).  Over
\(\Qbar\), the companion paper excludes their simultaneous algebraicity by
the additional identity \(e=\delta/E_1(1)\) and the transcendence of \(e\).
\end{remark}

\begin{corollary}[Witnesses exclude every relative class]
\label{cor:witness-all-K-trans}
If one member of the witness set \(\mathcal W\) in
\eqref{eq:witness-set} is algebraic, then all six constants displayed in
\cref{thm:five-classes} are transcendental over \(K\).
\end{corollary}

\begin{proof}
The normalized vector of an algebraic witness is either \([q_*]\) or a
two-point vector with one positive and one negative coefficient as in
\eqref{eq:G-certificate-vector}.  Its support and sign pattern differ from
each of the five vectors in the proof of \cref{thm:five-classes}.  Certificate
uniqueness therefore excludes all five relative assertions.
\end{proof}

We next separate finite field membership from countable analytic limits.
Let
\begin{equation}\label{eq:G-field}
 \calG=K\bigl(\Ein(\alpha):\alpha\in\Qbar^\times\bigr).
\end{equation}
Every element of \(\calG\) uses only finitely many algebraic-point values,
although the generating family is countable.

\begin{proposition}[Relative trace descent]
\label{prop:relative-trace-descent}
For \(m\ge2\),
\[
 S_m=P_m(\gamma^{-1}),\qquad
 \calG(S_m)\subset\calG(\gamma),\qquad
 [\calG(\gamma):\calG(S_m)]\le m.
\]
Consequently
\begin{equation}\label{eq:relative-trdeg}
 \trdeg_{\calG}\calG(S_m)=\trdeg_{\calG}\calG(\gamma).
\end{equation}
Moreover
\begin{equation}\label{eq:G-S2-S3}
 \boxed{\calG(S_2,S_3)=\calG(S_2,S_4)=\calG(\gamma)},
\end{equation}
and hence
\begin{equation}\label{eq:joint-gamma-free}
 S_2,S_3\in\calG
 \quad\Longleftrightarrow\quad
 S_2,S_4\in\calG
 \quad\Longleftrightarrow\quad
 \gamma\in\calG.
\end{equation}
Individually,
\[
 S_2\in\calG\Longrightarrow[\calG(\gamma):\calG]\le2,
 \qquad
 S_3\in\calG\Longrightarrow[\calG(\gamma):\calG]\le3,
 \qquad
 S_4\in\calG\Longrightarrow[\calG(\gamma):\calG]\le4.
\]
\end{proposition}

\begin{proof}
The polynomial identity comes from \cref{thm:universal-trace} at the Euler
level.  Thus \(\gamma^{-1}\) is algebraic of degree at most \(m\) over
\(\calG(S_m)\), proving the first assertions and
\eqref{eq:relative-trdeg}.  Equations \eqref{eq:gamma-recovery} and
\eqref{eq:gamma-even-recovery} prove \eqref{eq:G-S2-S3}; the remaining
statements follow.
\end{proof}

The reverse of the first individual implication for \(S_2\) is false without
an additional norm--trace condition.

\begin{theorem}[Exact quadratic descent criterion]
\label{thm:quadratic-descent}
One has \(S_2\in\calG\) if and only if either \(\gamma\in\calG\), or
\[
 [\calG(\gamma):\calG]=2
 \quad\text{and}\quad
 \Norm_{\calG(\gamma)/\calG}(\gamma)
 =2\Tr_{\calG(\gamma)/\calG}(\gamma).
\]
In the quadratic case the conjugate is
\begin{equation}\label{eq:gamma-conjugate}
 \gamma'=-\frac{2\gamma}{2-\gamma}.
\end{equation}
\end{theorem}

\begin{proof}
Let \(f(x)=x^{-2}-(2x)^{-1}\).  If \(\gamma\notin\calG\) and
\(f(\gamma)=S_2\in\calG\), then \cref{prop:relative-trace-descent} gives a
quadratic extension.  If \(\gamma'\) is its nontrivial conjugate, then
\[
 f(\gamma)-f(\gamma')
 =\frac{(\gamma'-\gamma)
 \{2(\gamma+\gamma')-\gamma\gamma'\}}
 {2\gamma^2{\gamma'}^2}.
\]
Thus invariance of \(f(\gamma)\) is equivalent to
\(\gamma\gamma'=2(\gamma+\gamma')\), which is exactly the displayed
norm--trace condition and rearranges to \eqref{eq:gamma-conjugate}.
Conversely that condition makes \(f(\gamma)\) fixed by the quadratic Galois
involution and hence an element of \(\calG\).
\end{proof}

For instance, degree two alone is insufficient: with base field \(\Q\) and
\(x=\sqrt2\), one has
\(f(x)=1/2-\sqrt2/4\notin\Q\).

\part{Finite Faber Monodromy and Barrier Escape}
\section{Uniform monodromy of polyexponential trace polynomials}

\subsection{Scope and status}

This note proves an all-order structural theorem for the trace polynomials
attached to every integer-order polyexponential. The proof has four parts:

\begin{enumerate}
\def\labelenumi{\arabic{enumi}.}
\tightlist
\item
  normalized Faber rigidity turns a polynomial decomposition into a
  vanishing coefficient of a reciprocal series;
\item
  a uniform 2-adic argument proves that none of those reciprocal
  coefficients vanishes;
\item
  two centered affine invariants exclude the cyclic, dihedral, and three
  exceptional rational polynomial-monodromy cases in Müller's
  classification;
\item
  the discriminant separates the geometric and arithmetic conclusions.
\end{enumerate}

The result does \textbf{not} prove that every trace polynomial is Morse. It proves
indecomposability for every pair of indices and reduces its monodromy to the
alternating/symmetric dichotomy, with the precise conclusions stated below.

Throughout, \(q\geq1\) and \(m\geq2\) are integers, and

\[
 E_q(z):=\operatorname {Ein}_q(z)
 =\sum_{n\geq1}\frac{(-1)^{n-1}}{n!\,n^q}z^n,
\]

\[
 P_{q,m}(X):=-m[z^m]\log(1-XE_q(z)).
\]

The coefficient formula

\[
 [X^k]P_{q,m}(X)=\frac{m}{k}[z^m]E_q(z)^k
\]

shows that \(P_{q,m}\in\mathbf Q[X]\), that its constant term is zero, and
that it is monic of degree \(m\).

\begin{center}\rule{0.5\linewidth}{0.5pt}\end{center}

\subsection{1. A general normalized-Faber rigidity lemma}

We first work with an arbitrary normalized formal series

\[
 f(z)=z+\sum_{n\geq2}a_nz^n
\]

over a characteristic-zero field and define

\[
 P_m^{(f)}(X):=-m[z^m]\log(1-Xf(z)).
\]

Put

\[
 M(w):=\frac1{f(1/w)}
 =w+d_0+d_1w^{-1}+d_2w^{-2}+\cdots .
\]

Let \(\Phi_m\) be the standard Faber polynomial of \(M\), normalized by

\[
 \log\frac{M(w)-X}{w}
 =-\sum_{m\geq1}\frac{\Phi_m(X)}m w^{-m}.
\]

\subsubsection{Lemma 1 (normalized Faber identity)}

For every \(m\geq1\),

\[
 P_m^{(f)}(X)=\Phi_m(X)-\Phi_m(0).
\]

\paragraph{Proof}

Since \(f(1/w)=M(w)^{-1}\),

\[
 \log(1-Xf(1/w))
 =\log\frac{M(w)-X}{w}-\log\frac{M(w)}w.
\]

Comparing the coefficient of \(w^{-m}\) gives the formula. \(\square\)

\subsubsection{Lemma 2 (inner-factor rigidity)}

Suppose \(m=ab\) with \(a,b>1\). If

\[
 P_m^{(f)}=G\circ H
\]

over a characteristic-zero extension field, then the decomposition may be
normalized so that \(G,H\) are monic, \(G(0)=H(0)=0\), and under this
normalization

\[
 H=P_b^{(f)}.
\]

\paragraph{Proof}

The Faber identity gives

\[
 P_m^{(f)}(M(w))=w^m-\Phi_m(0)+O(w^{-1}).
\]

Let \(\Psi\) be the formal inverse branch of the monic degree-\(a\)
polynomial \(G\) at infinity. It has a Puiseux expansion

\[
 \Psi(U)=U^{1/a}+c_0+c_1U^{-1/a}+c_2U^{-2/a}+\cdots .
\]

Substitution of \(U=w^{ab}+O(1)\) shows

\[
 H(M(w))=w^b+\text{constant}+O(w^{-1});
\]

in particular, there are no powers \(w^{b-1},\ldots,w\). The defining
uniqueness of the degree-\(b\) Faber polynomial implies that \(H\) differs
from \(\Phi_b\) by a constant. The normalization \(H(0)=0\) then gives
\(H=\Phi_b-\Phi_b(0)=P_b^{(f)}\). \(\square\)

\subsubsection{Lemma 3 (reciprocal-coefficient obstruction)}
\label{sec:reciprocal-coefficient-obstruction}

With

\[
 \frac1{f(z)}=z^{-1}+d_0+d_1z+d_2z^2+\cdots,
\]

a decomposition of \(P_m^{(f)}\) having inner degree \(b\) forces

\[
 d_b=[z^b]\frac1{f(z)}=0.
\]

\paragraph{Proof}

The first negative Faber coefficient is

\[
 \Phi_b(M(w))=w^b+b d_bw^{-1}+O(w^{-2}).
\]

Here is a formal-residue verification. If \(L\) is the compositional
inverse of \(M\) at infinity, then \(\Phi_b(X)\) is the polynomial part of
\(L(X)^b\). Formal change of variables gives

\[
 [X^{-1}]L(X)^b
 =\operatorname {res}_X L(X)^b\,dX
 =\operatorname {res}_w w^bM'(w)\,dw
 =-bd_b.
\]

Substituting \(X=M(w)\) into the polynomial/negative-power decomposition of
\(L(X)^b\) therefore gives the displayed coefficient \(+bd_b\). Hence

\[
 P_b^{(f)}(M(w))
 =w^b-\Phi_b(0)+bd_bw^{-1}+O(w^{-2}).
\]

In \(G(P_b^{(f)}(M(w)))\), the coefficient at

\[
 w^{m-b-1}=w^{b(a-1)-1}
\]

has the unique contribution \(ab d_b=md_b\) from the leading power of
\(G\). Indeed, a constant selection creates deficit \(b\), a negative term
\(w^{-j}\) creates deficit \(b+j\), and there are no terms of degrees
\(b-1,\ldots,1\); neither another selection in the leading outer power nor
a lower outer power can reach deficit \(b+1\). On the other hand,
\(P_m^{(f)}(M(w))\) has no positive power below \(w^m\). Therefore
\(md_b=0\), and characteristic zero gives \(d_b=0\). \(\square\)

\subsubsection{Remark (generic jet coordinates)}

If \(c_k=[X^k]P_m^{(f)}\), then

\[
 c_k=ma_{m-k+1}+Q_k(a_2,\ldots,a_{m-k}).
\]

Thus

\[
 (a_2,\ldots,a_m)\longmapsto(c_{m-1},\ldots,c_1)
\]

is a triangular polynomial automorphism with diagonal entries \(m\). The
specialized polyexponential problem is therefore not caused by a defect in
the universal trace construction: the general \((m-1)\)-jet produces the
generic monic degree-\(m\), zero-constant polynomial.

\begin{center}\rule{0.5\linewidth}{0.5pt}\end{center}

\subsection{2. Uniform 2-adic nonvanishing}

Write

\[
 A_q(z):=\frac{E_q(z)}z
 =\sum_{n\geq0}\frac{(-1)^n}{n!(n+1)^{q+1}}z^n,
 \qquad
 B_q(z):=\frac1{A_q(z)}=\sum_{n\geq0}B_{q,n}z^n.
\]

\subsubsection{Theorem 4 (exact reciprocal valuations)}

For all \(q\geq1\) and \(n\geq0\),

\[
 v_2(B_{q,n})=-(q+1)n.
\]

In particular, every \(B_{q,n}\) is nonzero.

\paragraph{Proof}

Let \(s_2(n)\) be the sum of the binary digits of \(n\), and put
\(r=v_2(n+1)\). Legendre's formula \(v_2(n!)=n-s_2(n)\) gives

\[
 v_2\!\left(
 \frac{2^{(q+1)n}}{n!(n+1)^{q+1}}
 \right)
 =qn+s_2(n)-(q+1)r.
\]

This value is zero for \(n=0,1\). If \(n\geq2\) and \(r=0\), it is
positive. If \(r\geq1\), the final \(r\) binary digits of \(n\) are all
one, so \(s_2(n)\geq r\), and

\[
 qn+s_2(n)-(q+1)r
 =q(n-r)+(s_2(n)-r)>0.
\]

It follows that

\[
 A_q(2^{q+1}z)\equiv1+z\pmod2.
\]

Taking formal reciprocals in \(\mathbf Z_2[[z]]\),

\[
 B_q(2^{q+1}z)
 \equiv(1+z)^{-1}
 =\sum_{n\geq0}z^n\pmod2.
\]

Thus every \(2^{(q+1)n}B_{q,n}\) is a 2-adic unit, which is exactly the
claimed valuation. \(\square\)

Since

\[
 \frac1{E_q(z)}=z^{-1}B_q(z),
 \qquad
 [z^b]\frac1{E_q(z)}=B_{q,b+1},
\]

Lemmas 2--3 immediately give the central structural result.

\subsubsection{Theorem 5 (all-order indecomposability)}

For every \(q\geq1\) and \(m\geq2\), the polynomial \(P_{q,m}\) is
indecomposable over \(\overline{\mathbf Q}\), and hence over every
characteristic-zero subfield.

\begin{center}\rule{0.5\linewidth}{0.5pt}\end{center}

\subsection{3. Centered invariants and the cyclic/dihedral cases}

For a general series

\[
 f(z)=z+a_2z^2+a_3z^3+a_4z^4+\cdots,
\]

the first upper coefficients of \(P_m^{(f)}\) are

\[
\begin{aligned}
P_m^{(f)}(X)={}&X^m+ma_2X^{m-1}
 +m\left(a_3+\frac{m-3}{2}a_2^2\right)X^{m-2}\\
&+m\left(a_4+(m-4)a_2a_3
 +\frac{(m-4)(m-5)}6a_2^3\right)X^{m-3}+\cdots .
\end{aligned}
\]

The last displayed term and the corresponding formula below are understood
for \(m\ge4\); the \(Y^{m-2}\) formula is valid for \(m\ge3\).

Center it by

\[
 \widehat P_m(Y):=P_m^{(f)}(Y-a_2)-P_m^{(f)}(-a_2).
\]

A direct expansion yields

\[
 [Y^{m-2}]\widehat P_m=m(a_3-a_2^2),
\]

\[
 [Y^{m-3}]\widehat P_m=m(a_2^3-2a_2a_3+a_4).
\]

For \(f=E_q\),

\[
 a_2=-\frac1{2^{q+1}},\qquad
 a_3=\frac1{2\cdot3^{q+1}},\qquad
 a_4=-\frac1{3\cdot2^{2q+3}}.
\]

Define

\[
 N_A(q):=2^{2q+1}-3^{q+1},
 \qquad
 N_B(q):=2^{2q+2}-3^{q+1}-6^q.
\]

Then

\[
 a_3-a_2^2
 =\frac{N_A(q)}{2^{2q+2}3^{q+1}},
\]

\[
 a_2^3-2a_2a_3+a_4
 =\frac{N_B(q)}{2^{3q+3}3^{q+1}}.
\]

The first numerator is nonzero by unique factorization, and

\[
 N_B(q)\equiv1\pmod3.
\]

Thus both centered coefficients are nonzero. A polynomial with cyclic
monodromy is linearly equivalent to a power map, whose centered coefficient
of degree \(m-2\) vanishes, so the cyclic case is excluded for \(m\ge3\).
A polynomial with dihedral monodromy is linearly
equivalent to a Chebyshev polynomial, whose centered coefficient of degree
\(m-3\) vanishes, so the proper dihedral case is excluded for \(m\ge4\).
In degree \(3\), the dihedral group is \(D_3\cong S_3\), while the cyclic
case has already been excluded. Degree \(2\) is immediate.

\begin{center}\rule{0.5\linewidth}{0.5pt}\end{center}

\subsection{4. Excluding Müller's exceptional degrees 6, 9, and 10}

For a monic centered polynomial

\[
 F(Y)=Y^m+C_2Y^{m-2}+C_3Y^{m-3}+\cdots,
 \qquad C_2\ne0,
\]

the quantity

\[
 \mathcal I(F):=\frac{C_3^2}{C_2^3}
\]

is invariant under linear equivalence: after monic normalization, input
scaling has weights \(2\) and \(3\) on \(C_2,C_3\) (with reciprocal weights
under the opposite input-scaling convention), and output scaling is fixed by
making the polynomial monic. In either convention the displayed ratio is
unchanged.

For \(P_{q,m}\), the preceding formulas give

\[
 \mathcal I(P_{q,m})=\frac{R_q}{m},
 \qquad
 R_q:=3^{q+1}\frac{N_B(q)^2}{N_A(q)^3}.
\]

Müller's three exceptional rational polynomials have the following centered
invariants:

\begin{longtable}[]{@{}rlr@{}}
\toprule\noalign{}
degree & representative & \(\mathcal I\) \\
\midrule\noalign{}
\endhead
\bottomrule\noalign{}
\endlastfoot
6 & \(X^4(X^2+6X+25)\) & \(18/5\) \\
9 & \(9X^9+108X^7+72X^6+\cdots\) & \(1/27\) \\
10 & \((X^2-405)^4(X^2+50X+945)\) & \(-8/81\) \\
\end{longtable}

For completeness, the centered pairs \((C_2,C_3)\) are respectively

\[
 (10,-60),\qquad(12,8),\qquad(-1800,-24000).
\]

Both \(N_A(q)\) and \(N_B(q)\) are 3-adic units. Moreover \(R_1=-9\),
while \(R_q>0\) for \(q\geq2\): indeed \(N_A(2)=5>0\), and the ratio of
the two exponential terms increases by a factor \(4/3\) with \(q\). We can
now exclude each row:

\begin{itemize}
\tightlist
\item
  Degree \(10\) would require \(R_q=-80/81\). The sign excludes
  \(q\geq2\), and \(R_1=-9\ne-80/81\).
\item
  Degree \(9\) would require \(R_q=1/3\), impossible because
  \(v_3(R_q)=q+1\).
\item
  Degree \(6\) would require \(R_q=108/5\). Comparing 3-adic valuations
  forces \(q=2\); but \(N_A(2)=5\), \(N_B(2)=1\), and
  \(R_2=27/125\ne108/5\).
\end{itemize}

Therefore none of Müller's exceptional rational monodromy polynomials is
linearly equivalent to any \(P_{q,m}\).

\begin{center}\rule{0.5\linewidth}{0.5pt}\end{center}

\subsection{5. Monodromy theorem}

Let

\[
 G^{\mathrm{geom}}_{q,m}
 :=\operatorname {Gal}
 \bigl(P_{q,m}(X)-T\,/\,\overline{\mathbf Q}(T)\bigr),
\]

\[
 G^{\mathrm{arith}}_{q,m}
 :=\operatorname {Gal}
 \bigl(P_{q,m}(X)-T\,/\,\mathbf Q(T)\bigr),
\]

where \texttt{Gal} denotes the Galois group of the splitting field.

We use the rational-coefficient form of Müller's classification: the
geometric monodromy of an indecomposable polynomial in \(\mathbf Q[X]\) is
alternating or symmetric, cyclic, dihedral, or belongs to one of the three
explicit rational exceptional cases above. Linear equivalence is taken over
\(\mathbf C\), equivalently here over \(\overline{\mathbf Q}\). Theorems 5
and Sections 3--4 remove every case except alternating and symmetric.

\subsubsection{Theorem 6 (uniform alternating/symmetric reduction)}

For all \(q\geq1\) and \(m\geq3\),

\[
 G^{\mathrm{geom}}_{q,m}\in\{A_m,S_m\}.
\]

In fact \(G^{\mathrm{geom}}_{q,3}=S_3\) by the degree-3 observation in
Section 3. For \(m=2\), both geometric and arithmetic groups are \(S_2\).

The discriminant is

\[
 D_{q,m}(T):=\operatorname {Disc}_X(P_{q,m}(X)-T),
\]

and, because \(P_{q,m}\) is monic of degree \(m\),

\[
 \deg_TD_{q,m}=m-1,
\]

\[
 \operatorname {LC}_T D_{q,m}
 =(-1)^{m(m-1)/2+m-1}m^m.
\]

\subsubsection{Theorem 7 (precise unconditional conclusions)}

For every \(q\geq1\):

\begin{enumerate}
\def\labelenumi{\arabic{enumi}.}
\tightlist
\item
  If \(m\) is even, then
  \[
  G^{\mathrm{geom}}_{q,m}
  =G^{\mathrm{arith}}_{q,m}=S_m.
  \]
\item
  If \(m\) is odd and not a perfect square, then
  \[
  G^{\mathrm{arith}}_{q,m}=S_m,
  \]
  while for odd \(m\ge5\) geometrically the present argument leaves
  \(G^{\mathrm{geom}}_{q,m}=A_m\) or \(S_m\). Degree \(3\) is already
  \(S_3\).
\item
  If \(m\) is an odd perfect square, both the arithmetic square-class test
  and the geometric test require additional information.
\end{enumerate}

\paragraph{Proof}

For even \(m\), the inertia generator at infinity is an \(m\)-cycle and is
odd, so the geometric group cannot be \(A_m\). Equivalently,
\(D_{q,m}\) has odd degree \(m-1\) and cannot be a square. Thus both groups
are \(S_m\).

Now let \(m\) be odd. The square class of the leading coefficient is

\[
 (-1)^{(m-1)/2}m.
\]

If \(m\equiv3\pmod4\), this is negative and hence not a rational square. If
\(m\equiv1\pmod4\) but \(m\) is not a square, it is again not a rational
square. Therefore \(D_{q,m}(T)\) is not a square in \(\mathbf Q(T)\), so
the arithmetic group is not contained in \(A_m\). Since the geometric group
is \(A_m\) or \(S_m\), the arithmetic group is \(S_m\).

Finally, every odd square is \(1\pmod8\), so its displayed leading
coefficient is itself a square and \(m-1\) is even. The leading-term test
then gives no conclusion. \(\square\)

\begin{center}\rule{0.5\linewidth}{0.5pt}\end{center}

\subsection{6. The remaining boundary}

The exact unresolved boundary of the uniform proof is therefore:

\begin{itemize}
\tightlist
\item
  \textbf{arithmetic:} odd perfect-square degrees \(m\ge9\);
\item
  \textbf{geometric:} odd degrees \(m\ge5\), where one must distinguish \(A_m\)
  from \(S_m\).
\end{itemize}

These are different questions. A nonsquare rational constant multiplying a
square polynomial can make the arithmetic group \(S_m\) while the geometric
group remains \(A_m\). To prove geometric \(S_m\), one must show that
\(D_{q,m}(T)\) is not of the form

\[
 cR(T)^2,
 \qquad c\in\mathbf Q^\times, R\in\mathbf Q[T].
\]

For the original order \(q=1\), exact finite-field certificates currently
prove

\[
 G^{\mathrm{geom}}_{1,m}=G^{\mathrm{arith}}_{1,m}=S_m
 \qquad(2\leq m\leq401).
\]

For odd \(57\leq m\leq401\), the certificate uses two regular
specializations modulo \(409\) at which the discriminant has opposite
quadratic characters; this rules out \(cR(T)^2\). The reproducible code and
the lower-order exact-Morse certificates are in

\begin{itemize}
\tightlist
\item
  \texttt{computation/trace\_galois\_401.py},
\item
  \texttt{computation/trace\_morse\_certificates.py},
\item
  \texttt{computation/TRACE\_MORSE\_EXACT\_CERTIFICATES.md}.
\end{itemize}

This finite verification does not replace an all-\(m\) proof and should be
reported as a certified range.

\begin{center}\rule{0.5\linewidth}{0.5pt}\end{center}

\subsection{Reference}

P. Müller, ``Primitive monodromy groups of polynomials,'' in \emph{Recent
Developments in the Inverse Galois Problem}, Contemporary Mathematics \textbf{186}
(1995), 385--401.
\section{Prime degree: a complete Morse theorem}

Write
\[
A_q(z)=\frac{E_q(z)}z
=\sum_{n\ge0}a_nz^n,
\qquad
a_n=\frac{(-1)^n}{(n+1)!(n+1)^q},
\]
and
\[
P'_{q,m}(X)=\sum_{j=0}^{m-1}d_jX^j.
\]
Direct coefficient extraction gives
\begin{equation}
\boxed{d_j=m[z^{m-j-1}]A_q(z)^{j+1}.}
\label{eq:dj}
\end{equation}

\begin{theorem}[Prime-degree Morse theorem; \UPASS]\label{thm:prime}
Let $m\ge3$ be prime. For every $q\ge1$, $P_{q,m}$ is Morse and
\begin{equation}
\boxed{G^{\mathrm{geom}}_{q,m}=G^{\mathrm{arith}}_{q,m}=S_m.}
\label{eq:prime-mon}
\end{equation}
\end{theorem}

\begin{proof}
Work over $\Q_m$. Formula~\eqref{eq:dj} gives
\[
v_m(d_0)=-q,
\qquad
v_m(d_j)\ge1\quad(1\le j\le m-2),
\qquad
v_m(d_{m-1})=1.
\]
Thus the Newton polygon of $P'_{q,m}$ is the single edge
\[
(0,-q)\longrightarrow(m-1,1).
\]
Every critical point $\alpha_i$ has valuation
\[
v_m(\alpha_i)=-\frac{q+1}{m-1}.
\]
After a ramified base extension the edge residual is a binomial
$u+vY^{m-1}$. Its derivative is nonzero because $m\nmid m-1$, so its
$m-1$ roots are simple and have distinct leading residues $y_i$.

Let $c_k=[X^k]P_{q,m}$. Then $c_1=d_0$, $v_m(c_k)\ge1$ for
$2\le k\le m-1$, and $c_m=1$. At every critical point,
$\alpha_i^m$ is the unique term of smallest valuation in
$P_{q,m}(\alpha_i)$. Indeed,
\[
v_m(c_1\alpha_i)-v_m(\alpha_i^m)=1,
\]
and, for $2\le k\le m-1$,
\[
v_m(c_k\alpha_i^k)-v_m(\alpha_i^m)
\ge1+\frac{(m-k)(q+1)}{m-1}>0.
\]
The leading critical-value residues are therefore $y_i^m$. Since all
$y_i^{m-1}$ equal the same nonzero constant, $y_i^m=Cy_i$, and these
residues are pairwise distinct. All critical points are simple and all
critical values are distinct, so $P_{q,m}$ is Morse. Its finite branch
cycles are transpositions and generate $S_m$. The arithmetic group contains
the geometric one.
\end{proof}

This is an infinite, computation-free theorem: it settles every prime
degree simultaneously for every polyexponential order.

\section{A moving-prime isolated transposition}

Let $m\ge7$ be odd and choose a prime
\begin{equation}
\frac{m+1}{2}<p<m,
\qquad
r=m-p.
\label{eq:window}
\end{equation}
Then $r$ is even and $2\le r\le p-2$. The relevant valuations of the
coefficients of $A_q$ are
\[
v_p(a_n)=
\begin{cases}
0,&0\le n\le p-2,\\
-(q+1),&n=p-1,\\
-1,&p\le n\le m-1.
\end{cases}
\]
Because $m\le2p-2$, at most one exceptional factor can occur in any
composition contributing to~\eqref{eq:dj}.  The endpoints $j=1,r$ attain
valuation $-q-1$, while every intermediate point lies on or above the
horizontal segment.  Hence the lower Newton polygon has the three edges
\begin{equation}
(0,-1)\longrightarrow(1,-q-1)
\longrightarrow(r,-q-1)\longrightarrow(m-1,0).
\label{eq:threeedge}
\end{equation}
The critical points split into one point of valuation $q$, $r-1$ unit
points, and $p-1$ points of valuation $-(q+1)/(p-1)$. The first point
$\alpha$ is simple, and its critical value $b=P_{q,m}(\alpha)$ satisfies
\begin{equation}
\boxed{v_p(b)=q-1.}
\label{eq:bval}
\end{equation}

For completeness, the exact valuation requires excluding cancellation
between the first two terms.  Write
$P_{q,m}(X)=\sum_{k\ge1}c_kX^k$ and
$P'_{q,m}(X)=\sum_{j\ge0}d_jX^j$, so that $c_1=d_0$ and
$c_2=d_1/2$.  The first Newton edge gives
\[
 d_1\alpha=-d_0(1+\varepsilon),\qquad v_p(\varepsilon)\ge q.
\]
Therefore
\[
 c_1\alpha+c_2\alpha^2
 =d_0\alpha+\frac{d_1}{2}\alpha^2
 =\frac{d_0\alpha}{2}(1-\varepsilon),
\]
which has valuation exactly $q-1$, since $p$ is odd.  For $k\ge3$,
\[
 v_p(c_k\alpha^k)\ge(k-1)q-1\ge2q-1>q-1,
\]
so no later term can cancel the leading contribution.

Define the residual lower-row derivative
\[
R_{q,r}(X)=P'_{q,r}(X)\bmod p.
\]
Keeping the horizontal-edge principal terms gives, up to a nonzero scalar,
\begin{equation}
p^{q+1}P_{q,m}(X)\equiv\frac{X^2}{r}R_{q,r}(X)\pmod p,
\label{eq:lowerP}
\end{equation}
and differentiation gives
\begin{equation}
p^{q+1}P'_{q,m}(X)
\equiv\frac{X}{r}\bigl(2R_{q,r}(X)+XR'_{q,r}(X)\bigr)\pmod p.
\label{eq:lowerPd}
\end{equation}
The only potentially dangerous denominator when recovering $P$ from $P'$ is
removed by Frobenius:
\begin{equation}
d_{p-1}=m[z^r]A_q(z)^p
\equiv m[z^r]A_q(z^p)=0\pmod p
\qquad(0<r<p).
\label{eq:frob}
\end{equation}

\begin{theorem}[Window-prime criterion; \UPASS]\label{thm:window}
If
\begin{equation}
\gcd(P'_{q,r},P''_{q,r})=1
\quad\text{in }\F_p[X],
\label{eq:gcd}
\end{equation}
then
\begin{equation}
\boxed{G^{\mathrm{geom}}_{q,m}=G^{\mathrm{arith}}_{q,m}=S_m.}
\label{eq:window-mon}
\end{equation}
\end{theorem}

\begin{proof}
If a unit critical point $\beta$ had a critical value of elevated
valuation, \eqref{eq:lowerP}--\eqref{eq:lowerPd} would force
\[
R_{q,r}(\bar\beta)=R'_{q,r}(\bar\beta)=0,
\]
contradicting~\eqref{eq:gcd}. Thus all unit critical values have valuation
$-(q+1)$.

On the final edge the principal derivative is
$d_rX^r+mX^{m-1}$. At any of its critical points the principal critical
value is
\[
d_rX^{r+1}
\left(\frac1{r+1}-\frac1m\right)
=d_rX^{r+1}\frac{p-1}{m(r+1)},
\]
so its valuation is
\[
-\frac{m(q+1)}{p-1}<0.
\]
Hence $b$ is the unique critical value of valuation $q-1$, and it is
supported by the single simple point $\alpha$. Its inertia is a
transposition. The pre-existing indecomposability theorem gives primitivity,
and a primitive group containing a transposition is $S_m$. The arithmetic
conclusion follows.
\end{proof}

The old degree-$m$ discriminant problem has therefore descended to a gcd for
a derivative of degree $r-1=m-p-1$.

\begin{corollary}[An all-order infinite family; \UPASS]\label{cor:shifted}
For every prime $p\ge5$, every $q\ge1$, and $m=p+2$,
\begin{equation}
\boxed{G^{\mathrm{geom}}_{q,p+2}=G^{\mathrm{arith}}_{q,p+2}=S_{p+2}.}
\label{eq:shifted}
\end{equation}
\end{corollary}

Indeed $r=2$, so $P'_{q,2}\bmod p$ is linear and square-free. In
particular, the former first odd geometric gap $m=403=401+2$ and the former
first odd-square arithmetic gap $m=441=439+2$ are closed for every $q$,
without a full discriminant calculation.

\section{Exact finite rectangles}

For $p>r$, the reduction of $P_{q,r}$ depends only on $q\bmod(p-1)$. The
criterion can therefore certify infinitely many orders by a finite
residue-class computation.

\begin{theorem}[All orders through degree 1001; \XPASS]\label{thm:1001}
For every $q\ge1$ and $2\le m\le1001$,
\begin{equation}
\boxed{G^{\mathrm{geom}}_{q,m}=G^{\mathrm{arith}}_{q,m}=S_m.}
\label{eq:1001}
\end{equation}
\end{theorem}

The exact audit finds one window prime working for every residue of $q$ in
492 odd degrees. Prime-degree \cref{thm:prime} handles
$m=3,5,11,17,167,577$. Only two composite degrees require combined
certificates:

The finite search policy is itself part of the certificate.  The audit is
run with \texttt{--bound 1001 --candidate-count 10} and examines at most the
ten largest eligible strict-window primes in descending order.  The displayed
prime-degree fallback list is therefore tied to this fixed cutoff and search
order, rather than being an intrinsic classification of the degrees.
\begin{itemize}
\item At $m=27$, the primes $17,19,23$ leave the single progression
$q\equiv1075\pmod{1584}$. Modulo $28$, this progression occupies the seven
classes $3,7,11,15,19,23,27$. For each class, reduction modulo $29$ gives
two regular specializations of
$\operatorname{Disc}_X(P_{q,27}(X)-T)$ with opposite quadratic character:
\end{itemize}

\begin{center}
\small
\begin{tabular}{@{}rcc@{}}
\toprule
$q\bmod 28$ & square $(T,D(T))$ & nonsquare $(T,D(T))$\\
\midrule
3  & $(3,7)$  & $(0,17)$\\
7  & $(0,16)$ & $(4,11)$\\
11 & $(0,22)$ & $(5,27)$\\
15 & $(2,16)$ & $(1,12)$\\
19 & $(1,16)$ & $(0,18)$\\
23 & $(0,22)$ & $(1,15)$\\
27 & $(0,22)$ & $(1,3)$\\
\bottomrule
\end{tabular}
\end{center}

Hence none of the seven discriminant polynomials is a constant times a
square, supplying the required odd geometric branch packet.

\begin{itemize}
\item At $m=989$, the bad order residues for $p=919$ are odd, while those
for $p=929$ are even. Their simultaneous CRT locus is empty.
\end{itemize}

The script and captured stdout are:
\begin{itemize}
\item \nolinkurl{trace_allq_modp_audit.py}\\[-.15em]
\textcolor{bluegray}{\small SHA-256: }
\texttt{\small\seqsplit{36ba3fd004df1e5a56a7bb6c61ed49cf7542d6f7c934ba458f88c1ebe105c31c}};
\item \nolinkurl{trace_allq_modp_audit_1001.txt}\\[-.15em]
\textcolor{bluegray}{\small SHA-256: }
\texttt{\small\seqsplit{1d4867aa086b4315d45d70c52ef27090ad4d1c256cd956cb6b18d72a8be374be}}.
\end{itemize}

\begin{notebox}
This theorem asserts full symmetric monodromy. It does \emph{not} assert
that every polynomial in this rectangle is Morse.
\end{notebox}

\begin{theorem}[Original tower through degree one million; \XPASS]\label{thm:million}
For every $2\le m\le10^6$,
\begin{equation}
\boxed{G^{\mathrm{geom}}_{1,m}=G^{\mathrm{arith}}_{1,m}=S_m.}
\label{eq:million}
\end{equation}
\end{theorem}

For each odd $7\le m\le10^6$, an exact finite-field search tests the largest
strict-window prime and, only when necessary, the second largest. The largest
residual degree is $114$. The largest-prime choice fails only at
$(m,p,r)=(215,211,4)$ and $(69649,69623,26)$; the second prime works in both
cases. Even degrees are structural, and prime degrees also follow from
\cref{thm:prime} independently of the audit.

The script and captured stdout are:
\begin{itemize}
\item \nolinkurl{window_prime_selection_audit.py}\\[-.15em]
\textcolor{bluegray}{\small SHA-256: }
\texttt{\small\seqsplit{e2eb3b8a14f758d84c0d80a4ed3c0ed11a3cf7eb38e3935960be3122d8fdcbf0}};
\item \nolinkurl{window_prime_selection_audit_1000000.txt}\\[-.15em]
\textcolor{bluegray}{\small SHA-256: }
\texttt{\small\seqsplit{51d3a027e9f0ce42f4918b39058288cdfdf808e25ded994ed009a701bc7c44f5}}.
\end{itemize}

Again, this is a monodromy theorem, not an all-degree Morse theorem.

\section{Large order: complex regular polygons replace real critical points}

Put
\[
G_q(z)=\sum_{n\ge1}\frac{z^n}{n!n^q},
\qquad
h_{q;m,k}=[z^m]G_q(z)^k>0.
\]
Then
\begin{equation}
P_{q,m}(X)
=\sum_{k=1}^m(-1)^{m-k}\frac{m}{k}h_{q;m,k}X^k.
\label{eq:P-positive}
\end{equation}
If $n_1+\cdots+n_k=m$ with every $n_i\ge1$, the product of the parts
satisfies
\begin{equation}
n_1\cdots n_k\ge m-k+1.
\label{eq:parts}
\end{equation}
For $2\le r\le m-1$, define
\begin{equation}
\lambda_{m,r}=\frac{m^{(r-1)/(m-1)}}r,
\qquad
\theta_m=\max_{2\le r\le m-1}\lambda_{m,r}<1.
\label{eq:theta}
\end{equation}
The strict inequality follows from the strict decrease of
$x\mapsto\log x/(x-1)$ on $(1,\infty)$.

Let
\[
s_q=m^{-q/(m-1)},
\qquad
\widetilde P_{q,m}(y)
=m^{qm/(m-1)}P_{q,m}(s_qy).
\]
The endpoint terms $k=1,m$ remain, while
\eqref{eq:parts}--\eqref{eq:theta} suppress every intermediate term
exponentially.

\begin{theorem}[Effective large-order Morse theorem; \UPASS]\label{thm:largeq}
For every fixed $m\ge2$,
\begin{equation}
\boxed{
\widetilde P_{q,m}(y)
\longrightarrow
S_m(y):=y^m+\frac{(-1)^{m-1}}{(m-1)!}y
\quad(q\to\infty)}
\label{eq:limit}
\end{equation}
coefficientwise and locally uniformly, with error
$O_{m,R}(\theta_m^q)$ on $|y|\le R$. Consequently there is an effectively
computable $Q(m)$ such that, for every $q\ge Q(m)$,
\begin{equation}
\boxed{
P_{q,m}\text{ is Morse},
\qquad
G^{\mathrm{geom}}_{q,m}=G^{\mathrm{arith}}_{q,m}=S_m.}
\label{eq:largeq-mon}
\end{equation}
\end{theorem}

\begin{proof}
The derivative of the limit is
\[
S_m'(y)=my^{m-1}+\frac{(-1)^{m-1}}{(m-1)!}.
\]
Its roots satisfy
\[
y^{m-1}=\frac{(-1)^m}{m!},
\]
so they are the distinct vertices of a regular $(m-1)$-gon. At a critical
point $y_j$,
\[
S_m(y_j)=(-1)^{m-1}\frac{m-1}{m!}y_j,
\]
and these critical values are another set of pairwise distinct regular
polygon vertices. Thus $S_m$ is Morse. The Morse locus is Zariski open in
coefficient space; alternatively, disjoint Rouch\'e circles around these
explicit roots and their separated values give an effective bound. The
exponential coefficient estimate therefore proves~\eqref{eq:largeq-mon}.
\end{proof}

For fixed odd $m$, the limiting critical equation has no real root. Hence
sufficiently large $q$ has no real critical point at all. The smallest exact
example is $(q,m)=(2,3)$. A uniform real-transposition strategy is therefore
false; complex critical-value separation is the correct replacement.

\section{Local Faber transfer beyond the polyexponential tower}

The moving-prime mechanism is a finite-jet theorem, not a peculiarity of
the differential equation for $E_q$. Let $L$ be a complete discretely
valued characteristic-zero field with residue characteristic $p\ge5$,
uniformizer $\pi$, and residue field $k$. For
\[
f(z)=zA(z),\qquad A(z)=\sum_{n\ge0}a_nz^n,\qquad a_0=1,
\]
put $P_m^f=-m[z^m]\log(1-Xf(z))$. Suppose $m=p+r$,
$2\le r\le p-2$, and, for integers $h>s\ge0$,
\begin{equation}
\begin{array}{ll}
v(a_n)\ge0 &(0\le n\le p-2),\\
v(a_{r-1})=0,\quad v(a_{p-1})=-h,\\
v(a_n)\ge-s &(p\le n\le m-2),\\
v(a_{m-1})=-s.
\end{array}
\label{eq:spike}
\end{equation}
Let $\bar A_{<p}$ be the reduction of the integral truncation below
degree $p-1$, and define
\begin{equation}
R_r(X)=\overline{(P_r^f)'(X)}
=r\sum_{j=1}^r[z^{r-j}]\bar A_{<p}(z)^jX^{j-1}.
\label{eq:R-general}
\end{equation}

\begin{theorem}[Single-spike Faber transfer; \UPASS]\label{thm:transfer}
Under \eqref{eq:spike}, the derivative Newton polygon is exactly
\begin{equation}
(0,-s)\longrightarrow(1,-h)\longrightarrow(r,-h)
\longrightarrow(m-1,0).
\label{eq:spike-NP}
\end{equation}
If $\gcd(R_r,R_r')=1$ in $k[X]$, then $P_m^f$ has a critical value
supported by exactly one simple critical point. The three critical-point
clusters and their critical-value valuations are
\begin{equation}
\begin{array}{c|c|c}
\text{number}&v(\text{critical point})&v(\text{critical value})\\ \hline
1&h-s&h-2s\\
r-1&0&-h\\
p-1&-h/(p-1)&-mh/(p-1).
\end{array}
\label{eq:cluster-table}
\end{equation}
Thus its geometric inertia contains a transposition. If the coefficients
lie in a number field $K$, the local hypotheses hold in a completion $K_v$
or a valued extension, and $P_m^f$ is indecomposable over $\overline K$,
then both its geometric and arithmetic monodromy groups are $S_m$.
\end{theorem}

\begin{proof}
The universal identities
\begin{equation}
[X^k]P_m^f=\frac mk[z^{m-k}]A^k,
\qquad
[X^j](P_m^f)'=m[z^{m-j-1}]A^{j+1}
\label{eq:faber-universal}
\end{equation}
show that at most one copy of the unique deep coefficient $a_{p-1}$ can
enter. Exact endpoint terms give \eqref{eq:spike-NP}. If
$\eta=\overline{\pi^ha_{p-1}}$, the horizontal initial forms are
\begin{equation}
\operatorname{in}_{-h}(P_m^f)
=\frac{\bar m\eta}{r}X^2R_r(X),
\qquad
\operatorname{in}_{-h}((P_m^f)')
=\frac{\bar m\eta}{r}X(2R_r+XR_r').
\label{eq:initial-forms}
\end{equation}
The apparent coefficient $[X^p]P_m^f=d_{p-1}/p$ is integral because
$[z^r]A^p\in p\mathcal O_L$ for $0<r<p$. Consequently the gcd condition
prevents cancellation of every unit-cluster critical value. The length-one
edge gives one simple point $\alpha$, and its derivative equation yields
$v(P_m^f(\alpha))=h-2s$. The other two values in
\eqref{eq:cluster-table} follow from \eqref{eq:initial-forms} and the
final-edge binomial. Since
\[
-\frac{mh}{p-1}<-h<h-2s,
\]
the value at $\alpha$ is globally isolated. Primitive polynomial monodromy
containing a transposition is $S_m$.
\end{proof}

For $r=2$, $R_2$ is linear, so the gcd condition is automatic. This
immediately transfers the local branch construction, for integers
$\tau\ge1$ and $\sigma,\tau\ge1$, respectively, to
\begin{equation}
\operatorname{Li}_\tau(z)=\sum_{N\ge1}\frac{z^N}{N^\tau}
\quad\text{and}\quad
f_{\sigma,\tau}(z)=\sum_{N\ge1}\frac{z^N}{(N!)^\sigma N^\tau},
\label{eq:transfer-families}
\end{equation}
with depths $(h,s)=(\tau,0)$ and $(\sigma+\tau,\sigma)$ over $\Q_p$,
respectively. Both families therefore acquire a transposition.  The next
subsection supplies the missing indecomposability for every pure
polylogarithmic member; for $f_{\sigma,\tau}$ it remains a separate gate.

\subsection{Global polylogarithmic closure}

For an integer $\tau\ge1$, write
\begin{equation}
L_{\tau,m}(X):=P_m^{\operatorname{Li}_\tau}(X)
=-m[z^m]\log\!\left(1-X\operatorname{Li}_\tau(z)\right),
\qquad
A_\tau(z):=\frac{\operatorname{Li}_\tau(z)}z.
\label{eq:polylog-faber}
\end{equation}
The missing global input for this family is supplied by a strict form of
Kaluza's reciprocal-sign mechanism.

\begin{lemma}[Strict Kaluza sign for polylogarithms; \UPASS]
\label{lem:strict-kaluza}
If
\[
A_\tau(z)=\sum_{n\ge0}a_nz^n,
\qquad a_n=(n+1)^{-\tau},
\qquad
\frac1{A_\tau(z)}=\sum_{n\ge0}b_nz^n,
\]
then $b_0=1$ and $b_n<0$ for every $n\ge1$.
\end{lemma}

\begin{proof}
The positive sequence $(a_n)$ is strictly log-convex, since
\[
\rho_n:=\frac{a_n}{a_{n-1}}
=\left(\frac n{n+1}\right)^\tau
\]
is strictly increasing. Kaluza's sign criterion \cite{Kaluza1928} gives
$b_n\le0$ for $n\ge1$. Strictness can also be read directly from its
recurrence. Put
\[
\frac1{A_\tau(z)}=1-\sum_{n\ge1}c_nz^n.
\]
The reciprocal identity gives
\[
a_n=\sum_{k=1}^n c_ka_{n-k}\qquad(n\ge1).
\]
Thus $c_1=a_1>0$, and comparison with the preceding identity multiplied by
$\rho_n$ yields, for $n\ge2$,
\begin{equation}
c_n=\sum_{k=1}^{n-1}c_ka_{n-1-k}
\bigl(\rho_n-\rho_{n-k}\bigr)>0.
\label{eq:strict-kaluza-recurrence}
\end{equation}
Induction proves $b_n=-c_n<0$ in every positive degree.
\end{proof}

\begin{theorem}[Global polylogarithmic Faber theorem; \UPASS]
\label{thm:polylog-global}
For every integer $\tau\ge1$ the following hold.
\begin{enumerate}
\item $L_{\tau,m}$ is absolutely indecomposable for every $m\ge2$.
\item If $\ell\ge3$ is prime, then $L_{\tau,\ell}$ is Morse and
\[
\Gal(L_{\tau,\ell}(X)-T/\overline{\Q}(T))
=\Gal(L_{\tau,\ell}(X)-T/\Q(T))=S_\ell.
\]
\item If $p\ge5$ is prime, then
\[
\Gal(L_{\tau,p+2}(X)-T/\overline{\Q}(T))
=\Gal(L_{\tau,p+2}(X)-T/\Q(T))=S_{p+2}.
\]
\item In the degrees covered by (2) or (3), the off-diagonal collision
polynomial
\[
\frac{L_{\tau,m}(X)-L_{\tau,m}(Y)}{X-Y}
\]
is absolutely irreducible over $\Q$.
\end{enumerate}
\end{theorem}

\begin{proof}
Since
\[
\frac1{\operatorname{Li}_\tau(z)}
=z^{-1}\frac1{A_\tau(z)},
\]
\Cref{lem:strict-kaluza} gives
\[
[z^b]\frac1{\operatorname{Li}_\tau(z)}=b_{b+1}<0
\qquad(b\ge0).
\]
The reciprocal-coefficient obstruction in
Section~\ref{sec:reciprocal-coefficient-obstruction} says that a decomposition
with inner degree $b>1$ would force this coefficient to vanish. This proves
(1).

For (2), work over $\Q_\ell$ and write
\[
L'_{\tau,\ell}(X)=\sum_{j=0}^{\ell-1}u_jX^j,
\qquad
u_j=\ell[z^{\ell-j-1}]A_\tau(z)^{j+1}.
\]
Then
\[
v_\ell(u_0)=1-\tau,
\qquad v_\ell(u_j)\ge1\ (1\le j\le\ell-2),
\qquad v_\ell(u_{\ell-1})=1.
\]
Indeed, the intermediate coefficient extractions involve only denominators
$1,\ldots,\ell-1$, which are $\ell$-adic units. Hence the derivative Newton
polygon is the single edge
\[
(0,1-\tau)\longrightarrow(\ell-1,1).
\]
Its residual binomial has $\ell-1$ distinct roots. The critical points
$\alpha_i$ therefore have valuation $-\tau/(\ell-1)$ and distinct leading
residues $y_i$. If $c_k=[X^k]L_{\tau,\ell}$, then
\[
c_1=u_0,\qquad c_k=u_{k-1}/k\ (2\le k\le\ell-1),
\qquad c_\ell=1.
\]
Consequently
\[
v_\ell(c_1\alpha_i)-v_\ell(\alpha_i^\ell)=1,
\]
and, for $2\le k\le\ell-1$,
\[
v_\ell(c_k\alpha_i^k)-v_\ell(\alpha_i^\ell)
\ge1+\frac{(\ell-k)\tau}{\ell-1}>0.
\]
Thus $\alpha_i^\ell$ is the unique leading term of the critical value.
Because the edge equation gives $y_i^{\ell-1}=C\ne0$, the leading critical
value residues $y_i^\ell=Cy_i$ are pairwise distinct. Hence
$L_{\tau,\ell}$ is Morse, its finite branch cycles are transpositions, and
its geometric and arithmetic monodromy groups are $S_\ell$.

For (3), apply \Cref{thm:transfer} with $r=2$ and $(h,s)=(\tau,0)$.
The residual polynomial $R_2$ is linear, and (1) supplies absolute
indecomposability, so both monodromy groups are $S_{p+2}$. Finally, an
$S_m$ polynomial cover is two-transitive. The point stabilizer is transitive
on the other sheets, which is equivalent to absolute irreducibility of the
displayed off-diagonal fiber product.
\end{proof}

\begin{corollary}[Fixed-order polylogarithmic window; \XPASS]
\label{cor:polylog-window}
For every $\tau\in\{1,2,3,4\}$ and every odd $7\le m\le1001$,
\begin{equation}
\boxed{
\Gal(L_{\tau,m}(X)-T/\overline{\Q}(T))
=\Gal(L_{\tau,m}(X)-T/\Q(T))=S_m.}
\label{eq:polylog-window}
\end{equation}
\end{corollary}

\begin{proof}
The exact finite-field audit described in
\Cref{app:computational-certificates} finds a strict-window prime satisfying
$\gcd(R_r,R_r')=1$ for every displayed row except
$(\tau,m)=(4,17)$.  In each certified row, \Cref{thm:transfer} supplies a
transposition and \Cref{thm:polylog-global}(1) supplies primitivity.  The sole
window exception has prime degree and is therefore covered unconditionally
by \Cref{thm:polylog-global}(2).  Thus it is a window-method exception, not a
monodromy exception.
\end{proof}

The case $\tau=1$ is $\operatorname{Li}_1(z)=-\log(1-z)$ and has additional
enumerative structure. The independent identity
\[
\frac z{-\log(1-z)}=\int_0^1(1-z)^t\,dt
\]
also shows directly that every positive-degree reciprocal coefficient is
nonzero. For
$f(z)=-\log(1-z)$,
\begin{equation}
P_m^f(X)=\sum_{k=1}^m
\frac{(k-1)!\,|s(m,k)|}{(m-1)!}X^k.
\label{eq:stirling}
\end{equation}
Thus all the polynomials in \eqref{eq:stirling} are absolutely
indecomposable, and \Cref{thm:polylog-global} gives
\begin{equation}
\boxed{
\Gal(P_{p+2}^f(X)-T/\overline{\Q}(T))
=\Gal(P_{p+2}^f(X)-T/\Q(T))=S_{p+2}
\quad(p\ge5\text{ prime})}
\label{eq:stirling-mon}
\end{equation}
for the logarithmic--Stirling trace tower. No literature-priority claim is
made without a dedicated search.

\section{Shifted discriminants, partial separability, and conditional density one}

For the original order $q=1$, put
\begin{equation}
F_r(X)=(r!)^2P'_{1,r}(X)\in\mathbf Z[X],
\qquad
\Delta_r=\operatorname{Disc}(F_r).
\label{eq:Delta}
\end{equation}
The integrality follows termwise from $\prod n_i!\mid r!$ and
$\prod n_i\mid r!$ for every composition $n_1+\cdots+n_k=r$. If $p>r$,
the leading coefficient $r(r!)^2$ is a unit modulo $p$, so
\begin{equation}
\boxed{
P'_{1,r}\bmod p\text{ is not square-free}
\quad\Longleftrightarrow\quad p\mid\Delta_r.}
\label{eq:Delta-equivalence}
\end{equation}
Coefficient majorization by $(e^z-1)^k$, followed by Mahler's
discriminant bound, gives whenever $\Delta_r\ne0$
\begin{equation}
\log|\Delta_r|=O(r^2\log r),
\qquad
\omega(\Delta_r)=O(r^2\log r).
\label{eq:Delta-bound}
\end{equation}

\begin{notebox}
\textbf{Hypothesis DS (\OPEN).} Every even residual derivative
$P'_{1,r}$, $r\ge2$, is square-free over $\Q$, equivalently
$\Delta_r\ne0$.
\end{notebox}

\begin{theorem}[Unconditional partial separability; \UPASS]\label{thm:partial-sep}
Let $r=2s$. For every $q\ge1$, the derivative $P'_{q,r}$ has at least
$s$ simple roots. After the dyadic normalization
\[
c_q=2^{q+1},\qquad
\widehat P_{q,r}(X)=c_q^rP_{q,r}(X/c_q),
\]
these are exactly the unit roots, while every possible repeated root lies
in the strict positive-valuation disk. Moreover, for every odd prime $p$
and every $q\ge1$,
\begin{equation}
\boxed{P'_{q,p+1}(X)\text{ is square-free over }\Q.}
\label{eq:prime-successor}
\end{equation}
\end{theorem}

\begin{proof}
Put $f_q(z)=E_q(c_qz)/c_q$. Exact coefficient valuations give the
order-independent reduction
\[
f_q(z)\equiv z+z^2\pmod2.
\]
If $a=v_2(r)$, coefficient extraction then yields
\begin{equation}
2^{-a}\widehat P'_{q,2s}(X)\equiv
X^{s-1}\Lambda_{2s+1}(X)\pmod2,
\label{eq:lambda-reduction}
\end{equation}
\begin{equation}
\Lambda_{2s+1}(X)=
\sum_{j=0}^s\binom{s+j}{s-j}X^j.
\label{eq:lambda}
\end{equation}
The Lucas recurrence $L_0=0,L_1=X$,
$L_n=X(L_{n-1}+L_{n-2})$ gives
$L_{2s+1}=X^{s+1}\Lambda_{2s+1}$. Its nonzero roots are
$X=\zeta+\zeta^{-1}$ for nontrivial $(2s+1)$-st roots of unity, modulo
$\zeta\leftrightarrow\zeta^{-1}$. Hence $\Lambda_{2s+1}$ has $s$
distinct nonzero roots and simple Hensel lifts, proving the first assertion.

For $m=p+1$, the $p$-adic Newton polygon of $P'_{q,m}$ is
\begin{equation}
(0,-1)\longrightarrow(1,-q-1)\longrightarrow(p,0).
\label{eq:prime-successor-NP}
\end{equation}
The first edge is linear; after adjoining a uniformizer for the possibly
fractional second slope, its nonzero-root initial binomial has exponent gap
$p-1$, prime to $p$. Thus all roots are simple.
\end{proof}

\begin{corollary}[Exact residual rectangle; \XPASS]\label{cor:residual-5000}
For the original order, exact finite-field computation proves
\begin{equation}
\boxed{P'_{1,r}\text{ is square-free for every even }2\le r\le5000.}
\label{eq:residual-5000}
\end{equation}
\end{corollary}

At the single good prime $998244353$, exact convolution and polynomial gcd
arithmetic certify \eqref{eq:residual-5000}. Leading coefficients stay
nonzero modulo this prime; for odd $3\le s\le4999$, those of the polar and
its derivative are $s+1$ and $(s-1)(s+1)$, respectively. Thus the relevant
degrees are preserved. The same run, together with the trivial constant
case $s=1$, also certifies every odd polar $2H_s+XH_s'$ for
$1\le s\le4999$, where
$H_s=\sum_k[z^s](-E_1(-z))^kX^{k-1}$.

The exact audit has SHA-256 hashes
\begin{itemize}
\item \nolinkurl{even_derivative_separability_audit.cpp}\\[-.15em]
\textcolor{bluegray}{\small SHA-256: }
\texttt{\small\seqsplit{d9c6f7c6438a9e7bf5d0aed50bbce727f460efcacc720f520c97b444b3f83b4d}};
\item \nolinkurl{even_derivative_separability_audit_5000.txt}\\[-.15em]
\textcolor{bluegray}{\small SHA-256: }
\texttt{\small\seqsplit{1e4b4073527d7d2543c8a087c37f0888fdbd99c7ca2b7880a418e7c533e6617b}}.
\end{itemize}
This does not prove DS: in row $r=10$, a later inner Newton face already
has inseparable residual polynomial $1+Y^2$ over $\F_2$.

\begin{theorem}[Conditional density-one selection; \CPASS]\label{thm:density-one}
Assume DS. Let $\mathcal E(X)$ be the odd $m\le X$ for which no prime
$(m+1)/2<p<m$ passes the residual gcd criterion \eqref{eq:gcd}. Then, for
every $\varepsilon>0$,
\begin{equation}
\boxed{\#\mathcal E(X)\ll_\varepsilon X^{15/16+\varepsilon}.}
\label{eq:density-one}
\end{equation}
In particular, under DS the original trace tower has symmetric geometric
and arithmetic monodromy for a density-one set of odd degrees.
\end{theorem}

\begin{proof}
In a dyadic interval $[X,2X]$, let $p<m$ be the largest prime and write
$r=m-p$. Set $H=X^{5/16}$. If $r\le H$ and $m$ is exceptional, then
\eqref{eq:Delta-equivalence}--\eqref{eq:Delta-bound} give
\[
\#\{m:r\le H,\ p\mid\Delta_r\}
\ll\sum_{r\le H}r^2\log r
\ll H^3\log H.
\]
The remaining integers lie more than $H$ places into a prime gap. Maynard's
mean-square prime-gap theorem \cite[Theorem~3.1]{Maynard2012}
\[
\sum_{p_n\le x}(p_{n+1}-p_n)^2
\ll_\varepsilon x^{5/4+\varepsilon}
\]
implies that their number is
$O_\varepsilon(X^{5/4+\varepsilon}/H)$. Balancing the two estimates gives
\eqref{eq:density-one}; dyadic summation completes the proof.
\end{proof}

The exact equivalence \eqref{eq:Delta-equivalence} also explains why
prime-gap bounds alone do not prove selection in every degree: one must
exclude the correlated divisibilities $m-r\mid\Delta_r$ for several
offsets. Already at $m=215$, the largest window prime $211$ is bad in row
$r=4$, although the next prime works. Thus a canonical nearest-prime proof
is false even though the selectable-prime criterion survives.

\section{Exact negative results that prevent a false proof}

The failed routes are useful because they make the remaining gap precise.
\begin{enumerate}
\item A simple real critical point does not suffice: another critical point
can share its critical value. For $q=1$, every odd degree has a simple
positive critical point, but global value separation remains necessary.
\item The natural coefficient triangle is not totally positive. At $q=1$
it already has an exact negative $9\times9$ minor. A proof cannot cite global
$TP_\infty$.
\item The dyadic derivative analysis isolates a simple outer critical
cluster, but coefficient Newton polygons alone do not exclude an exact
collision with a shallower cluster.
\item The moving-prime criterion is sufficient, not automatic. For odd
composite $m$, Bertrand applied to $(m+1)/2$ supplies a prime
$(m+1)/2<p<m$; compositeness excludes the only possible upper endpoint
$p=m$. Bertrand says nothing, however, about residual square-freeness for
an eligible prime.
\end{enumerate}

\section{The remaining theorem gate}

The cleanest unresolved statement is now finite-field rather than
transcendental.

\begin{problem}[\OPEN]
Prove that for every $q\ge1$ and every sufficiently large odd composite $m$,
at least one prime $p=p(q,m)$
\[
\frac{m+1}{2}<p<m
\]
makes
\[
\gcd(P'_{q,m-p},P''_{q,m-p})=1\pmod p.
\]
\end{problem}

For the original order $q=1$, exact computation verifies this through
$m=10^6$, always using one of the two largest eligible primes. Under DS,
\cref{thm:density-one} proves it for a density-one set, but no argument
currently promotes that finite fact to all $m$. A proof would close the
remaining odd composite degree problem by \cref{thm:window}.

An alternative is to continue the polar expansion at a multiple residual
root and prove that its lifted critical value cannot meet the isolated value
of~\eqref{eq:bval}. This is the precise local collision problem left by the
present work.
\section{Logarithmic--Stirling and supernecklace applications}
\label{sec:supernecklace}

The single-spike theorem above transfers an isolated local transposition to
the polylogarithmic and factorially weighted families in
\eqref{eq:transfer-families}.  The two scopes now diverge.  The strict Kaluza
argument in \Cref{thm:polylog-global} proves absolute indecomposability for
every pure polylogarithmic member and yields global symmetric families.  For
the factorially weighted $f_{\sigma,\tau}$ family the transfer remains local:
without a separate indecomposability theorem it does not by itself imply full
symmetric monodromy.  The logarithmic member $\operatorname{Li}_1$ has the
additional enumerative refinement developed below.

Put
\[
 C_m(X):=(m-1)!P_m^{(-\log(1-z))}(X)
 =\sum_{k=1}^m (k-1)!\stirling{m}{k}X^k,
 \tag{SN-1}\label{eq:SN-1}
\]
where \(\stirling{m}{k}\) is the unsigned Stirling number of the first kind.

\begin{proposition}[Type-III supernecklace enumerator; \UPASS]
The exponential generating function of the polynomials \(C_m\) is
\[
 \boxed{\displaystyle
 \sum_{m\ge1}C_m(X)\frac{z^m}{m!}
 =\log\!\left(\frac{1}{1-X\log\frac1{1-z}}\right).}
 \tag{SN-2}
\]
Consequently \(C_m\) is the cycle-number enumerator for a labelled cycle of
labelled cycles (the type-III supernecklace construction).
\end{proposition}

\begin{proof}
Expand the outer logarithm and use the classical Stirling exponential
generating function:
\[
 \log\!\left(\frac1{1-X L(z)}\right)
 =\sum_{k\ge1}\frac{X^k}{k}L(z)^k,
 \qquad
 L(z)^k=k!\sum_{m\ge k}\stirling{m}{k}\frac{z^m}{m!},
\]
where \(L(z)=\log(1/(1-z))\).  Coefficient comparison gives
\eqref{eq:SN-1}.  The labelled product/cycle interpretation is precisely the
standard exponential-formula reading of the same identity.
\end{proof}

\begin{theorem}[Symmetric supernecklace family; \UPASS]
For every prime \(p\ge5\), with \(m=p+2\),
\[
 \Gal(C_m(X)-T/\overline{\Q}(T))
 =\Gal(C_m(X)-T/\Q(T))=S_m.
 \tag{SN-3}\label{eq:SN-3}
\]
Moreover the off-diagonal collision polynomial
\[
 \boxed{\displaystyle
 \frac{C_m(X)-C_m(Y)}{X-Y}}
 \tag{SN-4}
\]
is absolutely irreducible over \(\Q\).
\end{theorem}

\begin{proof}
Scaling a polynomial by the nonzero constant \((m-1)!\) does not change its
geometric or arithmetic monodromy.  Equation~\eqref{eq:stirling-mon} therefore
gives \eqref{eq:SN-3}.  For a polynomial cover, the irreducible factors of the
off-diagonal fiber product correspond to the orbits of a point stabilizer on
the remaining sheets.  The natural action of \(S_m\) is two-transitive, so
that stabilizer has one such orbit.  This proves absolute irreducibility.
\end{proof}

\begin{remark}[Boundary of the application]
Qualitative Hilbert irreducibility supplies many specializations with group
\(S_m\), but no quantitative exceptional-set exponent is asserted here.
Nor is a genus formula claimed for the collision curve: symmetric
monodromy gives absolute irreducibility, while smoothness and the precise
ramification divisor require a separate Morse analysis.
\end{remark}

\part{Maximal Inverse Monodromy and Global Trace Collapse}
\section{All-order inverse monodromy for the polyexponential tower}

Uniform incomplete-gamma remainders, contraction boxes, and the covering-space
argument are included so that the proof is logically self-contained apart from
the quoted standard large-variable remainder theorem.  In particular,
\textbf{global pairwise noncollision of all critical values is neither assumed
nor claimed}.

Put

\[
 F_q(z)=\operatorname{Ein}_q(z)
 =\sum_{n\geq1}\frac{(-1)^{n-1}z^n}{n!n^q},\qquad q\geq1,
\]

and put \(F_0(z)=1-e^{-z}\). Then

\[
 zF_q'(z)=F_{q-1}(z).
\tag{Inv-0.1}
\]

The principal result is

\begin{quote}
\textbf{Theorem A (all-order maximal inverse monodromy).} For every fixed
integer \(q\geq1\) and every regular value \(c_0\) (that is, \(c_0\) is
not a critical value),
\[
 \operatorname{Mon}_{c_0}(F_q^{-1})
 =\operatorname{FSym}(F_q^{-1}(c_0))
 \cong\operatorname{FSym}(\mathbb N).
\tag{Inv-0.2}
\]
\end{quote}

For \(q=1\), all critical points are simple and all critical values are
pairwise distinct. For \(q\geq2\), those two global statements are not
needed: the proof shows that degeneracies and critical-value collisions, if
any, form only a finite exceptional set. A finite-exception permutation
lemma then upgrades the tail transpositions to the full finitary symmetric
group.

\subsection{1. Polynomial plus exponential decomposition}

For small \(u\),

\[
 \begin{aligned}
 \Phi(u,z)
 &:=\sum_{q\geq1}F_q(z)u^{q-1}\\
 &=\frac{z^u\Gamma(1-u)-1}{u}+z^u\Gamma(-u,z).
 \end{aligned}
\tag{Inv-1.1}
\]

Write \(w=\Log z+\gamma\) on the principal sheet and define

\[
 Q_q(w)=[u^q]\exp\!\left(
 wu+\sum_{m\geq2}\frac{\zeta(m)}m u^m\right).
\tag{Inv-1.2}
\]

Thus

\[
 Q_0=1,\qquad Q_q'=Q_{q-1},\qquad
 Q_q(w)=\frac{w^q}{q!}+O_q(|w|^{q-2}),
\tag{Inv-1.3}
\]

with the obvious interpretation for \(q=1\), and

\[
 F_q(z)=Q_q(\Log z+\gamma)+R_q(z),\qquad
 R_q(z)=[u^{q-1}]z^u\Gamma(-u,z).
\tag{Inv-1.4}
\]

The upper-incomplete-gamma expansion, uniformly for \(|\arg z|\leq\pi\)
with upper and lower boundary values on the cut, gives

\[
 z^u\Gamma(-u,z)
 =\frac{e^{-z}}z\left(
 \sum_{k=0}^{N-1}\frac{(-1)^k(1+u)_k}{z^k}
 +O_N(|z|^{-N})\right).
\tag{Inv-1.5}
\]

Here and below the required uniformity has the following precise meaning.
Fix \(0<\eta<1\). The standard large-variable remainder theorem for the
upper incomplete gamma function, applied with \(a=-u\) in the compact set
\(|u|\leq\eta\), says that for every \(N\geq1\)

\[
 z^u\Gamma(-u,z)=\frac{e^{-z}}z
 \left(\sum_{k=0}^{N-1}\frac{(-1)^k(1+u)_k}{z^k}
       +E_N(u,z)\right),
 \qquad
 \sup_{|u|\leq\eta}|E_N(u,z)|\leq C_{N,\eta}|z|^{-N}.
\tag{Inv-1.5a}
\]

This holds on \(|\arg z|\leq\pi\), with the corresponding boundary branch
on either side of the negative axis. It is the compact-parameter form of
the usual fixed-parameter expansion (for example, NIST DLMF, \S8.11(i)); it
follows from the same contour-integral remainder after the parameter is
restricted to a compact set. The remainder is holomorphic in \(u\).
Consequently Cauchy's formula on \(|u|=\eta\) gives, for every fixed \(j\),

\[
 |[u^j]E_N(u,z)|\leq C_{N,\eta}\eta^{-j}|z|^{-N}.
\tag{Inv-1.5b}
\]

On each closed subsector strictly inside \(|\arg z|<3\pi/2\), the same
remainder theorem may be differentiated in \(z\), and, after enlarging the
constant,

\[
 \sup_{|u|\leq\eta}|\partial_zE_N(u,z)|
 \leq C_{N,\eta}|z|^{-N-1}.
\tag{Inv-1.5c}
\]

For completeness, (1.5c) also follows from (1.5a) without a separate
differentiated theorem: put the chosen closed subsector inside a slightly
larger one and apply Cauchy's derivative estimate on a circle of radius
\(\varepsilon|z|\). In particular it holds in the upper and lower sectors
about the imaginary axes used in Section 3. Cauchy's formula in \(u\)
applies to (1.5c) as well. These uniform statements justify both
coefficient extraction and the derivative estimate for the contraction maps
below.

Coefficient extraction yields

\[
 R_q(z)=(-1)^{q-1}\frac{e^{-z}}{z^q}
 \left(1-\frac{q(q+1)}{2z}+O_q(z^{-2})\right).
\tag{Inv-1.6}
\]

Indeed, \((1+u)_k\) has degree \(k\); its coefficient of \(u^{q-1}\) first
occurs at \(k=q-1\), where it equals \(1\), and at \(k=q\) that coefficient
is \(1+\cdots+q=q(q+1)/2\). Taking \(N\geq q+1\) in (1.5a)--(1.5b) proves
(1.6), including its uniform error. For later uniform notation set

\[
 Q_0=1,\qquad R_0(z)=-e^{-z},\qquad F_0=Q_0+R_0.
\tag{Inv-1.7}
\]

Formula (1.6), together with (1.7), is the analytic engine for everything
that follows.

\subsection{2. No finite asymptotic values in any order}

\begin{quote}
\textbf{Lemma 2.1.} For every \(q\geq1\) and \(M>0\), every connected
component of
\[
 \{z:|F_q(z)|<M\}
\]
is bounded. Consequently \(F_q\) has no finite asymptotic value.
\end{quote}

Let \(z=x+iy\), \(r=|z|\), \(L=\log r\),
\(\theta=\arg z\in[-\pi,\pi]\), and put \(w=\Log z+\gamma\). Uniformly
as \(r\to\infty\),

\[
 Q_q(w)=\frac{w^q}{q!}\left(1+O_q(L^{-2})\right),
 \qquad
 R_q(z)=(-1)^{q-1}e^{-z}z^{-q}\left(1+O_q(r^{-1})\right).
\tag{Inv-2.0}
\]

If \(|F_q(z)|\leq M\), then \(|Q_q(w)|\asymp_q L^q\), and (1.4) gives

\[
 \frac{R_q(z)}{Q_q(w)}=-1+O_{q,M}(L^{-q}).
\tag{Inv-2.0a}
\]

Taking logarithmic moduli in (2.0)--(2.0a), and using
\(\log|w|=\log L+O(L^{-1})\), gives the following estimate uniformly on
the bounded sublevel, not merely along a selected sequence:

\[
 x=-qL-q\log L+\log(q!)+O_{q,M}(L^{-1}).
\tag{Inv-2.1}
\]

Hence \(|y|\sim r\) and
\(\arg z=\operatorname{sgn}(y)\pi/2+O_q(L/r)\). Taking arguments in the
same cancellation gives

\[
 (q-1)\pi-y-q\theta
 \equiv \pi+\arg Q_q(w)+O_{q,M}(L^{-q})\pmod{2\pi}.
\]

Since \(\arg Q_q(w)=q\arg w+O_q(L^{-2})\) and
\(\arg w=\theta/L+O(L^{-2})\), it follows, separately in the upper and
lower half-planes, that

\[
 \operatorname{dist}\!\left(
 y,2\pi\mathbb Z+\operatorname{sgn}(y)\frac{q\pi}{2}
 \right)\leq \frac{C_{q,M}}L.
\tag{Inv-2.2}
\]

To justify explicitly the word ``uniformly'\,', if either (2.1) or (2.2)
failed outside every disk, a sequence of counterexamples tending to infinity
would contradict (2.0)--(2.0a). Choose \(R=R(q,M)\) so that outside
\(\overline{D(0,R)}\) the sublevel is contained in the horizontal strips of
half-width \(\pi/4\) centered at the lattices in (2.2). These strips are
pairwise disjoint after coincident strips are identified (the upper and
lower lattices either agree or differ by \(\pi\)).

The intersection of the sublevel with any one fixed strip is bounded.
Indeed, otherwise it would contain a sequence with bounded \(y\) and
\(|x|\to\infty\), whereas (2.1) says that \(x\) has only logarithmic size in
\(|x|\), an impossibility. Components of an open subset of \(\mathbb C\)
are path-connected. A component avoiding \(\overline{D(0,R)}\) is therefore
contained in one strip and is bounded. If a component meets the disk, then
every remote strip used by it meets \(\partial D(0,R)\): follow a path from
the disk to a point in that strip and take its last exit from the disk. Only
finitely many of the strips meet that circle, and the sublevel part in each
of them is bounded. The component is again bounded.

If \(F_q(\gamma(t))\to a\) along a curve \(\gamma(t)\to\infty\), then for
some \(M>|a|\) a tail of \(\gamma\) is a connected unbounded subset of
\(\{|F_q|<M\}\), hence lies in an unbounded component. This contradiction
proves the last assertion of the lemma.

Also, every \(F_q\) is surjective. Since \(F_0(x)\) has the sign of \(x\),
(0.1) and induction show that \(F_q'(x)>0\) for real \(x\). The positive
and negative real-axis versions of (1.4)--(1.6) give respectively
\(F_q(x)\to+\infty\) as \(x\to+\infty\) and
\(F_q(x)\to-\infty\) as \(x\to-\infty\). Thus no real value is omitted.
The function is transcendental entire and has real Taylor coefficients, so
a nonreal omitted value would force its distinct conjugate to be omitted as
well, contrary to Picard's theorem.

\subsection{3. Large zeros and eventual separation of critical values}

The following is the required Hardy-type zero lemma. Recall that
\(F_p(z)=z e(-z,1\mid p+1)\) in Hardy's polyexponential notation, so its
leading zero distribution goes back to Hardy. The derivative and
critical-value consequences below are what are needed for monodromy.

\begin{quote}
\textbf{Lemma 3.1 (zero boxes).} Fix \(p\geq1\). Apart from finitely many
zeros, the nonzero zeros of \(F_p\) are conjugate pairs
\(\zeta_{p,n},\overline{\zeta_{p,n}}\), \(n\geq n_0(p)\), with exactly
one upper zero in every corresponding large phase box. There is a
\(C^1\) (indeed real-analytic) parametrized fixed point
\(\zeta_p(t)\), \(t\geq t_0\),
whose integer samples are those zeros and for which
\[
 \begin{aligned}
 \zeta_p(t)
 &=2\pi it-p\log(2\pi t)-p\log\log(2\pi t)+O_p(1),\\
 \zeta_p'(t)&=2\pi i+O_p(t^{-1}),\\
 \Log\zeta_{p,n+1}-\Log\zeta_{p,n}
 &=\frac1n+O_p\!\left(\frac{\log n}{n^2}\right).
 \end{aligned}
\tag{Inv-3.1}
\]
Every sufficiently large zero is simple.
\end{quote}

Here is a direct proof with explicit boxes. Put

\[
 \mathcal B_p(z)=(-1)^{p-1}e^z z^pR_p(z)
 =1-\frac{p(p+1)}{2z}+O_p(z^{-2}).
\]

Fix a large positive integer \(n\), put \(z_n^0=2\pi in\) and
\(L_n=\log(2\pi n)\), and initially work on a rough disk

\[
 E_{p,n}=\{z:|z-z_n^0|\leq C_p\log n\},
\tag{Inv-3.2a}
\]

where \(C_p\) is a sufficiently large fixed constant. On this disk the
principal \(\Log z\) is defined and

\[
 w=\Log z+\gamma=L_n+\gamma+\frac{\pi i}{2}
    +O_p(\log n/n).
\]

Thus \(Q_p(w)\ne0\), and we choose its logarithm by

\[
 \log Q_p(w)=p\Log w-\log(p!)+O_p(w^{-2}).
\]

Also \(\mathcal B_p(z)\ne0\), and we choose
\(\log \mathcal B_p(z)\to0\). The differentiated
uniform remainder in Section 1 gives

\[
 \mathcal B_p'(z)=O_p(n^{-2})\qquad (z\in E_{p,n}).
\tag{Inv-3.2b}
\]

At a zero of \(F_p\), exponentiation gives
\(e^zz^pQ_p(w)=(-1)^p\mathcal B_p(z)\). With
\(\log((-1)^p)=i\pi p\), the corresponding fixed-point equation is

\[
 z=T_{p,n}(z):=2\pi in+i\pi p-p\Log z
 -\log Q_p(\Log z+\gamma)+\log \mathcal B_p(z).
\tag{Inv-3.2c}
\]

Using \(-i\pi p\) instead would merely replace \(n\) by \(n+p\); (3.2c)
fixes the phase index once and for all. Define

\[
 c_{p,n}=T_{p,n}(z_n^0),\qquad
 D_{p,n}=\{z:|z-c_{p,n}|\leq1\}.
\tag{Inv-3.2d}
\]

The selected logarithms give

\[
 \begin{aligned}
 \Re c_{p,n}
 &=-pL_n-p\log L_n+\log(p!)+O_p(L_n^{-1}),\\
 \Im c_{p,n}
 &=2\pi n+\frac{p\pi}{2}
   -\frac{p\pi}{2L_n}+O_p(L_n^{-2}).
 \end{aligned}
\tag{Inv-3.2e}
\]

After increasing \(C_p\), the fixed-radius disk \(D_{p,n}\) lies in
\(E_{p,n}\) for every sufficiently large \(n\). On the rough disk,

\[
 T_{p,n}'(z)=-\frac pz
 -\frac{Q_{p-1}(w)}{zQ_p(w)}
 +\frac{\mathcal B_p'(z)}{\mathcal B_p(z)}
 =O_p(n^{-1}).
\tag{Inv-3.2f}
\]

Consequently, for \(z\in D_{p,n}\), the line segment from \(z_n^0\) to
\(z\) lies in \(E_{p,n}\) and

\[
 |T_{p,n}(z)-c_{p,n}|
 \leq \frac{C_p}{n}|z-z_n^0|
 \leq \frac{C_p(1+\log n)}n<1,
\]

while \(\sup_{D_{p,n}}|T_{p,n}'|<1/2\). Thus \(T_{p,n}\) has exactly one
fixed point in \(D_{p,n}\), and exponentiating (3.2c) shows that it is a
zero of \(F_p\). The centers of consecutive disks differ by
\(2\pi i+o(1)\), so these disks are pairwise disjoint for large \(n\).

Conversely, (2.1)--(2.2), applied to the zeros with \(M=1\) and \(q=p\),
show that every
sufficiently large upper zero has a unique integer \(n\) for which

\[
 z=z_n^0+O_p(\log n),\qquad
 \Im z=2\pi n+\frac{p\pi}{2}+O_p(L_n^{-1}).
\]

With the logarithms selected above it satisfies precisely (3.2c), so (3.2f)
gives

\[
 |z-c_{p,n}|=|T_{p,n}(z)-T_{p,n}(z_n^0)|
 =O_p(\log n/n)<1.
\]

It is therefore the already constructed fixed point. This proves both
existence in every large phase box and exhaustion of all large upper zeros;
conjugation gives the lower zeros.

For real \(t\) sufficiently large, replace \(n\) by \(t\) throughout the
same construction. The estimates are locally uniform in \(t\), and
\(1-T_{p,t}'(\zeta_p(t))\ne0\). The real implicit-function theorem gives a
local \(C^1\), indeed real-analytic, fixed point; uniqueness patches these
local functions into a single \(\zeta_p(t)\). Differentiating (3.2c) yields

\[
 \zeta_p'(t)=\frac{2\pi i}{1-T_{p,t}'(\zeta_p(t))}
 =2\pi i+O_p(t^{-1}).
\]

Equation (3.2e) gives the first line of (3.1), and hence

\[
 \begin{aligned}
 \Log\zeta_p(n+1)-\Log\zeta_p(n)
 &=\int_n^{n+1}\frac{\zeta_p'(t)}{\zeta_p(t)}\,dt\\
 &=\int_n^{n+1}\left(\frac1t
    +O_p\left(\frac{\log t}{t^2}\right)\right)dt
 =\frac1n+O_p\left(\frac{\log n}{n^2}\right),
 \end{aligned}
\]

which completes (3.1).

Finally, consecutive-remainder division, with (1.7) when \(p=1\), gives

\[
 \frac{R_{p-1}(z)}{R_p(z)}
 =-z\left(1+\frac pz+O_p(z^{-2})\right).
\tag{Inv-3.2g}
\]

At a large zero of \(F_p\), \(R_p=-Q_p\), and therefore

\[
 \begin{aligned}
 F_{p-1}(z)
 &=Q_{p-1}(w)+R_{p-1}(z)\\
 &=zQ_p(w)\left(1+O_p(n^{-1})
   +O_p((n\log n)^{-1})\right)
 =zQ_p(w)(1+o(1)).
 \end{aligned}
\tag{Inv-3.3}
\]

Thus \(F_p'(z)=F_{p-1}(z)/z\ne0\) for every sufficiently large zero,
including \(p=1\). Lemma 3.1 follows.

Now let \(q\geq2\), put \(p=q-1\), and define the upper critical values

\[
 b_{q,n}=F_q(\zeta_{p,n}).
\]

At a zero of \(F_p\), the quotient of two consecutive remainders in (1.6)
is

\[
 \frac{R_{p+1}(z)}{R_p(z)}
 =-\frac1z\left(1-\frac{p+1}{z}+O_p(z^{-2})\right).
\]

Since \(R_p=-Q_p\) there,

\[
 b_{q,n}=Q_{p+1}(w_n)+\frac{Q_p(w_n)}{\zeta_{p,n}}
 +O_p\!\left(\frac{(\log n)^p}{n^2}\right),
 \qquad w_n=\Log\zeta_{p,n}+\gamma.
\tag{Inv-3.4}
\]

Using \(Q_{p+1}'=Q_p\) and (3.1),

\[
 \delta w_n:=w_{n+1}-w_n
 =\frac1n+O_p\!\left(\frac{\log n}{n^2}\right).
\]

Taylor's formula, uniformly for \(w\) on the segment between \(w_n\) and
\(w_{n+1}\), gives

\[
 Q_{p+1}(w_{n+1})-Q_{p+1}(w_n)
 =Q_p(w_n)\delta w_n+O_p((\log n)^{p-1}n^{-2}).
\]

Moreover

\[
 \frac{Q_p(w_{n+1})}{\zeta_{p,n+1}}
 -\frac{Q_p(w_n)}{\zeta_{p,n}}
 =O_p((\log n)^p n^{-2}).
\]

The difference of the two remainders in (3.4) is bounded by the sum of
their absolute bounds. Consequently

\[
 b_{q,n+1}-b_{q,n}
 =\frac{Q_p(w_n)}n
 +O_p\!\left(\frac{(\log n)^{p+1}}{n^2}\right).
\tag{Inv-3.5}
\]

Since

\[
 w_n=\log(2\pi n)+\gamma+\frac{\pi i}{2}
 +O_p\!\left(\frac{\log n}{n}\right),
\]

equation (3.5) gives

\[
 \begin{aligned}
 \Re(b_{q,n+1}-b_{q,n})
 &=\frac{(\log(2\pi n))^p}{p!\,n}(1+o(1))>0,\\
 \Im(b_{q,n+1}-b_{q,n})
 &=\frac{\pi(\log(2\pi n))^{p-1}}
 {2(p-1)!\,n}(1+o(1))>0.
 \end{aligned}
\tag{Inv-3.6}
\]

The same formulas also give the individual-value asymptotics

\[
 \begin{aligned}
 \Re b_{q,n}
 &=\frac{(\log(2\pi n))^{p+1}}{(p+1)!}(1+o(1)),\\
 \Im b_{q,n}
 &=\frac{\pi(\log(2\pi n))^p}{2p!}(1+o(1))>0.
 \end{aligned}
\tag{Inv-3.7}
\]

In particular, \(|b_{q,n}|\to\infty\), and the lower critical values are
their conjugates. We now separate all possible tail collisions. For two
upper indices \(m>n\), summing the eventually positive real increments in
(3.6) gives \(\Re b_{q,m}>\Re b_{q,n}\). Upper and lower tail values cannot
coincide because (3.7) puts them in opposite half-planes. Finally, the
critical points not belonging to the two tails form a finite set, whereas
the tail values tend to infinity; after increasing the tail threshold, no
tail value equals a value coming from that finite set. We have proved:

\begin{quote}
\textbf{Corollary 3.2.} For every fixed \(q\geq2\), all but finitely many
critical points of \(F_q\) are simple, all but finitely many critical
values arise from exactly one critical point, and distinct such tail
critical points have distinct critical values. The complete critical
value set is closed and discrete. Every critical value has only finitely
many critical points above it.
\end{quote}

Indeed, the critical points of \(F_q\) are precisely the nonzero zeros of
\(F_{q-1}\). Lemma 3.1 leaves only finitely many of them outside the two
tails. Thus the full critical-value set is a finite set together with two
sequences escaping to infinity, which proves closedness and discreteness.
For any fixed critical value, (3.7) permits only finitely many tail
preimages, and there are only finitely many exceptional critical points.

For completeness, the order \(q=1\) calculation used below is global. The
critical points are

\[
 \operatorname{Crit}(F_1)=\{2\pi ik:k\in\mathbb Z\setminus\{0\}\},
\]

and they are simple because
\(F_1''(2\pi ik)=F_0'(2\pi ik)/(2\pi ik)=1/(2\pi ik)\ne0\). For \(k>0\),

\[
 F_1(2\pi ik)=\operatorname{Cin}(2\pi k)
              +i\operatorname{Si}(2\pi k).
\tag{Inv-3.8}
\]

The function \(\operatorname{Cin}(x)=\int_0^x(1-\cos t)t^{-1}\,dt\) is
strictly increasing at the points \(2\pi k\). Also
\(\operatorname{Si}(2\pi k)>0\): pair the positive half-wave
\([2j\pi,(2j+1)\pi]\) with the following negative half-wave and use the
strict decrease of \(1/t\). Hence the upper values in (3.8) are pairwise
distinct and nonreal, and conjugation separates them from all lower values.
Finally (1.4)--(1.6) on the positive imaginary axis give
\(\operatorname{Cin}(2\pi k)=\log(2\pi k)+\gamma+O(k^{-1})\), so this critical-value set
is closed and discrete, and every value has exactly one critical point above
it.

For \(q=2\), in fact every critical point is simple: a common nonzero zero
of \(F_1\) and \(F_1'\) would have to be \(2\pi ik\), but
\(F_1(2\pi ik)=\operatorname{Cin}(2\pi k)+i\operatorname{Si}(2\pi k)\ne0\).
Global noncollision for \(F_2\) remains a finite, numerically verified but
unproved problem. It is not required in Theorem A.

\subsection{4. The finite-exception permutation lemma}

\begin{quote}
\textbf{Lemma 4.1.} Let \(\Omega\) be countably infinite and let
\(G\leq\operatorname{FSym}(\Omega)\) be transitive. Suppose \(G\) is
generated by a family of transpositions together with finitely many
additional finitary permutations. Then \(G=\operatorname{FSym}(\Omega)\).
\end{quote}

Let \(H\) be the group generated by the transpositions and let its orbits be
the connected components of the associated edge graph. On each component,
\(H\) contains the full finitary symmetric group. If \(H\) is transitive,
then already \(H=\operatorname{FSym}(\Omega)\) and there is nothing to
prove. Assume henceforth that it is not transitive. Every \(H\)-orbit must
meet the finite union of the supports of the exceptional generators;
otherwise that proper orbit is \(G\)-invariant. Thus there are finitely many
\(H\)-orbits, and at least one is infinite.

Start from an infinite orbit \(U\), on which \(G\) contains
\(\operatorname{FSym}(U)\). If an exceptional generator \(g\) sends
\(a\in U\) to \(b\notin U\), choose \(x\in U\) outside the finite support
of \(g\). Then

\[
 g(a\ x)g^{-1}=(b\ x)\in G.
\]

This cross-transposition adjoins the whole \(H\)-orbit of \(b\) to \(U\).
More explicitly, form a finite incidence graph whose vertices are the
\(H\)-orbits and in which two vertices are joined when some exceptional
generator sends a point of one to a point of the other. This graph is
connected: a union of its connected components would otherwise give a
nontrivial \(G\)-invariant union of \(H\)-orbits, contrary to transitivity.
Following a spanning tree from the infinite orbit and repeating the above
cross-transposition construction therefore absorbs every orbit. Hence every
finitary transposition belongs to \(G\), proving the lemma.

\subsection{5. Proof of Theorem A}

Let \(B_q\) be the critical-value set. Corollary 3.2 (and the explicit
\(q=1\) calculation) makes \(B_q\) closed and discrete. Set

\[
 X_q=\mathbb C\setminus F_q^{-1}(B_q),\qquad
 Y_q=\mathbb C\setminus B_q.
\]

The preimage \(F_q^{-1}(B_q)\) is closed and discrete. Indeed, near a point
whose image is not in \(B_q\) this follows from closedness of \(B_q\); near
a point mapping to \(b\in B_q\), choose a target disk containing no other
point of \(B_q\), so locally the preimage is the discrete zero set of
\(F_q-b\). The complement of a closed discrete subset of the plane is
path-connected (perturb a polygonal path around the finitely many points it
meets), hence \(X_q\) is path-connected.

Let \(\Delta=D(a,r)\) be a disk with \(\overline\Delta\subset Y_q\), and
let \(U\) be a component of \(F_q^{-1}(\Delta)\). For
\(M>|a|+r\), the connected set \(U\) lies in one component of
\(\{|F_q|<M\}\); Lemma 2.1 therefore makes \(U\) bounded. The restriction
\(F_q:U\to\Delta\) is proper. To see this without any boundary
assumption, let \(K\Subset\Delta\) and take a sequence
\(z_j\in U\cap F_q^{-1}(K)\). Boundedness gives a convergent subsequence
\(z_j\to z\). Then \(F_q(z)\in K\subset\Delta\); since
\(F_q^{-1}(\Delta)\) is open, a small connected neighborhood of \(z\) lies
in the same component \(U\), so \(z\in U\).

The map on \(U\) is unramified. A proper local biholomorphism has finite
fibers and is a covering onto its image; that image is nonempty, open, and
closed in \(\Delta\), hence is all of \(\Delta\). Since \(U\) is connected
and \(\Delta\) is simply connected, this covering has one sheet. We have
therefore proved that

\[
 F_q:X_q\longrightarrow Y_q
\]

is a connected covering and its monodromy is transitive.

The regular fiber is countably infinite. Indeed, if \(F_q-c\) had only
finitely many zeros, the order-one Hadamard factorization theorem (the order
is immediate from the Taylor coefficients) would give

\[
 F_q(z)-c=P(z)e^{az+b}
\]

with a polynomial \(P\). If \(a\ne0\), a ray on which
\(\Re(az)\to-\infty\) makes the right side tend to zero and gives the finite
asymptotic value \(c\), contrary to Lemma 2.1. If \(a=0\), then \(F_q\)
would be a polynomial. Thus every regular fiber is infinite; being a
discrete subset of \(\mathbb C\), it is countably infinite.

The finitary nature of a meridian requires the bounded-component conclusion
of Lemma 2.1, not merely finiteness of the critical fiber. We record the
precise local statement.

\begin{quote}
\textbf{Lemma 5.1 (finitary meridians).} Let \(f\) be entire and suppose every
component of the inverse image of every bounded disk is bounded. Let
\(b\) be an isolated critical value with only finitely many critical
points above it. Then a meridian around \(b\) has finite support on a
regular fiber. If exactly one critical point lies above \(b\) and it is
simple, the meridian is a single transposition.
\end{quote}

Choose a disk \(\Delta_b\) whose closure contains no critical value of
\(f\) other than \(b\). The compactness argument above,
which did not use absence of critical points, shows that the restriction of
\(f\) to every component \(U\) of \(f^{-1}(\Delta_b)\) is a proper finite
branched cover onto \(\Delta_b\). If \(U\) contains no critical point, the
map is a connected unramified cover of a disk, hence has degree one, and its
sheet is fixed by the meridian. Every critical point in such a component
maps to \(b\); hence only finitely many components contain critical points.
Each of those components has finite degree, so the meridian is supported on
the finite union of their finite fibers. Its nontrivial cycles are exactly
the local ramification cycles. A unique simple critical point therefore
gives one transposition. Transporting the local meridian to the chosen base
point only conjugates this permutation and preserves its finite support and
cycle type. This proves the lemma.

Lemma 5.1 applies to \(F_q\) by Lemma 2.1 and Corollary 3.2 (or (3.8) when
\(q=1\)). Thus every meridian is finitary and all but finitely many
meridians are single transpositions. Finally, every loop in \(Y_q\) is a
finite word in meridians: its compact image lies in a sufficiently large
disk whose boundary avoids \(B_q\), and that disk contains only finitely
many points of the closed discrete set \(B_q\). It follows that the entire
monodromy group lies in \(\operatorname{FSym}\) and is generated by the
tail transpositions together with finitely many exceptional finitary
permutations. It is transitive and acts on a countably infinite fiber, so
Lemma 4.1 gives (0.2).

\subsection{6. Algebraic freedom of branches and first jets}

\begin{quote}
\textbf{Theorem B (all branches are ordinarily algebraically free).} On an
evenly covered simply connected disk, every finite family of distinct
inverse branches of \(F_q\) is algebraically independent over
\(\mathbb C(c)\).
\end{quote}

Suppose a relation exists and clear its coefficients' denominators to obtain
a nonzero polynomial

\[
 P(c;X_1,\ldots,X_n)\in\mathbb C[c,X_1,\ldots,X_n]
\]

which vanishes on the given branches. Choose \(c_*\) in the original disk
so that \(P(c_*;\mathbf X)\) is not the zero polynomial; only finitely many
values are excluded by its nonzero coefficient polynomials. Clearing the
denominators has removed all coefficient poles; equivalently, meridian
representatives could be perturbed away from their finite pole set. Analytic
continuation of the identity around loops based at \(c_*\), together with
the \(n\)-transitivity of \(\operatorname{FSym}\), makes
\(P(c_*;\mathbf X)\) vanish on every ordered distinct \(n\)-tuple in the
infinite fiber. Those tuples are Zariski dense in \(\mathbb C^n\): fix all
but one coordinate, omit the finitely many previously used fiber points, and
induct. This contradicts the choice of \(c_*\) and proves Theorem B.

There is a sharper jet statement.

\begin{quote}
\textbf{Theorem C (all-order first-jet freedom).} For every fixed \(q\geq1\),
distinct inverse branches \(\rho_1,\ldots,\rho_n\) satisfy
\[
 \{\rho_i,\rho_i':1\leq i\leq n\}
 \quad\text{algebraically independent over }\mathbb C(c).
\tag{Inv-6.1}
\]
\end{quote}

Fix an arbitrary regular value \(c_*\). We first write out the fiber version of the
contraction, rather than appeal to it implicitly. On the rough disk
\(E_{q,n}\) from (3.2a), \(Q_q(w)-c_*\ne0\) for large \(n\), and we choose
the logarithm asymptotic to \(q\Log w-\log(q!)\). The equation
\(F_q(z)=c_*\) is equivalent to

\[
 z=T^{c_*}_{q,n}(z):=2\pi in+i\pi q-q\Log z
 -\log(Q_q(\Log z+\gamma)-c_*)+\log\mathcal B_q(z).
\tag{Inv-6.2}
\]

Define \(c^{c_*}_{q,n}=T^{c_*}_{q,n}(2\pi in)\) and

\[
 D^{c_*}_{q,n}=\{z:|z-c^{c_*}_{q,n}|\leq1\}.
\]

The same direct expansion used in (3.2e), now with a fixed lower-order term
\(-c_*\), gives
\(|c^{c_*}_{q,n}-2\pi in|=O_{q,c_*}(\log n)\); hence this disk lies in
\(E_{q,n}\). Moreover,

\[
 (T^{c_*}_{q,n})'(z)
 =-\frac qz-
 \frac{Q_{q-1}(w)}{z(Q_q(w)-c_*)}
 +\frac{\mathcal B_q'(z)}{\mathcal B_q(z)}
 =O_{q,c_*}(n^{-1}).
\tag{Inv-6.3}
\]

Thus, for all sufficiently large \(n\),
\[
 |T^{c_*}_{q,n}(z)-c^{c_*}_{q,n}|
 \leq \frac{C_{q,c_*}}n|z-2\pi in|
 \leq \frac{C_{q,c_*}(1+\log n)}n<1
 \qquad(z\in D^{c_*}_{q,n}),
\]

and \(\sup_{D^{c_*}_{q,n}}|(T^{c_*}_{q,n})'|<1/2\). Thus
\(T^{c_*}_{q,n}(D^{c_*}_{q,n})\subset D^{c_*}_{q,n}\).
Exponentiating (6.2) identifies its unique
fixed point \(z_n(c_*)\) with a point of \(F_q^{-1}(c_*)\). This supplies
an infinite sequence \(|z_n(c_*)|\to\infty\), since the consecutive centers
again differ by \(2\pi i+o(1)\).

At such a point \(R_q=c_*-Q_q\). Consecutive-remainder division, including
\(q=1\) by (1.7), gives

\[
 \frac{R_{q-1}}{R_q}
 =-z\left(1+\frac qz+O_q(z^{-2})\right),
\]

and hence

\[
 F_{q-1}(z)=Q_{q-1}(w)
 +z(Q_q(w)-c_*)\left(1+\frac qz+O_q(z^{-2})\right).
\tag{Inv-6.4}
\]

Using (0.1) and
\(Q_{q-1}(w)/(z(Q_q(w)-c_*))=O_{q,c_*}((z\Log z)^{-1})\), we obtain

\[
 p(z):=\frac1{F_q'(z)}
 =\frac1{Q_q(\Log z+\gamma)-c_*}
  \left(1+O_{q,c_*}(z^{-1})\right)
 \sim\frac{q!}{(\Log z)^q}.
\tag{Inv-6.5}
\]

The jet set

\[
 J_{q,c_*}=\{(z,p(z)):F_q(z)=c_*\}
\]

is Zariski dense in \(\mathbb C^2\). Indeed, for a nonzero polynomial
\(A(z,p)\), choose its largest \(z\)-degree \(d\), and within that degree
choose the smallest occurring \(p\)-degree \(j\). Along the sequence above,
the term \(a_{d,j}z^dp^j\) dominates terms of the same \(z\)-degree and
larger \(p\)-degree because \(p\to0\). If \(N\) is the largest difference
of \(p\)-degrees that occurs, it dominates every lower \(z\)-degree term
because

\[
 |z|^{-1}|p|^{-N}=O((\log|z|)^{qN}/|z|)\longrightarrow0.
\]

Thus \(A\) cannot vanish on all of \(J_{q,c_*}\).

We use the following elementary product fact. If an infinite subset
\(S\subset\mathbb C^2\) is Zariski dense, then \(S\) with finitely many
points removed is still Zariski dense: otherwise the irreducible variety
\(\mathbb C^2\) would be the union of a proper closed set and finitely many
points. It follows by induction that

\[
 S^{[n]}=\{(s_1,\ldots,s_n)\in S^n:s_i\ne s_j\ (i\ne j)\}
\]

is Zariski dense in \((\mathbb C^2)^n\). Indeed, after fixing
\(s_1,\ldots,s_{n-1}\), a polynomial vanishing on \(S^{[n]}\) vanishes as a
polynomial in its last pair of variables on
\(S\setminus\{s_1,\ldots,s_{n-1}\}\); its coefficients then vanish on the
ordered distinct \((n-1)\)-fold product.

Since the construction and density proof apply to every regular \(c_*\), we
may now specialize a hypothetical relation at a convenient value. Suppose
that the \(2n\) first-jet functions satisfy an algebraic
relation. Clearing denominators gives a nonzero

\[
 P(c;Z_1,W_1,\ldots,Z_n,W_n)\in
 \mathbb C[c,Z_1,W_1,\ldots,Z_n,W_n].
\]

Choose \(c_*\) in the original regular disk so that its specialization in
the jet variables is nonzero. Clearing denominators removes coefficient
poles (or, equivalently, the finitely many poles can be avoided by the loop
representatives). Under continuation around a loop, a sheet
\(z\) and its derivative move together as
\((z,1/F_q'(z))\). The \(n\)-transitivity from Theorem A therefore makes
the specialized polynomial vanish on the ordered distinct \(n\)-fold
product of \(J_{q,c_*}\), which is Zariski dense by the preceding paragraph.
This contradicts nonzero specialization and proves (6.1).

\subsection{7. Ordinary freedom versus differential algebra}

If \(\rho(c)\) is an inverse branch and

\[
 \mathcal L_\rho=\frac{\rho}{\rho'}\frac d{dc},
\]

then (0.1) gives

\[
 H_q:=\mathcal L_\rho^q(c)=1-e^{-\rho},\qquad
 H_q'=\rho'(1-H_q).
\tag{Inv-7.1}
\]

After clearing powers of \(\rho'\), (7.1) is a differential-polynomial
equation of order at most \(q+1\). Hence every inverse branch is
differential-algebraic, while Theorems B--C show extreme ordinary algebraic
freedom.

For \(q=1\), (7.1) is exactly

\[
 \rho\rho''+(\rho')^3-(\rho+1)(\rho')^2=0.
\tag{Inv-7.2}
\]

Theorem C shows that no order-zero or order-one differential equation can
hold, so the differential order is exactly two. Differentiating (7.2)
expresses every higher derivative rationally in \(\rho,\rho'\). Therefore,
for any \(n\) distinct branches,

\[
 \operatorname{trdeg}_{\mathbb C(c)}
 \mathbb C(c)\langle\rho_1,\ldots,\rho_n\rangle=2n,
 \qquad
 \operatorname{dtrdeg}_{\mathbb C(c)}=0.
\tag{Inv-7.3}
\]

This is the sharpest ordinary/differential separation presently available.

\subsection{8. Coupling to exact two-trace recovery}

For a nonzero level \(c\), and the zero divisor of \(F_q(z)-c\), let

\[
 T_m^{(q)}(c)=\sum_{F_q(\rho)=c}\rho^{-m},\qquad m\geq2.
\]

The series is absolutely convergent because \(F_q-c\) has order one. We
recall the trace calculation, with the same normalization as the finite
Faber-monodromy section above. Hadamard factorization gives

\[
 F_q(z)-c=-c\,e^{\alpha z}
 \prod_{F_q(\rho)=c}\left(1-\frac z\rho\right)e^{z/\rho},
\]

with zeros repeated according to multiplicity. Comparing the coefficient
of \(z^m\), \(m\geq2\), in the logarithm at \(z=0\) gives

\[
 T_m^{(q)}(c)=P_{q,m}(c^{-1}),\qquad
 P_{q,m}(X):=-m[z^m]\log(1-XF_q(z)).
\tag{Inv-8.0}
\]

Indeed,

\[
 P_{q,m}(X)=\sum_{k=1}^m\frac{m}{k}[z^m]F_q(z)^kX^k,
\]

so \(P_{q,m}\in\mathbb Q[X]\) is monic of degree \(m\). Writing
\(X=c^{-1}\) and using the first four Taylor coefficients of \(F_q\) gives

\[
 \begin{aligned}
 T_2^{(q)}&=X^2-\frac{X}{2^q},\\
 T_4^{(q)}&=X^4-\frac{2X^3}{2^q}
 +\left(\frac2{3^{q+1}}+\frac1{2^{2q+1}}\right)X^2
 -\frac{X}{6\cdot4^q}.
 \end{aligned}
\tag{Inv-8.0a}
\]

In particular, if

\[
 A_q=\frac2{3^{q+1}}-\frac1{2^{2q+1}},
\]

and

\[
 C_q=-\frac1{6\cdot4^q}+\frac2{3\cdot6^q}
     -\frac1{2^{3q+1}}\ne0,
\]

(vanishing would give \(4^{q+1}=6^q+3^{q+1}\), impossible modulo \(3\)),

then (8.0a) gives

\[
 T_4^{(q)}-(T_2^{(q)})^2-A_qT_2^{(q)}=C_qX,
\]

and hence

\[
 c^{-1}=\frac{T_4^{(q)}-(T_2^{(q)})^2-A_qT_2^{(q)}}{C_q}.
\tag{Inv-8.1}
\]

Consequently

\[
 \mathbb C(c)=\mathbb C(T_2^{(q)},T_4^{(q)}).
\tag{Inv-8.2}
\]

Combining (8.2) with Theorems A--C produces the all-order
\textbf{local-freedom/global-collapse dichotomy}:

\begin{itemize}
\tightlist
\item
  every finite collection of local inverse branches, and even their first
  jets, is algebraically free;
\item
  the monodromy of all local sheets is the maximal finitary symmetric group;
\item
  nevertheless, all global reciprocal Newton traces of the infinite divisor
  are functions of two traces, and those two traces recover the level exactly.
\end{itemize}

This synthesis is substantially stronger than either the monodromy theorem or
the trace calculation in isolation.

\subsection{9. Literature boundary and attribution}

\begin{itemize}
\tightlist
\item
  NIST Digital Library of Mathematical Functions, \S8.11(i), equations
  8.11.2--8.11.3 (with the parameter restricted to a compact set), supplies
  the large-variable incomplete-gamma remainder used in (1.5a). The
  contour-remainder proof gives compact-parameter uniformity, and Cauchy's
  estimates give (1.5b)--(1.5c).
\item
  G. H. Hardy, \emph{On the Zeroes of Certain Classes of Integral Taylor
  Series, Part II}, Proc. London Math. Soc. (2) \textbf{2} (1905), 401--431,
  DOI: \texttt{10.1112/plms/s2-2.1.401}, studied the zeros and asymptotics of
  \(e(x,a\mid s)\). Since \(F_p(z)=z e(-z,1\mid p+1)\), the leading
  zero-box asymptotic must be credited to this line.
\item
  K. N. Boyadzhiev, \emph{Polyexponentials}, arXiv:0710.1332, surveys the
  defining identities and Hardy asymptotics.
\item
  L. Zelenko, \emph{Generic monodromy group of Riemann surfaces for inverses to
  entire functions of finite order}, arXiv:2105.14015, proves
  \(\operatorname{FSym}(\mathbb N)\) for a generic/``typical'\,' class with
  globally simple critical points and globally pairwise-distinct critical
  values.
\end{itemize}

Thus \(\operatorname{FSym}\) as an abstract phenomenon is not new. The
potentially new content here is the \textbf{explicit all-order polyexponential
theorem}, the proof that only finitely many exceptional critical events can
occur, the finite-exception permutation upgrade (which avoids proving global
critical-value noncollision), and the all-order branch/first-jet algebraic
freedom. A full submission still requires a systematic MathSciNet/zbMATH
priority search.

\subsection{10. What remains open}

\begin{enumerate}
\def\labelenumi{\arabic{enumi}.}
\tightlist
\item
  Prove or disprove global pairwise noncollision of the critical values of
  \(F_q\), already open in the finite remainder for \(q=2\).
\item
  Determine whether the inverse differential equation in (7.1) has minimal
  order \(q+1\) for every \(q\geq2\). Theorem C gives only the lower bound
  two.
\item
  Determine algebraic independence of higher inverse jets and the exact
  ordinary transcendence degree of the differential field generated by
  finitely many branches for \(q\geq2\).
\item
  Determine whether the local-freedom/global-collapse dichotomy has an
  analogue for other iterated-integral towers beyond the polyexponentials.
\end{enumerate}

\part{Ramanujan--Gamma Transmutations and Euler Continued Fractions}
\section{An explicit solution of a 2026 Euler recurrence}
\label{sec:ramanujan-recurrence}

The Ramanujan Challenge for AI \cite{RamanujanChallenge2026} asks, in its
Euler-constant problem, for a proof that the quotient of two distinguished
solutions of the following four-term recurrence tends to \(\gamma\).  We
give a self-contained solution.  The finite sums below belong to the
multiple-Laguerre framework developed by Wolfs--Van Assche
\cite{WolfsVanAssche2025}.  What is proved here is the adjacent
specialization, explicit all-index recurrence certificates, an elementary
positive concentration proof, and the flat matrix field.

Put
\begin{align*}
 A_n&=8n^3+51n^2+105n+68,\\
 B_n&=24n^5+337n^4+1833n^3+4818n^2+6092n+2928,\\
 C_n&=(n+2)(n+3)
 (24n^5+273n^4+1150n^3+2154n^2+1635n+268),\\
 D_n&=(n+1)(n+2)^4(n+3)
 (8n^3+75n^2+231n+232),
\end{align*}
and
\begin{equation}\label{eq:challenge-recurrence}
 A_nu_n=B_nu_{n-1}-C_nu_{n-2}+D_nu_{n-3}
 \qquad(n\ge0).
\end{equation}

\begin{theorem}[Closed solution of the Euler recurrence]
\label{thm:challenge-solution}
For \(n\ge-3\), set \(a=n+3\), \(b=n+4=a+1\), and
\begin{equation}\label{eq:Wab}
 W_{a,b}(k)=
 \frac{(a!b!)^2}{(k!)^3(a-k)!(b-k)!}
 \qquad(0\le k\le a).
\end{equation}
Then the two sequences determined by
\[
 (p_{-3},p_{-2},p_{-1})=(0,7,179),
 \qquad
 (q_{-3},q_{-2},q_{-1})=(1,12,306)
\]
and \eqref{eq:challenge-recurrence} have the exact formulas
\begin{align}
 q_n&=b\sum_{k=0}^{a}W_{a,b}(k),
 \label{eq:qn-closed}\\
 p_n&=b\sum_{k=0}^{a}W_{a,b}(k)
 \bigl(3H_k-H_{a-k}-H_{b-k}\bigr)+a!.
 \label{eq:pn-closed}
\end{align}
In particular,
\begin{equation}\label{eq:challenge-limit}
 \boxed{\displaystyle\lim_{n\to\infty}\frac{p_n}{q_n}=\gamma.}
\end{equation}
\end{theorem}

\begin{proof}
Substitution of \((a,b)=(0,1),(1,2),(2,3)\) gives the six initial
values.  The two symbolic telescoping identities in
\cref{app:telescopers} prove \eqref{eq:challenge-recurrence} for both
closed forms at every \(n\ge0\).  We prove the limit directly.

After cancelling factors independent of \(k\), define a probability law by
\begin{equation}\label{eq:challenge-pmf}
 \Pr(K_n=k)=
 \frac{\binom ak\binom bk/k!}
 {\sum_{j=0}^{a}\binom aj\binom bj/j!}.
\end{equation}
Then
\begin{equation}\label{eq:challenge-expectation}
 \frac{p_n}{q_n}
 =\mathbb E\bigl(3H_{K_n}-H_{a-K_n}-H_{b-K_n}\bigr)
 +\frac{a!}{q_n}.
\end{equation}
The ratio of consecutive unnormalized masses is
\begin{equation}\label{eq:challenge-ratio}
 r_k=\frac{(a-k)(b-k)}{(k+1)^3},
\end{equation}
which decreases strictly in \(k\).  Put \(x_n=(ab)^{1/3}\).  For each
fixed \(\eta>0\), the ratios outside
\([(1-\eta)x_n,(1+\eta)x_n]\) are uniformly separated from one once
\(n\) is large.  Multiplying the ratios on either side of the mode gives
\begin{equation}\label{eq:challenge-concentration}
 \Pr\bigl(|K_n/x_n-1|>\eta\bigr)
 =O\bigl(e^{-c_\eta x_n}\bigr).
\end{equation}
Hence \(K_n/x_n\to1\) in probability.

On the event \(|K_n/x_n-1|\le\eta\), the estimate
\(H_m=\log m+\gamma+O(m^{-1})\) yields
\begin{align*}
 3H_k-H_{a-k}-H_{b-k}
 &=\gamma+\log\frac{k^3}{(a-k)(b-k)}+o(1)\\
 &=\gamma+O(\eta)+o(1).
\end{align*}
The expression is \(O(\log n)\) on the full support, while the exceptional
probability in \eqref{eq:challenge-concentration} is exponentially small in
\(x_n\).  First let \(n\to\infty\), and then let \(\eta\downarrow0\); its
expectation therefore tends to \(\gamma\).  Finally the
\(k=0\) summand in \eqref{eq:qn-closed} gives
\[
 0\le\frac{a!}{q_n}\le\frac1{b\,b!}\longrightarrow0.
\]
Equation \eqref{eq:challenge-expectation} proves
\eqref{eq:challenge-limit}.
\end{proof}

\begin{remark}[What has and has not been solved]
The theorem gives an independent proof of the publicly stated limit problem.
The official proof and code are now public \cite{RamanujanOfficial2026}; the
sentence in Version 2.0 saying otherwise is superseded.  No priority claim is
made for the official Meijer--\(G\) trajectory, Miller selection, or the broad
multiple-Laguerre family.  The comparison with the terminating state is
established in \cref{sec:meijer-duality}.
The convergence \(p_n/q_n\to\gamma\) does not by itself imply the
irrationality of \(\gamma\).
\end{remark}

The same recurrence has a second nontrivial Ap\'ery limit.  Let \(w_n\) be
the solution with
\begin{equation}\label{eq:w-initial}
 (w_{-3},w_{-2},w_{-1})=(1,18,440),
\end{equation}
and let \(\mathfrak h\) be the positive constant in
\eqref{eq:h-constant} below.

\begin{theorem}[Complete initial-value limit classification]
\label{thm:ram-full-solution}
One has
\begin{equation}\label{eq:w-limit}
 \boxed{\displaystyle\frac{w_n}{q_n}\longrightarrow\mathfrak h
 =1.4322057346532244148\ldots.}
\end{equation}
For arbitrary rational initial data
\((x,y,z)=(u_{-3},u_{-2},u_{-1})\), the corresponding solution satisfies
\begin{align}
 \lim_{n\to\infty}\frac{u_n}{q_n}
 =\frac1{136}\bigl[&(-6x+179y-7z)\mathfrak h\notag\\
 &+(-228x-134y+6z)\gamma
 +(142x-179y+7z)\bigr].
 \label{eq:all-initial-limits}
\end{align}
Consequently the set of all rational-initial-data limits is exactly
\begin{equation}\label{eq:attainable-limit-span}
 \boxed{\Q+\Q\gamma+\Q\mathfrak h.}
\end{equation}
No assertion that this span has dimension three is made.
\end{theorem}

\begin{proof}
The official cyclic gauge at zero has rows
\begin{equation}\label{eq:U0-rows}
 \mathsf U(0)=
 \begin{pmatrix}1&18&440\\0&7&179\\0&1&45\end{pmatrix}.
\end{equation}
They generate solutions \(w,p,s\), respectively, and
\(q=w-p+s\).  The official Miller-selection theorem, transported through
the exact gauge of \cref{thm:ram-gauge}, gives
\begin{equation}\label{eq:Miller-limit-vector}
 \frac1{q_n}(w_n,p_n,s_n)
 \longrightarrow
 (\mathfrak h,\gamma,1+\gamma-\mathfrak h);
\end{equation}
see \cite[\S5.2]{RamanujanOfficial2026} and
\cref{thm:meijer-probability}.  The identity \(q=w-p+s\) normalizes the
right side to one.  Solving the initial vector
\((x,y,z)\) in the basis \((q,p,w)\) gives the three coefficients in
\eqref{eq:all-initial-limits}.  Conversely \(q,p,w\) realize
\(1,\gamma,\mathfrak h\), proving \eqref{eq:attainable-limit-span}.
\end{proof}

\begin{theorem}[Complete integer initial lattice]
\label{thm:ram-integer-lattice}
For a solution of \eqref{eq:challenge-recurrence},
\begin{equation}\label{eq:integer-lattice}
 \boxed{\displaystyle
 u_n\in\Z\ \text{for every }n\ge-3
 \Longleftrightarrow
 x,y,z\in\Z\ \text{and }13x+6y+z\equiv0\pmod{17}.}
\end{equation}
\end{theorem}

\begin{proof}
The first new term is
\begin{equation}\label{eq:u0-congruence}
 u_0=\frac{2784x-402y+732z}{17},
\end{equation}
so integrality forces the stated congruence.  Conversely its kernel in
\(\Z^3\) has basis
\[
 (1,0,4),\qquad(0,1,11),\qquad(0,0,17).
\]
The finite companion-gauge certificate in
\cref{app:integer-lattice-certificate} proves that the three corresponding
solutions are integral at every index.  Every integral vector satisfying the
congruence is their integral combination.
\end{proof}

\section{A flat conservative matrix field and exterior arithmetic}
\label{sec:cmf}

For integers \(a,b\ge0\), define
\begin{equation}\label{eq:Fab}
 F_{a,b}(z)={}_{2}F_{2}
 \!\left(\begin{matrix}-a,-b\\1,1\end{matrix};z\right),
 \qquad
 \theta=z\frac d{dz},
 \qquad
 \mathbf Y_{a,b}=
 \begin{pmatrix}F_{a,b}\\ \theta F_{a,b}\\ \theta^2F_{a,b}\end{pmatrix}.
\end{equation}
It satisfies
\begin{equation}\label{eq:Fab-ode}
 \theta^3F_{a,b}=z(\theta-a)(\theta-b)F_{a,b}.
\end{equation}
At \(z=1\), its terminating series gives the exact link
\begin{equation}\label{eq:q-Fab-link}
 F_{a,b}(1)=\sum_{k=0}^{a}\binom ak\binom bk\frac1{k!},
 \qquad q_n=b\,a!b!\,F_{a,b}(1)\quad(a=n+3,\ b=n+4).
\end{equation}

\begin{theorem}[Flat \(SL_3\) conservative matrix field]
\label{thm:flat-cmf}
For \(a,b>0\), put
\begin{align}
 M_a(a,b;z)&=\frac1a
 \begin{pmatrix}
 a&-1&0\\
 0&a&-1\\
 -zab&z(a+b)&a-z
 \end{pmatrix},\label{eq:Ma}\\
 M_b(a,b;z)&=\frac1b
 \begin{pmatrix}
 b&-1&0\\
 0&b&-1\\
 -zab&z(a+b)&b-z
 \end{pmatrix}.\label{eq:Mb}
\end{align}
Then
\[
 \mathbf Y_{a-1,b}=M_a(a,b;z)\mathbf Y_{a,b},
 \qquad
 \mathbf Y_{a,b-1}=M_b(a,b;z)\mathbf Y_{a,b},
\]
and
\begin{align}
 \det M_a&=\det M_b=1,\label{eq:cmf-det}\\
 M_b(a-1,b;z)M_a(a,b;z)
 &=M_a(a,b-1;z)M_b(a,b;z).
 \label{eq:cmf-flatness}
\end{align}
Thus the denominator system of \cref{thm:challenge-solution} is the
adjacent-diagonal path \((a,b)=(n+3,n+4)\) in an explicit flat \(SL_3\)
conservative matrix field.
\end{theorem}

\begin{proof}
The two contiguous relations follow by comparing the terminating
hypergeometric coefficients and using \eqref{eq:Fab-ode} for the third row.
Expansion of the determinants gives \(1\).  Direct matrix multiplication
gives the polynomial identity \eqref{eq:cmf-flatness}; an entrywise symbolic
certificate is also included with the source package.
\end{proof}

\begin{remark}[Path rigidity]
Flatness says that changing a lattice path while keeping the same endpoint
\((a,b)\) leaves \(F_{a,b}\), and hence the endpoint-defined denominator,
unchanged.  The numerator and quotient in \cref{thm:challenge-solution} are
also endpoint-defined by their closed formulas.  Rerouting a path therefore
does not change these particular values.  Changing the endpoint curve is a
different problem and is not excluded by this observation.
\end{remark}

The numerator is the first Frobenius jet of the same flat connection.  Let
\begin{equation}\label{eq:Phi-epsilon}
 \Phi_{a,b}(\epsilon,z)=
 \frac{a!b!}{(1+\epsilon)_a(1+\epsilon)_b}
 {}_3F_3\!\left(
 \begin{matrix}-a-\epsilon,-b-\epsilon,1\\
 1-\epsilon,1-\epsilon,1-\epsilon
 \end{matrix};z\right).
\end{equation}

\begin{theorem}[Flat first-jet \(SL_6\) lift]
\label{thm:ram-sl6}
Fix positive integers \(a,b\), expand the analytic \(\epsilon\)-germ at
zero modulo \(\epsilon^2\), and put
\(\mathbf Y_{a,b}(\epsilon)
=(\Phi_{a,b},\theta\Phi_{a,b},\theta^2\Phi_{a,b})^{\mathsf T}\).
Then
\begin{equation}\label{eq:epsilon-contiguous}
 \mathbf Y_{a-1,b}(\epsilon)=\mathscr M_a(\epsilon)\mathbf Y_{a,b}(\epsilon),
 \qquad
 \mathbf Y_{a,b-1}(\epsilon)=\mathscr M_b(\epsilon)\mathbf Y_{a,b}(\epsilon),
\end{equation}
where
\begin{equation}\label{eq:epsilon-Ma}
 \mathscr M_a(\epsilon)=\frac1a
 \begin{pmatrix}
 a+\epsilon&-1&0\\
 0&a+\epsilon&-1\\
 -zab-\epsilon z(a+b)&z(a+b)+2\epsilon z&a-z-2\epsilon
 \end{pmatrix},
\end{equation}
and
\begin{equation}\label{eq:epsilon-Mb}
 \mathscr M_b(\epsilon)=\frac1b
 \begin{pmatrix}
 b+\epsilon&-1&0\\
 0&b+\epsilon&-1\\
 -zab-\epsilon z(a+b)&z(a+b)+2\epsilon z&b-z-2\epsilon
 \end{pmatrix}.
\end{equation}
In
\(\Q(a,b,z)[\epsilon]/(\epsilon^2)\), these matrices have determinant one
and satisfy the same zero-curvature identity as
\eqref{eq:cmf-flatness}.

Writing
\(\mathscr M_a=\mathscr M_{a,0}+\epsilon\mathscr M_{a,1}\), the block matrix
\begin{equation}\label{eq:SL6-block}
 \begin{pmatrix}
 \mathscr M_{a,0}&0\\
 \mathscr M_{a,1}&\mathscr M_{a,0}
 \end{pmatrix}\in SL_6
\end{equation}
transports
\(\binom{\mathbf Y_{a,b}(0)}
{\partial_\epsilon\mathbf Y_{a,b}(0)}\), and similarly in the
\(b\)-direction.  Thus denominator and numerator path-independence are the
zero- and first-jet parts of one flat connection.

At \(b=a+1,z=1\),
\begin{equation}\label{eq:Phi-pq-jets}
 b\,a!b!\Phi_{a,b}(0,1)=q_{a-3},\qquad
 b\,a!b!\partial_\epsilon\Phi_{a,b}(0,1)=p_{a-3}.
\end{equation}
The boundary \(a=0,b=1\) in \eqref{eq:Phi-pq-jets} is verified directly.
\end{theorem}

\begin{proof}
Modulo \(\epsilon^2\), the Frobenius deformation satisfies
\begin{equation}\label{eq:Phi-dual-ODE}
 (\theta-\epsilon)^3\Phi_{a,b}
 =z(\theta-a-\epsilon)(\theta-b-\epsilon)\Phi_{a,b}.
\end{equation}
Coefficient comparison gives the first two rows of
\eqref{eq:epsilon-contiguous}, while \eqref{eq:Phi-dual-ODE} gives the third.
Direct determinant expansion and multiplication give determinant one and
flatness in the dual-number ring.  Expanding a
dual-number matrix action gives the block form \eqref{eq:SL6-block}.
At \(\epsilon=0\), \(\Phi_{a,b}=F_{a,b}\).  The derivative of the terminating
part gives the harmonic weight in \eqref{eq:pn-closed}; the regularized term
at \(k=b=a+1\) is \(1/(bb!)\) before multiplication by \(b\,a!b!\), and
hence gives exactly the endpoint \(+a!\).
\end{proof}

\begin{proposition}[First-jet exactness barrier]
\label{prop:first-jet-barrier}
Let \(\mathscr K_a\) be the rational carrier defined in
\eqref{eq:Kas-def}.  For \(b=a+1,z=1\),
\begin{equation}\label{eq:Phi-K-contact}
 \Phi_{a,b}(\epsilon,1)-\mathscr K_a(\epsilon)=O(\epsilon^2).
\end{equation}
Return here to the full analytic \(\epsilon\)-germ.  For every \(r\ge1\),
the term with \(k=b+r\) contributes
\begin{equation}\label{eq:E-tail-second-jet}
 [\epsilon^2]\,\Phi_{a,b}(\epsilon,1)\big|_{k=b+r}
 =-\frac{a!b!\,r!(r-1)!}{(b+r)!^3}\ne0.
\end{equation}
The total second-order difference cannot cancel:
\begin{equation}\label{eq:second-jet-total-positive}
 [\epsilon^2]\{\Phi_{a,b}(\epsilon,1)-\mathscr K_a(\epsilon)\}
 =
 \frac1{b^2a!}\left\{
 3H_b-\sum_{r\ge1}\frac{r!(r-1)!}{(b+1)_r^3}\right\}>0.
\end{equation}
Indeed \((b+1)_r\ge(r+1)!\), so the sum is less than one, while
\(3H_b\ge3\).
Thus the first derivative is exactly terminating after its resonant endpoint,
whereas the second derivative develops a genuine infinite hypergeometric
tail.  This is a barrier for this Frobenius deformation, not a universal
no-go for all higher-jet constructions.
\end{proposition}

\begin{proof}
For \(0\le k\le a\), the hypergeometric summand is identically the
corresponding summand of \(\mathscr J_a(\epsilon)\).  At the resonant index
\(k=b\), its analytic continuation is
\[
 \frac{\epsilon}{b^2a!}
 \prod_{j=1}^b\left(1-\frac{\epsilon}{j}\right)^{-3};
\]
its linear coefficient is \(1/(bb!)\), exactly the correction in
\(\mathscr K_a\), and its quadratic coefficient is
\(3H_b/(b^2a!)\).  For \(k=b+r\), \(r\ge1\), direct expansion gives
\eqref{eq:E-tail-second-jet}.  Summing those coefficients gives
\eqref{eq:second-jet-total-positive}; the displayed comparison proves its
strict positivity.  Hence the zero and first jets agree and the second
jets do not.
\end{proof}

\begin{proposition}[Forced divisor and exact Casoratian]
\label{prop:challenge-arithmetic}
Let \(d_b=\lcm(1,\ldots,b)\).  For every \(n\ge-3\),
\begin{equation}\label{eq:challenge-gcd}
 \boxed{\displaystyle
 \frac{b!}{d_b}\gcd\!\left(b,\frac{d_b}{b}\right)
 \mid\gcd(p_n,q_n).}
\end{equation}
If \(u^{(1)},u^{(2)},u^{(3)}\) are three solutions of
\eqref{eq:challenge-recurrence}, put, for \(n\ge-1\),
\[
 \mathcal W_n=\det\bigl(u^{(i)}_{n-j}\bigr)_{1\le i\le3,\,0\le j\le2},
\]
so in particular
\(\mathcal W_{-1}=\det(u^{(i)}_{-1-j})_{1\le i\le3,\,0\le j\le2}\).
Then, for every \(n\ge0\),
\begin{equation}\label{eq:challenge-casoratian}
 \boxed{\displaystyle
 \mathcal W_n=\mathcal W_{-1}
 \frac{A_{n+1}}{136}
 (n+1)!(n+2)!^4(n+3)!.}
\end{equation}
\end{proposition}

\begin{proof}
Write
\[
 W_{a,b}(k)=a!b!\binom ak\binom bk\frac1{k!}.
\]
Put \(g_b=\gcd(b,d_b/b)\).  After division by \(b!/d_b\), each summand of
\(q_n\) is \(b\) times an integer, because \(a!/k!\in\Z\).  The harmonic
part of \(p_n\) is likewise \(b\) times an integer, since
\(d_b(3H_k-H_{a-k}-H_{b-k})\in\Z\).  Its endpoint becomes \(d_b/b\).
Hence every term of both quotients is divisible by \(g_b\), proving
\eqref{eq:challenge-gcd}.

The companion determinant of \eqref{eq:challenge-recurrence} is
\(D_n/A_n\).  Since
\[
 D_n=(n+1)(n+2)^4(n+3)A_{n+1},
\]
iteration from \(-1\) to \(n\) telescopes and gives
\eqref{eq:challenge-casoratian}.
\end{proof}

\begin{remark}[Exterior-volume barrier]
The flat base field has determinant one, whereas scalar elimination together
with its factorial gauge produces the factorial-six Casoratian
\eqref{eq:challenge-casoratian}.  The forced divisor
\eqref{eq:challenge-gcd} is a lower bound for the common divisor, not an
upper bound.  These formulas record the factorial cost but do not rule out
further common-divisor cancellation; no small primitive integer linear form
is obtained here.
\end{remark}

\section{Exact Meijer--\texorpdfstring{\(G\)}{G}/hypergeometric duality}
\label{sec:meijer-duality}

The official solution of the Challenge is now public
\cite{RamanujanOfficial2026}.  It selects a recessive Meijer--\(G\) state by
Miller's algorithm, whereas \cref{sec:ramanujan-recurrence,sec:cmf} uses a
terminating \({}_2F_2\) state and a positive saddle.  The terminating module
is rationally gauge-equivalent to the contragredient of the normalized
official rank-three module.

Write
\[
 F_a(z)=F_{a,a+1}(z),\qquad
 \mathbf Y_a=
 \begin{pmatrix}F_a(1)\\ \theta F_a(1)\\ \theta^2F_a(1)\end{pmatrix}.
\]
The adjacent transition obtained by composing the two contiguous matrices
of \cref{thm:flat-cmf} is
\begin{equation}\label{eq:Delta-adjacent}
 \mathbf Y_a=\mathsf D_a\mathbf Y_{a+1},
\end{equation}
where
\begin{equation}\label{eq:D-adjacent}
\mathsf D_a=
\begin{pmatrix}
1&-\dfrac{2a+3}{(a+1)(a+2)}&\dfrac1{(a+1)(a+2)}\\[2mm]
1&\dfrac{a^2+a-1}{(a+1)(a+2)}&-\dfrac2{a+2}\\[2mm]
-(2a+1)&\dfrac{5a^2+11a+5}{(a+1)(a+2)}&\dfrac{a-2}{a+2}
\end{pmatrix}.
\end{equation}
For \(a\ge1\) this is the composite contiguous transfer.  The boundary
identity \(\mathbf Y_0=\mathsf D_0\mathbf Y_1\) is checked directly from the
terminating series.

Define
\begin{equation}\label{eq:official-Meijer-function}
 \mathcal H_a^G(z)=
 G_{2,3}^{3,2}\!\left(
 z\,\middle|\begin{matrix}1,1\\a+1,a+1,a+1\end{matrix}\right)
\end{equation}
with the Mellin--Barnes normalization of
\cite[\S5.2]{RamanujanOfficial2026}, and let
\[
 \mathbf S_a^G=
 \bigl(\mathcal H_a^G(1),\theta\mathcal H_a^G(1),
       \theta^2\mathcal H_a^G(1)\bigr)
\]
be the official row state, with the normalization of
\cite[\S5.2]{RamanujanOfficial2026}.  With \(b=a+1\), its transfer matrix is
\begin{equation}\label{eq:official-transfer}
 \mathbf S_{a+1}^G=\mathbf S_a^G\mathsf M_a^G,\qquad
 \mathsf M_a^G=
 \begin{pmatrix}
 0&b^3&b^3(3a+4)\\
 0&-3b^2&-b^2(8a+11)\\
 1&3b&6a^2+12a+5
 \end{pmatrix}.
\end{equation}
Put
\begin{equation}\label{eq:official-normalized-gauge}
 \mathsf P_a^G=b^{-2}\mathsf M_a^G,\qquad
 \mathsf B_a=
 \begin{pmatrix}
 b^2&-(2a+1)&1\\
 0&a/b&-1/b\\
 0&1/b&0
 \end{pmatrix}.
\end{equation}

\begin{theorem}[Exact rational \(SL_3\) coboundary]
\label{thm:ram-gauge}
For every \(a\ge0\),
\begin{equation}\label{eq:ram-determinants}
 \det\mathsf D_a=\det\mathsf B_a=\det\mathsf P_a^G=1,
 \qquad \det\mathsf M_a^G=(a+1)^6,
\end{equation}
and the exact gauge identity is
\begin{equation}\label{eq:ram-coboundary}
 \boxed{\displaystyle
 \mathsf D_a^{-\mathsf T}\mathsf B_{a+1}
 =\mathsf B_a(\mathsf P_a^G)^{-\mathsf T}.}
\end{equation}

If \(\mathbf V_a=\mathbf Y_a^{\mathsf T}\mathsf B_a\) and
\(\Pi_a^G=\mathsf M_0^G\cdots\mathsf M_{a-1}^G\), then
\begin{align}
 \mathbf V_{a+1}&=\mathbf V_a(\mathsf P_a^G)^{-\mathsf T},
 \label{eq:dual-trajectory}\\
 \boxed{\displaystyle
 \mathbf Y_a^{\mathsf T}\mathsf B_a
 =(a!)^2(1,-1,1)(\Pi_a^G)^{-\mathsf T}.}
 \label{eq:dual-trajectory-product}
\end{align}
The denominator coordinate has the precise factorial gauge
\begin{equation}\label{eq:q-dual-coordinate}
 \boxed{\displaystyle
 q_{a-3}=(a!)^2\mathbf V_ae_1.}
\end{equation}
Equivalently, \(\widehat{\mathbf V}_a=(a!)^2\mathbf V_a\) has first
coordinate \(q_{a-3}\) and evolves by the denominator-cleared dual matrix
\begin{equation}\label{eq:denominator-cleared-dual}
 \widehat{\mathbf V}_{a+1}
 =\widehat{\mathbf V}_a\widehat{\mathsf M}_a^G,\qquad
 \widehat{\mathsf M}_a^G=b^4(\mathsf M_a^G)^{-\mathsf T},\qquad
 \det\widehat{\mathsf M}_a^G=b^6.
\end{equation}
\end{theorem}

\begin{proof}
For \(a\ge1\), equation \eqref{eq:Delta-adjacent} is the product
\(M_b(a,a+2;1)M_a(a+1,a+2;1)\) from
\cref{thm:flat-cmf}; multiplication gives \eqref{eq:D-adjacent}.  At
\(a=0\), direct evaluation of the two terminating polynomials gives the
same identity.
The determinants in \eqref{eq:ram-determinants} are direct expansions.
Substitution of \cref{eq:D-adjacent,eq:official-normalized-gauge} into
\eqref{eq:ram-coboundary} reduces its nine entries to rational identities in
\(a\).

Rearranging \eqref{eq:ram-coboundary} gives
\(\mathsf D_a^{\mathsf T}\mathsf B_a
=\mathsf B_{a+1}(\mathsf P_a^G)^{\mathsf T}\), which proves
\eqref{eq:dual-trajectory}.  Since
\(\mathbf Y_0=(1,0,0)^{\mathsf T}\) and
\(\mathbf Y_0^{\mathsf T}\mathsf B_0=(1,-1,1)\), iteration gives
\eqref{eq:dual-trajectory-product}.  Finally,
\(\mathbf V_ae_1=b^2F_a(1)\), whereas
\eqref{eq:q-Fab-link} gives
\(q_{a-3}=(a!)^2b^2F_a(1)\).  This proves
\eqref{eq:q-dual-coordinate}; the last assertions follow by rescaling.
\end{proof}

The gauge fixes the finite connection pairing, not only its projective
direction.

\begin{theorem}[Finite Meijer--\(G\)/\({}_2F_2\) bilinear identity]
\label{thm:ram-bilinear}
For every \(a\ge0\),
\begin{equation}\label{eq:ram-bilinear-matrix}
 \boxed{\displaystyle
 \mathbf S_a^G\mathsf B_a^{\mathsf T}\mathbf Y_a=(a!)^2.}
\end{equation}
Writing \(b=a+1\), \(F=F_a(1)\), and
\(H=\mathcal H_a^G(1)\), this is
\begin{align}
 b^2HF
 &+(\theta H)\left(-(2a+1)F+\frac ab\theta F
                         +\frac1b\theta^2F\right)\notag\\
 &+(\theta^2H)\left(F-\frac1b\theta F\right)=(a!)^2.
 \label{eq:ram-bilinear-expanded}
\end{align}
Thus the normalized conserved primal--dual pairing is exactly one.
\end{theorem}

\begin{proof}
The official normalization gives
\(\mathbf S_0^G(1,-1,1)^{\mathsf T}=1\).  Use
\(\mathbf S_a^G=\mathbf S_0^G\Pi_a^G\) and
\eqref{eq:dual-trajectory-product}; the factors \(\Pi_a^G\) cancel.
Expanding \(\mathsf B_a^{\mathsf T}\mathbf Y_a\) gives
\eqref{eq:ram-bilinear-expanded}.
\end{proof}

The determinant-six phenomenon in \cref{prop:challenge-arithmetic} can now
be located exactly.  If \(\mathsf U(a)\) is the cyclic gauge of the official
solution, then
\begin{equation}\label{eq:cyclic-gauge-det}
 \det\mathsf U(a)=(a+1)^5(a+2)A_a,\qquad \det\mathsf U(0)=136.
\end{equation}
The coboundary transformation
\begin{equation}\label{eq:companion-gauge-definition}
 \mathsf C^{\mathrm{comp}}(a)=
 \mathsf U(a)^{-1}\widehat{\mathsf M}_a^G\mathsf U(a+1)
\end{equation}
gives the scalar companion transfer.  Consequently
\begin{align}
 \det\bigl(\mathsf C^{\mathrm{comp}}(0)\cdots
                 \mathsf C^{\mathrm{comp}}(n)\bigr)
 &=((n+1)!)^6\frac{\det\mathsf U(n+1)}{136}\notag\\
 &=\frac{A_{n+1}}{136}(n+1)!(n+2)!^4(n+3)!,
 \label{eq:casoratian-gauge-origin}
\end{align}
which is the multiplier of \(\mathcal W_{-1}\) in
\eqref{eq:challenge-casoratian}.  The sixth factorial power comes from the
denominator-cleared dual determinant; the shifted factorials and
\(A_{n+1}/136\) are the endpoint cyclic-gauge determinant.

\section{Euler--Beta--log transmutation and the offset ladder}
\label{sec:beta-log}

The adjacent sums are an exact integral transform of the cubic Pilehrood
polynomials, not merely an analogous family.  Define
\begin{align}
 \mathcal Q_m(t)&=\sum_{j=0}^m\binom mj^3j!\,t^j,
 \label{eq:cubic-Q}\\
 \mathcal P_m(t)&=\sum_{j=0}^m\binom mj^3j!
   \bigl(3H_{m-j}-2H_j\bigr)t^j.
 \label{eq:cubic-P}
\end{align}
These are the cubic polynomials occurring in the Euler-constant construction
of Hessami Pilehrood--Hessami Pilehrood
\cite{Pilehrood2013,PilehroodBell2012}.  For \(d\ge1\), put
\begin{equation}\label{eq:Beta-operator}
 \mathcal B_d[f]=\frac1{(d-1)!}\int_0^1(1-t)^{d-1}f(t)\,dt
\end{equation}
and
\begin{equation}\label{eq:transmutation-operator}
 \mathcal T_{m,d}[f](z)=
 \frac{(m+d)!}{m!^2}z^m\mathcal B_d[f(t/z)].
\end{equation}

\begin{theorem}[Exact Beta--log transmutation]
\label{thm:beta-log-transform}
For every \(m\ge0\) and \(d\ge1\), the first identity below holds for
\(z\ne0\) and extends polynomially to \(z=0\); at \(z=1\), the second
identity also holds:
\begin{align}
 F_{m,m+d}(z)&=\mathcal T_{m,d}[\mathcal Q_m](z),
 \label{eq:offset-F-transform}\\
 \mathcal G_{m,d}(1)&=
 \mathcal T_{m,d}[\mathcal P_m+\mathcal Q_m\log t](1),
 \label{eq:offset-G-transform}
\end{align}
where
\begin{equation}\label{eq:Gmd-def}
 \mathcal G_{m,d}(z)=
 \sum_{k=0}^m\binom mk\binom{m+d}k\frac{z^k}{k!}
 \bigl(3H_k-H_{m-k}-H_{m+d-k}\bigr).
\end{equation}
Moreover the transform intertwines Euler derivations:
\begin{equation}\label{eq:Beta-intertwining}
 \theta_z\mathcal T_{m,d}
 =\mathcal T_{m,d}(m-\theta_t).
\end{equation}

For the Challenge specialization \(m=a\), \(d=1\), \(z=1\),
\begin{align}
 q_{a-3}&=(a+1)^3a!\int_0^1\mathcal Q_a(t)\,dt,
 \label{eq:q-Beta-transform}\\
 p_{a-3}&=(a+1)^3a!
 \int_0^1\bigl(\mathcal P_a(t)+\mathcal Q_a(t)\log t\bigr)\,dt+a!.
 \label{eq:p-Beta-transform}
\end{align}
\end{theorem}

\begin{proof}
Put \(j=m-k\).  The beta integrals
\begin{align*}
 \mathcal B_d[t^j]&=\frac{j!}{(j+d)!},\\
 \mathcal B_d[t^j\log t]&=\frac{j!}{(j+d)!}(H_j-H_{j+d})
\end{align*}
turn the coefficient of \(z^{m-j}=z^k\) on the right of
\eqref{eq:offset-F-transform} into
\[
 \frac{m!(m+d)!}{k!^3(m-k)!(m+d-k)!},
\]
which is the coefficient of \(F_{m,m+d}\).  Adding
\(3H_{m-j}-2H_j\) and \(H_j-H_{j+d}\) gives exactly
\(3H_k-H_{m-k}-H_{m+d-k}\), proving
\eqref{eq:offset-G-transform} at \(z=1\).  Differentiating
\(z^mf(t/z)\) gives
\eqref{eq:Beta-intertwining}.  Finally
\cref{eq:q-Beta-transform,eq:p-Beta-transform} follow from
\cref{eq:qn-closed,eq:pn-closed}.
\end{proof}

\begin{corollary}[Fixed-offset Euler ladder]
\label{cor:offset-ladder}
For every fixed \(d\ge1\),
\begin{equation}\label{eq:offset-gamma-limit}
 \boxed{\displaystyle
 \frac{\mathcal G_{m,d}(1)}{F_{m,m+d}(1)}\longrightarrow\gamma.}
\end{equation}
\end{corollary}

\begin{proof}
Use the probability weights
\[
 \Prob(K=k)\propto\binom mk\binom{m+d}k\frac1{k!}.
\]
Their consecutive ratio is
\((m-k)(m+d-k)/(k+1)^3\), so, with
\(s_m=(m(m+d))^{1/3}\), the same ratio argument as in
\cref{thm:challenge-solution} gives \(K/s_m\to1\) exponentially in
probability.  On the window \(|K/s_m-1|\le\eta\),
\[
 3H_K-H_{m-K}-H_{m+d-K}
 =\gamma+\log\frac{K^3}{(m-K)(m+d-K)}+o(1)
 =\gamma+O(\eta)+o(1).
\]
The logarithmic growth on the complement is absorbed by the exponential
tail.  First let \(m\to\infty\), and then let \(\eta\downarrow0\).
Taking expectations proves \eqref{eq:offset-gamma-limit}.
\end{proof}

The multiple-Laguerre background and general approximation families are
already present in the cited literature and in
\cite{WolfsVanAssche2025}.  The assertion here is the exact transmutation
between that cubic system and the 2026 Challenge normalization, together
with its full fixed-offset form.

\section{Rational Gamma jets and contact polynomials}
\label{sec:gamma-jets}

The numerator and denominator are the first two jets of one rational
deformation.  Continue to write \(b=a+1\), and define
\begin{align}
 \mathscr J_a(s)&=a!b!\sum_{k=0}^a
 \frac1{(1-s)_k^3(1+s)_{a-k}(1+s)_{b-k}},
 \label{eq:Jas-def}\\
 \mathscr K_a(s)&=\mathscr J_a(s)+\frac{s}{b\,b!}.
 \label{eq:Kas-def}
\end{align}
The small final term is essential: it records the resonant endpoint
correction \(+a!\) in \eqref{eq:pn-closed}.

\begin{theorem}[Master Gamma-jet transform]
\label{thm:gamma-jet}
For every \(a\ge0\),
\begin{equation}\label{eq:zero-first-jets}
 \boxed{\displaystyle
 q_{a-3}=b\,a!b!\,\mathscr K_a(0),\qquad
 p_{a-3}=b\,a!b!\,\mathscr K_a'(0),}
\end{equation}
and hence
\begin{equation}\label{eq:gamma-first-log-jet}
 \frac{p_{a-3}}{q_{a-3}}=(\log\mathscr K_a)'(0).
\end{equation}
As \(a\to\infty\), locally uniformly on the strip \(|\Re s|<1\),
\begin{equation}\label{eq:Gamma-master-limit}
 \boxed{\displaystyle
 \frac{\mathscr J_a(s)}{\mathscr J_a(0)},\quad
 \frac{\mathscr K_a(s)}{\mathscr K_a(0)}
 \longrightarrow
 \Gamma(1-s)^3\Gamma(1+s)^2.}
\end{equation}
Consequently, for every \(m\ge2\),
\begin{equation}\label{eq:zeta-log-jets}
 \boxed{\displaystyle
 \frac{\left.\dfrac{d^m}{ds^m}\log\mathscr K_a(s)\right|_{s=0}}
 {(3+2(-1)^m)(m-1)!}\longrightarrow\zeta(m),}
\end{equation}
while the first logarithmic jet tends to \(\gamma\).
Every finite-stage logarithmic jet in \eqref{eq:zeta-log-jets} is rational.
\end{theorem}

\begin{proof}
Differentiating a reciprocal Pochhammer symbol at zero gives
\begin{align*}
 \left.\frac d{ds}\log(1-s)_k^{-3}\right|_{s=0}&=3H_k,\\
 \left.\frac d{ds}\log(1+s)_r^{-1}\right|_{s=0}&=-H_r.
\end{align*}
Equations \cref{eq:qn-closed,eq:pn-closed} now give
\eqref{eq:zero-first-jets}; the derivative of \(s/(bb!)\) produces exactly
the endpoint \(a!\).

Normalize the summands of \(\mathscr J_a(0)\) to the probability law
\eqref{eq:challenge-pmf}.  Uniformly for \(s\) in a compact subset of the
strip \(|\Re s|<1\), the ratio of a deformed summand to its value at zero is
 \begin{align}
 &\left(\frac{k!}{(1-s)_k}\right)^3
 \frac{(a-k)!}{(1+s)_{a-k}}
 \frac{(b-k)!}{(1+s)_{b-k}}\notag\\
 &\quad=\Gamma(1-s)^3\Gamma(1+s)^2
 k^{3s}(a-k)^{-s}(b-k)^{-s}(1+o(1)).
 \label{eq:uniform-gamma-ratio}
 \end{align}
Choose a shrinking central window
\(|k/(ab)^{1/3}-1|\le\eta_a\), where
\(\eta_a\downarrow0\) and \(\eta_a^2(ab)^{1/3}\to\infty\).
The concentration estimate \eqref{eq:challenge-concentration} makes its
complement exponentially small.  On the window,
\(k=(ab)^{1/3}(1+o(1))\), \(a-k=a(1+o(1))\), and
\(b-k=b(1+o(1))\); hence the power factor in
\eqref{eq:uniform-gamma-ratio} tends uniformly to one.  Uniform Gamma-ratio
bounds on each compact substrip give at most a fixed polynomial in
\(k+1,a-k+1,b-k+1\), which the exponential tail absorbs.  This proves the
locally uniform limit for \(\mathscr J_a\).  The extra term in
\(\mathscr K_a\) is \(O((bb!)^{-1})\), whereas
\(\mathscr J_a(0)\ge1\), so the same limit holds for \(\mathscr K_a\).

Finally,
\begin{equation}\label{eq:Gamma-log-series}
 \log\{\Gamma(1-s)^3\Gamma(1+s)^2\}
 =\gamma s+\sum_{m\ge2}\frac{3+2(-1)^m}{m}\zeta(m)s^m.
\end{equation}
Local uniformity permits termwise differentiation of logarithms near zero
and proves \eqref{eq:zeta-log-jets}.  Rationality at finite \(a\) follows
because \(\mathscr K_a\in\Q(s)\) and \(\mathscr K_a(0)\ne0\).
\end{proof}

For example, as \(a\to\infty\), the correctly normalized first cases are
\begin{equation}\label{eq:first-zeta-jets}
 \frac{(\log\mathscr K_a)''(0)}5\to\zeta(2),\qquad
 \frac{(\log\mathscr K_a)^{(3)}(0)}2\to\zeta(3),\qquad
 \frac{(\log\mathscr K_a)^{(4)}(0)}{30}\to\zeta(4).
\end{equation}
The raw derivative \(\mathscr K_a''(0)\) cannot replace the second
logarithmic derivative.

\begin{corollary}[Exact parity separation]
\label{cor:gamma-jet-parity}
As \(a\to\infty\), locally uniformly on \(|\Re s|<1\),
\begin{align}
 \frac{\mathscr K_a(s)\mathscr K_a(-s)}{\mathscr K_a(0)^2}
 &\longrightarrow
 \left(\frac{\pi s}{\sin\pi s}\right)^5,
 \label{eq:even-jet-limit}\\
 \frac{\mathscr K_a(s)}{\mathscr K_a(-s)}
 &\longrightarrow\frac{\Gamma(1-s)}{\Gamma(1+s)}.
 \label{eq:odd-jet-limit}
\end{align}
Thus the logarithm of the first quotient generates only even zeta values,
while the logarithm of the second generates \(\gamma\) and the odd zeta
values.
\end{corollary}

\begin{proof}
Multiply and divide the two limits in \eqref{eq:Gamma-master-limit}, and use
\(\Gamma(1-s)\Gamma(1+s)=\pi s/\sin\pi s\).
\end{proof}

The Bell-polynomial zeta tower itself is classical
\cite{PilehroodBell2012}.  The point of
\cref{thm:gamma-jet,cor:gamma-jet-parity} is its exact finite connection to
the individual 2026 Challenge approximants, including their endpoint term.

The same deformation has a finite integer spectral avatar.  Put
\begin{equation}\label{eq:contact-denominator}
 \mathscr D_a(s)=
 \prod_{j=1}^a(s-j)^3(s+j)^2(s+a+1).
\end{equation}

\begin{theorem}[Monic integer contact polynomial]
\label{thm:contact-polynomial}
There is a unique polynomial \(\mathscr N_a\in\Z[s]\) such that
\begin{equation}\label{eq:contact-rational-form}
 \mathscr K_a(s)=
 \frac{\mathscr N_a(s)}{a!(a+1)^2\mathscr D_a(s)}.
\end{equation}
It is monic of degree \(5a+2\), and
\begin{equation}\label{eq:contact-degree}
 \deg\{\mathscr N_a-s\mathscr D_a\}=3a.
\end{equation}
Consequently \(\mathscr N_a\) and \(s\mathscr D_a\) have the same first
\(2a+1\) Newton power sums.  Explicitly, if the zeros of
\(\mathscr N_a\) are counted with multiplicity, then for
\(1\le m\le2a+1\),
\begin{equation}\label{eq:contact-Newton}
 \sum_{\mathscr N_a(\nu)=0}\nu^m
 =3\sum_{j=1}^aj^m+2\sum_{j=1}^a(-j)^m+(-a-1)^m.
\end{equation}
At the opposite end of the polynomial,
\begin{align}
 \mathscr N_a(0)
 &=(-1)^a(a+1)(a!)^4q_{a-3},
 \label{eq:contact-N0}\\
 \mathscr N_a'(0)
 &=(-1)^a(a+1)(a!)^4
 \left[p_{a-3}-q_{a-3}\left(H_a-\frac1{a+1}\right)\right].
 \label{eq:contact-N1}
\end{align}
\end{theorem}

\begin{proof}
Multiplication of \eqref{eq:Kas-def} by
\(a!(a+1)^2\mathscr D_a(s)\) clears every Pochhammer denominator.  Each
resulting quotient is in \(\Z[s]\); the endpoint term becomes exactly
\(s\mathscr D_a(s)\).  The latter is monic of degree \(5a+2\), while the
\(k\)-th summand of the remaining part has degree \(3a-k\).  The \(k=0\)
term has nonzero leading coefficient, proving monicity and
\eqref{eq:contact-degree}.

Two monic degree-\(5a+2\) polynomials whose difference has degree \(3a\)
share their first \(2a+1\) coefficients below the leading term.  Newton's
identities give \eqref{eq:contact-Newton}; the roots of
\(s\mathscr D_a\) are \(0\), three copies of \(1,\ldots,a\), two copies of
\(-1,\ldots,-a\), and \(-a-1\).  Finally evaluate
\eqref{eq:contact-rational-form} and its logarithmic derivative at zero,
use \eqref{eq:zero-first-jets}, and compute
\[
 \frac{\mathscr D_a'(0)}{\mathscr D_a(0)}
 =-H_a+\frac1{a+1}.
\]
This gives \cref{eq:contact-N0,eq:contact-N1}.
\end{proof}

\section{Probabilistic realizations and zero geometry}
\label{sec:ram-probability}

The positive law in \eqref{eq:challenge-pmf} has two complementary
realizations: exponential maxima at finite index and independent Bernoulli
variables through the zeros of its generating polynomial.

For \(m\ge1\), let \(M_m\) be the maximum of \(m\) independent
\(\operatorname{Exp}(1)\) variables, put \(M_0=0\), and use independent
copies whenever superscripts are present.  Then
\begin{equation}\label{eq:maxima-mgf}
 \mathbb Ee^{tM_m}=\prod_{j=1}^m\frac j{j-t}
 =\frac{m!}{(1-t)_m}\qquad(\Re t<1).
\end{equation}

\begin{theorem}[Exact finite maxima model]
\label{thm:ram-probability}
For \(a\ge0\), let \(J_a\) have the law
\begin{equation}\label{eq:Ja-law}
 \Prob(J_a=k)=
 \frac{\binom ak\binom{a+1}k/k!}{F_{a,a+1}(1)}
 \qquad(0\le k\le a),
\end{equation}
and, conditionally on \(J_a\), define
\begin{equation}\label{eq:Za0}
 Z_a^{(0)}=
 M_{J_a}^{(1)}+M_{J_a}^{(2)}+M_{J_a}^{(3)}
 -M_{a-J_a}^{(4)}-M_{a+1-J_a}^{(5)}.
\end{equation}
Then
\begin{equation}\label{eq:exact-maxima-mgf}
 \boxed{\displaystyle
 \mathbb Ee^{tZ_a^{(0)}}
 =\frac{\mathscr J_a(t)}{\mathscr J_a(0)}}
 \qquad(|\Re t|<1).
\end{equation}
If
\begin{equation}\label{eq:Za-correction}
 c_a=\frac{a!}{q_{a-3}},\qquad Z_a=Z_a^{(0)}+c_a,
\end{equation}
then
\begin{equation}\label{eq:finite-mean-approximant}
 \boxed{\displaystyle\mathbb EZ_a=\frac{p_{a-3}}{q_{a-3}}.}
\end{equation}
Moreover, as \(a\to\infty\), with independent standard Gumbel variables
\(G_1,\ldots,G_5\),
\begin{equation}\label{eq:Za-all-moments}
 Z_a\longrightarrow X_0:=G_1+G_2+G_3-G_4-G_5
\end{equation}
in distribution and with convergence of every moment.  Equivalently,
locally uniformly for \(|\Re t|<1\),
\begin{equation}\label{eq:Za-mgf-limit}
 \mathbb Ee^{tZ_a}\longrightarrow
 \Gamma(1-t)^3\Gamma(1+t)^2.
\end{equation}
Thus
\begin{align}
 \kappa_1(Z_a)&\longrightarrow\gamma,\notag\\
 \kappa_m(Z_a)&\longrightarrow
 (m-1)!\{3+2(-1)^m\}\zeta(m)\qquad(m\ge2).
 \label{eq:Za-cumulants}
\end{align}
\end{theorem}

\begin{proof}
Multiply the weight in \eqref{eq:Ja-law} by the five conditional moment
generating functions from \eqref{eq:maxima-mgf}.  All ordinary factorials
cancel, leaving exactly the summand of \(\mathscr J_a(t)\); normalization at
zero proves \eqref{eq:exact-maxima-mgf}.  Taking the first derivative and
using \eqref{eq:zero-first-jets} gives
\eqref{eq:finite-mean-approximant}.  Since \(c_a\to0\),
\cref{thm:gamma-jet} gives \eqref{eq:Za-mgf-limit} on a neighborhood of the
origin.  Uniform exponential moments on every smaller strip imply
distributional and all-moment convergence.  Finally differentiate the
logarithm and use \eqref{eq:Gamma-log-series}.
\end{proof}

The distinction between the two deformations is important:
\begin{equation}\label{eq:mgf-not-K}
 \mathbb Ee^{tZ_a}
 =e^{c_at}\frac{\mathscr J_a(t)}{\mathscr J_a(0)}
 \ne\frac{\mathscr K_a(t)}{\mathscr K_a(0)}
\end{equation}
at finite \(a\).  The function \(\mathscr K_a\) is the exact rational
zero/first-jet carrier; \(\mathscr J_a\) is the exact probability carrier.
They have the same Gamma limit.

We next identify the full official Meijer state.  Let
\begin{equation}\label{eq:Xn-def}
 A_{n,i}\sim\Gamma(n+1,1)\quad(1\le i\le3),\qquad
 E_4,E_5\sim\operatorname{Exp}(1)
\end{equation}
be independent, and put
\begin{equation}\label{eq:Xn-ratio}
 X_n=\log\frac{E_4E_5}{A_{n,1}A_{n,2}A_{n,3}}.
\end{equation}

\begin{theorem}[Stop-loss interpretation of the Meijer state]
\label{thm:meijer-probability}
For the official normalization in \cref{sec:meijer-duality},
\begin{equation}\label{eq:Meijer-stop-loss}
 \boxed{\displaystyle
 \mathcal H_n^G(e^{-x})=(n!)^3\mathbb E(X_n-x)_+}
 \qquad(x\in\R).
\end{equation}
Consequently
\begin{equation}\label{eq:Meijer-state-probability}
 \boxed{\displaystyle
 \frac{\bigl(\mathcal H_n^G(1),\theta\mathcal H_n^G(1),
 \theta^2\mathcal H_n^G(1)\bigr)}{(n!)^3}
 =\bigl(\mathbb E(X_n)_+,\Prob(X_n>0),f_{X_n}(0)\bigr).}
\end{equation}
The functional shift is
\begin{equation}\label{eq:Meijer-functional-shift}
 \mathcal H_{n+1}^G(z)=z\theta^2\mathcal H_n^G(z),
\end{equation}
and hence
\begin{equation}\label{eq:probability-conservation}
 f_{X_n}(x)=e^x(n+1)^3\mathbb E(X_{n+1}-x)_+.
\end{equation}

At \(n=0\),
\begin{equation}\label{eq:h-constant}
 \mathfrak h:=\mathbb E(X_0)_+
 =\int_0^\infty e^{-x}\frac{\log x}{x-1}\,dx
 =1.4322057346532244148\ldots,
\end{equation}
whereas
\begin{equation}\label{eq:Powers-gamma}
 \mathbb EX_0=\Prob(X_0>0)=\gamma.
\end{equation}
In particular \(\mathfrak h\ne\gamma\).
\end{theorem}

\begin{proof}
The Mellin transform of \(X_n\) is
\begin{equation}\label{eq:Xn-mgf}
 \mathbb Ee^{sX_n}
 =\frac{\Gamma(1+s)^2\Gamma(n+1-s)^3}{(n!)^3},
 \qquad -1<\Re s<n+1.
\end{equation}
Perron's inverse Mellin formula on a line \(0<c<n+1\) gives
\eqref{eq:Meijer-stop-loss}; differentiation in \(x\), using
\(\theta=-d/dx\) at \(z=e^{-x}\), gives
\eqref{eq:Meijer-state-probability}.  Shifting the three gamma factors in
the Mellin--Barnes integral proves \eqref{eq:Meijer-functional-shift}, and
comparison with \eqref{eq:Meijer-stop-loss} gives
\eqref{eq:probability-conservation}.  At \(n=0\), substitute the five
exponential densities into \(\mathbb E(X_0)_+\), integrate four variables
first, and apply Frullani's identity; the remaining integral is
\[
 \int_0^\infty e^{-x}\frac{\log x}{x-1}\,dx,
\]
which proves \eqref{eq:h-constant}.
The probability identity in \eqref{eq:Powers-gamma} is due to Powers
\cite{Powers2023}; the expectation identity is immediate from
\(\mathbb EG_j=\gamma\).
\end{proof}

The Bernoulli realization comes from real-rootedness.  For integers
\(b\ge a\ge1\), write the zeros of \(F_{a,b}\), with multiplicities, as
\(-\rho_{a,b,1},\ldots,-\rho_{a,b,a}\).

\begin{theorem}[Poisson--binomial law and zero-density limit]
\label{thm:poisson-binomial}
All zeros of \(F_{a,b}\) are negative real numbers, and
\begin{equation}\label{eq:Poisson-binomial-factor}
 \frac{F_{a,b}(z)}{F_{a,b}(1)}
 =\prod_{j=1}^a(1-p_j+p_jz),
 \qquad p_j=\frac1{1+\rho_{a,b,j}}.
\end{equation}
Thus the law with weights proportional to
\(\binom ak\binom bk/k!\) is exactly a sum of independent Bernoulli
variables.  Its mass sequence is strictly log-concave and ultra-log-concave,
and it is monotone in each of \(a,b\) in the monotone-likelihood-ratio order.

Fix \(d\ge0\), put \(b=a+d\) and \(s_a=(ab)^{1/3}\).  As \(a\to\infty\),
if \(J_a\) has these weights, then
\begin{align}
 \mathbb EJ_a&\sim s_a,&
 \Var(J_a)&\sim\frac{s_a}{3},
 \label{eq:Ja-mean-var}\\
 \frac{J_a-\mathbb EJ_a}{\sqrt{\Var(J_a)}}&\Longrightarrow N(0,1).
 \label{eq:Ja-CLT}
\end{align}
The corresponding local central limit theorem holds, and \(J_a/s_a\)
satisfies a large-deviation principle with speed \(s_a\) and rate
\begin{equation}\label{eq:Ja-rate}
 I(x)=
 \begin{cases}
 3(x\log x-x+1),&x\ge0,\\
 +\infty,&x<0,
 \end{cases}
 \qquad(0\log0:=0).
\end{equation}

Finally, the locally finite zero measures converge vaguely on
\([0,\infty)\):
\begin{equation}\label{eq:zero-density-limit}
 \boxed{\displaystyle
 \frac1{s_a}\sum_{j=1}^a\delta_{\rho_{a,b,j}}
 \Longrightarrow
 \frac{\sqrt3}{2\pi}\rho^{-2/3}\,d\rho.}
\end{equation}
\end{theorem}

\begin{proof}
The polynomial \(L_a(-z)=\sum_k\binom akz^k/k!\) has only negative real
zeros.  The coefficient sequence \(\{\binom bk\}_{k\ge0}\) is a classical
P\'olya--Schur multiplier sequence because its exponential generating
function is \(L_b(-z)\); see \cite{PolyaSchur1914}.  Applying the multiplier
theorem to \(L_a(-z)\)
proves real-rootedness, and positive coefficients make every root negative.
Factorization and normalization at \(z=1\) give
\eqref{eq:Poisson-binomial-factor}.

The consecutive mass ratio
\[
 \frac{w_{k+1}}{w_k}=\frac{(a-k)(b-k)}{(k+1)^3}
\]
is strictly decreasing, proving strict log-concavity.  Ultra-log-concavity
follows because \(w_k/\binom ak=\binom bk/k!\) has decreasing consecutive
ratio.  Finally
\[
 \frac{w_{a+1,b;k}}{w_{a,b;k}}=\frac{a+1}{a+1-k}
\]
is increasing in \(k\), and similarly for \(b\), proving the likelihood
ratio assertion.

A Taylor expansion about the unique real saddle
\(\kappa_a\), determined by
\[
 (a-\kappa_a)(b-\kappa_a)=(\kappa_a+1)^3,
\]
shows that the discrete law has one mode or two adjacent modes neighboring
\(\kappa_a\).  Its logarithmic curvature is
\[
 \frac1{a-\kappa_a}+\frac1{b-\kappa_a}
 +\frac3{\kappa_a+1}\sim\frac3{s_a},
\]
which gives \eqref{eq:Ja-mean-var}.  The Bernoulli representation,
\(\Var(J_a)\to\infty\), and the standard Poisson--binomial local limit
theorem \cite[Chapter~VII]{Petrov1975} give the central and local limit
assertions.  Stirling's
formula at \(k=xs_a\), uniformly on compact \(x\)-intervals, gives
\eqref{eq:Ja-rate} and exponential tightness gives the full large-deviation
principle.

For \(z>0\), tilt the weights by \(z^k\).  Their ratio has the extra factor
\(z\), so \(J_{a,z}/s_a\to z^{1/3}\).  Therefore
\begin{equation}\label{eq:Stieltjes-transform-limit}
 \frac1{s_a}\frac{F'_{a,b}(z)}{F_{a,b}(z)}
 =\frac{\mathbb E_zJ_{a,z}}{s_az}\longrightarrow z^{-2/3}.
\end{equation}
The left side is the Stieltjes transform of the measure on the left of
\eqref{eq:zero-density-limit}, whereas
\[
 \int_0^\infty\frac1{\rho+z}
 \frac{\sqrt3}{2\pi}\rho^{-2/3}\,d\rho=z^{-2/3}.
\]
Uniqueness and the standard continuity theorem for Stieltjes transforms
\cite[Chapter~I]{Akhiezer1965} give the vague limit.
\end{proof}

No simplicity or strict-interlacing assertion is needed above.  Those
stronger statements require a separate strict multiplier or
multiple-orthogonality input and are not claimed here.

\begin{remark}[Priority boundary for the Ramanujan package]
The official solution already proves Question~2 through a rank-three
Meijer--\(G\) trajectory, a Mellin--Barnes recessive estimate, a rational
companion gauge, and inverse-transpose Miller selection
\cite{RamanujanOfficial2026}.  The cubic sums and their Bell--zeta
approximants occur in \cite{PilehroodBell2012,Pilehrood2013}, while the
five-exponential ratio law, its signed-Gumbel logarithm, and the associated
Gamma-product transform occur in \cite{Powers2023}.  The general
multiple-orthogonal and multiple-Laguerre background is developed in
\cite{VanAsscheWolfs2023,WolfsVanAssche2025}.  The contribution of the
present sections is the explicit finite-level transmutation among the
Challenge recurrence, its terminating \({}_2F_2\) dual, the official
Meijer--\(G\) trajectory, and the cubic Bell family, together with the exact
finite probability carrier and the all-order logarithmic-jet limits in the
Challenge normalization.
\end{remark}

\section{The regularized Gamma determinant}
\label{sec:gamma-determinant}

The positive trace-class operator of \cref{sec:spectral} encodes the single
Euler trace \(S_2\).  A different, signed Hilbert--Schmidt operator encodes
the complete limiting Gamma jet.  On
\(\ell^2(\N)^{\oplus5}\), define
\begin{equation}\label{eq:A-Gamma}
 \mathsf A_\Gamma=\diag\left(
 \frac1n,\frac1n,\frac1n,-\frac1n,-\frac1n:n\ge1\right).
\end{equation}
Let \(\mathsf A_{\Gamma,N}\) be its compression to the first \(N\)
coordinates in each of the five summands.  We use the standard
Hilbert--Schmidt determinant
\[
 \det{}_2(I-T)=\det\bigl((I-T)e^T\bigr).
\]

\begin{theorem}[Regularized determinant and zeta traces]
\label{thm:gamma-det2}
One has
\(\mathsf A_\Gamma\in\mathcal S_2\setminus\mathcal S_1\), and
\begin{equation}\label{eq:FP-trace-gamma}
 \boxed{\displaystyle
 \FPTr\mathsf A_\Gamma
 :=\lim_{N\to\infty}
 \left(\Tr\mathsf A_{\Gamma,N}-\log N\right)=\gamma.}
\end{equation}
For every \(m\ge2\),
\begin{equation}\label{eq:Gamma-operator-traces}
 \boxed{\displaystyle
 \Tr\mathsf A_\Gamma^m
 =\{3+2(-1)^m\}\zeta(m).}
\end{equation}
Equivalently,
\begin{equation}\label{eq:Gamma-det2}
 \boxed{\displaystyle
 \Gamma(1-t)^3\Gamma(1+t)^2
 =e^{\gamma t}\det{}_2(I-t\mathsf A_\Gamma)^{-1}.}
\end{equation}
In particular
\begin{equation}\label{eq:Gamma-even-odd-traces}
 \Tr\mathsf A_\Gamma^{2r}=5\zeta(2r),\qquad
 \Tr\mathsf A_\Gamma^{2r+1}=\zeta(2r+1)\qquad(r\ge1).
\end{equation}
\end{theorem}

\begin{proof}
Square summability and failure of absolute summability are immediate from
the spectrum.  The truncated trace is \(H_N\), so
\eqref{eq:FP-trace-gamma} is Euler's definition.  Taking powers of the five
eigenvalue strings gives \eqref{eq:Gamma-operator-traces}.  Finally, for a
Hilbert--Schmidt operator,
\[
 \log\det{}_2(I-t\mathsf A_\Gamma)
 =-\sum_{m\ge2}\frac{t^m}{m}\Tr\mathsf A_\Gamma^m.
\]
Compare with \eqref{eq:Gamma-log-series}; this first proves
\eqref{eq:Gamma-det2} for \(|t|<1\), and meromorphic continuation gives the
identity on \(\C\).
\end{proof}

This regularized determinant is an explicit diagonal spectral realization of the
Gamma product, but it does not furnish a new irrationality result for odd
zeta values.  It must also not be confused with the ordinary Fredholm
determinant \(\Delta_2\) of the positive trace-class operator in
\cref{thm:Delta2}.

\section{Cross-place Euler--Pad\'e and continued-fraction convergence}
\label{sec:cross-place-pade}

The Pad\'e system for Euler's factorial series gives a separate arithmetic
transfer.  Put
\begin{equation}\label{eq:Euler-Pade-Q}
 Q_n(T)=\sum_{h=0}^nh!\binom nh^2(-T)^h,
\end{equation}
and let \(P_n(T)\), of degree at most \(n-1\), be the Pad\'e numerator
characterized by
\begin{equation}\label{eq:Euler-Pade-condition}
 Q_n(T)\sum_{m\ge0}m!T^m-P_n(T)=O(T^{2n})
\end{equation}
as a formal series.  For a prime \(p\), the series
\begin{equation}\label{eq:p-adic-Euler-series}
 E_p(t)=\sum_{m\ge0}m!t^m
\end{equation}
converges for \(|t|_p\le1\).  At a real negative argument put
\begin{equation}\label{eq:real-Euler-function}
 E_\infty(t)=\int_0^\infty\frac{e^{-x}}{1-tx}\,dx.
\end{equation}
The standard Pad\'e remainder gives, for \(t\in\Z_p\),
\begin{equation}\label{eq:p-adic-remainder-bound}
 v_p\!\left(E_p(t)-\frac{P_n(t)}{Q_n(t)}\right)
 \ge2v_p(n!)-v_p(Q_n(t)),
\end{equation}
whenever \(Q_n(t)\ne0\); see
\cite{ErnvallMatalaSeppala2022}.

\begin{theorem}[Unit-boundary residue criterion]
\label{thm:p-adic-unit-criterion}
Let \(p\) be prime and \(t\in\Z_p^\times\).  For
\(0\le r<p\), define
\begin{equation}\label{eq:q-rp}
 q_{r,p}(X)=\sum_{h=0}^rh!\binom rh^2(-X)^h\in\mathbb F_p[X].
\end{equation}
If
\begin{equation}\label{eq:good-residue-condition}
 q_{r,p}(\bar t)\ne0\qquad(0\le r<p),
\end{equation}
then every \(Q_n(t)\) is a \(p\)-adic unit and
\begin{equation}\label{eq:full-p-adic-convergence}
 \boxed{\displaystyle
 \frac{P_n(t)}{Q_n(t)}\longrightarrow E_p(t).}
\end{equation}
The following are sample good nonzero residue classes:
\begin{equation}\label{eq:good-residue-table}
 \begin{array}{c|c}
 p&t\pmod p\\ \hline
 3&2\\
 5&2,3,4\\
 7&4,5\\
 11&2,3,5,7\\
 13&6,9,11,12
 \end{array}.
\end{equation}
In particular, at \(t=-1\), the full convergent sequence tends to
\(E_p(-1)\) for \(p=2,3,5,13\).
\end{theorem}

\begin{proof}
Modulo \(p\), all terms with \(h\ge p\) vanish and, for \(h<p\),
\(\binom nh\bmod p\) depends only on \(n\bmod p\).  Hence
\begin{equation}\label{eq:Q-period-p}
 Q_n(t)\equiv q_{r,p}(\bar t)\pmod p
 \qquad(n\equiv r\pmod p).
\end{equation}
Condition \eqref{eq:good-residue-condition} makes every denominator a unit;
then \eqref{eq:p-adic-remainder-bound} tends to infinity.  The table is a
finite calculation in the indicated residue fields.

For \(p=2,t=-1\), an elementary calculation modulo four gives
\begin{equation}\label{eq:Q-v2}
 v_2(Q_n(-1))=
 \begin{cases}0,&n\text{ even},\\1,&n\text{ odd}.
 \end{cases}
\end{equation}
The same remainder estimate therefore proves full convergence also at two.
For \(p=3,5,13\), the assertion is contained in the table.
\end{proof}

The congruence behind a simultaneous subsequence is stronger than periodicity.

\begin{theorem}[A common real and all-prime subsequence]
\label{thm:common-placewise-subsequence}
For all integers \(n\ge0\) and \(k\ge1\),
\begin{equation}\label{eq:Carlitz-congruence}
 Q_{n+k}(T)\equiv Q_n(T)\pmod{k\Z[T]}.
\end{equation}
Let \(N_j=\lcm(1,\ldots,j)\).  For every fixed negative integer \(t\), the
single rational sequence
\begin{equation}\label{eq:common-rational-sequence}
 x_j(t)=\frac{P_{N_j}(t)}{Q_{N_j}(t)}
\end{equation}
satisfies
\begin{equation}\label{eq:common-placewise-limits}
 \boxed{\displaystyle
 x_j(t)\longrightarrow E_\infty(t)\quad\text{in }\R,
 \qquad
 x_j(t)\longrightarrow E_p(t)\quad\text{in }\Q_p
 \text{ for every fixed }p.}
\end{equation}
At \(t=-1\), the real limit is the Euler--Gompertz constant
\(\delta=eE_1(1)\), while the \(p\)-adic limit is \(E_p(-1)\).
\end{theorem}

\begin{proof}
For each \(h\), the polynomial
\(h!\binom Xh=X(X-1)\cdots(X-h+1)\) belongs to \(\Z[X]\).  Hence
\[
 h!\left\{\binom{n+k}h^2-\binom nh^2\right\}
 =\left\{h!\binom{n+k}h-h!\binom nh\right\}
  \left\{\binom{n+k}h+\binom nh\right\}
\]
is divisible by \(k\), coefficient by coefficient in
\eqref{eq:Euler-Pade-Q}.  This proves \eqref{eq:Carlitz-congruence}.
Taking \(n=0,k=N_j\) gives
\begin{equation}\label{eq:QN-unit-all-p}
 Q_{N_j}(T)\equiv1\pmod{N_j\Z[T]}.
\end{equation}
Thus for every fixed \(p\) and all \(j\ge p\), \(Q_{N_j}(t)\) is a
\(p\)-adic unit, and \eqref{eq:p-adic-remainder-bound} proves the \(p\)-adic
limit.  The real Stieltjes--Pad\'e convergence to
\eqref{eq:real-Euler-function} holds for every \(t<0\), and therefore along
the same subsequence.
\end{proof}

\begin{proposition}[\(p\)-adic denominator interpolation]
\label{prop:p-adic-denominator-interpolation}
For every prime \(p\), the series
\begin{equation}\label{eq:q-p-interpolation}
 q_p(x)=\sum_{h\ge0}h!\binom xh^2\qquad(x\in\Z_p)
\end{equation}
converges uniformly and defines a \(1\)-Lipschitz function satisfying
\(q_p(n)=Q_n(-1)\) for \(n\in\N_0\).  Moreover
\begin{equation}\label{eq:q-p-root-criterion}
 \boxed{\displaystyle
 \sup_{n\ge0}v_p(Q_n(-1))=\infty
 \Longleftrightarrow q_p\text{ has a zero in }\Z_p.}
\end{equation}
In particular, if \(q_p\) is rootless then the full Euler convergent
sequence at \(t=-1\) tends to \(E_p(-1)\).
\end{proposition}

\begin{proof}
Uniform convergence follows from \(v_p(h!)\to\infty\) and
\(\binom xh\in\Z_p\).  For each summand,
\[
 h!\left\{\binom xh^2-\binom yh^2\right\}
 =\left(x^{\underline h}-y^{\underline h}\right)
  \left\{\binom xh+\binom yh\right\},
\]
which is divisible by \(x-y\); passing to the uniform limit gives the
Lipschitz bound.  The interpolation identity follows because
\(\binom nh=0\) for \(h>n\).  If the valuations are unbounded, compactness
of \(\Z_p\) supplies a convergent subsequence of the corresponding integers,
and continuity gives a zero.  Conversely, approximate a zero by
nonnegative integers and use continuity.  If there is no zero, compactness
bounds \(v_p(q_p(x))\), and \eqref{eq:p-adic-remainder-bound} proves the
last assertion.
\end{proof}

The word ``simultaneous'' in \cref{thm:common-placewise-subsequence} is
placewise: no convergence in the restricted adelic product topology is
asserted.  Nor does the theorem prove irrationality of any \(E_p(-1)\).
At the residue \(n=p-1\),
\begin{equation}\label{eq:Kurepa-contact}
 Q_{p-1}(-1)\equiv\sum_{h=0}^{p-1}h!\pmod p,
\end{equation}
so Kurepa's left-factorial problem appears as one boundary congruence, but is
not solved here.  Full convergence at \(t=-1\) for every prime remains open;
the observed growth bound \(v_p(Q_n(-1))=O(\log n)\) is not presently a
theorem.  The existence of a zero of \(q_p\) means only that denominator
valuations are unbounded; it is not equivalent to failure of convergence,
which would require comparison with \(v_p(n!)\).

\part{Positive Spectra, Jacobi Towers, and Gamma Closures}
\section{Order-axis polyexponential total positivity and the Jacobi tower}

\subsection{Standing input}

Let

\[
K=\overline{\mathbf Q}\bigl(e^\beta:\beta\in\overline{\mathbf Q}\bigr),\qquad
\Ein_q(\alpha)=\sum_{n\ge1}\frac{(-1)^{n-1}\alpha^n}{n^q n!}.
\]

The only arithmetic input is {[}001, Corollary 8.4{]}: the values
\(\Ein_q(\alpha)\), jointly for integers \(q\ge1\) and nonzero algebraic
\(\alpha\), are algebraically independent over \(K\).

\subsection{1. Master strict-total-positivity theorem}

For real \(\alpha,q>0\), set

\[
\mathcal E(\alpha,q)=\Gamma(q)\Ein_q(\alpha).
\]

Then

\[
\mathcal E(\alpha,q)
=\int_0^\infty y^{q-1}\bigl(1-e^{-\alpha e^{-y}}\bigr)\,dy. \tag{TP-1}
\]

\subsubsection{Theorem}

The kernel \((\alpha,q)\mapsto\mathcal E(\alpha,q)\) is strictly totally
positive of infinite order on \((0,\infty)^2\). Thus, for every

\[
0<\alpha_1<\cdots<\alpha_r,\qquad 0<q_1<\cdots<q_r,
\]

one has

\[
\det[\mathcal E(\alpha_i,q_j)]_{i,j=1}^r>0. \tag{TP-2}
\]

\subsubsection{Composition proof, including orientation and strictness}

Put

\[
B(\alpha,y)=1-e^{-\alpha e^{-y}}
=e^{-y}\int_0^\alpha e^{-u e^{-y}}\,du
=\int_0^\infty \mathbf1_{u<\alpha}\,e^{-y}e^{-u e^{-y}}\,du. \tag{TP-3}
\]

For increasing \(\alpha_i\) and increasing \(u_j\), the Volterra kernel
\(\mathbf1_{u<\alpha}\) is totally nonnegative. For increasing \(u_i\) and
increasing \(y_j\), the numbers \(x_j=e^{-y_j}\) are decreasing, and
\([e^{-u_i x_j}]\) has positive determinant: this is the exponential kernel
\(e^{uv}\) with both \(u_i\) and \(v_j=-x_j\) increasing. Positive column
factors \(e^{-y_j}\) do not change signs. Continuous Cauchy--Binet applied
to (3) gives total nonnegativity of \(B\). It is strict because the cell

\[
0<u_1<\alpha_1<u_2<\alpha_2<\cdots<u_r<\alpha_r
\]

has positive measure; on it the Volterra determinant equals one and the
exponential determinant is positive. Hence \(B(\alpha,y)\) is STP\(_\infty\)
for \(\alpha,y\) in increasing order.

The kernel

\[
C(q,y)=y^{q-1}=e^{(q-1)\log y}
\]

is STP\(_\infty\) for increasing \(q,y>0\). A second continuous
Cauchy--Binet/Andreief composition applied to (1) proves (2). All integrals
are finite because the integrand is bounded near zero and is
\(O(y^{q-1}e^{-y})\) at infinity.

In particular, every finite ordered family
\(\alpha\mapsto\mathcal E(\alpha,q_j)\) is a strict Chebyshev system:
every nonzero linear combination of \(r\) such functions has at most
\(r-1\) distinct zeros. The transform is also variation diminishing in the
standard strict-total-positivity sense. Separately,
for every fixed \(q>0\), \(\alpha\mapsto\mathcal E(\alpha,q)\) is a
Bernstein function, since

\[
(-1)^{m-1}\partial_\alpha^m\mathcal E(\alpha,q)>0\qquad(m\ge1).
\]

\subsection{2. Arithmetic-generic totally positive matrices and positive geometry}

Take algebraic \(0<\alpha_1<\cdots<\alpha_r\) and positive integers
\(q_1<\cdots<q_s\), and put

\[
M_{ij}=\Gamma(q_j)\Ein_{q_j}(\alpha_i).
\]

Then every minor of \(M\) is positive, while its \(rs\) entries are jointly
algebraically independent over \(K\). Hence \(M\) is an explicit
\(K\)-generic point of the open totally positive matrix cone.

Consequences:

\begin{enumerate}
\def\labelenumi{\arabic{enumi}.}
\tightlist
\item
  Every minor is a positive number transcendental over \(K\).
\item
  For \(r<s\), the maximal minors have exactly the classical Pluecker
  relations and no accidental relation; their field has transcendence
  degree \(r(s-r)+1\). After projectivization, the row span is a generic
  point of the positive Grassmannian and has transcendence degree
  \(r(s-r)\).
\item
  For a square \(n\)-by-\(n\) sample, the raw nested Pluecker coordinates
  first form the corresponding multicone.  After separate projectivization,
  the nested first-row spans give a generic point of the totally positive
  full flag variety.
  After projectivizing separately in every rank, its transcendence degree
  is \(n(n-1)/2\), and its relation ideal is the usual flag
  Pluecker/incidence ideal.
\item
  For any infinite increasing algebraic sequence \(\alpha_i\) and increasing
  integer sequence \(q_i\), the leading minors
  \(D_n=\det[M_{ij}]_{1\le i,j\le n}\) are jointly algebraically independent
  over \(K\). Expand \(D_n\) in its private variable \(M_{nn}\); its
  coefficient is \(D_{n-1}\ne0\), while earlier \(D_j\) do not contain
  \(M_{nn}\).
\end{enumerate}

\subsubsection{Oscillation and Galois corollary}

For square \(M\), strict total positivity makes \(M\) oscillatory. Its
eigenvalues are positive and simple, and their eigenvectors have the usual
exact sign-variation pattern. The characteristic coefficients are
algebraically independent over \(K\): a universal relation would remain
after specializing a generic matrix to a diagonal matrix and would give a
relation among elementary symmetric functions. Therefore the \(n\)
eigenvalues are jointly algebraically independent over \(K\). Over the field
generated by its characteristic coefficients, the characteristic polynomial
is a generic monic polynomial and has Galois group \(S_n\).

\subsection{3. The order-axis Stieltjes moment problem}

For algebraic \(\alpha>0\), define

\[
\mu_{\alpha,k}=k!\Ein_{k+1}(\alpha)
=\int_0^\infty y^k w_\alpha(y)\,dy,
\qquad w_\alpha(y)=1-e^{-\alpha e^{-y}}. \tag{TP-4}
\]

All \(\mu_{\alpha,k}\), jointly in positive algebraic \(\alpha\) and
\(k\ge0\), are algebraically independent over \(K\).

After normalizing the measure, the change of coordinates
\(m_{\alpha,k}=\mu_{\alpha,k}/\mu_{\alpha,0}\) is birational. Hence
\(\mu_{\alpha,0}\) together with all nonconstant probability moments
\(m_{\alpha,k}\), jointly in \(\alpha\), is algebraically independent.
The triangular moment--cumulant transform gives the same assertion for all
probability cumulants. In particular, the explicit mean and variance

\[
\frac{\Ein_2(\alpha)}{\Ein_1(\alpha)},\qquad
\frac{2\Ein_3(\alpha)}{\Ein_1(\alpha)}
-\left(\frac{\Ein_2(\alpha)}{\Ein_1(\alpha)}\right)^2
\]

belong to a joint algebraically independent family over \(K\).

The Stieltjes problem is determinate. Indeed,

\[
\mu_{\alpha,k}=\alpha k!\left(1-\frac{\alpha}{2^{k+2}}
+O_\alpha(3^{-k})\right), \tag{TP-5}
\]

so \(\sum_k\mu_{\alpha,2k}^{-1/(2k)}=\infty\).

The infinite Hankel matrix \([\mu_{\alpha,i+j}]_{i,j\ge0}\) is strictly
totally positive. Every one of its finite minors is positive and
transcendental over \(K\). If

\[
H_{\alpha,d}=\det[\mu_{\alpha,i+j}]_{0\le i,j<d},
\]

then \(\{H_{\alpha,d}:\alpha\in\overline{\mathbf Q}_{>0},d\ge1\}\) is
jointly algebraically independent over \(K\), since

\[
\partial H_{\alpha,d}/\partial\mu_{\alpha,2d-2}=H_{\alpha,d-1}.
\]

\subsection{4. Jacobi operator, finite inverse spectra, and Galois groups}

Let the monic orthogonal polynomials satisfy

\[
p_{n+1}(y)=(y-a_{\alpha,n})p_n(y)-b_{\alpha,n}p_{n-1}(y)
\quad(n\ge1),\qquad b_{\alpha,n}>0,
\]

with \(p_0=1\) and \(p_1=y-a_{\alpha,0}\).

Finite Stieltjes--Favard elimination is birational. Consequently the full
family

\[
\{\mu_{\alpha,0},a_{\alpha,n},b_{\alpha,n+1}:
\alpha\in\overline{\mathbf Q}_{>0},\ n\ge0\}
\]

is jointly algebraically independent over \(K\).

Carleman determinacy makes the minimal Jacobi matrix essentially
self-adjoint. Its closure is unitarily equivalent to multiplication by \(y\)
on

\[
L^2\!\left([0,\infty),\frac{w_\alpha(y)}{\Ein(\alpha)}dy\right).
\]

It therefore has simple purely absolutely continuous spectrum \([0,\infty)\)
of multiplicity one, with the displayed measure as the spectral measure of
the first basis vector.

For the leading \(N\)-by-\(N\) Jacobi truncation \(J_{\alpha,N}\):

\begin{itemize}
\tightlist
\item
  its \(N\) eigenvalues are positive, simple, and jointly algebraically
  independent over \(K\);
\item
  the spectra of \(J_{\alpha,N-1}\) and \(J_{\alpha,N}\) strictly interlace
  and their combined \(2N-1\) points are jointly algebraically independent;
\item
  the complete spectral tower through size \(N\) has transcendence degree
  exactly \(2N-1\), with the top consecutive pair a transcendence basis up
  to algebraic extension;
\item
  the characteristic polynomial of \(J_{\alpha,N}\), over its coefficient
  field, has generic Galois group \(S_N\); the pair of consecutive
  characteristic polynomials has group \(S_N\times S_{N-1}\) over their
  joint coefficient field.
\end{itemize}

All assertions remain joint for finitely many distinct algebraic \(\alpha\),
with one chosen consecutive pair at each parameter.

If \(\lambda_{N,j}\) and \(\omega_{N,j}>0\) are the Gaussian nodes and
unnormalized quadrature masses, then

\[
\mu_{\alpha,k}=\sum_{j=1}^N\omega_{N,j}\lambda_{N,j}^k
\quad(0\le k\le2N-1).
\]

The \(2N\) numbers \(\lambda_{N,1},\ldots,\lambda_{N,N}, \omega_{N,1},\ldots,\omega_{N,N}\) are jointly algebraically independent
over \(K\). This follows because Gaussian quadrature is birational up to the
finite splitting extension between these data and
\(\mu_0,\ldots,\mu_{2N-1}\).

\subsection{5. Exact probability law and spectral flow in alpha}

The normalized law \(Y_\alpha\) has density

\[
f_\alpha(y)=\frac{1-e^{-\alpha e^{-y}}}{\Ein(\alpha)},\qquad y\ge0,
\]

and exact survival function

\[
\Pr(Y_\alpha>y)=\frac{\Ein(\alpha e^{-y})}{\Ein(\alpha)}. \tag{TP-6}
\]

Its moment generating function for \(\Re s<1\) is

\[
\mathbf E e^{sY_\alpha}
=\frac{\alpha^s\gamma(1-s,\alpha)-(1-e^{-\alpha})}
{s\Ein(\alpha)}
=\frac1{\Ein(\alpha)}\sum_{n\ge1}
\frac{(-1)^{n-1}\alpha^n}{n!(n-s)}. \tag{TP-7}
\]

The meromorphic continuation has a nonzero simple pole at every positive
integer; hence this EGF is not D-finite and \((\mu_{\alpha,k})\) is not
P-recursive (this analytic statement holds for every real \(\alpha>0\)).

The density ratio for \(\alpha_2>\alpha_1\) is strictly increasing in \(y\),
so the family is increasing in monotone-likelihood-ratio order. Moreover

\[
\partial_\alpha\log w_\alpha(y)
=\frac{e^{-y}}{e^{\alpha e^{-y}}-1}
\]

is strictly increasing in \(y\). Markov's parameter theorem for orthogonal
polynomials therefore gives strict spectral flow: every zero of every fixed
\(p_N(\cdot;\alpha)\) is strictly increasing in \(\alpha\).

As \(\alpha\downarrow0\), \(Y_\alpha\) converges with all moments to
\(\mathrm{Exp}(1)\); for fixed \(n\),

\[
a_{\alpha,n}\to2n+1,\qquad b_{\alpha,n}\to n^2,
\]

and finite spectra tend to Laguerre zeros.

As \(\alpha\to\infty\), with \(L=\log\alpha\),

\[
Y_\alpha/L\Longrightarrow\mathrm{Unif}[0,1]
\]

with convergence of every moment. Hence, for fixed \(n\),

\[
\frac{a_{\alpha,n}}L\to\frac12,
\qquad
\frac{b_{\alpha,n}}{L^2}\to\frac{n^2}{4(4n^2-1)},
\]

and the finite spectra divided by \(L\) tend to shifted-Legendre zeros.

The Hankel endpoints are

\[
\frac{H_{\alpha,d}}{\alpha^d}\to\prod_{j=0}^{d-1}(j!)^2
\quad(\alpha\downarrow0),
\]

and

\[
\frac{H_{\alpha,d}}{(\log\alpha)^{d^2}}
\to\det\left[\frac1{i+j+1}\right]_{0\le i,j<d}
=\prod_{j=0}^{d-1}\frac{(j!)^4}{(2j)!(2j+1)!}
\quad(\alpha\to\infty).
\]

\subsection{6. Large-degree law and determinant free energy}

For fixed \(\alpha>0\), concavity of \(1-e^{-\alpha x}\) gives

\[
(1-e^{-\alpha})e^{-y}\le w_\alpha(y)\le\alpha e^{-y}. \tag{TP-8}
\]

Thus the weight is a bounded positive perturbation of the Laguerre weight.
Its Gram determinant obeys

\[
(1-e^{-\alpha})^d\prod_{j=0}^{d-1}(j!)^2
\le H_{\alpha,d}\le
\alpha^d\prod_{j=0}^{d-1}(j!)^2. \tag{TP-9}
\]

In particular

\[
\lim_{d\to\infty}\frac{H_{\alpha,d}^{1/d^2}}d=e^{-3/2}. \tag{TP-10}
\]

Standard norm-comparison/weighted-potential theory for Laguerre weights
then gives the scaled zero law

\[
\frac1N\sum_{j=1}^N\delta_{\lambda_{N,j}/N}
\Longrightarrow
\frac1{2\pi}\sqrt{\frac{4-x}{x}}\,\mathbf1_{(0,4)}(x)\,dx,
\]

the Marchenko--Pastur law of parameter one.

\subsection{7. Related literature and exact scope}

The exact polyexponential kernel \(\Gamma(q)\Ein_q(\alpha)\), the associated
order-axis Jacobi measure, and the arithmetic-generic positive
Grassmannian/oscillatory-matrix consequences are not supplied by the sources
cited below.  The analytic ingredients are
classical: exponential-kernel total positivity and Karlin's basic composition
formula. Close precedents are:

\begin{itemize}
\tightlist
\item
  D. St.~P. Richards, \emph{Totally positive kernels, Polya frequency functions,
  and generalized hypergeometric series}, LAA 137/138 (1990), 467--478,
  DOI 10.1016/0024-3795(90)90139-4.
\item
  V. Buchstaber and A. Glutsyuk, \emph{Total positivity, Grassmannian and modified
  Bessel functions}, arXiv:1708.02154 (integer-index Bessel TP matrices).
\item
  D. S. P. Salazar, \emph{Strict Total Positivity from Spectral Darboux and Toeplitz
  Smoothing Mechanisms}, arXiv:2607.02778 (2026), which solves the real-order
  Bessel question but contains no Ein/polyexponential result.
\item
  A. Postnikov, \emph{Total positivity, Grassmannians, and networks},
  arXiv:math/0609764.
\end{itemize}

The resulting package is an explicit \textbf{arithmetic-generic STP kernel}
simultaneously realizing generic positive Grassmannians, oscillatory matrices
with generic Galois groups, and a determinate Jacobi spectral flow.  Its
arithmetic conclusions are relative to [001], while the positivity and
spectral deductions use classical composition and moment machinery.
\section{Order-axis Hankel double scaling: exact ramp equilibrium, free energy, and arithmetic phase transition}

\subsection{0. Scope and normalization}

For \(\alpha>0\), put

\[
w_\alpha(y)=1-e^{-\alpha e^{-y}},\qquad
\mu_{\alpha,k}=\int_0^\infty y^k w_\alpha(y)\,dy
=k!\,\Ein_{k+1}(\alpha),
\]

and define the size-\(N\) order-axis Hankel determinant

\[
H_N(\alpha)=\det[\mu_{\alpha,i+j}]_{i,j=0}^{N-1}.
\]

The purpose of this note is to close the simultaneous limit

\[
\frac{\log\alpha_N}{N}\longrightarrow c\in\mathbf R
\]

at the level of the equilibrium measure, zero distribution, determinant
free energy, and the transition from Marchenko--Pastur to arcsine behavior.
The final section records the arithmetic strengthening supplied by {[}001{]}.

Throughout,

\[
\mathcal I_c(\mu)
=\int_0^\infty (x-c)_+\,d\mu(x)
-\iint\log|x-y|\,d\mu(x)d\mu(y),
\]

and \(\mathcal F(c)=\inf_{\mu\in\mathcal P([0,\infty))}\mathcal I_c(\mu)\).

\begin{center}\rule{0.5\linewidth}{0.5pt}\end{center}

\subsection{1. Master theorem}

\subsubsection{Theorem 1 (exact order-axis double-scaling law)}

Let \(\alpha_N>0\) and assume

\[
c_N:=\frac{\log\alpha_N}{N}\longrightarrow c\in\mathbf R.
\]

Then:

\begin{enumerate}
\def\labelenumi{\arabic{enumi}.}
\item
  The effective fields after \(y=Nx\),

  \[
  Q_N(x)=-\frac1N\log w_{\alpha_N}(Nx),
  \]

  converge uniformly on \([0,\infty)\) to

  \[
  V_c(x)=(x-c)_+.
  \]
\item
  The normalized zero-counting measures of the degree-\(N\) monic
  orthogonal polynomials for \(w_{\alpha_N}(y)dy\), after division of their
  zeros by \(N\), converge weakly to the unique minimizer \(\mu_c\) of
  \(\mathcal I_c\).
\item
  The Hankel determinant satisfies

  \[
  \boxed{\ \log H_N(\alpha_N)
  =N^2\log N-N^2\mathcal F(c)+o(N^2).\ }
  \]
\item
  For \(c\le0\),

  \[
  d\mu_c(x)=\frac1{2\pi}\sqrt{\frac{4-x}{x}}\,
  \mathbf1_{(0,4)}(x)\,dx,
  \qquad
  \mathcal F(c)=\frac32-c.
  \]
\item
  For \(c>0\), let \(\phi=\phi(c)\in(0,\pi/2)\) be the unique solution

  \[
  \boxed{\ c=\frac{2\pi\cos^2\phi}
  {\phi+\sin\phi\cos\phi}.\ } \tag{DS-1.1}
  \]

  Set

  \[
  D(\phi)=\phi+\sin\phi\cos\phi,
  \qquad
  b=b(c)=\frac{2\pi}{D(\phi)}.
  \tag{DS-1.2}
  \]

  Thus \(c=b\cos^2\phi\) and \(b-c=b\sin^2\phi\). The support is exactly
  \([0,b]\), and on \(0<x<b\), \(x\ne c\), its density is

  \[
  \boxed{
  \begin{aligned}
  \rho_c(x)
  &=\frac{\phi}{\pi^2}\sqrt{\frac{b-x}{x}}\\
  &\quad+\frac1{2\pi^2}
  \log\left|
  \frac{\sqrt{c(b-x)}+\sqrt{(b-c)x}}
       {\sqrt{c(b-x)}-\sqrt{(b-c)x}}
  \right|.
  \end{aligned}} \tag{DS-1.3}
  \]

  Equivalently,

  \[
  \rho_c(x)=\frac1{2\pi^2}\sqrt{\frac{b-x}{x}}
  \left[
  2\phi+\sqrt{\frac{x}{b-x}}
  \log\left|
  \frac{\cot\phi+\sqrt{x/(b-x)}}
       {\cot\phi-\sqrt{x/(b-x)}}
  \right|
  \right]. \tag{DS-1.4}
  \]
\item
  The mass lying in the linearly confined part of the potential is

  \[
  \boxed{\ p(c):=\mu_c([c,b])
  =\frac{b\phi^2}{\pi^2}
  =\frac{2\phi^2}{\pi D(\phi)}.\ } \tag{DS-1.5}
  \]
\item
  The equilibrium free energy is

  \[
  \boxed{
  \mathcal F(c)=
  \frac32+\log\frac{2D(\phi)}\pi
  -\frac{2\phi\cos^2\phi
  \bigl(2\phi+\sin\phi\cos\phi\bigr)}{D(\phi)^2},
  \qquad c>0.} \tag{DS-1.6}
  \]

  It obeys the envelope identity

  \[
  \mathcal F'(c)=-p(c). \tag{DS-1.7}
  \]
\end{enumerate}

\begin{center}\rule{0.5\linewidth}{0.5pt}\end{center}

\subsection{2. Uniform reduction to the ramp potential}

For every \(t>0\),

\[
(1-e^{-1})\min(1,t)\le1-e^{-t}\le\min(1,t).
\]

Taking \(t=\alpha_Ne^{-Nx}\) gives the uniform bounds

\[
(x-c_N)_+
\le Q_N(x)
\le(x-c_N)_++\frac{-\log(1-e^{-1})}{N}. \tag{DS-2.1}
\]

Since

\[
\sup_{x\ge0}|(x-c_N)_+-(x-c)_+|\le|c_N-c|,
\]

the convergence \(Q_N\to V_c\) is uniform on the whole half-line, not
merely locally away from the kink.

Heine's identity and \(y_i=Nx_i\) give

\[
H_N(\alpha_N)
=\frac{N^{N^2}}{N!}
\int_{[0,\infty)^N}\Delta(x)^2
\prod_{j=1}^N e^{-NQ_N(x_j)}\,dx_j. \tag{DS-2.2}
\]

The logarithmic-gas Laplace principle, together with (2.1), yields

\[
\frac1{N^2}\log\frac{H_N(\alpha_N)}{N^{N^2}}
\longrightarrow-\mathcal F(c). \tag{DS-2.3}
\]

Here is the short triangular-array reduction implicit in this step. Put

\[
\delta_N=\lVert Q_N-V_c\rVert_{L^\infty([0,\infty))}=o(1)
\]

and let \(Z_N[Q]\) denote the \(N\)-particle integral in (2.2), with
\(Q_N\) replaced by \(Q\). Direct comparison of the integrands gives

\[
e^{-N^2\delta_N}Z_N[V_c]
\le Z_N[Q_N]\le
e^{N^2\delta_N}Z_N[V_c]. \tag{DS-2.4}
\]

Thus the partition-function limit for the triangular array follows from
the classical fixed-field limit for \(V_c\). The same reduction applies to
zeros. Write
\(\lVert p\rVert^2_{Q,N}=\int_0^\infty |p(x)|^2e^{-NQ(x)}\,dx\).
If \(p_N^{(N)}\) is the monic degree-\(N\) extremal polynomial for
\(Q_N\), and \(p_N^{(c)}\) is the corresponding extremal polynomial for
\(V_c\), then

\[
\lVert p_N^{(N)}\rVert^2_{V_c,N}
\le e^{2N\delta_N}
\lVert p_N^{(c)}\rVert^2_{V_c,N}. \tag{DS-2.5}
\]

Consequently \(p_N^{(N)}\) is asymptotically extremal for the fixed
admissible field \(V_c\), and the standard weighted zero-distribution
theorem gives the same equilibrium limit. Equations (2.4)--(2.5) isolate
all dependence on the varying sequence \(Q_N\).

For precision, the standard varying-weight result used here is the following
Gonchar--Rakhmanov/Mhaskar--Saff form. If \(Q_N\to Q\) uniformly on a closed
real set, \(Q\) is continuous and admissible
\(\bigl(Q(x)-2\log(1+|x|)\to+\infty\bigr)\), and the reference measure has
positive density almost everywhere, then the normalized zeros of the monic
degree-\(N\) \(L^2(e^{-NQ_N}dx)\)-extremal polynomials converge to the
weighted equilibrium measure for \(Q\). The corresponding log-gas partition
functions obey the energy Laplace principle (2.3). On an unbounded set this
is reduced to the compact theorem by the classical restricted-range
inequality.

All hypotheses are explicit here. Equation (2.1) gives uniform convergence,
and, because \(c_N\) is eventually bounded,

\[
Q_N(x)\ge x-C\qquad(x\ge C+1)
\]

with one constant \(C\). Thus admissibility, exponential tightness, and the
restricted-range reduction are uniform in \(N\). Lebesgue measure has
positive density everywhere. This proves both the zero-counting statement
and (2.3). This is the only place where general weighted-potential theory is
used; all remaining formulas are obtained explicitly below.

For \(c\le0\), \(V_c(x)=x-c\) on \([0,\infty)\). Adding the constant \(-c\)
does not alter the minimizer, so it is the Laguerre/Marchenko--Pastur
equilibrium measure. Its first moment is one and its logarithmic energy gives
\(\mathcal F(0)=3/2\): explicitly, the MP law has
\(\int x\,d\mu_0=1\), \(\int\log x\,d\mu_0=-1\), and evaluating the
Euler--Lagrange constant at the hard edge gives \(\ell=-2\), so
\(\mathcal F(0)=(1-\ell)/2=3/2\). Hence
\(\mathcal F(c)=3/2-c\).

\begin{center}\rule{0.5\linewidth}{0.5pt}\end{center}

\subsection{\texorpdfstring{3. Solution of the ramp equilibrium problem for \(c>0\)}{3. Solution of the ramp equilibrium problem for c\textgreater0}}

\subsubsection{3.1 Resolvent and endpoint}

Suppose the support is \([0,b]\), with \(b>c\). Let

\[
M_c(z)=\int_0^b\frac{d\mu_c(x)}{z-x},
\qquad
R(z)=\sqrt{\frac{z-b}{z}},\qquad R(z)\to1\quad(z\to\infty).
\]

The differentiated Euler--Lagrange condition is

\[
M_{c,+}(x)+M_{c,-}(x)=V_c'(x)
=\mathbf1_{(c,b)}(x). \tag{DS-3.1}
\]

The scalar Riemann--Hilbert solution is

\[
M_c(z)=\frac{R(z)}{2\pi}
\int_c^b\frac{1}{z-s}\sqrt{\frac{s}{b-s}}\,ds. \tag{DS-3.2}
\]

The condition \(M_c(z)=z^{-1}+O(z^{-2})\) is precisely

\[
\int_c^b\sqrt{\frac{s}{b-s}}\,ds=2\pi. \tag{DS-3.3}
\]

Put \(c=b\cos^2\phi\), \(b-c=b\sin^2\phi\). Evaluating (3.3) gives

\[
b\bigl(\phi+\sin\phi\cos\phi\bigr)=2\pi,
\]

which is (1.1)--(1.2). The map

\[
c(\phi)=\frac{2\pi\cos^2\phi}{D(\phi)}
\]

is a decreasing bijection from \((0,\pi/2)\) onto \((0,\infty)\), with
\(c(\phi)\to\infty\) as \(\phi\downarrow0\) and \(c(\phi)\to0\) as
\(\phi\uparrow\pi/2\), since

\[
c'(\phi)
=-\frac{4\pi\cos\phi(\cos\phi+\phi\sin\phi)}{D(\phi)^2}<0. \tag{DS-3.4}
\]

Thus the endpoint is unique.

The integral in (3.2) is elementary. With consistent branches one may also
write

\[
\boxed{
M_c(z)=\frac1\pi\left[
\arctan\left(\tan\phi\sqrt{\frac{z}{z-b}}\right)
-\phi\sqrt{\frac{z-b}{z}}
\right].} \tag{DS-3.5}
\]

\subsubsection{3.2 Density and Frostman verification}

Taking the jump of (3.2) gives

\[
\rho_c(x)=\frac1{2\pi^2}\sqrt{\frac{b-x}{x}}
\operatorname{PV}\!\int_c^b
\frac1{s-x}\sqrt{\frac{s}{b-s}}\,ds. \tag{DS-3.6}
\]

If

\[
a=\sqrt{\frac{x}{b-x}},\qquad u_c=\sqrt{\frac{c}{b-c}}=\cot\phi,
\]

then direct integration yields

\[
\operatorname{PV}\!\int_c^b
\frac1{s-x}\sqrt{\frac{s}{b-s}}\,ds
=2\phi+a\log\left|\frac{u_c+a}{u_c-a}\right|. \tag{DS-3.7}
\]

The right side is strictly positive both for \(a<u_c\) and \(a>u_c\).
Hence (1.3) defines a positive probability density on \((0,b)\).

Equation (3.1) makes the Frostman expression constant on \([0,b]\). For
\(x>b\), the derivative of

\[
V_c(x)-2\int\log|x-y|\,d\mu_c(y)
\]

is \(1-2M_c(x)>0\): the Stieltjes transform decreases from
\(M_c(b+)=1/2\). Thus the required Frostman inequality holds off the support.
This proves that (1.3) is the unique equilibrium measure, rather than only a
formal solution of the singular-integral equation.

The three distinguished local behaviors are

\[
\rho_c(x)\sim\frac{\phi\sqrt b}{\pi^2\sqrt x}
\quad(x\downarrow0), \tag{DS-3.8}
\]

\[
\rho_c(x)=\frac1{2\pi^2}
\log\frac{4c(b-c)}{b|x-c|}+O(1)
\quad(x\to c), \tag{DS-3.9}
\]

and

\[
\rho_c(x)\sim
\frac{\phi+\cot\phi}{\pi^2\sqrt b}\sqrt{b-x}
\quad(x\uparrow b). \tag{DS-3.10}
\]

Thus the hard edge survives, the ramp kink produces an integrable logarithmic
bulk singularity, and the right endpoint is soft.

\begin{center}\rule{0.5\linewidth}{0.5pt}\end{center}

\subsection{4. Tail mass and exact free energy}

Under \(x=b\sin^2\theta\), the kink corresponds to
\(\theta=\pi/2-\phi\). Integrating (1.4) over \([c,b]\) gives

\[
p(c)=\frac{b}{\pi^2}
\int_0^{\tan\phi}
\left[
\frac{2\phi t^2}{(1+t^2)^2}
+\frac{t}{(1+t^2)^2}
\log\frac{\tan\phi+t}{\tan\phi-t}
\right]dt. \tag{DS-4.1}
\]

The two integrals are

\[
\phi(\phi-\sin\phi\cos\phi),
\qquad
\phi\sin\phi\cos\phi,
\]

respectively. Hence (1.5) follows.

Since \(\mu_c\) has no atom at \(c\), the envelope theorem gives

\[
\mathcal F'(c)
=\int\partial_c(x-c)_+\,d\mu_c(x)
=-\mu_c((c,\infty))=-p(c). \tag{DS-4.2}
\]

Using \(c=2\pi\cos^2\phi/D(\phi)\), one checks

\[
\frac{d}{d\phi}\left[
\log D(\phi)
-\frac{2\phi\cos^2\phi
(2\phi+\sin\phi\cos\phi)}{D(\phi)^2}
\right]
=-p(c(\phi))c'(\phi). \tag{DS-4.3}
\]

At \(\phi=\pi/2\), \(c=0\), \(D=\pi/2\), and
\(\mathcal F(0)=3/2\). Integrating (4.3) proves (1.6).

\begin{center}\rule{0.5\linewidth}{0.5pt}\end{center}

\subsection{5. The MP-to-arcsine interpolation and its phase transition}

\subsubsection{\texorpdfstring{5.1 Onset at \(c=0\)}{5.1 Onset at c=0}}

Write \(\phi=\pi/2-\varepsilon\). Then

\[
c=4\varepsilon^2+O(\varepsilon^4),
\qquad
b=4+\frac{2}{3\pi}c^{3/2}+O(c^2), \tag{DS-5.1}
\]

\[
p(c)=1-\frac{2}{\pi}\sqrt c+O(c), \tag{DS-5.2}
\]

and

\[
\boxed{
\mathcal F(c)=\frac32-c+\frac{4}{3\pi}c^{3/2}+O(c^2)
\qquad(c\downarrow0).} \tag{DS-5.3}
\]

For \(c<0\), \(\mathcal F(c)=3/2-c\). Therefore \(\mathcal F\) is
\(C^1\) but not \(C^2\) at \(c=0\); more precisely,

\[
\mathcal F''(c)\sim\frac1{\pi\sqrt c}\qquad(c\downarrow0).
\]

This is a genuine continuous hard-edge onset transition with singular
exponent \(3/2\). Calling it a third-order transition would be inaccurate
under the usual derivative-based convention.

\subsubsection{5.2 Arcsine limit}

As \(c\to\infty\),

\[
\phi=\frac\pi c-\frac{2\pi^3}{3c^3}+O(c^{-5}),
\qquad
b=c+\frac{\pi^2}{c}+O(c^{-3}), \tag{DS-5.4}
\]

\[
p(c)=\frac1c-\frac{\pi^2}{3c^3}+O(c^{-5}), \tag{DS-5.5}
\]

and

\[
\mathcal F(c)=\log\frac4c-\frac{\pi^2}{6c^2}+O(c^{-4}). \tag{DS-5.6}
\]

Let \(\nu_c\) be the pushforward of \(\mu_c\) by \(x\mapsto x/c\). Then

\[
\boxed{
\nu_c\Longrightarrow
\frac{\mathbf1_{(0,1)}(u)}{\pi\sqrt{u(1-u)}}\,du
\qquad(c\to\infty).} \tag{DS-5.7}
\]

Indeed, for fixed \(0<u<1\), (1.3) gives

\[
c\rho_c(cu)\longrightarrow\frac1{\pi\sqrt{u(1-u)}},
\]

while \(b/c\to1\) and the tail mass above the kink is \(p(c)\to0\).
Thus this single explicit equilibrium family starts at the parameter-one
Marchenko--Pastur law and ends, after its natural macroscopic rescaling, at
the arcsine equilibrium law of a compact interval.

\begin{center}\rule{0.5\linewidth}{0.5pt}\end{center}

\subsection{6. Arithmetic strengthening from {[}001{]}}

Let

\[
K=\overline{\mathbf Q}
\bigl(e^\beta:\beta\in\overline{\mathbf Q}\bigr).
\]

The arithmetic input {[}001, Corollary 8.4{]} says that all
\(\Ein_q(\alpha)\), jointly for \(q\ge1\) and nonzero algebraic \(\alpha\),
are algebraically independent over \(K\).

\subsubsection{Theorem 2 (an algebraically independent phase-transition sequence)}

For every sequence \(\alpha_N\in\overline{\mathbf Q}_{>0}\), the countable
family

\[
\{H_N(\alpha_N):N\ge1\}
\]

is algebraically independent over \(K\), in the finite-subset sense.

\paragraph{Proof}

For fixed \(N\), the coordinate

\[
\mu_{\alpha_N,2N-2}=(2N-2)!\Ein_{2N-1}(\alpha_N)
\]

occurs only in the bottom-right entry of the \(N\)-by-\(N\) Hankel matrix.
Expansion in this coordinate gives

\[
H_N(\alpha_N)
=H_{N-1}(\alpha_N)\mu_{\alpha_N,2N-2}+B_N, \tag{DS-6.1}
\]

where \(B_N\) does not involve \(\mu_{\alpha_N,2N-2}\), and
\(H_{N-1}(\alpha_N)\ne0\) by strict positivity. No earlier determinant
\(H_j(\alpha_j)\), \(j<N\), uses the coordinate
\((\alpha_N,2N-1)\). Since all underlying polyexponential coordinates are
jointly algebraically independent, (6.1) proves the assertion inductively.
For completeness, there is no hidden cancellation in this pivot argument.
If a nonzero relation \(R(H_1,\ldots,H_N)=0\) had degree \(d\) in its last
variable, induction would make its leading coefficient
\(R_d(H_1,\ldots,H_{N-1})\) nonzero. After substituting (6.1) and viewing
the result as a polynomial in the new coordinate
\(\mu_{\alpha_N,2N-2}\), its leading coefficient would be

\[
R_d(H_1,\ldots,H_{N-1})H_{N-1}(\alpha_N)^d\ne0.
\]

This contradicts the algebraic independence of that coordinate from all
coordinates occurring earlier. The reasoning also covers repeated values
among the \(\alpha_N\), because the order \(2N-1\) coordinate cannot occur
in a determinant of smaller size.

For every real \(c\), one can choose positive rational \(\alpha_N\) with
\(\log\alpha_N/N\to c\). Combining Theorems 1 and 2 therefore produces,
for every \(c\in\mathbf R\), an explicit sequence of positive, jointly
algebraically independent transcendental Hankel determinants satisfying

\[
\log H_N(\alpha_N)
=N^2\log N-N^2\mathcal F(c)+o(N^2),
\]

with the nonanalytic \(3/2\)-onset at \(c=0\) and the MP-to-arcsine
spectral interpolation. This is the cleanest direct bridge between the
arithmetic theorem of {[}001{]} and a macroscopic random-matrix phase diagram.

\begin{center}\rule{0.5\linewidth}{0.5pt}\end{center}

\subsection{7. Literature boundary and honest level assessment}

The general mechanisms used here are classical:

\begin{itemize}
\tightlist
\item
  logarithmic equilibrium and varying-weight zero asymptotics go back to
  Gonchar--Rakhmanov and weighted-potential theory;
\item
  the Laguerre endpoint \(c\le0\) is the parameter-one
  Marchenko--Pastur law;
\item
  the \(c\to\infty\) endpoint is the classical equilibrium/arcsine law on
  an interval;
\item
  large-\(N\) asymptotics for analytic one-cut varying weights were
  developed by Deift--Kriecherbauer--McLaughlin--Venakides--Zhou.
\end{itemize}

The cited results do not directly treat the exact weight
\(1-e^{-\alpha e^{-y}}\), the exact ramp field \((x-c)_+\) with formulas
(1.1)--(1.7), or its coupling to the polyexponential Hankel tower.  The
remaining local problem is the analysis of the logarithmic kink (3.9): the
mean spacing there is of
order \((N\log N)^{-1}\) in the scaled \(x\)-coordinate, while the original
Gumbel smoothing of the weight occurs on the larger \(N^{-1}\) scale.
A rigorous local kernel, recurrence-coefficient asymptotics, and a uniform
critical expansion across \(c=0\) would constitute a substantially deeper
Riemann--Hilbert problem.

\subsubsection{Primary background located}

\begin{itemize}
\tightlist
\item
  A. A. Gonchar and E. A. Rakhmanov, \href{https://ui.adsabs.harvard.edu/abs/1986SbMat..53..119G}{\emph{Equilibrium measure and the
  distribution of zeros of extremal polynomials}},
  Math. USSR-Sb. 53 (1986), 119--130.
\item
  P. Deift, T. Kriecherbauer, K. T.-R. McLaughlin, S. Venakides and
  X. Zhou, \href{https://doi.org/10.1002/(SICI)1097-0312(199911)52:11\%3C1335::AID-CPA1\%3E3.0.CO;2-1}{\emph{Uniform asymptotics for polynomials orthogonal with respect to
  varying exponential weights and applications to universality questions
  in random matrix theory}},
  Comm. Pure Appl. Math. 52 (1999), 1335--1425.
\item
  A. B. J. Kuijlaars and P. M. J. Tibboel, \emph{The asymptotic behaviour of
  recurrence coefficients for orthogonal polynomials with varying
  exponential weights}, \href{https://arxiv.org/abs/0708.3956}{arXiv:0708.3956}.
\end{itemize}
\section{Fixed-cut incomplete Gamma tail: differential transcendence, zero density, and the Euler-motive barrier}

\subsection{Status and scope}

This section combines the all-order algebraic-point theorem of companion
paper \textbf{[001]} with the universal inverse-trace identities developed
above.  Accordingly, every arithmetic-independence or transcendence statement
in this section has status \RPASS{}.  Analytic identities are unconditional.
No claim about the individual irrationality or transcendence of \(\gamma\)
or \(\delta\), and no comparison-injectivity claim, is made or used.

Put

\[
 K:=\overline{\mathbf Q}\bigl(e^\alpha:\alpha\in\overline{\mathbf Q}\bigr),
 \qquad E:=E_1(1)=\frac{\delta}{e},
 \qquad \eta:=\Ein(1)=\gamma+E .
\]

Write the Laurent expansion and the upper-tail moments as

\[
 \Gamma(s)=\frac1s+\sum_{k\ge0}g_ks^k,
 \qquad
 T_k:=\frac1{k!}\int_1^\infty e^{-u}\frac{(\log u)^k}{u}\,du .
\]

The split at one proved in {[}001{]} is

\[
 c_k:=g_k-T_k=(-1)^{k+1}\Ein_{k+1}(1),                         \tag{GT-1}
\]

and the family \((c_k)_{k\ge0}\) is jointly algebraically independent over \(K\).

\begin{center}\rule{0.5\linewidth}{0.5pt}\end{center}

\subsection{1. Exact finite gamma-jet field and the tail transcendence-degree bound}

\subsubsection{Lemma 1 (exact gamma-jet field)}

For \(N=0\),

\[
 K(g_0)=K(\gamma).
\]

For every \(N\ge1\),

\[
 \boxed{
 K(g_0,\ldots,g_N)
 =K\!\left(\gamma,\pi^2,\zeta(3),\zeta(5),\ldots,
 \zeta\!\left(2\left\lfloor\frac N2\right\rfloor+1\right)\right).}       \tag{GT-2}
\]

Consequently,

\[
 \operatorname{trdeg}_K K(g_0,\ldots,g_N)
 \le 2+\left\lfloor\frac N2\right\rfloor \qquad(N\ge1).                 \tag{GT-3}
\]

\paragraph{Proof}

The standard identity

\[
 \Gamma(1+s)=
 \exp\!\left(-\gamma s+\sum_{j\ge2}\frac{(-1)^j\zeta(j)}j s^j\right)
\]

shows triangularly that \(g_0,\ldots,g_N\) and
\(\gamma,\zeta(2),\ldots,\zeta(N+1)\) generate the same field: taking the
formal logarithm gives the inverse triangular transformation. The even zeta
values all lie in \(\mathbf Q(\pi^2)\), while the odd indices in
\([3,N+1]\) are precisely

\[
 3,5,\ldots,2\left\lfloor\frac N2\right\rfloor+1,
\]

of which there are \(\lfloor N/2\rfloor\). This proves (2) and (3). Notice
that (3) is only an upper bound: no algebraic independence of \(\gamma\),
\(\pi^2\), or odd zeta values over \(K\) is being assumed. \(\square\)

\subsubsection{Theorem 2 (half-rank lower bound for the fixed-cut tail)}

For every \(N\ge1\),

\[
 \boxed{
 \operatorname{trdeg}_K K(T_0,\ldots,T_N)
 \ge \left\lceil\frac N2\right\rceil-1.}                                  \tag{GT-4}
\]

For \(N=0\), the corresponding lower bound is \(0\). In particular,

\[
 \operatorname{trdeg}_K K(T_0,T_1,\ldots)=\infty.                          \tag{GT-5}
\]

\paragraph{Proof}

By (1),

\[
 K(c_0,\ldots,c_N)\subset
 K(g_0,\ldots,g_N,T_0,\ldots,T_N).
\]

The left side has transcendence degree \(N+1\) over \(K\). Subadditivity
and (3) therefore give

\[
\begin{aligned}
N+1
&\le \operatorname{trdeg}_K K(g_0,\ldots,g_N,T_0,\ldots,T_N)\\
&\le 2+\left\lfloor\frac N2\right\rfloor
 +\operatorname{trdeg}_K K(T_0,\ldots,T_N),
\end{aligned}
\]

which is (4). \(\square\)

The bound is collective. It implies that infinitely many tail moments are
transcendental over \(K\), with lower density at least \(1/2\) among the
indices in the weak counting sense supplied by (4), but it does \textbf{not} name
\(T_0=\delta/e\) as one of them.

\begin{center}\rule{0.5\linewidth}{0.5pt}\end{center}

\subsection{2. Mahler's coefficient-field theorem and a fixed-cut DT theorem}

\subsubsection{The exact Mahler input}

The needed result is the formal-power-series theorem in K. Mahler,
``A remark on algebraic differential equations,'' \emph{Rend. Accad. Naz. Lincei},
Ser. VIII \textbf{50} (1971), 402--412, especially §§3--5
(\href{https://carmamaths.org/resources/mahler/docs/176.pdf}{scan}).

\begin{quote}
\textbf{Mahler's theorem (form used here).} Let \(L\) be any field of
characteristic zero and let \(f(z)=\sum_{n\ge0}f_nz^n\in L[[z]]\). If
\(f\) satisfies a nonzero algebraic differential equation
\[
F(z,f,f',\ldots,f^{(m)})=0,
\qquad F\in L[z,Y_0,\ldots,Y_m],
\]
then the coefficient field \(\mathbf Q(f_0,f_1,\ldots)\) is a finitely
generated field extension of \(\mathbf Q\); in particular, it has finite
transcendence degree over \(\mathbf Q\).
\end{quote}

Rational functions in \(z\) are allowed after clearing denominators. The
statement is formal and includes singular formal solutions. Mahler first
constructs a defining equation over the coefficient field. The
Hurwitz--Popken recurrence in §4 has the eventual form

\[
 A(n)f_{n+1}=B_n(f_0,\ldots,f_n),                                           \tag{GT-6}
\]

where \(A\) is a nonzero polynomial and the coefficients used in all
\(B_n\)'s belong to the finite set of coefficients of the defining equation.
Since a nonzero polynomial has only finitely many integer zeros, all
sufficiently late coefficients are rational functions of finitely many early
coefficients and finitely many equation coefficients. Section 5 records the
result as

\[
 \mathbf Q(f_0,f_1,\ldots)=\mathbf Q(s_1,\ldots,s_r,t),
\]

with \(s_1,\ldots,s_r\) algebraically independent and \(t\) algebraic over
\(\mathbf Q(s_1,\ldots,s_r)\).

\subsubsection{Theorem 3 (fixed-cut upper incomplete Gamma is differentially transcendental)}

Let

\[
 U(s):=\Gamma(s,1)=\int_1^\infty e^{-u}u^{s-1}\,du .
\]

Then

\[
 \boxed{U(s)\text{ is differentially transcendental over }\mathbf C(s).}    \tag{GT-7}
\]

Equivalently, no nonzero polynomial in \(s,U,U',\ldots,U^{(m)}\), with
coefficients in \(\mathbf C\), vanishes identically.

\paragraph{Proof}

Uniform domination on compact subsets of the \(s\)-plane permits
differentiation under the integral, so \(U\) is entire and

\[
 U(s)=\sum_{k\ge0}T_ks^k.                                                   \tag{GT-8}
\]

If \(U\) were differentially algebraic over \(\mathbf C(s)\), clearing
denominators and applying Mahler's theorem with \(L=\mathbf C\) would make
\(\mathbf Q(T_0,T_1,\ldots)\) finite-transcendence-degree over
\(\mathbf Q\). But (4) tends to infinity. Algebraic independence over
\(K\) implies algebraic independence over \(\mathbf Q\), so
\(\operatorname{trdeg}_{\mathbf Q}\mathbf Q(T_0,T_1,\ldots)=\infty\), a
contradiction. \(\square\)

This strictly strengthens the safe consequences ``not D-finite'' and ``the
coefficient sequence is not P-recursive.'' It also implies that \(U\) belongs
to no finite-type ordinary Picard--Vessiot extension of \(\mathbf C(s)\),
because every element of such an extension is differentially algebraic over
the base differential field.

\subsubsection{Spectral formal-series corollary}

Let \(m_k\ge2\) be arbitrary and put

\[
 \sigma_k:=P_{m_k}(T_k^{-1}),
\]

where \(P_m\) is the universal inverse-trace polynomial of v2.3. For every
finite index set \(I\),

\[
 \operatorname{trdeg}_K K(\sigma_k:k\in I)
 =\operatorname{trdeg}_K K(T_k:k\in I),                                    \tag{GT-9}
\]

because \(\sigma_k\in K(T_k)\) and \(T_k\) is algebraic over
\(K(\sigma_k)\). Hence the formal series \(\sum_{k\ge0}\sigma_kz^k\) is
also differentially transcendental over \(\mathbf C(z)\). Convergence is
not needed for this statement: Mahler's theorem is formal.

\subsubsection{Exact overlap with Arreche (2013)}

C. E. Arreche, ``A Galois-theoretic proof of the differential transcendence
of the incomplete Gamma function,'' \emph{J. Algebra} \textbf{389} (2013), 119--127
(\href{https://arxiv.org/abs/1304.1917}{arXiv:1304.1917}), proves

\[
 \gamma(t,x)=\int_0^x u^{t-1}e^{-u}\,du
 \quad\text{is }\partial_t\text{-transcendental over }\mathbf C(t,x).       \tag{GT-10}
\]

Its setting is parameterized Picard--Vessiot theory: \(x\) is the main
differential variable, \(t\) is the parameter, and

\[
 \partial_x^2Y-\frac{t-1-x}{x}\partial_xY=0.                               \tag{GT-11}
\]

Arreche's criterion in fact applies to any nonconstant solution of (11), so
it also covers the bivariate upper incomplete Gamma after placing that
solution in the ambient differential extension.

The overlap and non-overlap are therefore precise:

\begin{itemize}
\tightlist
\item
  Both results assert differential transcendence in the \textbf{parameter/order}
  direction.
\item
  Arreche's theorem is unconditional and bivariate over \(\mathbf C(t,x)\),
  obtained from a parameterized Galois group.
\item
  Theorem 3 above first fixes the cut at \(x=1\) and then proves that the
  one-variable entire function \(s\mapsto\Gamma(s,1)\) is DT over
  \(\mathbf C(s)\), using the arithmetic independence theorem of {[}001{]} and
  Mahler's coefficient theorem.
\item
  The bivariate assertion (10) does \textbf{not formally specialize} to the
  fixed-cut assertion (7): differential transcendence need not survive
  specialization of a second variable. Conversely, (7) says nothing about
  the full PPV group in the \(x\)-direction.
\end{itemize}

Thus (7) is not a corollary of the cited Arreche theorem as stated. A final
priority claim would still require a dedicated search of the broader
shift-difference/hypertranscendence literature, because \(U\) also satisfies
\(U(s+1)=sU(s)+e^{-1}\).

\begin{center}\rule{0.5\linewidth}{0.5pt}\end{center}

\subsection{3. A half-density theorem for wholly transcendental tail divisors}

For \(k\ge0\), set

\[
 F_k(z):=\Ein(z)-T_k.
\]

\subsubsection{Lemma 4 (unique real tail root)}

There is a unique \(r_k\in(0,1)\) with

\[
 \Ein(r_k)=T_k,                                                             \tag{GT-12}
\]

and the sequence \((r_k)\) is strictly decreasing.

\paragraph{Proof}

With \(u=e^y\), put

\[
 I_k:=k!T_k=\int_0^\infty e^{-e^y}y^k\,dy.
\]

Integration by parts gives

\[
 (k+1)I_k
 =\int_0^\infty e^y e^{-e^y}y^{k+1}\,dy>I_{k+1},
\]

so \(T_{k+1}<T_k\). Also \(0<T_k\le T_0=E<\eta=\Ein(1)\).
The real function \(\Ein\) is strictly increasing because
\((1-e^{-x})/x>0\) for every real \(x\ne0\). This proves existence,
uniqueness, and strict decrease. \(\square\)

\subsubsection{Lemma 5 (only possible algebraic zero)}

The only possible algebraic zero of \(F_k\) is \(r_k\).

\paragraph{Proof}

There is only one real zero by strict monotonicity. If a nonreal algebraic
zero \(z\) existed, then \(\bar z\ne z\) would also be algebraic and

\[
 \Ein(z)=T_k=\Ein(\bar z).
\]

This is a polynomial relation between two distinct coordinates of the
jointly algebraically independent family in {[}001, Cor. 8.4{]}, a contradiction.
\(\square\)

\subsubsection{Theorem 6 (half-density all-zero theorem)}

Let

\[
 A_N:=\#\{0\le k\le N:r_k\in\overline{\mathbf Q}\}.
\]

Then for every \(N\ge1\),

\[
 \boxed{A_N\le2+\left\lfloor\frac N2\right\rfloor.}                         \tag{GT-13}
\]

Consequently at least

\[
 \boxed{\left\lceil\frac N2\right\rceil-1}                                 \tag{GT-14}
\]

of the divisors \(F_0,\ldots,F_N\) have \textbf{all} their zeros transcendental.
In density notation,

\[
 \liminf_{N\to\infty}
 \frac{\#\{0\le k\le N:\text{every zero of }F_k\text{ is transcendental}\}}
 {N+1}\ge\frac12.                                                           \tag{GT-15}
\]

\paragraph{Proof}

Let \(I\subset\{0,\ldots,N\}\) be the set of indices with algebraic
\(r_k\). The \(2|I|\) resonant coordinates

\[
 \bigl\{\Ein_{k+1}(1),\ \Ein_1(r_k):k\in I\bigr\}
\]

are pairwise distinct and jointly algebraically independent over \(K\) by
{[}001, Cor. 8.4{]}. Via (1) and \(g_k=c_k+T_k\), the invertible blockwise
linear transformation

\[
 (c_k,T_k)\longleftrightarrow(c_k,g_k)
\]

shows that the selected \(g_k\)'s are algebraically independent over
\(K\). Hence

\[
 |I|\le\operatorname{trdeg}_K K(g_0,\ldots,g_N)
 \le2+\left\lfloor\frac N2\right\rfloor,
\]

which proves (13). Lemma 5 then gives (14) and (15). \(\square\)

\subsubsection{A sharp level-zero fork}

The unique real zero at level zero is

\[
 r_0=\Ein^{-1}(E_1(1))\approx0.2321964241289549.
\]

Exactly one of the following holds:

\begin{enumerate}
\def\labelenumi{\arabic{enumi}.}
\tightlist
\item
  every zero of \(\Ein(z)-E_1(1)\) is transcendental;
\item
  \(r_0\) is algebraic, and \(\gamma\) and \(\delta\) are algebraically
  independent over \(K\).
\end{enumerate}

Indeed, in the second case \(\Ein(1)=\eta\) and \(\Ein(r_0)=E\) are jointly
algebraically independent. The invertible linear/scalar change

\[
 (\eta,E)\longleftrightarrow(\gamma=\eta-E,\delta=eE)
\]

proves the assertion. This is a genuine zero criterion, but not an
unconditional irrationality result: the arithmetic nature of \(r_0\) is
unknown.

\subsubsection{An Euler--Gompertz pencil strengthening}

For real algebraic \(t\in[0,1]\), put

\[
 c_t:=(1-t)\gamma+tE,
 \qquad \Ein(q_t)=c_t,\quad q_t\in(0,1).                                    \tag{GT-16}
\]

The elementary bounds \(E<e^{-1}<1/2<\gamma\) show that the \(c_t\)'s,
and hence the \(q_t\)'s, are distinct as \(t\) varies. Each divisor
\(\Ein-c_t\) has at most one algebraic zero, namely \(q_t\), and among the
entire algebraic pencil there is at most one algebraic \(q_t\). For if
\(q_s,q_t\) were algebraic with \(s\ne t\), the three distinct resonant
values \(\eta,c_s,c_t\) would be jointly algebraically independent, whereas

\[
 (t-s)\eta+(1-2t)c_s-(1-2s)c_t=0.                                          \tag{GT-17}
\]

Moreover \(q_{1/2}\) is unconditionally transcendental because
\(2\Ein(q_{1/2})-\Ein(1)=0\). If the one possible exceptional \(q_t\) is
algebraic, then \(t\ne1/2\), and the matrix taking
\((\gamma,E)\) to \((\eta,c_t)\) has determinant \(2t-1\ne0\). Therefore
that exception would again imply algebraic independence of \(\gamma\) and
\(\delta\) over \(K\).

\begin{center}\rule{0.5\linewidth}{0.5pt}\end{center}

\subsection{4. What the inverse traces do and do not add to the Euler motive}

Let \(P_m\in\mathbf Q[X]\) be the monic universal trace polynomial of v2.3,

\[
 \mathcal T_m(c)=P_m(c^{-1}),\qquad m\ge2,
\]

and recall the exact recovery identity

\[
 c^{-1}=144\mathcal T_4(c)-144\mathcal T_2(c)^2-14\mathcal T_2(c).          \tag{GT-18}
\]

\subsubsection{Exact spectral reformulations}

For distinct algebraic \(s,t\in[0,1]\) and arbitrary \(m,n\ge2\),

\[
 \boxed{
 \operatorname{trdeg}_K
 K\bigl(\mathcal T_m(c_s),\mathcal T_n(c_t)\bigr)
 =\operatorname{trdeg}_K K(\gamma,\delta)\in\{1,2\}.}                     \tag{GT-19}
\]

Each level is algebraic over any one of its nonconstant trace coordinates,
and the two distinct linear pencil levels recover \(\gamma,E\). Thus
\(\gamma,\delta\) are algebraically independent over \(K\) if and only if
any such pair of distinct pencil traces is algebraically independent. This
is an exact algebraic-matroid/spectral criterion, not a proof of either side.

For the Euler exponential motive with period matrix

\[
 P_\gamma=\begin{pmatrix}1&\gamma\\0&2\pi i\end{pmatrix},
\]

full comparison injectivity of

\[
 \overline{\mathbf Q}[a,a^{-1},b]\longrightarrow\mathbf C,
 \qquad a\mapsto2\pi i,\quad b\mapsto\gamma,                              \tag{GT-20}
\]

is equivalent to

\[
 \operatorname{trdeg}_{\overline{\mathbf Q}}
 \overline{\mathbf Q}(2\pi i,\gamma)=2.                                   \tag{GT-21}
\]

If \(S_m=P_m(\gamma^{-1})\), (18) gives

\[
 \boxed{
 (20)\text{ injective}
 \iff
 \operatorname{trdeg}_{\overline{\mathbf Q}}
 \overline{\mathbf Q}(2\pi i,S_2,S_4)=2.}                                 \tag{GT-22}
\]

Again, (22) is a birational re-encoding of the same unknown period
statement.

\subsubsection{The exact barriers}

\begin{enumerate}
\def\labelenumi{\arabic{enumi}.}
\item
  \textbf{Global DT does not isolate the constant term.} Equation (7) only says
  that infinitely much arithmetic complexity occurs throughout the tail.
  It is compatible with \(T_0=\delta/e\) being algebraic. Likewise (4)
  can be carried entirely by higher \(T_k\)'s.
\item
  \textbf{The gamma-jet upper field has enough unknown capacity.} The odd
  zeta values and \(\pi^2\) in (2) can absorb roughly half of the finite
  transcendence rank. No argument above forces \(\gamma\) to contribute
  a transcendence degree.
\item
  \textbf{Inverse traces are not automatically motivic periods.} The period
  algebra of a Tannakian category is not generally closed under inversion
  of an arbitrary numerical period coordinate. Formula (18) analytically
  recovers \(\gamma\), but it does not construct \(S_2,S_4\) as matrix
  coefficients of a tensor construction on \(M_\gamma\). Consequently it
  supplies no new comparison injectivity by itself.
\item
  \textbf{Degree-one fullness is already the hard \(\gamma\) problem.} The
  Euler motive is known to be a non-split extension
  \(0\to\mathbf Q(0)\to M_\gamma\to\mathbf Q(-1)\to0\). In this fixed
  rank-two setting, the assertion ``every numerical relation
  \(\lambda\gamma+\beta=0\) is motivic and forces splitting'' is logically
  equivalent to transcendence of \(\gamma\): non-splitting makes fullness
  imply transcendence, while transcendence leaves no such relation to
  lift.
\item
  \textbf{Bare Frobenius does not retain the extension.} The unfiltered
  rank-two Frobenius object splits automatically in the regime proved in
  {[}001{]}. Any successful arithmetic comparison must retain filtered,
  syntomic, or coherent adelic norm data and prove a numerical-to-motivic
  lifting theorem. The new inverse-root or Hankel identities do not supply
  that missing functorial comparison.
\item
  \textbf{Function-level differential Galois rigidity is a different axis.}
  The conclusion that \(U(s)\) lies in no finite Picard--Vessiot extension
  concerns variation in the parameter \(s\). It does not turn a numerical
  relation at \(s=0\) into a morphism of the Euler motive and therefore does
  not imply irrationality of \(\delta\), transcendence of \(\gamma\), or
  period-map injectivity.
\end{enumerate}

\subsubsection{One genuine motivic-size consequence}

No single fixed finite-dimensional Tannakian object with finite-type motivic
Galois group can contain the entire numerical tail tower \((T_k)_{k\ge0}\)
in the field generated by its period-matrix entries: that field has finite
transcendence degree, whereas (5) has infinite transcendence degree. Any
motivic realization of all tail jets must therefore use an increasing tower
or an infinite/pro-object. This is a structural consequence, not a new
statement about \(M_\gamma\)'s rank-two comparison map.

\begin{center}\rule{0.5\linewidth}{0.5pt}\end{center}

\subsection{5. Consolidated classification}

\small
\begin{longtable}[]{@{}
  >{\raggedright\arraybackslash}p{(\columnwidth - 4\tabcolsep) * \real{0.2727}}
  >{\raggedleft\arraybackslash}p{(\columnwidth - 4\tabcolsep) * \real{0.3636}}
  >{\raggedleft\arraybackslash}p{(\columnwidth - 4\tabcolsep) * \real{0.3636}}@{}}
\toprule\noalign{}
\begin{minipage}[b]{\linewidth}\raggedright
Claim
\end{minipage} & \begin{minipage}[b]{\linewidth}\raggedleft
Status from {[}001{]} + v2.3
\end{minipage} & \begin{minipage}[b]{\linewidth}\raggedleft
Arithmetic force on \(\gamma,\delta\)
\end{minipage} \\
\midrule\noalign{}
\endhead
\bottomrule\noalign{}
\endlastfoot
\(\operatorname{trdeg}_K K(T_0,\ldots,T_N)\ge\lceil N/2\rceil-1\) & proved & none individually \\
\(\Gamma(s,1)\) DT over \(\mathbf C(s)\) & proved via Mahler & none individually \\
no finite ordinary PV extension contains \(\Gamma(s,1)\) & proved & none individually \\
lower density \(\ge1/2\) of tail divisors with all zeros transcendental & proved & none individually \\
algebraic \(r_0\Rightarrow \gamma,\delta\) jointly AI over \(K\) & proved conditional fork & very strong if trigger is settled \\
at most one algebraic root in the Euler--Gompertz pencil & proved & an exception forces joint AI \\
two trace coordinates recover an Euler/pencil level & proved & exact criterion only \\
full Euler period-map injectivity & open & equivalent to \(\pi,\gamma\) algebraic independence \\
\(\gamma\) transcendental from motivic non-splitting alone & not proved & needs numerical-to-motivic fullness \\
\(\delta\) irrational/transcendental & not proved & tail rank does not select \(T_0\) \\
\end{longtable}
\normalsize

The strongest transfer package is therefore the combination of (i) a
\textbf{fixed-cut differential-transcendence theorem}, (ii) a
\textbf{half-density all-zero theorem}, and (iii) sharp \textbf{conditional zero/trace
criteria}. None closes Euler's or Gompertz's individual arithmetic problem
without one further trigger: an algebraic exceptional inverse root, or a
numerical-to-motivic comparison theorem.
\section{A finite-rank Gamma-jet spectral sieve across all polyexponential orders}

\begin{notebox}
\textbf{Status and scope.}  The arithmetic statements in this section have
status \RPASS{} and use the joint algebraic-independence theorem of
\textbf{[001]}.  The analytic fiber statements use the all-order inverse-surface
theorem proved above.  No unconditional irrationality or transcendence
statement about Euler's constant is asserted.
\end{notebox}

Put

\[
 K=\overline{\mathbf Q}
   (e^\xi:\xi\in\overline{\mathbf Q})
\]

and assume the arithmetic input

\[
 \bigl(\operatorname{Ein}_q(\beta)\bigr)_
 {(q,\beta)\in\mathbf N_{\geq1}\times
 \overline{\mathbf Q}^{\times}}
 \quad\text{is jointly algebraically independent over }K        \tag{GS-0.1}
\]

in the finite-subset sense; the coordinate labels are the ordered pairs
\((q,\beta)\).

For positive algebraic \(\alpha\) and \(k\geq0\), define

\[
 U_{\alpha,k}=[s^k]\Gamma(s,\alpha)
 =\frac1{k!}\int_\alpha^\infty
 e^{-u}\frac{(\log u)^k}{u}\,du.                              \tag{GS-0.2}
\]

Writing

\[
 \Gamma(s)=s^{-1}+\sum_{k\geq0}g_ks^k,
\]

the lower/upper incomplete-Gamma splitting gives

\[
 U_{\alpha,k}=g_k+R_{\alpha,k},                                \tag{GS-0.3}
\]

where, with \(L_\alpha=\log\alpha\),

\[
 R_{\alpha,k}=-\frac{L_\alpha^{k+1}}{(k+1)!}
 +\sum_{j=0}^k
 \frac{(-1)^jL_\alpha^{k-j}}{(k-j)!}
 \operatorname{Ein}_{j+1}(\alpha).                            \tag{GS-0.4}
\]

\subsection{1. The all-order spectral sieve}

Let \(\Lambda\subset\overline{\mathbf Q}_{\geq1}\) consist of distinct
positive algebraic numbers whose multiplicative span has rank \(r<\infty\).
For \(\alpha\in\Lambda\), \(q\geq1\), and fixed \(k\geq0\), set

\[
 \mathcal D_{\alpha,q,k}
 =\{z\in\mathbf C:
 \operatorname{Ein}_q(z)=U_{\alpha,k}\}.                      \tag{GS-1.1}
\]

\begin{quote}
\textbf{Theorem 1 (finite-rank all-order spectral sieve).} Among all pairs
\((\alpha,q)\in\Lambda\times\mathbf N\), at most \(r+1\) divisors
\(\mathcal D_{\alpha,q,k}\) contain an algebraic point. More precisely:

\begin{enumerate}
\def\labelenumi{\arabic{enumi}.}
\tightlist
\item
  if a divisor contains an algebraic point, that point is its unique real
  point \(\beta\) in \((0,1)\); in particular it contains no nonreal
  algebraic point;
\item
  for any fixed \(\alpha\), at most one outer order \(q\) can give such an
  algebraic point;
\item
  every divisor outside the at most \(r+1\) exceptional pairs is infinite
  and consists wholly of transcendental numbers.
\end{enumerate}
\end{quote}

\subsubsection{Proof}

For \(\alpha\geq1\), every integrand in (0.2) is nonnegative and

\[
 \sum_{k\geq0}U_{\alpha,k}
 =\int_\alpha^\infty e^{-u}\,du=e^{-\alpha}.
\]

Consequently

\[
 0<U_{\alpha,k}<e^{-\alpha}\leq e^{-1}.                       \tag{GS-1.2}
\]

On the other hand, the alternating series gives, for every \(q\geq1\),

\[
 \operatorname{Ein}_q(1)
 >1-\frac1{2^{q+1}}\geq\frac34>e^{-1}.                        \tag{GS-1.3}
\]

The real restriction of \(\operatorname{Ein}_q\) is strictly increasing,
vanishes at zero, and is onto. Therefore the equation
\(\operatorname{Ein}_q(x)=U_{\alpha,k}\) has exactly one real solution,
and (1.2)--(1.3) put it in \((0,1)\).

If a nonreal algebraic \(\beta\) belonged to the divisor, then its distinct
conjugate would also belong to it because \(\operatorname{Ein}_q\) has real
coefficients. This would give

\[
 \operatorname{Ein}_q(\beta)
 =\operatorname{Ein}_q(\overline\beta),
\]

contradicting (0.1). Thus any algebraic point is the unique real point just
described.

For fixed \(k\), differentiation under the integral sign gives

\[
 \frac{\partial}{\partial\alpha}U_{\alpha,k}
 =-\frac{e^{-\alpha}(\log\alpha)^k}{k!\alpha}<0
 \quad(\alpha>1),                                             \tag{GS-1.4}
\]

with strict decrease from \(\alpha=1\) as well. Hence distinct cuts give
distinct levels. Also, if the same cut \(\alpha\) had algebraic inverse
points \(\beta_1,\beta_2\) at two distinct outer orders \(q_1,q_2\), then

\[
 \operatorname{Ein}_{q_1}(\beta_1)
 =U_{\alpha,k}
 =\operatorname{Ein}_{q_2}(\beta_2),
\]

again contradicting (0.1). Thus bad pairs use distinct cuts.

Suppose now that \(s\) distinct bad pairs
\((\alpha_i,q_i)\) exist, with algebraic inverse points
\(\beta_i\in(0,1)\). Let \(X\) be the collection of all source coordinates

\[
 \operatorname{Ein}_j(\alpha_i),
 \qquad 1\leq i\leq s,\quad1\leq j\leq k+1,
\]

and put \(M=|X|\). The new coordinates

\[
 Y_i=\operatorname{Ein}_{q_i}(\beta_i)=U_{\alpha_i,k}
\]

have mutually distinct coordinate labels, and their labels are disjoint
from those of \(X\), because \(\beta_i\in(0,1)\) while
\(\alpha_i\geq1\). Assumption (0.1) therefore
gives

\[
 \operatorname{trdeg}_K K(X,Y_1,\ldots,Y_s)=M+s.              \tag{GS-1.5}
\]

Choose logarithms \(\lambda_1,\ldots,\lambda_{r'}\), with
\(r'\leq r\), spanning the logarithms of the selected cuts over
\(\mathbf Q\). Equations (0.3)--(0.4) show that

\[
 K(X,Y_1,\ldots,Y_s)
 \subset K(X,\lambda_1,\ldots,\lambda_{r'},g_k).
\]

The right side has transcendence degree at most \(M+r'+1\). Comparison
with (1.5) yields

\[
 s\leq r'+1\leq r+1.                                         \tag{GS-1.6}
\]

Finally, the all-order inverse-surface theorem shows that every fiber of
\(\operatorname{Ein}_q\) is infinite. This proves the theorem. \(\square\)

\subsection{2. The rank-one cut tower}

Fix a positive algebraic \(A>1\) and put \(\alpha_n=A^n\), \(n\geq1\).
Its multiplicative rank is one.

\begin{quote}
\textbf{Corollary 2 (two-exception theorem across the whole \((n,q)\)-grid).}
For each fixed \(k\geq0\), among all pairs \((n,q)\in\mathbf N^2\), at
most two divisors
\[
 \operatorname{Ein}_q(z)-U_{A^n,k}
\]
contain any algebraic zero. Every other divisor is infinite and all of
its zeros are transcendental.
\end{quote}

This is stronger than a density-one statement in each row: the total number
of exceptional entries in the entire countable two-dimensional grid is at
most two.

At outer order \(q=1\), every positive real level is regular, because the
critical values of \(\operatorname{Ein}\) are nonreal. Hence, for fixed
\(k\), all but at most two of

\[
 \operatorname{Ein}(z)-U_{A^n,k}
\]

are infinite, simple, pairwise disjoint divisors consisting wholly of
transcendental points. Their inverse monodromy is the full finitary
symmetric group.

The most classical row is \(k=0\):

\[
 U_{A^n,0}=E_1(A^n).
\]

Thus, for example, among every divisor

\[
 \operatorname{Ein}_q(z)-E_1(2^n),
 \qquad n,q\geq1,
\]

at most two can contain an algebraic point.

\subsection{3. A cofinite algebraically independent spectral forest}

The fixed-order upper-jet defect theorem from the Moving-Level manuscript
says that, after deleting at most \(r+1\) cuts, the levels
\(U_{\alpha,k}\) are algebraically independent over \(K\). Delete also the
at most \(r+1\) cuts occurring in Theorem 1. We obtain:

\begin{quote}
\textbf{Corollary 3 (cofinite all-transcendental trace forest).} There is a
subset \(E\subset\Lambda\), \(|E|\leq2r+2\), such that:

\begin{itemize}
\tightlist
\item
  the levels \(\{U_{\alpha,k}:\alpha\in\Lambda\setminus E\}\) are jointly
  algebraically independent over \(K\);
\item
  for every \(\alpha\notin E\) and every \(q\geq1\), the divisor
  \(\mathcal D_{\alpha,q,k}\) consists wholly of transcendental points;
\item
  if one chooses arbitrary orders \(q_\alpha\geq1\) and
  \(m_\alpha\geq2\), then the reciprocal traces
  \[
   T_{q_\alpha,m_\alpha}(U_{\alpha,k})
   =P_{q_\alpha,m_\alpha}(U_{\alpha,k}^{-1})
  \]
  are jointly algebraically independent over \(K\).
\end{itemize}
\end{quote}

For the tower \(A^n\), one may take \(|E|\leq4\). At \(q=1\) this gives a
cofinite family of pairwise disjoint, simple, wholly transcendental spectral
divisors whose arbitrary one-trace transversal is jointly algebraically
independent.

The last assertion follows because every
\(P_{q,m}(X)\in\mathbf Q[X]\) is a nonconstant monic polynomial: adjoining
\(P_{q,m}(U^{-1})\) makes \(U\) algebraic, so separate independent level
coordinates remain independent after the coordinatewise polynomial maps.

\subsection{4. Extremal rigidity}

The defect bound also records exactly what two exceptional algebraic roots
would cost.

\begin{quote}
\textbf{Corollary 4 (two exceptions force Euler--logarithm independence).}
Let \(A>1\) be algebraic. If, at one fixed \(k\), two distinct pairs
\((n_1,q_1)\), \((n_2,q_2)\) have algebraic inverse points, then
\[
 \log A\quad\text{and}\quad g_k
\]
are algebraically independent over \(K\). For \(k=0\), since
\(g_0=-\gamma\), this forces
\[
 \boxed{\ \gamma\text{ and }\log A
 \text{ jointly algebraically independent over }K.\ }
\]
\end{quote}

Indeed, with two bad pairs, equality holds at both ends of the
transcendence-degree comparison:

\[
 M+2
 =\operatorname{trdeg}_K K(X,Y_1,Y_2)
 \leq\operatorname{trdeg}_K K(X,\log A,g_k)
 \leq M+2.
\]

Thus \(\log A,g_k\) are algebraically independent even over \(K(X)\).
This is a conditional fork, not a proof of the transcendence of
\(\gamma\).

\subsection{5. Significance and boundary}

The result is a Diophantine finiteness theorem for algebraic inverse images
of explicit upper incomplete-Gamma jets. Its distinctive point is the
quantifier order: the \(r+1\) bound holds \textbf{simultaneously over every outer
polyexponential order \(q\)}. In rank one, a two-parameter infinite grid
has at most two algebraically contaminated divisors.

The bound is the strongest conclusion available from (0.1) without new
information on the \(r\) logarithmic coordinates and the single Gamma
coefficient \(g_k\). Removing the last \(r+1\) possibilities would require
additional arithmetic input; the extremal rigidity corollary explains
why Euler's constant already appears at \(k=0\).
\section{Arithmetic Transfer Closures from the Companion Independence Theorem}
\label{sec:transfer-closures}

Put
\[
 K=\Qbar\bigl(e^\xi:\xi\in\Qbar\bigr),
 \qquad
 \Ein_q(z)=\sum_{n\ge1}\frac{(-1)^{n-1}z^n}{n^q n!}
 \quad(q\ge1).
\]
Thus \(\Ein_1=\Ein\).  This section combines two inputs already isolated
elsewhere in the monograph:

\begin{enumerate}
\item the companion theorem [001] that the coordinates
\[
 \{\Ein_q(\alpha):q\ge1,\ \alpha\in\Qbar^\times\}
\]
are countably algebraically independent over \(K\); and
\item the universal trace identity
\[
 T_m(c)=P_m(c^{-1})\quad(m\ge2),
 \qquad
 c^{-1}=144T_4(c)-144T_2(c)^2-14T_2(c).
\]
\end{enumerate}

\begin{statusbox}
Every conclusion in this section that uses the first input is marked
\RPASS{} and is relative to [001].  Formal trace identities, positivity,
analytic non-holonomy, and the general \(p\)-adic criteria are marked
\UPASS{}.  Finite residue-field certificates are marked \XPASS{}.  None of
the results below determines the arithmetic nature of an individual
classical connection constant.
\end{statusbox}

\subsection{Algebraic inputs in signed polyexponential levels}

\begin{theorem}[Signed algebraic-input classification (\RPASS)]
\label{thm:transfer-signed-classification}
Let \(\alpha,\beta\in\Qbar^\times\), \(q\ge1\), and
\(a\in\Qbar^\times\).  Then
\[
 \Ein(\beta)=a\Ein_q(\alpha)
 \quad\Longleftrightarrow\quad
 (q,\beta,a)=(1,\alpha,1).
\]
Consequently, every solution of
\(\Ein(z)=a\Ein_q(\alpha)\) is transcendental except for the possible
algebraic solution \(z=\alpha\) in the tautological case \(q=1\), \(a=1\).
\end{theorem}

\begin{proof}
If \((1,\beta)\ne(q,\alpha)\), the displayed equality is a nonzero
\(K\)-linear relation between two distinct coordinates in the [001]
family.  If the coordinates coincide, it becomes
\((1-a)\Ein(\alpha)=0\).  Algebraic independence implies
\(\Ein(\alpha)\ne0\), and hence \(a=1\).  Finally, \(z=0\) cannot solve
the equation because \(a\Ein_q(\alpha)\ne0\); applying the classification
to any nonzero algebraic solution proves the last assertion.
\end{proof}

For example, if \(\alpha>0\) is algebraic and \(q\ge2\), strict
monotonicity and unboundedness of \(\Ein\) on the positive real axis give
a unique \(x_{q,\alpha}>0\) such that
\[
 \Ein(x_{q,\alpha})=\Ein_q(\alpha).
\]
The theorem says that \(x_{q,\alpha}\) is transcendental over \(\Qbar\).

\subsection{Gamma-jet levels and their inverse spectra}

Write
\[
 \Gamma(s)=\frac1s+\sum_{k\ge0}g_ks^k,
 \qquad
 U_k=\frac1{k!}\int_1^\infty
 e^{-u}(\log u)^k\frac{du}{u},
\]
and set
\[
 c_k=g_k-U_k=(-1)^{k+1}\Ein_{k+1}(1).
\]
For the zero multiset
\(Z_k=\{\rho\in\C:\Ein(\rho)=c_k\}\), put
\[
 \tau_{m,k}=\sum_{\rho\in Z_k}\rho^{-m}=P_m(c_k^{-1})
 \qquad(m\ge2).
\]

\begin{theorem}[Gamma-jet spectral divisors (\RPASS)]
\label{thm:transfer-gamma-divisors}
The following statements hold.
\begin{enumerate}
\item Every point of every divisor \(Z_k\), \(k\ge0\), is
transcendental.
\item Each \(\tau_{m,k}\) is transcendental over \(K\).
\item For an arbitrary selection of orders \(m_k\ge2\), the family
\[
 \{\tau_{m_k,k}:k\ge0\}
\]
is countably algebraically independent over \(K\).
\end{enumerate}
\end{theorem}

\begin{proof}
The levels \(c_k\) are nonzero algebraic multiples of distinct members of
the [001] family, and hence are jointly algebraically independent.  For
\(k=0\) the signed level is \(-\Ein_1(1)\); for \(k\ge1\) its order is
\(k+1\ge2\).  In either case
\cref{thm:transfer-signed-classification} excludes an algebraic zero.

If \(P_m(c_k^{-1})\) were algebraic over \(K\), then \(c_k^{-1}\) would
be algebraic over \(K\), because it is a root of a nonconstant polynomial
over the algebraic closure of \(K\).  This is impossible.  For a finite
set of indices, each \(c_k\) is algebraic over
\(K(\tau_{m_k,k})\).  Equality of transcendence degrees therefore transfers
the algebraic independence of the levels to the selected traces.
\end{proof}

\subsection{Positive-integer incomplete-Gamma derivatives}

Let
\[
 \Gamma_{<}(s,1)=\int_0^1e^{-t}t^{s-1}\,dt,
 \qquad
 D_{r,k}=\left.\frac{\partial^k}{\partial s^k}
 \Gamma_{<}(s,1)\right|_{s=r}
 \quad(r\ge1,\ k\ge1).
\]

\begin{theorem}[Fixed-parameter Gamma derivatives (\RPASS)]
\label{thm:transfer-incomplete-gamma-derivatives}
For every fixed positive integer \(r\), the countable family
\[
 \{D_{r,k}:k\ge1\}
\]
is algebraically independent over \(K\).  In particular every one of its
members is transcendental over \(K\).  At \(r=1\),
\[
 D_{1,k}=(-1)^k k!\Ein_k(1).
\]
No assertion of independence is made after combining different values of
\(r\), since the incomplete-Gamma recurrence creates relations between
those rows.
\end{theorem}

\begin{proof}
Expand
\[
 \mathscr R(u)=\frac1u-\Gamma_{<}(u,1)
 =\sum_{j\ge0}R_ju^j,
 \qquad R_j=(-1)^j\Ein_{j+1}(1).
\]
Iterating
\(\Gamma_{<}(s+1,1)=s\Gamma_{<}(s,1)-e^{-1}\) gives a polynomial
\(Q_r\in\Z[u]\) such that
\[
 \Gamma_{<}(u+r,1)
 =\frac{(u)_r}{u}-(u)_r\mathscr R(u)-e^{-1}Q_r(u),
 \qquad
 (u)_r=u(u+1)\cdots(u+r-1).
\]
Comparison of the coefficient of \(u^k\), for \(k\ge1\), yields
\[
 \frac{D_{r,k}}{k!}
 =-(r-1)!R_{k-1}
 +A_{r,k}(R_0,\ldots,R_{k-2}),
\]
where \(A_{r,k}\) is affine over \(K\).  This is an invertible affine
lower-triangular transformation on every finite initial segment.  The [001]
independence of \(R_0,R_1,\ldots\) proves the claim.  The formula at
\(r=1\) follows directly by differentiation under the integral sign.
\end{proof}

\subsection{The order-axis Stieltjes--Hankel tower}

For \(\alpha\in\Qbar_{>0}\) and \(k\ge0\), define
\[
 \mu_{\alpha,k}=k!\Ein_{k+1}(\alpha).
\]
Mellin inversion of \(n^{-k-1}\) gives the two positive representations
\begin{equation}
\label{eq:transfer-order-moments}
 \mu_{\alpha,k}
 =\int_0^1\frac{1-e^{-\alpha t}}{t}(-\log t)^k\,dt
 =\int_0^\infty y^k\bigl(1-e^{-\alpha e^{-y}}\bigr)\,dy.
\end{equation}
Thus \((\mu_{\alpha,k})_{k\ge0}\) is a strict Stieltjes moment sequence.
Put
\[
 H_{\alpha,d}=\det(\mu_{\alpha,i+j})_{0\le i,j<d}
 \qquad(d\ge1).
\]

\begin{theorem}[Order-axis Hankel tower (\UPASS{} positivity;
\RPASS{} independence)]
\label{thm:transfer-order-hankel}
One has
\[
 H_{\alpha,d}>0
 \qquad(\alpha\in\Qbar_{>0},\ d\ge1),
\]
and the complete family
\[
 \{H_{\alpha,d}:\alpha\in\Qbar_{>0},\ d\ge1\}
\]
is algebraically independent over \(K\).  Consequently, if one order
\(m_{\alpha,d}\ge2\) is selected for every pair \((\alpha,d)\), then
\[
 \{T_{m_{\alpha,d}}(H_{\alpha,d}):
   \alpha\in\Qbar_{>0},\ d\ge1\}
\]
is algebraically independent over \(K\).
\end{theorem}

\begin{proof}
The second integral in \eqref{eq:transfer-order-moments} is the moment
functional of a positive density on an interval with infinite support.
Its Gram determinants are therefore strictly positive.

The [001] theorem makes all variables \(\mu_{\alpha,k}\), over all
\(\alpha\) and \(k\), algebraically independent over \(K\).  Within one
\(\alpha\)-block,
\[
 \frac{\partial H_{\alpha,d}}
      {\partial\mu_{\alpha,2d-2}}=H_{\alpha,d-1},
 \qquad H_{\alpha,0}=1.
\]
The Jacobian of \((H_{\alpha,1},\ldots,H_{\alpha,D})\) with respect to
\((\mu_{\alpha,0},\mu_{\alpha,2},\ldots,
\mu_{\alpha,2D-2})\) is triangular with nonzero diagonal.  Distinct
\(\alpha\)-blocks use disjoint variables.  Hence every finite family of
the Hankel determinants is algebraically independent.  The final assertion
follows because each nonconstant trace polynomial leaves its level algebraic
over the generated one-coordinate trace field.
\end{proof}

\subsection{Explicit relative transcendental combinations}

Let
\[
 S_m=T_m(\gamma)=P_m(\gamma^{-1})\qquad(m\ge2)
\]
and let
\[
 X\in\{\delta,E_1(1),\Ei(1),\Ci(1),\Chi(1)\}.
\]
All constants in this display are real.

\begin{theorem}[Complex separation of relative classes (\RPASS)]
\label{thm:transfer-complex-combinations}
For every \(a\in\Qbar\setminus\R\),
\[
 S_m+aX
\]
is transcendental over \(K\).  Moreover \(e\) and \(S_m+aX\) are
algebraically independent over \(\Q\).  In particular, each number
\(S_m+iX\) is an explicit member of this family.
\end{theorem}

\begin{proof}
The field \(K\) is stable under complex conjugation.  If
\(z=S_m+aX\) were algebraic over \(K\), then so would be
\(\overline z=S_m+\overline aX\), and hence
\[
 X=\frac{z-\overline z}{a-\overline a}
 \quad\hbox{and}\quad S_m=z-aX
\]
would both be algebraic over \(K\).  Since
\(S_m=P_m(\gamma^{-1})\) with \(P_m\) nonconstant, this would also make
\(\gamma\) algebraic over \(K\).  The five-relative-class theorem, which
is itself a consequence of [001], excludes the simultaneous algebraicity
of the class of \(\gamma\) and the class containing \(X\).  Thus \(z\) is
transcendental over \(K\).

If \(e\) and \(z\) satisfied a nonzero polynomial over \(\Q\), the
transcendence of \(e\) would make \(z\) algebraic over \(\Q(e)\), hence
over \(K\), a contradiction.
\end{proof}

\begin{remark}
The theorem is a collective separation statement relative to [001].  It
does not imply an irrationality or transcendence result for either summand
separately.
\end{remark}

\subsection{Defect at most one after adjoining two connection constants}

Set
\[
 F=K(\gamma,\delta),
 \qquad d=\trdeg_KF.
\]
Because \(\Ein(1)=\gamma+\delta/e\) is transcendental over \(K\) under
[001], one has \(d\in\{1,2\}\).  Write
\(Y_{\alpha,q}=\Ein_q(\alpha)\) and exclude the anchor
\((\alpha,q)=(1,1)\).

\begin{theorem}[Two-constant defect bound (\RPASS)]
\label{thm:transfer-two-constant-defect}
For every finite set \(I\) of non-anchor indices,
\[
 \trdeg_FF(Y_{\alpha,q}:(\alpha,q)\in I)
 \ge |I|+1-d.
\]
Thus the loss of rank over \(F\) is at most one.  More precisely:
\begin{enumerate}
\item if \(d=1\), all non-anchor coordinates remain algebraically
independent over \(F\);
\item if one non-anchor coordinate is algebraic over \(F\), then \(d=2\)
and all the remaining non-anchor coordinates are algebraically independent
over \(F\).
\end{enumerate}
The same statements hold after replacing each selected coordinate by one
nonconstant inverse trace \(P_m(Y_{\alpha,q}^{-1})\).
\end{theorem}

\begin{proof}
The anchor belongs to \(F\), while [001] gives
\[
 |I|+1
 =\trdeg_KK\bigl(Y_{1,1},Y_{\alpha,q}:(\alpha,q)\in I\bigr)
 \le d+\trdeg_FF(Y_{\alpha,q}:(\alpha,q)\in I).
\]
This proves the inequality and the first refinement.  If a coordinate
\(Y_0\) is algebraic over \(F\), applying the inequality to
\(I\cup\{0\}\), with \(0\notin I\), gives
\[
 \trdeg_FF(Y_{\alpha,q}:(\alpha,q)\in I)\ge |I|+2-d.
\]
The case \(I=\varnothing\) forces \(d=2\), and the general case then gives
full rank \(|I|\).  Finally, a nonconstant one-variable trace coordinate
and its level generate fields differing only by an algebraic extension, so
all transcendence degrees are preserved.
\end{proof}

\subsection{The complete trace relation ideal and spectral rigidity}

The next result is formal and does not use [001].  Let \(L\) be a field of
characteristic zero, let \(M\subset\{2,3,\ldots\}\) be finite with
\(2,4\in M\), and introduce variables \(X_{i,m}\) for
\(1\le i\le n\), \(m\in M\).  Put
\[
 R_i=144X_{i,4}-144X_{i,2}^2-14X_{i,2}
\]
and
\[
 J_M=\left\langle
 X_{i,2}-P_2(R_i),\quad
 X_{i,m}-P_m(R_i):m\in M\setminus\{2,4\},\ 1\le i\le n
 \right\rangle.
\]

\begin{proposition}[Complete universal relation ideal (\UPASS)]
\label{prop:transfer-trace-ideal}
For the homomorphism
\[
 \phi:L[X_{i,m}:1\le i\le n,\ m\in M]
 \longrightarrow L[Z_1,\ldots,Z_n],
 \qquad X_{i,m}\longmapsto P_m(Z_i),
\]
one has \(\ker\phi=J_M\), and \(\phi\) is surjective.  Consequently
\[
 L[X_{i,m}]/J_M\simeq L[Z_1,\ldots,Z_n].
\]
In particular, \(J_M\) is a prime complete-intersection ideal and its
variety is smooth, normal, factorial, and isomorphic to \(\mathbb A_L^n\).
\end{proposition}

\begin{proof}
The recovery identity gives
\[
 144P_4(Z_i)-144P_2(Z_i)^2-14P_2(Z_i)=Z_i,
\]
so every \(Z_i\) lies in the image.  In the quotient by \(J_M\), all
coordinates except \(X_{i,4}\) are polynomials in \(R_i\), while the
definition of \(R_i\) expresses \(X_{i,4}\) as a polynomial in
\(R_i\) and \(X_{i,2}=P_2(R_i)\).  Hence the quotient is generated freely
by \(R_1,\ldots,R_n\), and the induced inverse sends \(Z_i\) to \(R_i\).
This proves the asserted isomorphism.  The geometric properties follow
from the polynomial coordinate ring; the displayed generators have the
expected codimension and form a regular sequence.
\end{proof}

\begin{corollary}[No hidden cross-level relations (\UPASS{} formally;
\RPASS{} for [001] levels)]
\label{cor:transfer-no-hidden-relations}
If the levels \(c_1,\ldots,c_n\) are algebraically independent over \(L\),
then \(J_M\) is the complete ideal of algebraic relations among the numbers
\(T_m(c_i)\).  In particular this applies to every finite family of
distinct [001] coordinates and to the derived independent levels above.
\end{corollary}

\begin{proof}
The elements \(c_i^{-1}\) are algebraically independent, so evaluation
\(L[Z_1,\ldots,Z_n]\to L[c_1^{-1},\ldots,c_n^{-1}]\) is injective.
Compose this map with \(\phi\).
\end{proof}

\begin{proposition}[Minimal two-moment rigidity (\UPASS)]
\label{prop:transfer-two-moment-rigidity}
The pair \((T_2(c),T_4(c))\) recovers \(c\) and therefore the complete
divisor of \(\Ein-c\).  No single trace \(T_m\) is globally injective as
a function of \(c\in\C^\times\), since \(P_m\) has degree \(m\ge2\).
Moreover the only nontrivial collision of \((T_2,T_3)\) is
\[
 \{c,d\}=\{6-2\sqrt3,6+2\sqrt3\},
 \qquad
 (T_2,T_3)=\left(-\frac1{24},\frac1{96}\right).
\]
Thus trace order four is the smallest maximal order that gives a globally
rigid trace pair in this family.
\end{proposition}

\begin{proof}
Recovery from \((T_2,T_4)\) is the displayed universal identity.  A
nonconstant degree-\(m\) polynomial cannot be injective on \(\C\), and
the point \(c=0\), which is outside the level domain, does not remove all
generic fibers.  The collision calculation follows by eliminating
\(x=c^{-1}\) and \(y=d^{-1}\) from \(P_2(x)=P_2(y)\) and
\(P_3(x)=P_3(y)\); the non-diagonal branch has
\(x+y=1/2\), \(xy=1/24\), giving the stated reciprocal levels.
\end{proof}

\subsection{Analytic non-holonomy of the moving tower}

Recall the dyadic levels
\[
 Y_N=(N+1)\Ein(2^N)-N\Ein(2^{N+1})
 =\gamma+\varepsilon_N,
 \qquad
 \varepsilon_N\sim\frac{N+1}{2^N}e^{-2^N},
\]
and put \(\tau_{m,N}=P_m(Y_N^{-1})\).

\begin{theorem}[All-order moving-index non-holonomy (\UPASS)]
\label{thm:transfer-nonholonomy}
For every fixed \(m\ge2\), neither \((\tau_{m,N})_{N\ge1}\) nor its
signed increment sequence
\((\tau_{m,N+1}-\tau_{m,N})_{N\ge1}\) is P-recursive over \(\C\).
Equivalently, neither ordinary generating function is D-finite.
\end{theorem}

\begin{proof}
Let \(g_m(x)=P_m(x^{-1})\).  This is nonconstant and analytic at
\(x=\gamma\).  If \(r\ge1\) is the first nonzero Taylor order of
\(g_m(\gamma+u)-g_m(\gamma)\), then
\[
 \tau_{m,N}-T_m(\gamma)=b_m\varepsilon_N^r(1+o(1)),
 \qquad b_m\ne0.
\]
For each fixed \(j>0\),
\[
 \frac{\varepsilon_{N+j}}{\varepsilon_N}
 =\exp\bigl(-(2^j-1)2^N+O(N)\bigr),
\]
which tends to zero faster than every rational function of \(N\).
In a hypothetical polynomial-coefficient recurrence for the centered
sequence, divide by the term having the least shift with nonzero
coefficient.  Every later-shift ratio vanishes, leaving \(1=0\).  The
increment has the same leading scale, with coefficient \(-b_m\), so the
same argument applies.  P-recursive coefficient sequences are equivalent
to D-finite ordinary generating functions.
\end{proof}

There is a second, independent analytic obstruction.  The meromorphic
function
\[
 \mathscr R(u)=u^{-1}-\Gamma_{<}(u,1)
\]
has poles at every negative integer.  A D-finite meromorphic function can
have finite-plane singularities only among the finitely many zeros of the
leading polynomial of its differential equation.  Therefore
\(\mathscr R\) is not D-finite over \(\C(u)\).  This analytic statement is
\UPASS{} and does not depend on any arithmetic independence assertion.

\subsection{\texorpdfstring{Exact \(p\)-adic closures and their remaining
barriers}{Exact p-adic closures and their remaining barriers}}

For the Euler factorial Pad\'e system, recall
\[
 Q_n(t)=\sum_{h=0}^nh!\binom nh^2(-t)^h.
\]
For a prime \(p\), \(t\in\Z_p^\times\), and \(0\le r<p\), set
\[
 q_{r,p}(X)=\sum_{h=0}^rh!\binom rh^2(-X)^h\in\mathbb F_p[X].
\]

\begin{theorem}[Unit-boundary residue closure (\UPASS{} with \XPASS{}
certificates)]
\label{thm:transfer-padic-residue}
If \(q_{r,p}(\bar t)\ne0\) for every \(0\le r<p\), then every
\(Q_n(t)\) is a \(p\)-adic unit and the full Pad\'e sequence converges to
\[
 E_p(t)=\sum_{n\ge0}n!t^n.
\]
Exact finite residue-field calculation gives the sample classes
\[
\begin{array}{c|l}
p& t\pmod p\\ \hline
3&2\\
5&2,3,4\\
7&4,5\\
11&2,3,5,7\\
13&6,9,11,12.
\end{array}
\]
Together with the separate modulo-four calculation at \(p=2\), this proves
full convergence at \(t=-1\) for \(p=2,3,5,13\).  In particular, at
\(p=13\),
\[
 v_{13}\!\left(E_{13}(-1)-\frac{P_n(-1)}{Q_n(-1)}\right)
 \ge2v_{13}(n!)=\frac{n-s_{13}(n)}6.
\]
\end{theorem}

\begin{proof}
Modulo \(p\), terms with \(h\ge p\) vanish and the remaining binomial
coefficients depend only on \(n\bmod p\).  Hence
\(Q_n(t)\equiv q_{r,p}(\bar t)\pmod p\) when \(n\equiv r\pmod p\).
The unit conclusion and the standard Pad\'e remainder bound then give
convergence.  The listed classes are exact finite certificates.  For
\(p=2,t=-1\), one has \(v_2(Q_n(-1))=0\) for even \(n\) and \(1\) for odd
\(n\).  Legendre's formula gives the last equality.
\end{proof}

\begin{theorem}[One rational subsequence at all fixed places (\UPASS)]
\label{thm:transfer-common-subsequence}
Let \(N_j=\lcm(1,\ldots,j)\).  For every fixed negative integer \(t\),
the rational numbers
\[
 \frac{P_{N_j}(t)}{Q_{N_j}(t)}
\]
converge at the real place to the Stieltjes value
\(\int_0^\infty e^{-x}(1-tx)^{-1}\,dx\), and converge in \(\Q_p\) to
\(E_p(t)\) for every fixed prime \(p\).
\end{theorem}

\begin{proof}
The Carlitz congruence
\(Q_{n+k}(T)\equiv Q_n(T)\pmod{k\Z[T]}\) gives
\(Q_{N_j}(T)\equiv1\pmod{N_j\Z[T]}\).  Thus the denominators are eventually
units at each fixed prime, and the Pad\'e remainder tends to zero.  Real
Stieltjes--Pad\'e convergence holds for every \(t<0\), hence also along
this subsequence.
\end{proof}

Finally, the uniformly convergent interpolation
\[
 q_p(x)=\sum_{h\ge0}h!\binom xh^2\qquad(x\in\Z_p)
\]
is \(1\)-Lipschitz, satisfies \(q_p(n)=Q_n(-1)\), and gives the exact
\UPASS{} equivalence
\[
 \sup_{n\ge0}v_p(Q_n(-1))=\infty
 \quad\Longleftrightarrow\quad
 q_p\text{ has a zero in }\Z_p.
\]

\begin{statusbox}
\textbf{Open \(p\)-adic gates (\OPEN).}
Full convergence at \(t=-1\) is not known for every prime.  A zero of
\(q_p\) records unbounded denominator valuation; by itself it does not
imply failure of convergence, which requires comparison with
\(v_p(n!)\).  The fixed-prime valuation gain supplied here is only linear
in \(n\), whereas the relevant global denominator height has an
\(n\log n\) term, so these estimates do not furnish an irrationality
criterion for \(E_p(-1)\).  The common subsequence theorem is placewise and
does not assert convergence in a restricted adelic product topology.
\end{statusbox}

\subsection{Transfer boundary}

The results of this section are exact transports: [001] supplies the
arithmetic rank, and the trace, Gamma, Hankel, or conjugation maps preserve
enough of that rank to yield the displayed \RPASS{} conclusions.  The
formal geometry and analytic statements marked \UPASS{} require no such
input.  None of these mechanisms isolates an individual boundary constant
from an infinite limit, improves the fixed-prime height budget to the
required global scale, or converts a trace coordinate into a new comparison
theorem.  Those are the remaining gates rather than suppressed hypotheses.
\section{Multigenerator Jossen-type reductions for the resonant polyexponential algebra}

\begin{notebox}
\textbf{Scope.}  This section proves a rigorous reduction and two complete
non-cyclic subclasses.  It does \textbf{not} claim the full multigenerator
Jossen problem for the algebra below.  Statements using the algebraic-point
independence theorem of \textbf{[001]} have status \RPASS{}.
\end{notebox}

\subsection{1. The algebra and the arithmetic generic-point theorem}

Put

\[
  E_{q,\lambda}(z)=\operatorname{Ein}_q(\lambda z),\qquad
  \mathscr A=\overline{\mathbf Q}
  [E_{q,\lambda}:q\ge 1,\ \lambda\in\overline{\mathbf Q}^{\times}],
\]

and let

\[
  K_{\exp}=\overline{\mathbf Q}
  (e^\alpha:\alpha\in\overline{\mathbf Q}).
\]

Only finitely many generators occur in any element of \(\mathscr A\).

\subsubsection{Theorem 1.1 (polynomiality and algebraic-time genericity)}

After identifying repeated pairs \((q,\lambda)\), the functions
\(E_{q,\lambda}\) are algebraically independent over
\(\overline{\mathbf Q}\). Hence

\[
  \mathscr A\simeq
  \overline{\mathbf Q}[X_{q,\lambda}]
\]

is a polynomial UFD in the finite-subset sense.

For every \(\beta\in\overline{\mathbf Q}^{\times}\), evaluation at \(z=\beta\)
is injective on \(\mathscr A\), and the point

\[
  \Phi(\beta)=
  (\operatorname{Ein}_q(\lambda\beta))_{q,\lambda}
\]

is a generic point over \(K_{\exp}\): every finite set of its distinct
coordinates is algebraically independent over \(K_{\exp}\).

Consequently, if \(F\in\mathscr A\) is nonconstant, then

\[
  F(\beta)\notin K_{\exp}
  \quad(\beta\in\overline{\mathbf Q}^{\times}).                 \tag{J-1.1}
\]

More generally, \(\Phi(\beta)\) misses every proper algebraic subvariety
defined over \(K_{\exp}\).

\paragraph{Proof}

Apply Theorem 8.3 and Corollary 8.4 of 001 to the finitely many pairs
occurring in a proposed relation, with the arguments \(\lambda\beta\).
They are distinct exactly when the corresponding \(\lambda\)'s are distinct,
and the theorem simultaneously includes all selected orders \(q\). A
functional polynomial relation would specialize at any fixed nonzero
algebraic \(\beta\), contradicting the value theorem. The remaining claims
are immediate.

This is stronger than ordinary Zariski density: \textbf{every} nonzero algebraic
time is generic. It does not say anything comparable at a transcendental
time.

\subsubsection{Corollary 1.2 (complete algebraic-zero exclusion)}

For nonconstant \(F\in\mathscr A\) and \(a\in K_{\exp}\), every zero of
\(F(z)-a\), other than a possible zero at \(z=0\), is transcendental.
In particular, no nonzero element of \(\mathscr A\) has a nonzero algebraic
zero.

\subsubsection{Proposition 1.3 (every fiber is infinite)}

For every nonconstant \(F\in\mathscr A\) and every \(a\in\mathbf C\), the
fiber

\[
  \{z\in\mathbf C:F(z)=a\}
\]

is infinite.

\paragraph{Proof}

Every element of \(\mathscr A\) is an E-function and is of exponential type,
hence has order at most one. If \(F-a\) had only finitely many zeros,
Hadamard factorization would give

\[
  F(z)-a=e^{Az+B}R(z)                                      \tag{J-1.2}
\]

with \(R\in\mathbf C[z]\).

The functional proof of 001, Lemmas 8.1--8.2, is unchanged after adjoining
one extra exponential \(e^{Az}\): take a \(\mathbf Z\)-basis of the finitely
generated additive group spanned by \(A\) and all slopes \(\lambda\) occurring
in \(F\), and repeat the Laurent-character/Darboux-polynomial argument.
Thus the selected \(E_{q,\lambda}\)'s remain algebraically independent over

\[
  \mathbf C(z,e^{Az},e^{\lambda z}:\lambda\text{ selected}).
\]

Equation (1.2) is a nontrivial polynomial relation over this field, a
contradiction.

\subsection{2. Exact reduction of both parts of Jossen's conjecture}

Let \(F,G\in\mathscr A\setminus\{0\}\). Since only finitely many variables
occur, their gcd \(D=\gcd_{\mathscr A}(F,G)\) is well-defined. Write

\[
  F=DF_0,\qquad G=DG_0,\qquad \gcd_{\mathscr A}(F_0,G_0)=1.    \tag{J-2.1}
\]

\subsubsection{Theorem 2.1 (the exact remaining common-zero obstruction)}

\begin{enumerate}
\def\labelenumi{\arabic{enumi}.}
\item
  Neither \(F_0\) nor \(G_0\) can vanish at a nonzero algebraic point.
\item
  If \(F_0(0)=G_0(0)=0\), the common zero at the origin is explained by
  the nonunit E-function \(h(z)=z\): both \(F_0/z\) and \(G_0/z\) are
  E-functions.
\item
  Therefore the only unresolved case of Jossen part (ii) for
  \(\mathscr A\) is a \textbf{nonzero transcendental} point \(\xi\) satisfying

  \[
    F_0(\xi)=G_0(\xi)=0.                                  \tag{J-2.2}
  \]
\end{enumerate}

Equivalently, for a finite coordinate map \(\Phi:\mathbf C\to\mathbf A^n\),
the hard locus is the transcendental-time pullback

\[
  \Phi^{-1}(V(P,Q)),\qquad \operatorname{codim}V(P,Q)\ge2,  \tag{J-2.3}
\]

for coprime \(P,Q\in\overline{\mathbf Q}[X_1,\ldots,X_n]\).

The factor \(z\) in item 2 is standard E-function zero removal at the
algebraic point zero.  The contribution of [001] is to remove all other
algebraic times simultaneously.

\subsubsection{Theorem 2.2 (an entire-quotient counterexample would be infinite)}

Assume \(F/G\) is entire. With (2.1), exactly one of the following occurs.

\begin{itemize}
\tightlist
\item
  \(G_0\) is constant, in which case \(F/G=F_0/G_0\in\mathscr A\), so
  Jossen part (i) holds inside \(\mathscr A\).
\item
  \(G_0\) is nonconstant. Then \(G_0\) has infinitely many zeros by
  Proposition 1.3, every one of them is a zero of \(F_0\), and all but a
  possible zero at the origin are transcendental. Thus a failure of
  divisibility in \(\mathscr A\) would force an \textbf{infinite transcendental
  codimension-two intersection carrying the complete divisor of \(G_0\)}.
\end{itemize}

This is a useful sharpening of the gate: one isolated mysterious root is the
obstruction for part (ii), whereas part (i) requires an infinite divisor-sized
exceptional intersection.

\subsubsection{Origin contact and transversality}

For a finite list \(E_i=E_{q_i,\lambda_i}\),

\[
  E_i(z)=\lambda_i z+O(z^2).
\]

Let \(P_d\) be the first nonzero homogeneous part of
\(P\in\overline{\mathbf Q}[X_1,\ldots,X_n]\). If

\[
  P_d(\lambda_1,\ldots,\lambda_n)\ne0,
\]

then

\[
  \operatorname{ord}_{z=0}P(E_1(z),\ldots,E_n(z))=d.          \tag{J-2.4}
\]

In particular, the origin is a simple zero precisely under the familiar
first-order transversality condition

\[
  \sum_i\lambda_i\partial_iP(0)\ne0.                         \tag{J-2.5}
\]

If the expression in (2.5) vanishes, higher jets give an effective valuation
algorithm. This completely settles the algebraic-time Jacobian question but
does not control (2.2).

\subsection{2A. A general affine E-surface theorem}

There is a substantially broader two-dimensional closure that does not use
any special identity of \(\operatorname{Ein}\).

\subsubsection{Theorem 2A.1 (factorial E-surface criterion)}

Let \(B\) be a finitely generated factorial
\(\overline{\mathbf Q}\)-domain of Krull dimension two, embedded as a
subring of the ring of E-functions. Assume that \(B\) contains at least one
nonconstant polynomial coordinate

\[
  \tau(z)\in\overline{\mathbf Q}[z].
\]

Then both parts of Jossen's conjecture hold for every nonzero
\(f,g\in B\) (so in particular the quotient below has \(g\ne0\)):

\begin{enumerate}
\def\labelenumi{\arabic{enumi}.}
\tightlist
\item
  if \(f/g\) is entire, then \(f/g\) is an E-function;
\item
  if \(f\) and \(g\) have a common zero, they have a common nonunit
  E-function divisor.
\end{enumerate}

\paragraph{Proof of the common-zero assertion}

Take a gcd in the UFD \(B\):

\[
  f=df_0,\qquad g=dg_0,\qquad \gcd_B(f_0,g_0)=1.             \tag{J-2A.1}
\]

Let \(\xi\) be a common zero. If \(d(\xi)=0\), then \(d\) itself is a
nonunit E-function divisor of \(f\) and \(g\).

Assume \(d(\xi)\ne0\). Then \(f_0(\xi)=g_0(\xi)=0\). Since \(f_0,g_0\)
are nonzero and \(B\) is a factorial domain, neither the height-zero prime nor
any height-one prime contains both of them. Hence every prime over
\((f_0,g_0)\) has height two. Because \(\dim B=2\),

\[
  B/(f_0,g_0)
\]

is a zero-dimensional affine \(\overline{\mathbf Q}\)-algebra. The kernel
of evaluation at \(\xi\) is a prime over \((f_0,g_0)\), hence is maximal.
By Zariski's lemma its residue field is algebraic over
\(\overline{\mathbf Q}\), hence equals \(\overline{\mathbf Q}\). In particular
\(\tau(\xi)\in\overline{\mathbf Q}\). Since \(\xi\) is a root of the
nonzero polynomial \(\tau(T)-\tau(\xi)\), this proves
\(\xi\in\overline{\mathbf Q}\).

If \(\xi\ne0\), Beukers's algebraic-point zero-removal proposition gives

\[
  \frac{f_0(z)}{z-\xi},\qquad
  \frac{g_0(z)}{z-\xi}
\]

as E-functions; for \(\xi=0\), the same conclusion follows directly by
shifting the Taylor coefficients of an E-function that vanishes at the
origin. Thus \(z-\xi\) is the required common nonunit E-function
divisor (after multiplying the displayed quotients by \(d\)).

\paragraph{Proof of the entire-quotient assertion}

Cancel \(d\) as in (2A.1); the cancelled quotient extends to the same entire
function across the zeros of \(d\). If \(f_0/g_0\) is entire, every zero of \(g_0\)
is a common zero of \(f_0,g_0\). The zero-dimensional scheme above has only
finitely many closed points. Each such point fixes an algebraic value of
\(\tau\), and a nonconstant polynomial has finite fibers. Thus \(g_0\) has
finitely many zeros, all algebraic. Let

\[
  r(z)=\prod_\xi(z-\xi)^{\operatorname{ord}_\xi g_0}
  \in\overline{\mathbf Q}[z].                               \tag{J-2A.2}
\]

Repeated algebraic-point zero removal shows that \(g_0/r\) is a zero-free
E-function. A zero-free entire function of order at most one is
\(Ce^{az}\) by Hadamard factorization. Here

\[
  C=(g_0/r)(0)\in\overline{\mathbf Q}^{\times},\qquad
  a=(g_0/r)'(0)/(g_0/r)(0)\in\overline{\mathbf Q},
\]

so \(g_0/r=Ce^{az}\) is a unit E-function. Entirety also gives
\(\operatorname{ord}_\xi f_0\ge\operatorname{ord}_\xi g_0\), hence
\(f_0/r\) is an E-function. Finally,

\[
  \frac fg=\frac{f_0/r}{Ce^{az}}
\]

is an E-function.

\subsubsection{\texorpdfstring{Corollary 2A.2 (one arbitrary E-function plus \(z\))}{Corollary 2A.2 (one arbitrary E-function plus z)}}

Let \(F\) be any nonpolynomial E-function. Then \(z\) and \(F(z)\) are
algebraically independent: an entire function algebraic over
\(\overline{\mathbf Q}(z)\) is a polynomial. Consequently

\[
  B_F=\overline{\mathbf Q}[z,F(z)]
  \simeq\overline{\mathbf Q}[X,Y]
\]

is a factorial affine E-surface, and \textbf{both parts of Jossen's conjecture hold
unconditionally throughout \(B_F\)}.

In direct plane language, if coprime \(P,Q\in\overline{\mathbf Q}[X,Y]\)
have a common pullback zero \(\xi\), then
\((\xi,F(\xi))\in V(P,Q)\). The plane intersection is finite and defined
over \(\overline{\mathbf Q}\), so \(\xi\) is algebraic and \(z-\xi\) is the
common E-factor. For an entire quotient, the same finite intersection traps
all denominator zeros and the preceding Hadamard argument completes the
proof.

Equivalently, for arbitrary \(P,Q\) and
\(D=\gcd(P,Q)\), one has the exact set-theoretic decomposition

\[
 \begin{aligned}
 Z(P(z,F))\cap Z(Q(z,F))
  ={}& Z(D(z,F))\\
    &{}\cup S_{P,Q,F},
 \end{aligned}                                                \tag{J-2A.3}
\]

where \(S_{P,Q,F}\) is a finite set of algebraic numbers. Thus every
transcendental common zero is automatically carried by the explicit common
E-factor \(D(z,F)\); only finitely many algebraic residual points require the
linear factors \(z-\xi\).

The polynomial case reduces to the univariate algebra
\(\overline{\mathbf Q}[z]\) and is elementary. Thus the corollary holds for
every E-function \(F\), with the nonpolynomial case being the genuinely
two-dimensional one.

More generally, if \(\tau\in\overline{\mathbf Q}[z]\) is any nonconstant
polynomial and \(F\) is a nonpolynomial E-function, then \(\tau(z),F(z)\)
are algebraically independent (otherwise \(F\) would be algebraic over
\(\overline{\mathbf Q}(z)\)). Therefore

\[
  \overline{\mathbf Q}[\tau(z),F(z)]
\]

also satisfies both parts of Jossen.

For \(F=e^{\beta z}\), this recovers the familiar one-support exponential-
polynomial situation. The point of Corollary 2A.2 is that no exponential,
hypergeometric, differential-order, or zero-geometry assumption on \(F\) is
needed beyond the E-function axioms.

In particular, every nonconstant \(F\in\mathscr A\) is nonpolynomial (by the
functional independence in 001), and hence

\[
  \boxed{\overline{\mathbf Q}[z,F]\text{ satisfies both parts of Jossen}}
                                                                  \tag{J-2A.4}
\]

for \textbf{every single polyexponential polynomial observable} \(F\). This is a
strict extension of the cyclic result
\(\overline{\mathbf Q}[F]\): the numerator and denominator may depend
independently on the physical variable \(z\) and on \(F(z)\).

This surface theorem explains precisely why adding a second independent
polyexponential beyond \(z,F\) changes the problem: in Krull dimension three,
two coprime equations can leave a positive-dimensional intersection, and the
argument forcing \(\xi\) algebraic disappears.

\subsubsection{Corollary 2A.3 (holonomic transfer)}

The proof is not intrinsically restricted to Siegel E-functions. Let
\(\mathcal H_{\overline{\mathbf Q}}^{\mathrm{fin}}\) be the ring of
finite-order entire functions that are D-finite over
\(\overline{\mathbf Q}(z)\) and have algebraic Taylor coefficients. It is
closed under sums and products and under division by \(z-\xi\) at an
algebraic zero, with finite order preserved. A zero-free finite-order entire
function \(h\) is \(Ce^{A(z)}\) with \(A\) a polynomial. Normalize this
representation by taking \(A(0)=0\) and \(C=h(0)\). Then
\(C\in\overline{\mathbf Q}^{\times}\), and the formal logarithm of \(h/C\)
shows that \(A\in\overline{\mathbf Q}[z]\). Hence \(h^{-1}=C^{-1}e^{-A(z)}\)
is again D-finite of finite order with algebraic Taylor coefficients.

Therefore Theorem 2A.1, with ``E-function'' replaced by ``finite-order entire
D-finite function with algebraic Taylor coefficients,'' remains valid. In
particular, every factorial affine holonomic surface containing a nonconstant
polynomial coordinate is closed under its entire quotients and has the same
common-zero factor property. This is an elementary algebraic--analytic
transfer to the holonomic setting.

\subsection{\texorpdfstring{3. A complete genuinely bivariate blow-up subclass inside \(\mathscr A\)}{3. A complete genuinely bivariate blow-up subclass inside \textbackslash mathscr A}}

The next result is not merely cyclic in one generator. It exhibits the
geometric mechanism by which a codimension-two point becomes an E-divisor.

Choose

\[
  H(z)=\operatorname{Ein}(\lambda z),\qquad
  R(z)=\operatorname{Ein}_p(\mu z),
\]

with \((p,\mu)\ne(1,\lambda)\), and fix
\(a\in\overline{\mathbf Q}\), \(m\ge1\). Theorem 1.1 makes \(H,R\)
algebraically independent. Put

\[
  U=a+H^mR.                                                   \tag{J-3.1}
\]

Then \(H,U\) are also algebraically independent, so
\(\mathscr B=\overline{\mathbf Q}[H,U]\) is a genuine polynomial
two-generator subalgebra of \(\mathscr A\).

For \(0\ne P\in\overline{\mathbf Q}[X,Y]\), expand

\[
  P(X,a+Y)=\sum_{i,j\ge0}c_{ij}X^iY^j
\]

and define the weighted exceptional valuation

\[
  \nu_m(P)=\min\{i+mj:c_{ij}\ne0\}.                          \tag{J-3.2}
\]

\subsubsection{Theorem 3.1 (exact pullback gcd)}

If \(P,Q\in\overline{\mathbf Q}[X,Y]\) are coprime, then, up to a nonzero
algebraic constant,

\[
 \gcd_{\mathscr A}\bigl(P(H,U),Q(H,U)\bigr)
 =H^{\min(\nu_m(P),\nu_m(Q))}.                               \tag{J-3.3}
\]

\paragraph{Proof}

Substitution gives

\[
  P(H,U)=\sum c_{ij}H^{i+mj}R^j
         =H^{\nu_m(P)}P^\sharp(H,R),
\]

where \(P^\sharp\) is not divisible by \(H\). After inverting \(H\), the
substitution is the isomorphism

\[
 \overline{\mathbf Q}[X,Y,X^{-1}]
 \xrightarrow{\sim}
 \overline{\mathbf Q}[H,R,H^{-1}],
 \quad X\mapsto H,\quad Y\mapsto a+H^mR,
\]

whose inverse sends \(R\) to \((Y-a)/X^m\). Coprimality therefore forbids
every common irreducible factor other than \(H\), and (3.3) follows.

\subsubsection{Corollary 3.2 (complete Jossen common-zero theorem at the contracted point)}

Assume

\[
  \sqrt{(P,Q)}=(X,Y-a),                                      \tag{J-3.4}
\]

so the affine complete intersection of \(P\) and \(Q\) is exactly the point
\((0,a)\). Then

\[
 \{P(H,U)=Q(H,U)=0\}=\{H=0\}.                               \tag{J-3.5}
\]

The right side is infinite; zero is its only algebraic point, and all its
other points are transcendental. Every zero is simple. All of these common
zeros, including the transcendental ones, are explained by the explicit
nonunit E-function factor in (3.3). Thus Jossen part (ii) holds completely
for this bivariate class.

Indeed,

\[
 H'(z)=\frac{1-e^{-\lambda z}}z.
\]

A multiple nonzero zero would have \(\lambda z=2\pi ik\). But

\[
 \Re\operatorname{Ein}(2\pi ik)
 =\int_0^1\frac{1-\cos(2\pi kt)}t\,dt>0\qquad(k\ne0),         \tag{J-3.6}
\]

so this is impossible. The zero at zero is simple because \(H'(0)=\lambda\).
There are infinitely many zeros since otherwise Hadamard factorization would
make \(H=e^{Az+B}S(z)\) and hence \(H'\) would have only finitely many zeros,
whereas \(H'(2\pi ik/\lambda)=0\) for every \(k\ne0\). Their transcendence
follows from Theorem 1.1.

A basic example is

\[
  P=X,\qquad Q=Y-a,
\]

for which the coprime bivariate pair pulls back to

\[
  H,\qquad H^mR.
\]

It has infinitely many transcendental common zeros, but the common factor is
exactly \(H\). Geometrically, (3.1) contracts the divisor \(H=0\) to the
codimension-two point \((0,a)\); Jossen's factor is the exceptional divisor.

\subsubsection{Theorem 3.3 (an exact non-cyclic entire-quotient window)}

Let \(1\le\nu\le m\) and \(P\in\overline{\mathbf Q}[X,Y]\setminus\{0\}\).
Then

\[
 \frac{P(H,U)}{H^\nu}\text{ is entire}
 \quad\Longleftrightarrow\quad
 \partial_X^jP(0,a)=0\quad(0\le j<\nu)                     \tag{J-3.7}
\]

and, whenever these equivalent conditions hold,

\[
  \frac{P(H,U)}{H^\nu}\in\mathscr A.                        \tag{J-3.8}
\]

Hence Jossen part (i) holds for this whole bivariate numerator class and the
denominators \(H^\nu\), \(1\le\nu\le m\).

\paragraph{Proof}

For weights strictly below \(\nu\le m\), only the pure terms \(c_{i0}X^i\)
can occur in (3.2); every term involving \(Y-a\) acquires a factor \(H^m\).
Thus the derivative conditions in (3.7) are equivalent to
\(\nu_m(P)\ge\nu\), which makes the quotient in (3.8) a polynomial in
\(H,R\). If one derivative does not vanish, the first offending pure
\(H^i\)-term produces a pole at every simple zero of \(H\). This proves the
reverse implication.

For example,

\[
  \frac{U-a}{H^m}=R
\]

recovers an independent polyexponential coordinate as an entire E-function
quotient; this is not contained in the cyclic algebra
\(\overline{\mathbf Q}[H]\).

\subsection{4. A natural differential two-generator subclass}

Let

\[
  H=\operatorname{Ein}(\lambda z),\qquad
  \Theta H=zH'=1-e^{-\lambda z}.
\]

The pair \((H,\Theta H)\) is algebraically independent over
\(\overline{\mathbf Q}\), by the functional exponential/polyexponential
theorem in 001. Equation (3.6) proves

\[
  \{H=0\}\cap\{\Theta H=0\}=\{0\}.                          \tag{J-4.1}
\]

Consequently, for any \(P,Q\in\overline{\mathbf Q}[X,Y]\) satisfying

\[
  \sqrt{(P,Q)}=(X,Y),                                       \tag{J-4.2}
\]

the common-zero set of \(P(H,\Theta H)\) and
\(Q(H,\Theta H)\) is exactly \(\{0\}\), and the common zero is explained by
the E-function factor \(z\). This is a complete non-cyclic common-zero
theorem in the first differential closure of \(\mathscr A\).

There is also a broad imported quotient theorem. Fischler--Rivoal prove that
if \(g^r\) is an E-function and \(L\in\overline{\mathbf Q}(z)[d/dz]\) is such
that \(L(g)/g\) is entire, then \(L(g)/g\) is a polynomial. Thus Jossen part
(i) holds for every pair \((L(G),G)\) lying in \(\mathscr A^2\). For example,
\(\Theta\) maps

\[
 \overline{\mathbf Q}[\operatorname{Ein}_q(\lambda z):q\ge2]
 \quad\text{into}\quad\mathscr A.
\]

For a nonconstant \(G\) in this subalgebra, \(\Theta G/G\) is in fact never
entire: Proposition 1.3 supplies a nonzero zero of \(G\), and the logarithmic
derivative has a pole there.

\subsection{5. Why the full multigenerator statement is still open}

\subsubsection{5.1 Algebraic independence is not codimension-two avoidance}

An algebraically nondegenerate finite-order entire curve can meet a
codimension-two target infinitely often. The elementary map

\[
  z\longmapsto(\sin z,z\sin z)
\]

is algebraically nondegenerate and meets \((0,0)\) at every \(z\in\pi\mathbf Z\).
Here the hidden divisor \(\sin z\) explains the intersection. This is exactly
the sort of divisor that Jossen asks one to recover inside the E-function
ring. Therefore Theorem 1.1 alone cannot settle (2.3).

\subsubsection{5.2 Why standard Nevanlinna gcd estimates do not immediately close it}

Nevanlinna gcd theorems on algebraic tori give
\(o(T)\)-type bounds for codimension-two intersections under zero-free or
small-zero hypotheses on the coordinate functions. The polyexponential
coordinates have infinite divisors of their own, so those hypotheses are not
automatic. Even an \(o(T)\) upper bound would need a matching lower bound for
the zero counting function of every possible denominator \(Q(\Phi)\) before
it could prove Jossen part (i).

Thus a usable Nevanlinna route requires two new ingredients:

\begin{enumerate}
\def\labelenumi{\arabic{enumi}.}
\tightlist
\item
  a gcd estimate adapted to the additive/iterated-integral differential
  group containing the \(\operatorname{Ein}_q\)'s, rather than a torus; and
\item
  a uniform lower zero-density theorem for every nonconstant element of
  \(\mathscr A\), stronger than Proposition 1.3.
\end{enumerate}

\subsubsection{5.3 The exact remaining structural gate}

Any one of the following would close a substantial part of the remaining
problem.

\begin{itemize}
\tightlist
\item
  \textbf{Transcendental codimension-two saturation:} if coprime
  \(P,Q\in\mathscr A\) have a nonzero common root, prove that their pullbacks
  share a nonunit E-function factor (necessarily outside \(\mathscr A\)).
\item
  \textbf{Strong avoidance:} prove that a full-\(\mathscr A\)-coprime pair has no
  nonzero common root at all. This is stronger than Jossen and may be false;
  it should not be asserted without a new zero theorem.
\item
  \textbf{Divisor-sized gcd estimate:} for coprime \(P,Q\), show that the common-zero
  counting function is too small to contain the complete divisor of
  \(Q(\Phi)\). Together with Theorem 2.2 this proves the entire-quotient part.
\item
  \textbf{Bourget-type seed:} already the assertion that two distinct primitive
  coordinates \(\operatorname{Ein}_q(\lambda z)\) and
  \(\operatorname{Ein}_p(\mu z)\) have no common nonzero zero is not supplied
  by 001. Proving it uniformly would be a natural first step.
\end{itemize}

\subsection{6. Literature boundary and exact limitations}

Fischler--Rivoal state Jossen's conjecture as follows:

\begin{enumerate}
\def\labelenumi{\arabic{enumi}.}
\tightlist
\item
  if E-functions \(f,g\) have an entire quotient \(f/g\), then that quotient
  is an E-function;
\item
  if two E-functions have a common root, then they have a common nonunit
  E-function divisor.
\end{enumerate}

They record that part (ii) is known when the common root is algebraic, part
(i) is known for polynomial denominators, and prove the differential-operator
case used above. Under Schanuel they obtain much stronger common-zero results
for exponential polynomials. The transcendental-root case for general
E-functions remains the hard boundary.

Primary sources:

\begin{itemize}
\tightlist
\item
  S. Fischler and T. Rivoal, \href{https://doi.org/10.1007/s12215-026-01475-x}{\emph{Zeros of E-functions and of exponential
  polynomials defined over} \(\overline{\mathbf Q}\)},
  \emph{Rend. Circ. Mat. Palermo (2)} \textbf{75} (2026), Art. 158; see also the
  \href{https://arxiv.org/abs/2503.20345}{2025 preprint}, especially Conjecture 1.1,
  Theorem 1.2, and Sections 5--6.
\item
  F. Beukers, \emph{A refined version of the Siegel--Shidlovskii theorem},
  Ann. of Math. 163 (2006), 369--379, especially the algebraic-point
  zero-removal proposition used in Theorem 2A.1.
\item
  A. Levin and J. T.-Y. Wang, \href{https://arxiv.org/abs/1903.03876}{\emph{Greatest common divisors of analytic functions
  and Nevanlinna theory on algebraic tori}}.
\item
  001, Theorem 8.3 and Corollary 8.4, for the resonant value and functional
  independence used throughout.
\end{itemize}

The precise remaining obstruction is the transcendental codimension-two
saturation gate of Section 5.3.  None of the preceding subclasses is used to
claim the full multigenerator conjecture.

\part{Arithmetic Barriers and Remaining Gates}
\section{The critical moving-point barrier}
\label{sec:barrier}

We first separate the coefficient of \(\gamma\) from the gamma-free tail.
Fix \(q>1\), put
\[
 f_q(t)=E_1(q^t),
\]
and consider a finitely supported coefficient vector \((c_n)\).  The two
linear statistics
\[
 \omega(c)=\sum_nc_n,
 \qquad
 \mu(c)=\sum_nnc_n
\]
give the exact decomposition
\begin{equation}\label{eq:augmentation-decomposition}
 \sum_nc_n\Ein(q^n)
 =\omega(c)\gamma+\mu(c)\log q+\sum_nc_nf_q(n).
\end{equation}
We call \(\omega=1,\mu=0\) the augmentation-one, log-balanced stratum.
The dyadic level belongs to this stratum.  The next theorem completely solves
its natural cumulative-positive cone.

Direct differentiation gives
\begin{equation}\label{eq:fq-convex}
 f_q'(t)=-(\log q)e^{-q^t}<0,
 \qquad
 f_q''(t)=(\log q)^2q^te^{-q^t}>0.
\end{equation}

\begin{theorem}[Sharp cumulative-positive extremality]
\label{thm:augmentation-extremality}
Let \(M\ge2\), and suppose that real coefficients \(c_0,\ldots,c_M\)
satisfy
\[
 \sum_{n=0}^Mc_n=1,
 \qquad
 \sum_{n=0}^Mnc_n=0.
\]
Put \(C_k=\sum_{n=0}^kc_n\), and assume \(C_k\ge0\) for \(0\le k<M\).
Then the complete sharp range is
\begin{equation}\label{eq:augmentation-range}
 \begin{split}
 M f_q(M-1)-(M-1)f_q(M)
 &\le \sum_{n=0}^Mc_nf_q(n)\\
 &\le f_q(M)+M\{f_q(0)-f_q(1)\}.
 \end{split}
\end{equation}
The lower endpoint is attained uniquely by
\begin{equation}\label{eq:augmentation-minimizer}
 c_{M-1}=M,\qquad c_M=-(M-1),
 \qquad c_0=\cdots=c_{M-2}=0.
\end{equation}
The upper endpoint is attained uniquely by
\(c_0=M,c_1=-M,c_M=1\), with the other coefficients zero.
Consequently the moving level \(Y_{M-1}(q)-\gamma\) is the unique minimum
tail in this cone, and
\[
 \min\sum_nc_nf_q(n)
 \sim \frac{M}{q^{M-1}}e^{-q^{M-1}}.
\]
\end{theorem}

\begin{proof}
Abel summation and the moment normalization give
\begin{align*}
 \sum_{n=0}^Mc_nf_q(n)
 &=f_q(M)+\sum_{k=0}^{M-1}C_k\{f_q(k)-f_q(k+1)\},\\
 \sum_{k=0}^{M-1}C_k&=M.
\end{align*}
By strict convexity, the positive differences
\(f_q(k)-f_q(k+1)\) decrease strictly with \(k\).  The displayed sum is
therefore minimized by putting all mass \(M\) at \(k=M-1\), and maximized
by putting it at \(k=0\).  Unwinding the partial sums gives the two unique
coefficient vectors.  The feasible \(C\)-vectors form the simplex
\(\{C_k\ge0:\sum C_k=M\}\); its continuous linear image is the whole
interval between these two endpoints, which proves the asserted complete
range.  The asymptotic is \eqref{eq:E1-asymptotic}.
\end{proof}

The same doubly exponential scale survives when the two levels use
multiplicatively independent bases.

\begin{proposition}[Two-base moving-level obstruction]
\label{prop:two-base-obstruction}
For \(q>1\), write
\[
 Y_N(q)=(N+1)\Ein(q^N)-N\Ein(q^{N+1})
       =\gamma+\varepsilon_N(q).
\]
Let
\[
 M_N=\left\lfloor\frac{N\log2}{\log3}\right\rfloor+1,
 \qquad W_N=Y_N(2)-Y_{M_N}(3).
\]
Then \(W_N>0\) for all sufficiently large \(N\), and
\begin{equation}\label{eq:two-base-asymptotic}
 \boxed{\displaystyle
 -\log W_N
 =2^N+N\log2-\log(N+1)+O(1).}
\end{equation}
In particular \(-\log W_N=2^N+O(N)\).  Thus changing from one
multiplicative orbit to two independent ones does not remove the
moving-point scale of the fixed-form obstruction.
\end{proposition}

\begin{proof}
Put \(x=2^N\) and \(y=3^{M_N}\).  Since
\(N\log2/\log3\) is not an integer, \(y>x\); both are integers, so
\(y\ge x+1\).  The standard bounds
\[
 \frac{e^{-u}}{u+1}<E_1(u)<\frac{e^{-u}}u
\]
show, uniformly here, that
\[
 \varepsilon_N(2)
 =\frac{N+1}{x}e^{-x}\{1+O(x^{-1})\}.
\]
Moreover \(M_N/N\to\log2/\log3<1\), and the same bounds give
\[
 \frac{\varepsilon_{M_N}(3)}{\varepsilon_N(2)}
 \le
 \left(\frac{\log2}{\log3}+o(1)\right)e^{-1}<1.
\]
Hence \(W_N\) is positive and differs from \(\varepsilon_N(2)\) by a
factor bounded above and below away from zero.  Taking logarithms yields
\eqref{eq:two-base-asymptotic}.
\end{proof}

Outside this cone, augmentation one alone has no sign content.

\begin{proposition}[Three-point flexibility]
\label{prop:three-point-flexibility}
Let \(0<a<b<c\) be positive integers.  Put
\[
 c^{(0)}=\left(\frac b{b-a},-\frac a{b-a},0\right),
 \qquad
 v=(c-b,a-c,b-a).
\]
For \(t\in\R\), the vector \(c(t)=c^{(0)}+tv\), supported at \(a,b,c\),
is augmentation one and log balanced.  Its tail is \(A+tD\), where
\[
 A>0,
 \qquad
 D=(b-a)(c-b)(c-a)f_q[a,b,c]>0.
\]
Hence rational admissible coefficients give a dense set of tail values in
\(\R\); in particular both signs and arbitrarily small nonzero values occur.
For example, at \(q=2\) and exponents \(1,2,3\),
\[
 -E_1(2)+5E_1(4)-3E_1(8)<0.
\]
\end{proposition}

\begin{proof}
The two normalizations are immediate.  The slope \(D\) is a positive second
divided difference by \eqref{eq:fq-convex}.  Since \(A+D\Q\) is dense in
\(\R\), all qualitative assertions follow.  For the explicit example use
\[
 E_1(x)>\frac{e^{-x}}{x+1},
 \qquad
 E_1(x)<\frac{e^{-x}}x:
\]
the displayed value is less than
\(-e^{-2}/3+5e^{-4}/4<0\).
\end{proof}

The next estimate identifies the coefficient height at which such
cancellation can begin.  Write a rational coefficient vector primitively as
\(c_j=a_j/D\), with \(D\ge1\), and put
\(H(c)=\max(D,|a_1|,\ldots,|a_r|)\).

\begin{proposition}[First-tail dominance]
\label{prop:first-tail-dominance}
Let \(n_1<\cdots<n_r\) be integers with \(a_1\ne0\).  If
\begin{equation}\label{eq:height-subcritical}
 (r-1)H(c)\frac{1+q^{-n_1}}q
 e^{-(q-1)q^{n_1}}\le\frac12,
\end{equation}
then the tail has the sign of \(a_1\) and
\begin{equation}\label{eq:first-tail-lower}
 \boxed{\displaystyle
 \left|\sum_{j=1}^rc_jE_1(q^{n_j})\right|
 \ge\frac{e^{-q^{n_1}}}{2H(c)(q^{n_1}+1)}.}
\end{equation}
Thus, for fixed \(q,r\), cancellation of the first active tail requires
\[
 \log H(c)\ge(q-1)q^{n_1}+O(1).
\]
\end{proposition}

\begin{proof}
The reverse triangle inequality gives
\[
 \left|\sum_jc_jE_1(q^{n_j})\right|
 \ge\frac{E_1(q^{n_1})-(r-1)H(c)E_1(q^{n_1+1})}{H(c)}.
\]
The elementary bounds used in the preceding proof imply
\[
 \frac{E_1(q^{n+1})}{E_1(q^n)}
 \le\frac{1+q^{-n}}q e^{-(q-1)q^n}.
\]
Now apply \eqref{eq:height-subcritical}.
\end{proof}

For consecutive triples the transition is exact.  Let
\[
 F_{j,N}=E_1(q^{N+j}),
 \quad
 A_N=(N+2)F_{1,N}-(N+1)F_{2,N},
 \quad
 B_N=F_{0,N}-2F_{1,N}+F_{2,N},
\]
and \(\rho_N=B_N/A_N\).  Both \(A_N,B_N\) are positive.

\begin{theorem}[Exact denominator phase transition]
\label{thm:denominator-transition}
Every rational augmentation-one, log-balanced coefficient vector on
\(N,N+1,N+2\) has the form
\[
 c(s)=(s,N+2-2s,s-N-1).
\]
For \(s=a/D\) with \(D\ge1\), its tail \(\varepsilon_N(s)\) satisfies
\begin{equation}\label{eq:denominator-exact}
 \boxed{\displaystyle
 |\varepsilon_N(a/D)|
 =\frac{B_N}{D}\left|a+\frac D{\rho_N}\right|,}
\end{equation}
and hence
\[
 \min_{a\in\Z}|\varepsilon_N(a/D)|
 =\frac{B_N}{D}\left\|\frac D{\rho_N}\right\|.
\]
In particular, \(D\le\rho_N/2\) implies
\(|\varepsilon_N(a/D)|\ge A_N\) for every integer \(a\).  Conversely,
whenever \(\rho_N\ge\tfrac12\) (hence for all sufficiently large \(N\)),
choosing a positive integer \(D\) nearest to \(\rho_N\) and taking
\(a=-1\) gives either an exact zero or
\begin{equation}\label{eq:denominator-upper}
 0<|\varepsilon_N(-1/D)|\le\frac{A_N}{2D}.
\end{equation}
Furthermore, with \(X=q^N\),
\begin{equation}\label{eq:rho-asymptotic}
 \rho_N\sim\frac q{N+2}e^{(q-1)X}.
\end{equation}
Thus the denominator threshold is exactly
\(e^{(q-1)q^N}/N\) up to a constant factor.
\end{theorem}

\begin{proof}
Substitution gives \(\varepsilon_N(s)=A_N+B_Ns\), proving
\eqref{eq:denominator-exact}.  If \(D/\rho_N\le1/2\), its nearest integer
is zero.  If \(|D-\rho_N|\le1/2\), then
\(|1-D/\rho_N|\le1/(2\rho_N)\), which gives
\eqref{eq:denominator-upper}.  Finally \eqref{eq:E1-asymptotic} yields
\[
 B_N\sim E_1(X),
 \qquad
 A_N\sim(N+2)E_1(qX),
\]
and hence \eqref{eq:rho-asymptotic}.
\end{proof}

When \(q\) is algebraic, exact zero can occur for at most one \(N\); by
\cref{thm:uN-defect}, any such zero would already force the algebraic
independence of \(\gamma,\delta\) over \(K\).  Discarding that possible one
index, the nearest-denominator construction gives normalized coefficient
height
\[
 H_N=q e^{(q-1)X}(1+o(1))
\]
and
\begin{equation}\label{eq:critical-height-upper}
 0<|\varepsilon_N|
 \le\frac{(N+2)^2}{2q^2X}
 e^{-(2q-1)X}(1+o(1)).
\end{equation}
Consequently a uniform lower bound, on normalized nonzero rational
augmentation-one triples, of the form
\[
 |\varepsilon|\ge C e^{-\beta X}H^{-\kappa}X^{-\lambda}
\]
requires
\begin{equation}\label{eq:height-frontier}
 \boxed{\beta+\kappa(q-1)\ge2q-1.}
\end{equation}
At equality it also requires \(\lambda\ge1\).  In particular, at the natural
point cost \(\beta=1\), one must have \(\kappa\ge2\).  This exponent refers
to the affine normalization \(\sum c_j=1\); clearing denominators multiplies
the value and changes the corresponding homogeneous height exponent.

The preceding sections construct exact finite spectral data converging very
rapidly to Euler-level traces.  To extract the arithmetic nature of the
limit, one would need a lower bound uniform in a moving algebraic sample
point.  We first show that sublinear point cost is impossible.

Let
\[
 F_1(z)=\Ein(z),\qquad
 F_2(z)=\Ein(2z),\qquad
 F_3(z)=\Ein(4z),
\]
and let \(P_0=X_1-2X_2+X_3\).

\begin{theorem}[A fixed-form critical-scale no-go]
\label{thm:sublinear-no-go}
For \(x\ge1\),
\begin{equation}\label{eq:R-def}
 R(x):=P_0(F_1(x),F_2(x),F_3(x))
 =E_1(x)-2E_1(2x)+E_1(4x)>0,
\end{equation}
and
\begin{equation}\label{eq:R-asymptotic}
 R(x)=\frac{e^{-x}}x
 \left(1+O(x^{-1})+O(e^{-x})\right).
\end{equation}
Consequently, for every \(A>0\) and every \(\vartheta<1\),
\begin{equation}\label{eq:R-sublinear}
 0<R(n)<e^{-An^\vartheta}
\end{equation}
for all sufficiently large integers \(n\).

More sharply, for every \(C_0>0\), a lower bound
\begin{equation}\label{eq:critical-polynomial-template}
 R(n)\ge C_0e^{-\sigma n}n^{-\lambda}
\end{equation}
cannot hold for all large \(n\) if either \(\sigma<1\), or if
\(\sigma=1\) and \(\lambda<1\).

In particular, no lower bound
\begin{equation}\label{eq:uniform-template}
 |P(F(\xi))|\ge
 \exp\!\left[-C(1+|\xi|)^\vartheta
 \Phi(\deg P,\log H(P))\right]
\end{equation}
can hold uniformly for all positive integral \(\xi\), all nonzero values and
all fixed \(E\)-function systems if \(\vartheta<1\), when \(C\) and \(\Phi\)
are independent of \(\xi\).
\end{theorem}

\begin{proof}
The connection formula \eqref{eq:Ein-connection} cancels \(\gamma\),
\(\log x\), and the scale logarithms because
\[
 1-2+1=0,
 \qquad
 \log x-2\log(2x)+\log(4x)=0.
\]
This proves the identity in \eqref{eq:R-def}.  Using
\(E_1(x)=\int_1^\infty e^{-xt}\,dt/t\), its integrand numerator is, with
\(u=e^{-xt}\),
\[
 u-2u^2+u^4=u(1-u)(1-u-u^2)>0,
\]
because \(0<u\le e^{-1}<(\sqrt5-1)/2\).  This proves positivity.
Asymptotic formula \eqref{eq:R-asymptotic} follows from
\eqref{eq:E1-asymptotic}; the terms at \(2x\) and \(4x\) are exponentially
smaller.  Hence
\[
 -\log R(n)=n+\log n+o(1).
\]
Since \(n^\vartheta=o(n)\), \eqref{eq:R-sublinear} follows.  In
\eqref{eq:uniform-template}, \(P_0\) has fixed degree one and fixed height
two, so its right-hand exponent is \(O(n^\vartheta)\), contradicting
\eqref{eq:R-sublinear}.
Finally, \eqref{eq:R-asymptotic} gives
\[
 \frac{R(n)}{C_0e^{-\sigma n}n^{-\lambda}}
 =C_0^{-1}e^{-(1-\sigma)n}n^{\lambda-1}(1+o(1)),
\]
which tends to zero in each of the two stated cases and contradicts
\eqref{eq:critical-polynomial-template}.
\end{proof}

\begin{remark}[Exact scope of the no-go]
The theorem gives the necessary condition \(\vartheta\ge1\) under the
normalization \eqref{eq:uniform-template}, and at the critical exponential
scale it forces at least the polynomial loss \(n^{-1}\).  It does not prove
the existence or sufficiency of such a bound.  If arbitrary point-dependent
constants are allowed, or if the function system itself is changed with the
sample point, the exponent \(\vartheta\) alone is not an invariant.
\end{remark}

\begin{remark}[Augmentation zero versus one]
The fixed form \(P_0=(1,-2,1)\) has \(\omega(P_0)=0\) and zero logarithmic
moment, so it lies in the gamma-free stratum of
\eqref{eq:augmentation-decomposition}.  It disproves bounds that treat all
coefficient directions uniformly; by itself it does not disprove a theorem
restricted to \(\omega=1\).
\Cref{thm:augmentation-extremality,thm:denominator-transition}
describe what survives, and what fails, after that restriction.
\end{remark}

The same fixed second difference is built into the dyadic tower.

\begin{proposition}[Increment identity]
\label{prop:increment-identity}
For every \(N\ge1\),
\begin{equation}\label{eq:YN-increment}
 \boxed{\displaystyle
 Y_N-Y_{N+1}=(N+1)R(2^N).}
\end{equation}
Consequently
\begin{equation}\label{eq:epsilon-tail-R}
 \varepsilon_N=\sum_{k\ge N}(k+1)R(2^k).
\end{equation}
\end{proposition}

\begin{proof}
Expanding \eqref{eq:YN} gives
\[
 Y_N-Y_{N+1}
 =(N+1)\{\Ein(2^N)-2\Ein(2^{N+1})+\Ein(2^{N+2})\},
\]
which is \eqref{eq:YN-increment}.  Telescoping and
\(Y_N\to\gamma\) give \eqref{eq:epsilon-tail-R}.
\end{proof}

Thus the obstruction to a sublinear measure is not merely analogous to the
tower: it is its exact level increment.

We now formulate the remaining arithmetic question with a coefficient that
is fixed across every level.  For \(s\in\C\), define
\begin{equation}\label{eq:LN-def}
 L_N(s)=2sY_N^2+Y_N-2.
\end{equation}

\begin{theorem}[Common-coefficient critical-slope criterion]
\label{thm:critical-slope}
For \(s\in\C\), the limit
\[
 \lambda(s)=\lim_{N\to\infty}
 \frac{-\log|L_N(s)|}{2^N}
\]
exists whenever the logarithm is defined for all sufficiently large \(N\),
and
\begin{equation}\label{eq:lambda-dichotomy}
 \lambda(s)=
 \begin{cases}
 1,&s=S_2,\\
 0,&s\ne S_2.
 \end{cases}
\end{equation}
For \(s=S_2\), more precisely,
\begin{align}
 L_N(S_2)
 &=\frac{4-\gamma}{\gamma}\varepsilon_N
 +2S_2\varepsilon_N^2>0,\label{eq:LN-exact}\\
 -\log L_N(S_2)
 &=2^N+N\log2-\log(N+1)+O(1).
 \label{eq:LN-log-asymptotic}
\end{align}
Hence
\begin{equation}\label{eq:S2-slope-equivalence}
 \boxed{\displaystyle
 S_2\in\Qbar
 \quad\Longleftrightarrow\quad
 \exists s\in\Qbar\text{ with }\lambda(s)=1.}
\end{equation}
If \(s\in\Qbar\), the family \(\{L_N(s):N\ge1\}\) is countably
algebraically independent over \(K\).
\end{theorem}

\begin{proof}
Since \(2S_2\gamma^2+\gamma-2=0\), substituting
\(Y_N=\gamma+\varepsilon_N\) gives \eqref{eq:LN-exact}; combine this with
\eqref{eq:epsilon-properties} to get \eqref{eq:LN-log-asymptotic}.  If
\(s\ne S_2\), then
\[
 L_N(s)\longrightarrow2\gamma^2(s-S_2)\ne0,
\]
so the normalized logarithm tends to zero.  This proves
\eqref{eq:lambda-dichotomy} and \eqref{eq:S2-slope-equivalence}.

Finally, for algebraic \(s\), one has \(s\in K\) and
\[
 2sY_N^2+Y_N-(L_N(s)+2)=0.
\]
Thus \(K(Y_N:N\in I)\) is algebraic over
\(K(L_N(s):N\in I)\) for every finite set \(I\); when \(s=0\), the relation
is linear.  Algebraic independence of the \(Y_N\) therefore transfers to the
\(L_N(s)\).
\end{proof}

There is an equivalent bounded-distortion form that uses only finite-stage
gamma-free data.

\begin{corollary}[Bounded-distortion criterion]
\label{cor:bounded-distortion}
One has
\begin{equation}\label{eq:bounded-distortion-limit}
 \frac{L_N(S_2)}{Y_N-Y_{N+1}}
 \longrightarrow\frac{4-\gamma}{\gamma},
\end{equation}
whereas for \(s\ne S_2\),
\[
 \left|\frac{L_N(s)}{Y_N-Y_{N+1}}\right|\longrightarrow\infty.
\]
Consequently
\begin{equation}\label{eq:bounded-equivalence}
 S_2\in\Qbar
 \quad\Longleftrightarrow\quad
 \exists s\in\Qbar:\quad
 \left\{\frac{L_N(s)}{Y_N-Y_{N+1}}\right\}_{N\gg1}
 \text{ is bounded}.
\end{equation}
\end{corollary}

\begin{proof}
By \cref{prop:increment-identity} and \eqref{eq:R-asymptotic},
\[
 Y_N-Y_{N+1}\sim\frac{N+1}{2^N}e^{-2^N}
 \sim\varepsilon_N.
\]
Now use \eqref{eq:LN-exact}; if \(s\ne S_2\), the numerator tends to a
nonzero constant while the denominator tends to zero.
\end{proof}

\begin{remark}[Why a generic critical measure is not enough]
The fixed-form obstruction has size \(R(x)\sim e^{-x}/x\).  Under the
algebraicity hypothesis \(s=S_2\), the target at \(x=2^N\) has size
\[
 L_N(S_2)\sim C_\gamma(\log x)\frac{e^{-x}}x.
\]
It is therefore larger, by a logarithmic factor, than an unconditional
fixed-form example already permits.  A lower bound constrained only by this
worst-case pointwise scale cannot distinguish the two.  A surviving route
must use additional structure, for example the repetition of the \emph{same}
algebraic coefficient \(s\) over all levels.
\end{remark}

Existing quantitative \(E\)-function theorems cited here are fixed-point
theorems: the constants in the algebraic-independence measure depend on the
chosen nonzero algebraic evaluation point; see
Adamczewski--Faverjon \cite{AdamczewskiFaverjon2025} and
Fischler--Rivoal \cite{FischlerRivoal2024}.  They do not directly provide
the moving uniformity required by \cref{thm:critical-slope}.  This statement
is only about the scope of those theorems, not an impossibility claim for all
future methods.

\section{An exact asymmetric Sondow integral}
\label{sec:sondow}

One of the discarded barrier heuristics attempted to assign a universal
entropy floor to an asymmetric two-variable Sondow integral.  The correct
replacement is an exact finite formula.  For integers \(A\ge0\) and
\(B\ge1\), put
\begin{equation}\label{eq:Sondow-AB}
 \mathcal I_{A,B}=
 \int_0^1\!\!\int_0^1
 \frac{x^A(1-x)^B y^A(1-y)^B}
 {(1-xy)(-\log(xy))}\,dx\,dy.
\end{equation}

\begin{theorem}[Exact asymmetric evaluation]
\label{thm:Sondow-asymmetric}
One has
\begin{align}
 \mathcal I_{A,B}
 ={}&\binom{2B}{B}\gamma\notag\\
 &+2\sum_{j=0}^B\binom Bj^2
 (H_j-H_{B-j})\log((A+j)!)\notag\\
 &-\sum_{j=0}^B\binom Bj^2H_{A+j}.
 \label{eq:Sondow-exact}
\end{align}
In particular,
\begin{equation}\label{eq:Sondow-example}
 \mathcal I_{1,2}=6\gamma+3\log6-\frac{53}{6}.
\end{equation}
\end{theorem}

\begin{proof}
Use Tonelli's theorem,
\[
 -\frac1{\log(xy)}=\int_0^\infty(xy)^u\,du,
 \qquad
 \frac1{1-xy}=\sum_{k\ge0}(xy)^k,
\]
and the beta integral.  With
\[
 R_B(t)=\frac{B!}{t(t+1)\cdots(t+B)},
\]
this gives
\[
 \mathcal I_{A,B}
 =\sum_{m=A+1}^{\infty}\int_m^\infty R_B(t)^2\,dt.
\]
The rational kernel has the partial fraction identity
\begin{equation}\label{eq:Sondow-partial-fraction}
 \left\{\frac{B!}{t(t+1)\cdots(t+B)}\right\}^{\!2}
 =\sum_{j=0}^B\binom Bj^2
 \left\{\frac1{(t+j)^2}
 +\frac{2(H_j-H_{B-j})}{t+j}\right\}.
\end{equation}
Put \(a_j=\binom Bj^2\) and \(h_j=H_j-H_{B-j}\).  Integrating and summing
only up to \(m=M\) gives the finite identity
\begin{align*}
 \mathcal I_{A,B}^{(M)}
 ={}&\sum_{j=0}^Ba_j(H_{M+j}-H_{A+j})\\
 &-2\sum_{j=0}^Ba_jh_j
   \log\frac{(M+j)!}{(A+j)!}.
\end{align*}
The potentially divergent pieces cancel by
\[
 \sum_ja_j=\binom{2B}{B},\qquad
 \sum_ja_jh_j=0,
\]
together with the derivative of Vandermonde's identity
\[
 2\sum_jja_jh_j=\binom{2B}{B}.
\]
Indeed \(H_{M+j}=\log M+\gamma+o(1)\), while Stirling's formula gives
\[
 \log(M+j)!
 =M(\log M-1)+(j+\tfrac12)\log M
  +\tfrac12\log(2\pi)+o(1).
\]
Letting \(M\to\infty\) in the finite identity gives
\[
 \binom{2B}{B}\gamma
 -\sum_ja_jH_{A+j}
 +2\sum_ja_jh_j\log((A+j)!),
\]
which is \eqref{eq:Sondow-exact}.  Direct substitution gives
\eqref{eq:Sondow-example}.
\end{proof}

Formula \eqref{eq:Sondow-exact} replaces, rather than supports, any fixed
numerical entropy floor for this asymmetric family.  The remaining
Diophantine bottleneck is the transverse height of the prime-exponent vector
in the logarithmic kernel.  We make no novelty claim beyond this exact
evaluation without a separate priority comparison with the broader Sondow
literature \cite{Sondow2003}.

\section{The Pr\'evost grid and residual-prime saturation}
\label{sec:prevost-grid}

For \(0\le m\le n\), Pr\'evost's family \cite{Prevost2008} can be written
\begin{align}
 J_{n,m}&=\gamma+L_{n,m}-2H_n,\label{eq:prevost-form}\\
 L_{n,m}&=\sum_{k=0}^n(-1)^{n+k}
 \binom nk\binom{n+k}k\log(n-m+k+1).
 \label{eq:prevost-log}
\end{align}
Let \(d_n=\lcm(1,\ldots,n)\) and \(j_{n,m}=d_nJ_{n,m}\).  If
\(m/n\to\alpha\in(0,1)\), its saddle rate is
\begin{equation}\label{eq:prevost-rate}
 |J_{n,m}|=\exp\{(\log r(\alpha)+o(1))n\},
 \qquad
 r(\alpha)=
 \alpha^\alpha(1-\alpha)^{2-2\alpha}(2-\alpha)^{\alpha-2},
\end{equation}
up to a polynomial factor.  The unique optimum is
\begin{equation}\label{eq:prevost-optimum}
 \alpha_*=1-\frac1{\sqrt2},
 \qquad
 r(\alpha_*)=3-2\sqrt2,
 \qquad
 -\log r(\alpha_*)=\log(3+2\sqrt2).
\end{equation}

We first isolate an exact linear-algebra obstruction that is independent of
the saddle estimate.  Fix a finite set \(I\) of cells and collect the
normalized forms into
\begin{equation}\label{eq:prevost-vector}
 \mathbf j=\gamma\mathbf B-\mathbf A
 +\sum_{\ell\in\mathcal P_I}(\log\ell)\mathbf V_\ell\in\R^I,
\end{equation}
where \(\mathcal P_I\) is the finite set of primes occurring in the
logarithmic coordinates.  The vectors \(\mathbf B,\mathbf A,\mathbf V_\ell\)
are integral.  We write \(B(x)=\langle\mathbf B,x\rangle\), and similarly
for the other coordinate functionals.

\begin{theorem}[Residual-prime Hilbert saturation]
\label{thm:residual-prime-saturation}
Let \(\mathcal S\subset\mathcal P_I\), let \(a/b\in\Q\) be reduced with
\(b>0\), and put
\[
 E_{\mathcal S,a/b}
 =\bigcap_{\ell\notin\mathcal S}\ker V_\ell
  \cap\ker(aB-bA).
\]
If \(P\) is orthogonal projection onto this subspace, then
\begin{equation}\label{eq:residual-projection}
 \boxed{\displaystyle
 bP\mathbf j=(b\gamma-a)P\mathbf B
 +b\sum_{p\in\mathcal S}(\log p)P\mathbf V_p.}
\end{equation}

Let \(c=(c_p)_{p\in\mathcal S}\in\Z^{\mathcal S}\) be primitive and
nonzero.  Choose \(h_p\in\Z\) with \(\sum h_pc_p=1\), set
\[
 \mathbf T=\sum_{p\in\mathcal S}h_p\mathbf V_p,
 \qquad
 \alpha_c=\sum_{p\in\mathcal S}c_p\log p,
\]
write \(T(x)=\langle\mathbf T,x\rangle\), and define
\[
 E_{c,a/b}=E_{\mathcal S,a/b}
 \cap\bigcap_{p\in\mathcal S}\ker(V_p-c_pT).
\]
Denote its orthogonal projection by \(P_c\).
If \(\gamma=a/b\), then
\begin{equation}\label{eq:residual-direction}
 \boxed{P_c\mathbf j=\alpha_cP_c\mathbf T.}
\end{equation}
If \(\Lambda=E_{c,a/b}\cap\Z^I\) and \(T(\Lambda)=g\Z\), \(g>0\), then
the real quotient minimum
\[
 \lambda_\R=\inf\{\|x\|:x\in E_{c,a/b},\ T(x)=g\}
\]
satisfies
\begin{equation}\label{eq:residual-saturation}
 \boxed{\displaystyle
 \lambda_\R=\frac g{\|P_c\mathbf T\|},
 \qquad
 \lambda_\R\|P_c\mathbf j\|=g|\alpha_c|.}
\end{equation}
\end{theorem}

\begin{proof}
Project \eqref{eq:prevost-vector}.  All primes outside \(\mathcal S\)
vanish and \(aP\mathbf B=bP\mathbf A\), proving
\eqref{eq:residual-projection}.  The fixed-direction equations give
\(P_c\mathbf V_p=c_pP_c\mathbf T\), and
\eqref{eq:residual-direction} follows under \(\gamma=a/b\).
Finally, the restriction of \(T\) to \(E_{c,a/b}\) is represented by
\(P_c\mathbf T\).  Cauchy--Schwarz is sharp in its direction and gives
the first identity in \eqref{eq:residual-saturation}; the second follows
from \eqref{eq:residual-direction}.
\end{proof}

\begin{corollary}[Phase-blind quotient barrier]
\label{cor:phase-blind-quotient}
Under the hypotheses of the second part of
\cref{thm:residual-prime-saturation}, and under \(\gamma=a/b\),
\[
 \lambda_\R\|\mathbf j\|\ge g|\alpha_c|.
\]
Thus a strict inequality in the opposite direction cannot follow from this
real Euclidean quotient identity and Cauchy--Schwarz alone.  Integral
divisibility or genuinely lattice-theoretic information is not excluded.
Any strict reverse inequality would already imply \(\gamma\ne a/b\).
\end{corollary}

\begin{remark}[The one-prime shear]
Let \(P\) now denote orthogonal projection onto
\(\bigcap_{\ell\ne2}\ker V_\ell\), with no candidate rationality
hyperplane imposed, and suppose \(x=P\mathbf V_2\ne0\).  Put
\[
 z=P(\gamma\mathbf B-\mathbf A),
 \qquad \varepsilon=\|P\mathbf j\|.
\]
Then \(P\mathbf j=z+(\log2)x\), and hence
\begin{equation}\label{eq:log2-shear}
 \frac{\det\operatorname{Gram}(x,z)}{\|x\|^2}
 =\|P_{x^\perp}P\mathbf j\|^2\le\varepsilon^2,
 \qquad
 \frac{\langle x,z\rangle}{\|x\|^2}
 =-\log2+O\!\left(\frac\varepsilon{\|x\|}\right).
\end{equation}
Thus \(-\log2\) is the longitudinal shear coefficient, while
\(P_{x^\perp}P\mathbf j\) is the narrow transverse residual.  For an
external proxy \(\widetilde\gamma=a/b\ne\gamma\),
\[
 P(\widetilde\gamma\mathbf B-\mathbf A)
 =z+(\widetilde\gamma-\gamma)P\mathbf B.
\]
More explicitly, if an integral coefficient vector \(c\) kills the
odd-prime labels and has residual label \(k=V_2(c)\), then
\[
 (aB(c)-bA(c))+bk\log2
 =b\langle\mathbf j,c\rangle+(a-b\gamma)B(c).
\]
Thus a proxy form is not controlled by the analytic norm of \(P\mathbf j\)
unless \(a/b=\gamma\); finite-window diagnostics based on an external proxy
cannot be used as irrationality evidence.
\end{remark}

\subsection{The fixed-place Ridout deficit}

For a fixed finite set \(S\) of rational primes and \(0\ne r\in\Z\), write
\[
 r_S=\prod_{p\in S}p^{v_p(r)},
 \qquad r_{S^c}=|r|/r_S.
\]

\begin{proposition}[Fixed-\(S\) Ridout budget]
\label{prop:ridout-budget}
Let \(\xi\) be a real algebraic irrational number.  For every
\(\epsilon>0\), all but finitely many reduced fractions \(a/b\) with
\(b>0\), \(a\ne0\), and \(|\xi-a/b|<1\) satisfy
\begin{equation}\label{eq:ridout-linear-form}
 |b\xi-a|\ge\frac{a_S}{b_{S^c}b^\epsilon}.
\end{equation}
Consequently, suppose along a sequence that
\[
 \log b_N=hN+o(N),
 \quad
 \log(a_N)_S+\log(b_N)_S=\sigma N+o(N),
 \quad
 |b_N\xi-a_N|=e^{-(\beta+o(1))N}.
\]
Assume here that \(h>0\), that \(b_N\to\infty\), and that the reduced
fractions \(a_N/b_N\) take infinitely many distinct values.
If \(\beta>h-\sigma\), then \(\xi\) is not algebraic irrational.  If
\(\beta\le h-\sigma\), these rates alone do not enter the range excluded by
Ridout's theorem.
\end{proposition}

\begin{proof}
Ridout's theorem \cite{Ridout1958} says that
\[
 \prod_{p\in S}|ab|_p
 \min\{1,|\xi-a/b|\}<b^{-2-\epsilon}
\]
has only finitely many reduced solutions.  For a good approximant the left
side is \(|\xi-a/b|/(a_Sb_S)\).  Multiplying the complementary inequality
by \(b\) gives \eqref{eq:ridout-linear-form}; in the homogeneous-height
form of the theorem, \(|\xi-a/b|<1\) makes
\(\max(|a|,b)\ll_\xi b\), and applying the theorem with
\(\epsilon/2\) absorbs the fixed comparison constant for large \(b\).
Taking logarithms and then letting \(\epsilon\downarrow0\) proves the rate
assertion.
\end{proof}

\begin{theorem}[The Pr\'evost saddle--lcm ledger]
\label{thm:prevost-ridout-deficit}
For any integer sequence \(m_n\) with \(m_n/n\to\alpha_*\),
\[
 |d_nJ_{n,m_n}|=
 \exp\{-(\log(3+2\sqrt2)-1+o(1))n\}.
\]
For every fixed finite set \(S\),
\[
 \log(d_n)_S=|S|\log n+O_S(1)=o(n),
 \qquad
 \log d_n=n+o(n).
\]
Thus, in the favorable height-one ledger in which the lcm is the only
denominator cost and a fixed \(S\) contributes no additional exponential
divisibility, the analytic exponent falls short of the Ridout threshold by
\begin{equation}\label{eq:ridout-deficit}
 \boxed{\displaystyle
 \Delta_R=2-\log(3+2\sqrt2)
 =0.2372528259\ldots.}
\end{equation}
This is an exact saddle--denominator deficit, not yet a rational
approximation to \(\gamma\), because logarithmic coordinates remain in
\(J_{n,m_n}\).
\end{theorem}

\begin{proof}
Combine \eqref{eq:prevost-rate} with the prime-number theorem form
\(\log d_n=\psi(n)=n+o(n)\).  Differentiating \(\log r\) gives
\[
 (\log r)'(\alpha)
 =\log\frac{\alpha(2-\alpha)}{(1-\alpha)^2},
\]
which proves the unique optimum \eqref{eq:prevost-optimum}.  Finally
\(v_p(d_n)=\lfloor\log n/\log p\rfloor\), so every fixed \(S\) contributes
only \(o(n)\).  Comparing the resulting exponent
\(\log(3+2\sqrt2)-1\) with the height-one threshold (1) gives
\eqref{eq:ridout-deficit}.
\end{proof}

\begin{corollary}[Conditional rational-extraction budget]
Suppose an elimination of the logarithmic coordinates produces infinitely
many distinct reduced fractions \(a_n/b_n\) such that
\[
 \log b_n=n+o(n),\qquad
 \log(a_n)_S+\log(b_n)_S=o(n),
\]
and preserves the saddle rate
\[
 |b_n\gamma-a_n|
 =\exp\{-(\log(3+2\sqrt2)-1+o(1))n\}.
\]
Then \cref{prop:ridout-budget} does not exclude algebraic irrationality of
\(\gamma\).  Within this height-one template, crossing the Ridout threshold
requires an additional fixed-\(S\) valuation density greater than
\(\Delta_R\), an archimedean gain of the same size, or a corresponding
reduction in height.
\end{corollary}

\begin{remark}[Exact scope]
This is a budget theorem, not an exclusion of every linear combination of
Pr\'evost forms.  Further cancellation could increase the archimedean
exponent, and a new automatic divisibility theorem could increase the
\(S\)-part.  Neither input is presently available.  Enlarging a fixed finite
set \(S\) alone cannot help because \(\log(d_n)_S=o(n)\).
\end{remark}

\begin{proposition}[Bounded-degree logarithmic transfer]
\label{prop:bounded-degree-transfer}
After cancelling all odd-prime coordinates, write a resulting form as
\[
 J=B\gamma-A+k\log2,
 \qquad A,B,k\in\Z,
\]
and assume \(k\ne0\).  If \(\gamma\) is algebraic of degree \(d\), then
\[
 \beta=\frac{A-B\gamma}{k}
\]
has degree at most \(d\), naïve height at most
\(C_\gamma\max(1,|A|,|B|,|k|)^d\), and
\[
 |\log2-\beta|=\frac{|J|}{|k|}.
\]
Thus any such attack is bounded by the degree-\(d\) algebraic approximation
exponent of \(\log2\), with an additional factor \(d\) in the height
exponent.
\end{proposition}

\begin{proof}
If \(f\) is the minimal polynomial of \(\gamma\) and \(B\ne0\), then
\(\beta\) is a zero of
\[
 B^df\!\left(\frac{A-kX}{B}\right),
\]
and the affine transformation shows that its degree is exactly \(d\).
The displayed polynomial has coefficient height
\(O_\gamma(\max(1,|A|,|B|,|k|)^d)\); the standard factor-height bound
gives the same estimate, up to a constant depending on \(\gamma\), for the
minimal polynomial of \(\beta\).  The case \(B=0\) is rational, and the
approximation identity is immediate.
\end{proof}

\begin{remark}
The available bounds are
\[
 \mu(\log2)\le3.574553902525\ldots,
 \qquad
 \mu_2(\log2)\le12.841618132152\ldots;
\]
see \cite{Marcovecchio2009,Marcovecchio2025}.  Thus a rational branch must
beat the first exponent, while the degree-two height transfer makes
\(2\mu_2(\log2)=25.683236264304\ldots\) the corresponding coefficient-height
target.  These are target exponents once the growth of
\(\max(1,|A|,|B|,|k|)\) is known; no comparison with a projected Pr\'evost
family is claimed here without a separate height lemma.
\end{remark}

\section{Fixed-prime Rodrigues twists and universal saddle cancellation}
\label{sec:fixed-prime-rodrigues}

The deficit \eqref{eq:ridout-deficit} suggests inserting a factor
\(p^{cn}\), for one fixed prime \(p\), directly into the Pr\'evost kernel.
The following deformation preserves integrality, degree, and a linear number
of moments, and produces this factor in the endpoint coefficient.  Even if
the whole endpoint factor is favorably credited as nonarchimedean gain, the
real saddle cancels it in the precise sense of
\cref{thm:universal-fixed-p-cancellation}.

Fix a prime \(p\).  For \(n\ge1\), \(1\le\ell\le n\), put
\begin{equation}\label{eq:fixed-p-rodrigues}
 \mathcal P_{n,\ell,p}(x)=
 \frac{(-1)^\ell}{\ell!}\frac{d^\ell}{dx^\ell}
 \left[x^\ell(1-x)^\ell
 \bigl(1+(p-1)x\bigr)^{n-\ell}\right].
\end{equation}
Write \(\mathcal P_{n,\ell,p}(x)=\sum_{k=0}^nc_kx^k\).

\begin{proposition}[Arithmetic and orthogonality]
\label{prop:fixed-p-arithmetic}
One has \(\mathcal P_{n,\ell,p}\in\Z[x]\),
\(\deg\mathcal P_{n,\ell,p}=n\),
\[
 \int_0^1x^j\mathcal P_{n,\ell,p}(x)\,dx=0
 \qquad(0\le j<\ell),
\]
and
\begin{equation}\label{eq:fixed-p-endpoint}
 \boxed{\mathcal P_{n,\ell,p}(1)=p^{n-\ell}.}
\end{equation}
\end{proposition}

\begin{proof}
Expanding before differentiation gives
\[
 c_k=(-1)^\ell\binom{\ell+k}{\ell}
 \sum_{a+b=k}(-1)^a
 \binom\ell a\binom{n-\ell}b(p-1)^b,
\]
which proves integrality and the degree assertion.  Integration by parts
\(\ell\) times proves the moment identities; all boundary terms vanish.
At \(x=1\), only the term in which all derivatives fall on
\((1-x)^\ell\) survives, giving \eqref{eq:fixed-p-endpoint}.
\end{proof}

Put
\[
 W(x)=\frac1{\log x}+\frac1{1-x}
\]
and, for \(0\le m\le\ell\), define
\begin{equation}\label{eq:fixed-p-moment}
 \mathcal J_{n,\ell,m,p}=
 \int_0^1x^{\ell-m}\mathcal P_{n,\ell,p}(x)W(x)\,dx.
\end{equation}

\begin{proposition}[Exact logarithmic linear form]
\label{prop:fixed-p-linear-form}
Let
\[
 Q_{n,\ell,p}=p^{n-\ell},
 \qquad
 A_{n,\ell,p}=\sum_{k=0}^nc_kH_k,
\]
and
\[
 L_{n,\ell,m,p}=
 \sum_{k=0}^nc_k\log(\ell-m+k+1).
\]
Then
\begin{equation}\label{eq:fixed-p-exact-form}
 \boxed{\mathcal J_{n,\ell,m,p}
 =Q_{n,\ell,p}\gamma-A_{n,\ell,p}+L_{n,\ell,m,p},}
\end{equation}
and \(d_nA_{n,\ell,p}\in\Z\).
\end{proposition}

\begin{proof}
Use Euler's integral \(\gamma=\int_0^1W(x)\,dx\) and Frullani's identity
\[
 \int_0^1\frac{x^r-1}{\log x}\,dx=\log(r+1).
\]
The rational part is
\begin{align*}
 &\int_0^1\frac{\mathcal P(1)-x^{\ell-m}\mathcal P(x)}{1-x}\,dx\\
 &\quad=\int_0^1\frac{\mathcal P(1)-\mathcal P(x)}{1-x}\,dx
 +\int_0^1\mathcal P(x)\frac{1-x^{\ell-m}}{1-x}\,dx.
\end{align*}
The second integral vanishes by the moment conditions, and the first is
\(\sum c_kH_k\).  This proves \eqref{eq:fixed-p-exact-form}; its denominator
divides \(d_n\).
\end{proof}

Pr\'evost's positive Markov--Stieltjes representation is
\begin{equation}\label{eq:prevost-markov}
 W(u)=\int_0^1\frac{\omega(t)}{1-(1-u)t}\,dt,
 \qquad
 \omega(t)=\frac1{t\{\log^2(t^{-1}-1)+\pi^2\}}>0.
\end{equation}

\begin{proposition}[Positive double integral]
\label{prop:fixed-p-positive}
One has
\begin{align}
 (-1)^m\mathcal J_{n,\ell,m,p}
 &=\int_0^1\!\int_0^1
 \frac{u^\ell(1-u)^\ell
 (1+(p-1)u)^{n-\ell}
 (1-t)^{\ell-m}t^m}
 {(1-(1-u)t)^{\ell+1}}
 \omega(t)\,dt\,du>0.
 \label{eq:fixed-p-double}
\end{align}
\end{proposition}

\begin{proof}
Insert \eqref{eq:prevost-markov}, use Fubini, and integrate by parts
\(\ell\) times.  The polynomial part is annihilated, while
\[
 \frac{d^\ell}{du^\ell}
 \frac{u^{\ell-m}}{1-(1-u)t}
 =(-1)^m\ell!\frac{(1-t)^{\ell-m}t^m}
 {(1-(1-u)t)^{\ell+1}}.
\]
This gives \eqref{eq:fixed-p-double} and positivity.
\end{proof}

Let \(\theta=\ell/n\), \(\alpha=m/\ell\), and set
\begin{align}
 \phi_\alpha(u,t)&=
 \frac{u(1-u)(1-t)^{1-\alpha}t^\alpha}
 {1-(1-u)t},\label{eq:prevost-saddle-kernel}\\
 \Phi_{\theta,\alpha,p}(u,t)&=
 \phi_\alpha(u,t)^\theta
 (1+(p-1)u)^{1-\theta},\label{eq:fixed-p-saddle}\\
 R_{\theta,\alpha,p}&=-\log\max_{0\le u,t\le1}
 \Phi_{\theta,\alpha,p}(u,t).
\end{align}
We extend \(\phi_\alpha\), and hence
\(\Phi_{\theta,\alpha,p}\), continuously by zero at the boundary corner
\((u,t)=(0,1)\).  Indeed the quotient in
\eqref{eq:prevost-saddle-kernel} tends to zero there for
\(0<\alpha<1\).

For the finite parameters \(\theta_n=\ell/n\),
\(\alpha_n=m/\ell\), the positive double integral is
\[
 (-1)^m\mathcal J_{n,\ell,m,p}
 =\int_0^1\!\int_0^1
 \Phi_{\theta_n,\alpha_n,p}(u,t)^n
 \frac{\omega(t)}{1-(1-u)t}\,dt,du.
\]
The remaining measure has finite mass
\[
 \int_0^1\!\int_0^1
 \frac{\omega(t)}{1-(1-u)t}\,dt,du
 =\int_0^1W(u)\,du=\gamma.
\]
This gives the upper Laplace bound, while a small rectangle around an
interior maximizer gives the matching lower bound.  Consequently,
whenever \(\ell/n\to\theta\in(0,1]\) and
\(m/\ell\to\alpha\in(0,1)\),
\begin{equation}\label{eq:fixed-p-laplace}
 \lim_{n\to\infty}|\mathcal J_{n,\ell,m,p}|^{1/n}
 =e^{-R_{\theta,\alpha,p}}.
\end{equation}

\begin{theorem}[Universal fixed-prime saddle cancellation]
\label{thm:universal-fixed-p-cancellation}
For every prime \(p\), \(0<\theta\le1\), and \(0<\alpha<1\),
\begin{equation}\label{eq:universal-fixed-p-bound}
 \boxed{\displaystyle
 R_{\theta,\alpha,p}+(1-\theta)\log p
 \le\log(3+2\sqrt2).}
\end{equation}
Thus no member of \eqref{eq:fixed-p-rodrigues} improves the original
Pr\'evost saddle after the complete endpoint factor
\(p^{n-\ell}\) is favorably credited as fixed-\(p\) gain.
\end{theorem}

\begin{proof}
Since
\[
 e^{-\{R_{\theta,\alpha,p}+(1-\theta)\log p\}}
 =\max_{u,t}
 \left[
  \phi_\alpha(u,t)^\theta
  \left(\frac{1+(p-1)u}{p}\right)^{1-\theta}
 \right],
\]
it is enough to bound this normalized maximum from below.  Now
For \(0\le u\le1\),
\[
 \frac{1+(p-1)u}{p}=u+\frac{1-u}{p}\ge u.
\]
At \(t=1/2\), the dependence on \(\alpha\) disappears:
\[
 \phi_\alpha\!\left(u,\frac12\right)
 =\frac{u(1-u)}{1+u}.
\]
Consequently the normalized maximum is at least
\[
 \max_{0\le u\le1}
 u\left(\frac{1-u}{1+u}\right)^\theta
 \ge\max_{0\le u\le1}\frac{u(1-u)}{1+u}
 =3-2\sqrt2.
\]
Taking negative logarithms proves \eqref{eq:universal-fixed-p-bound}.
\end{proof}

\begin{corollary}[The fixed-place deficit survives]
\label{cor:fixed-p-deficit}
Even under the favorable assumptions that logarithmic coordinates are
removed at no exponential cost, \(d_n=e^{n+o(n)}\) is the only denominator
cost, and the full \(p^{n-\ell}\) factor survives primitive reduction, the
integerized exponent satisfies
\[
 \bigl(R_{\theta,\alpha,p}-1\bigr)+(1-\theta)\log p
 \le\log(3+2\sqrt2)-1<1.
\]
Thus, even after the entire endpoint factor is credited, this favorable
ledger does not cross the fixed-\(S\) Ridout threshold; its remaining gap is
at least \(\Delta_R\) from \eqref{eq:ridout-deficit}.
\end{corollary}

\begin{proposition}[Balanced factorial ratios have zero fixed-prime slope]
\label{prop:balanced-factorial-fixed-p}
Let
\[
 \mathcal R(n,k)=
 \frac{\prod_{i=1}^rL_i(n,k)!}{\prod_{j=1}^sM_j(n,k)!},
\]
where the integral linear forms are nonnegative and \(O(n)\) on the index
range, and suppose \(\sum_iL_i=\sum_jM_j\).  Then for every fixed prime
\(p\),
\[
 \lvert v_p(\mathcal R(n,k))\rvert=O(\log n)
\]
uniformly wherever the ratio is defined.  In particular its fixed-prime
valuation density is zero.
\end{proposition}

\begin{proof}
Legendre's formula gives a sum over \(p^\nu\).  At each level the linear
parts cancel by balance, leaving a bounded floor error; only
\(O(\log n)\) levels are nonzero.
\end{proof}

\begin{remark}
This covers the standard balanced Ap\'ery, Franel, Domb, and Christol
factorial ratios.  The valuation exponents supplied directly by the
digit-count and \(p\)-Lucas mechanisms of \cite{Delaygue2018} involve at
most \(1+\lfloor\log_p n\rfloor\) base-\(p\) digits and therefore provide
only an \(O(\log n)\) guaranteed gain.  This does not exclude additional
problem-specific cancellation inside a multisum.  If a factorial-balance
defect is \(\delta n+O(1)\), \(\delta\ne0\), Stirling's formula produces a
\(\delta n\log n\) term in the archimedean ledger.  A successful kernel
must therefore use information not already encoded by the balanced
factorial ratio or the positive endpoint weight twist analyzed here.
\end{remark}

\section{The Stokes-framing barrier}
\label{sec:stokes-framing}

The preceding barriers distinguish gamma-free directions from the Euler
extension.  The analogous distinction for an elementary differential
operator annihilating the Euler \(E\)-function is exact.

Put
\[
 \mathscr E(z)=\sum_{n\ge1}\frac{z^n}{n\,n!},
 \qquad F(x)=\mathscr E(-x)=-\Ein(x),
\]
and let \(\partial=\partial_x\).  Consider
\begin{equation}\label{eq:euler-E-operator}
 \mathcal L=(\partial+1)\circ\partial\circ x\circ\partial,
 \qquad
 \mathcal Ly=xy'''+(x+2)y''+y'.
\end{equation}

\begin{theorem}[Euler Stokes-framing barrier]
\label{thm:stokes-framing}
On \(\Omega=\C\setminus(-\infty,0]\), with principal branches,
\[
 \mathcal B_\infty=(1,\log x,E_1(x)),
 \qquad
 \mathcal B_0=(1,F(x),\log x)
\]
are solution bases of \(\mathcal L\), and
\begin{equation}\label{eq:C-gamma}
 \begin{pmatrix}1\\F(x)\\\log x\end{pmatrix}
 =C_\gamma
 \begin{pmatrix}1\\\log x\\E_1(x)\end{pmatrix},
 \qquad
 C_\gamma=
 \begin{pmatrix}
 1&0&0\\
 -\gamma&-1&-1\\
 0&1&0
 \end{pmatrix}.
\end{equation}
The complete upper/lower lateral transition across the negative real axis is
\begin{equation}\label{eq:stokes-transition}
 \begin{pmatrix}1\\\log_+x\\E_{1,+}(x)\end{pmatrix}
 =S
 \begin{pmatrix}1\\\log_-x\\E_{1,-}(x)\end{pmatrix},
 \qquad
 S=
 \begin{pmatrix}
 1&0&0\\
 2\pi i&1&0\\
 -2\pi i&0&1
 \end{pmatrix}.
\end{equation}
In particular \(S\) is independent of \(\gamma\), and the coefficient row
\(c_\gamma=(-\gamma,-1,-1)\) satisfies
\begin{equation}\label{eq:stokes-cocycle-blind}
 \boxed{c_\gamma S=c_\gamma}
\end{equation}
for every complex value of \(\gamma\).

For every \(t\in\Qbar\), \(F+t\) is an \(E\)-function solution of the same
operator.  With \(\mathcal B_0^{(t)}=(1,F+t,\log x)\), its connection
matrix is
\begin{equation}\label{eq:framing-shear}
 C_t=U_tC_\gamma,
 \qquad
 U_t=\begin{pmatrix}1&0&0\\t&1&0\\0&0&1\end{pmatrix}.
\end{equation}
Consequently, under \(\gamma\in\Qbar\), the choice \(t=\gamma\) removes
the Euler entry by an algebraic change of Frobenius framing without changing
the operator or any unframed Stokes conjugacy class.
\end{theorem}

\begin{proof}
Since \(xF'(x)=e^{-x}-1\),
\[
 (\partial+1)\partial(xF')=(\partial+1)(-e^{-x})=0.
\]
Moreover,
\[
 W(1,\log x,E_1(x))=\frac{e^{-x}}{x^2}\ne0,
\]
and \(\det C_\gamma=1\), so both displayed triples are indeed bases.
The connection identity
\[
 F(x)=-\gamma-\log x-E_1(x)
\]
proves \eqref{eq:C-gamma}.  For \(r>0\),
\[
 \log_+(-r)-\log_-(-r)=2\pi i,
 \qquad
 E_{1,+}(-r)-E_{1,-}(-r)=-2\pi i,
\]
which gives \eqref{eq:stokes-transition}.  Matrix multiplication proves
\eqref{eq:stokes-cocycle-blind}.  Adding the constant solution \(t\) adds
\(t\) times the first row of the connection matrix to its second row,
proving \eqref{eq:framing-shear}.
\end{proof}

\begin{corollary}[Exterior blindness]
\label{cor:stokes-exterior-blind}
One has
\[
 \det C_\gamma=1,
 \qquad \det S=1,
 \qquad \operatorname{tr}S=3.
\]
The displayed lateral transition and every invariant of its unframed
conjugacy class are independent of \(\gamma\).  Their determinant, trace,
and exterior-power characteristic data lose the off-diagonal Euler framing.
\end{corollary}

\begin{remark}[Strict factors and framed Stokes constants]
Fischler--Rivoal use ``Stokes constants'' for sector-dependent coefficients
of a distinguished \(E\)-function in a formal basis at infinity; in that
broader sense \(-\gamma\) is such a coefficient.  After the logarithmic
formal monodromy is separated, the strict factor is
\(E_1\mapsto E_1-2\pi i\) and is gamma-free.  Their arithmetic theorem is a
containment theorem for possible coefficients, not a numerical-to-motivic
separation theorem; see \cite{FischlerRivoal2016}.
\end{remark}

\begin{remark}[Class-group averaging does not erase \(\gamma\)]
\label{rem:Chowla-Selberg-correction}
Suppose a Kronecker--Chowla--Selberg family is written in class-group
components as \(A_C=\kappa\gamma+B_C\), with the same Euler coefficient in
each class.  Character projection gives
\[
 \sum_C\overline{\chi(C)}A_C
 =\kappa\gamma\sum_C\overline{\chi(C)}
 +\sum_C\overline{\chi(C)}B_C.
\]
Thus the trivial character, equivalently the class-group average, preserves
the Euler term.  It disappears only in nontrivial character components or
other coefficient-sum-zero directions.  In representation-theoretic
language, \(\gamma\) occupies the trivial isotypic component.  This corrects
the opposite averaging heuristic; it is not a new Chowla--Selberg theorem.
See \cite{GrossCS}.
\end{remark}

\section{The mixed-Gevrey aperture}
\label{sec:mixed-gevrey}

Fischler--Rivoal call a formal Laurent series mixed arithmetic-Gevrey when
its nonnegative part is an \(E\)-function and its negative part is an
arithmetic Gevrey series of order one
\cite{FischlerRivoalMixed2024}.  The Euler connection formula produces a
particularly transparent family in this class.

For \(\lambda\in\Qbar_{>0}\), set
\begin{equation}\label{eq:mixed-dilation-family}
 \widehat D_\lambda(z)=
 \sum_{n\ge0}(-1)^n n!\lambda^{-n-1}z^{-n-1},\qquad
 U_\lambda(z)=e^{\lambda z},\qquad
 W_\lambda(z)=e^{\lambda z}\Ein(\lambda z)-\widehat D_\lambda(z).
\end{equation}

\begin{proposition}[Dilation system and functional rank]
\label{prop:mixed-dilation-rank}
Direction zero is not anti-Stokes for \(\widehat D_\lambda\), and its
Borel sum is
\begin{equation}\label{eq:Dhat-Borel-sum}
 (\widehat D_\lambda)_0(z)=e^{\lambda z}E_1(\lambda z).
\end{equation}
Consequently
\begin{equation}\label{eq:Wlambda-sum}
 (W_\lambda)_0(z)=e^{\lambda z}
 (\gamma+\log\lambda+\log z),
\end{equation}
on the principal positive sector, and
\begin{equation}\label{eq:mixed-dilation-system}
 \binom{U_\lambda}{W_\lambda}'
 =\begin{pmatrix}\lambda&0\\z^{-1}&\lambda\end{pmatrix}
 \binom{U_\lambda}{W_\lambda}.
\end{equation}
For distinct \(\lambda_1,\ldots,\lambda_m\in\Qbar_{>0}\), the \(2m\)
mixed functions \(U_{\lambda_i},W_{\lambda_i}\) are linearly independent
over \(\C(z)\).
\end{proposition}

\begin{proof}
The Laplace integral of the formal negative part is
\[
 \int_0^\infty\frac{e^{-\lambda zt}}{1+t}\,dt
 =e^{\lambda z}E_1(\lambda z),
\]
which proves \eqref{eq:Dhat-Borel-sum}; its Borel singularity lies on the
negative direction.  Subtracting this identity from
\(e^{\lambda z}\Ein(\lambda z)\) and using
\eqref{eq:Ein-connection} proves \eqref{eq:Wlambda-sum}.  Direct
differentiation proves \eqref{eq:mixed-dilation-system}.

Suppose a rational-function linear relation among all the displayed
functions exists, and sum it in direction zero.  Analytic continuation once
around zero adds \(2\pi i e^{\lambda_i z}\) to
\((W_{\lambda_i})_0\) and fixes \(U_{\lambda_i}\).  Subtraction gives a
rational-function relation among exponentials with distinct exponents, so
all coefficients of the \(W\)-terms vanish.  The remaining exponential
relation then forces all \(U\)-coefficients to vanish.
\end{proof}

\begin{proposition}[Unconditional non-divisibility]
\label{prop:mixed-nondivisibility}
For every \(\lambda\in\Qbar_{>0}\) and \(a\in\Qbar\), the formal quotient
\begin{equation}\label{eq:mixed-quotient}
 \frac{W_\lambda-aU_\lambda}{z-1}
\end{equation}
is not a mixed arithmetic-Gevrey function.
\end{proposition}

\begin{proof}
Write the normalized positive part of the numerator as
\[
 F(z)=e^{\lambda z}\{\Ein(\lambda z)-a\}
\]
and its negative part as \(-\widehat D_\lambda(z)\), which has zero constant
term.  If the quotient were mixed, comparison of nonnegative and
nonpositive Laurent parts after multiplication by \(z-1\) would force
\(F(1)\in\Qbar\).  But
\[
 F(1)=e^\lambda\{\Ein(\lambda)-a\},
\]
which is transcendental because \(\Ein(\lambda)\) is algebraically
independent over \(K\supset\Qbar(e^\lambda)\) by
\cref{thm:input-independence}.  This is a contradiction.
\end{proof}

We can now state the exact conditional payoff.  Write
\(\mathbf E,\mathbf D\) for the rings of algebraic-point values of
\(E\)-functions and of directionally summed arithmetic Gevrey-one series,
respectively, in the notation of Fischler--Rivoal.

\begin{theorem}[Conditional simultaneous dilation theorem]
\label{thm:conditional-mixed-dilations}
Assume the intersection conjecture
\begin{equation}\label{eq:ED-intersection}
 \mathbf E\cap\mathbf D=\Qbar.
\end{equation}
For any distinct
\(\lambda_1,\ldots,\lambda_m\in\Qbar_{>0}\), the \(2m\) numbers
\begin{equation}\label{eq:mixed-dilation-values}
 \boxed{\displaystyle
 e^{\lambda_i},\qquad
 e^{\lambda_i}(\gamma+\log\lambda_i)
 \quad(1\le i\le m)}
\end{equation}
are linearly independent over \(\Qbar\).  In particular every
\(\gamma+\log\lambda_i\) is transcendental; the case \(\lambda=1\) gives
the conditional transcendence of \(\gamma\).
\end{theorem}

\begin{proof}
Assume that the values in \eqref{eq:mixed-dilation-values} satisfy a
linear relation over \(\Qbar\).  Using
\(W_\lambda=e^{\lambda z}\Ein(\lambda z)-\widehat D_\lambda\), move all
summed Gevrey terms to the right.  The left side is in \(\mathbf E\), the
right side is in \(\mathbf D\), so their common value is algebraic by
\eqref{eq:ED-intersection}.  The left side has the form
\[
 \sum_i a_i e^{\lambda_i}
 +\sum_i b_i e^{\lambda_i}\Ein(\lambda_i).
\]
The \(\Ein(\lambda_i)\) are algebraically independent over the field \(K\)
containing all \(e^{\lambda_i}\), so algebraicity forces every \(b_i=0\).
The Lindemann--Weierstrass theorem applied to
\(0,\lambda_1,\ldots,\lambda_m\) then forces every \(a_i=0\).  This proves
linear independence.  The final transcendence assertion follows because
an algebraic ratio would make the two values in its block linearly
dependent.
\end{proof}

The dilation construction has an all-order resonant extension.  Write
\(\stirling{n}{r}\) for the unsigned Stirling number of the first kind, and for
\(r\ge1\) put
\begin{equation}\label{eq:all-order-Dhat}
 \widehat D_{\lambda,r}(z)=
 \sum_{n\ge0}(-1)^n\stirling{n+1}{r}
 \lambda^{-n-1}z^{-n-1},
\end{equation}
and
\begin{equation}\label{eq:all-order-mixed-family}
 \mathcal W_{\lambda,0}(z)=e^{\lambda z},\qquad
 \mathcal W_{\lambda,r}(z)
 =e^{\lambda z}\Ein_r(\lambda z)-\widehat D_{\lambda,r}(z)
 \quad(r\ge1).
\end{equation}

\begin{proposition}[All-order mixed jet]
\label{prop:all-order-mixed-jet}
Every \(\mathcal W_{\lambda,r}\) is mixed arithmetic-Gevrey, direction zero
is non-anti-Stokes, and
\begin{equation}\label{eq:all-order-mixed-generating}
 \boxed{\displaystyle
 \sum_{r\ge0}(\mathcal W_{\lambda,r})_0(z)t^r
 =e^{\lambda z}\Gamma(1-t)(\lambda z)^t.}
\end{equation}
They satisfy the triangular differential system
\begin{equation}\label{eq:all-order-mixed-system}
 \mathcal W_{\lambda,0}'=\lambda\mathcal W_{\lambda,0},
 \qquad
 \mathcal W_{\lambda,r}'
 =\lambda\mathcal W_{\lambda,r}+z^{-1}\mathcal W_{\lambda,r-1}
 \quad(r\ge1).
\end{equation}
For distinct positive algebraic \(\lambda_i\), every finite family of the
functions \(\mathcal W_{\lambda_i,r}\) is linearly independent over
\(\C(z)\).
\end{proposition}

\begin{proof}
The normalized Borel transform of \eqref{eq:all-order-Dhat} is
\[
 \frac{(-1)^{r-1}}{\lambda(r-1)!}
 \frac{\log^{r-1}(1+\xi/\lambda)}{1+\xi/\lambda},
\]
a \(G\)-function whose singularities lie on the negative ray.  Hence the
formal series is arithmetic Gevrey of order one and direction zero is
admissible.  The elementary split
\[
 \Gamma(1-t)x^t
 =1+t\sum_{r\ge1}\Ein_r(x)t^{r-1}
   -t x^t\Gamma(-t,x)
\]
together with the large-
\(x\) expansion
\[
 e^x x^t\Gamma(-t,x)
 \sim\sum_{n\ge0}(-1)^n(t+1)_n x^{-n-1}
\]
proves \eqref{eq:all-order-mixed-generating}, since
\([t^{r-1}](t+1)_n=\stirling{n+1}{r}\).  Differentiating the generating
identity and comparing coefficients proves
\eqref{eq:all-order-mixed-system}.

The sectorial sum on the right of
\eqref{eq:all-order-mixed-generating} is \(e^{\lambda z}\) times a
polynomial in \(\log z\) of degree \(r\), with nonzero leading coefficient
\(1/r!\).  Repeated monodromy differences about zero strip off the highest
logarithmic degree.  At each step the linear independence of exponentials
with distinct exponents forces the corresponding rational coefficients to
vanish.  Descending induction on \(r\) proves the last assertion.
\end{proof}

\begin{theorem}[Conditional all-order Gamma-jet independence]
\label{thm:conditional-all-order-gamma}
Assume \eqref{eq:ED-intersection}.  Let
\(\lambda_1,\ldots,\lambda_m\in\Qbar_{>0}\) be distinct and choose
integers \(R_i\ge0\).  Then all numbers
\begin{equation}\label{eq:conditional-all-order-values}
 \boxed{\displaystyle
 e^{\lambda_i}[t^r]\{\Gamma(1-t)\lambda_i^t\},
 \qquad 1\le i\le m,\quad0\le r\le R_i,}
\end{equation}
are linearly independent over \(\Qbar\).  In particular
\begin{equation}\label{eq:conditional-gamma-derivative-LI}
 \boxed{1,\Gamma'(1),\Gamma''(1),\ldots}
\end{equation}
are linearly independent over \(\Qbar\).
\end{theorem}

\begin{proof}
In a putative relation, split every \(r\ge1\) value as the value of the
\(E\)-function \(e^{\lambda z}\Ein_r(\lambda z)\) minus the summed
Gevrey value \((\widehat D_{\lambda,r})_0\), and move all Gevrey terms to
one side.  By \eqref{eq:ED-intersection}, the common value is algebraic.
The other side is a linear form in the jointly algebraically independent
coordinates \(\Ein_r(\lambda_i)\) over \(K\), by
\eqref{eq:resonant-input}; therefore every coefficient with \(r\ge1\)
vanishes.  Lindemann--Weierstrass then removes the remaining \(r=0\)
exponential terms.  Formula \eqref{eq:all-order-mixed-generating} at
\(z=1\) proves \eqref{eq:conditional-all-order-values}.  Taking
\(\lambda=1\) and using
\([t^r]\Gamma(1-t)=(-1)^r\Gamma^{(r)}(1)/r!\) proves
\eqref{eq:conditional-gamma-derivative-LI}.
\end{proof}

\begin{remark}[Priority and scope]
Fischler--Rivoal already prove from \eqref{eq:ED-intersection} the
transcendence of Gamma derivatives at the relevant rational points; in
particular the one-point transcendence of \(\gamma\) is not claimed as new
here.  Their stronger mixed zero-removal Conjecture~3 implies
\eqref{eq:ED-intersection}; alternatively, their conditional linear
specialization theorem and \cref{prop:mixed-dilation-rank} give a second
proof of the displayed rank.  The contribution of
\cref{thm:conditional-mixed-dilations,thm:conditional-all-order-gamma} is
the explicit simultaneous dilation family, its resonant all-order lift, the
use of the companion multipoint theorem, and the full linear rank.  It
does not assert algebraic independence: products of mixed functions need
not be mixed, so the cited specialization theorem is deliberately linear.
The value \((W_1)_0(1)=e\gamma\) is the same Euler extension coordinate
that reappears in the framed motive below, but no implication between the
two conjectural routes is known.
\end{remark}

\section{The Euler--Tate matched-line target}
\label{sec:euler-tate}

The exponential-motive construction supplies the canonical framing absent
from the scalar Stokes conjugacy class.

\begin{theorem}[Euler--Mascheroni motive; Fres\'an--Jossen]
\label{thm:FJ-euler-motive}
There is a rank-two exponential motive \(M(\gamma)\) with adapted de Rham
and Betti bases \((d_0,d_1)\) and \((b_0,b_1)\) for which, with columns
indexed by the de Rham basis and rows by the Betti basis, the comparison
matrix is
\begin{equation}\label{eq:euler-period-matrix}
 P_\gamma=\begin{pmatrix}1&\gamma\\0&2\pi i\end{pmatrix}.
\end{equation}
Equivalently,
\[
 \operatorname{comp}(d_0)=b_0,
 \qquad
 \operatorname{comp}(d_1)=\gamma b_0+2\pi i\,b_1.
\]
It lies in a non-split exact sequence
\begin{equation}\label{eq:euler-motive-extension}
 0\longrightarrow\Q(0)\longrightarrow M(\gamma)
 \longrightarrow\Q(-1)\longrightarrow0,
\end{equation}
satisfies \(\operatorname{Hom}(M(\gamma),\Q(0))=0\), and has motivic
Galois group
\[
 G_{M(\gamma)}=
 \left\{\begin{pmatrix}1&b\\0&a\end{pmatrix}:
 a\in\mathbb G_m,\ b\in\mathbb G_a\right\}.
\]
\end{theorem}

\begin{proof}
This is the Euler--Mascheroni construction and Galois-group computation of
Fres\'an--Jossen \cite[\S12.8]{FresanJossen2024}.
\end{proof}

Its determinant is \(2\pi i\), the same as that of the semisimplification,
and is again blind to the extension coordinate.  The missing input can be
stated for this one object.

\begin{proposition}[Complete matched-line classification]
\label{prop:matched-line-classification}
Every de Rham line mapping isomorphically to the quotient is uniquely of
the form
\[
 W_{\mathrm{dR}}(x)=\Qbar(xd_0+d_1),\qquad x\in\Qbar,
\]
and every Betti line with the same property is uniquely of the form
\[
 W_{\mathrm B}(u)=\Qbar(ub_0+b_1),\qquad u\in\Qbar.
\]
The two lines are matched if and only if
\begin{equation}\label{eq:matched-line-equation}
 \boxed{x+\gamma=2\pi i\,u.}
\end{equation}
Consequently a matched pair exists if and only if
\begin{equation}\label{eq:matched-line-locus}
 \boxed{\gamma\in\Qbar+2\pi i\,\Qbar
              =\Qbar+\pi\Qbar.}
\end{equation}
\end{proposition}

\begin{proof}
The displayed parameterizations are the affine charts of lines transverse
to the respective subobjects.  Comparison sends
\[
 xd_0+d_1\longmapsto(x+\gamma)b_0+2\pi i\,b_1.
\]
This spans \(ub_0+b_1\) exactly when
\eqref{eq:matched-line-equation} holds.  Eliminating \(x,u\in\Qbar\)
gives \eqref{eq:matched-line-locus}; the two fields coincide because
\(2i\in\Qbar^\times\).
\end{proof}

\begin{problem}[Euler--Tate matched-line lifting]
\label{prob:euler-tate-fullness}
After extension to \(\Qbar\), suppose one-dimensional de Rham and Betti
subspaces both map isomorphically onto the corresponding realization of
\(\Q(-1)\), and comparison maps one line to the other.  Prove that the
matched pair is induced by a subobject
\[
 \Q(-1)_{\Qbar}\hookrightarrow M(\gamma)_{\Qbar}.
\]
\end{problem}

\begin{theorem}[Strengthened payoff of matched-line lifting]
\label{thm:matched-line-implies-gamma}
A positive solution of \cref{prob:euler-tate-fullness} implies that
\begin{equation}\label{eq:matched-line-payoff}
 \boxed{\gamma\notin\Qbar+\pi\Qbar.}
\end{equation}
Equivalently, \(1,\pi,\gamma\) are linearly independent over \(\Qbar\).
In particular \(\gamma\), \(\gamma/\pi\), and \(S_2\) are transcendental.
\end{theorem}

\begin{proof}
If \(\gamma\in\Qbar+\pi\Qbar\),
\cref{prop:matched-line-classification} supplies a matched pair.  The
proposed lifting would give a copy of \(\Q(-1)\) mapping nontrivially to the quotient
in \eqref{eq:euler-motive-extension}.  The composite is a nonzero scalar,
so rescaling gives a section and hence a splitting.  Moreover
\(\operatorname{Hom}(-,-)\) commutes here with extension from \(\Q\) to
\(\Qbar\), so the Hom-vanishing in \cref{thm:FJ-euler-motive} remains
zero.  This contradiction proves \eqref{eq:matched-line-payoff}.  Membership
in \(\Qbar+\pi\Qbar\) is exactly a nontrivial \(\Qbar\)-linear relation
among \(1,\pi,\gamma\).  Algebraicity of either \(\gamma\) or
\(\gamma/\pi\) would give such a relation.  Finally, if \(S_2\) were
algebraic, \(2S_2\gamma^2+\gamma-2=0\) would make \(\gamma\) algebraic.
\end{proof}

\begin{remark}[Rational and algebraic versions]
A matched-line lifting statement restricted to \(\Q\)-defined lines would
already exclude the corresponding rational matched locus.  The
\(\Qbar\)-coefficient formulation in \cref{prob:euler-tate-fullness} is the
version that yields the full linear-independence conclusion
\eqref{eq:matched-line-payoff}.
\end{remark}

\begin{remark}[Exact scope of current rapid-decay structures]
Snodgrass constructs cyclotomic rational structures on rapid-decay
cohomology under a hypothesis chosen so that Stokes constants do not enter
\cite[\S1.2]{Snodgrass2026}.  This supplies a realization and period pairing,
but not the numerical-to-motivic fullness asserted in
\cref{prob:euler-tate-fullness}.  The problem is a single-object rank-one
target, far more narrowly scoped than the full exponential period conjecture
but not presently proved.
\end{remark}

\begin{remark}[Three independent apertures]
The common-coefficient problem \cref{prob:common-coefficient} is quantitative
and uses one algebraic coefficient across infinitely many moving levels.
The matched-line problem is categorical and uses one canonical framed
extension.  The mixed-Gevrey route uses the intersection
\(\mathbf E\cap\mathbf D=\Qbar\).  Each would prove the transcendence of
\(\gamma\); matched-line lifting would additionally prove the
transcendence of \(\gamma/\pi\), while the mixed route gives an all-order
linear-independence family.  No implication among the three inputs is known.
\end{remark}

\section{Conclusions and remaining problems}
\label{sec:conclusion}

The completed results fall into the following thematic groups.
\begin{enumerate}[label=\textup{(\arabic*)}]
 \item Every divisor \(\Ein-c\) has a universal polynomial inverse-trace
 calculus.  Its full trace ring is exactly \(F[c^{-1}]\), generated by the
 two even traces \(T_2,T_4\), while every transversal moving selection
 preserves algebraic independence.  The only nontrivial collision of
 \((T_2,T_3)\) is the pair \(\{6-2\sqrt3,6+2\sqrt3\}\); at every order
 there are also algebraic levels with zero inverse trace and wholly
 transcendental spectral divisors.  Finite-rank multiplicative groups have
 defect at most their rank and admit cofinite jointly independent
 subfamilies; the complete gamma-free, log-free cancellation space is
 identified.  The tails \(E_1(\alpha)\), and jointly the four classical
 families \(E_1,\Ei,\Ci,\operatorname{Si}\), have defect at most one more
 than the multiplicative rank.  Regularized lower incomplete-Gamma jets
 have defect at most the rank, uniformly across all jet orders, while fixed
 upper-jet rows have one further possible defect.  Sine-integral graph
 ideals, Dickman log-Laplace relations, and characteristic-zero algebraic
 matroids give further realizations of this engine.  At every finite set of
 levels, the complete trace field is exactly the rational quotient by the
 one-dimensional gauge shift.
 \item The dyadic and fixed-power balanced levels approach \(\gamma\) through
 explicit algebraically independent finite-stage data.  Their order-two
 and order-four traces give positive algebraically independent spectral
 decompositions of \(S_2\) and \(S_4\); the two positive carrier traces
 recover \(\gamma\).  The order-two carrier also gives an order-zero
 Laguerre--P\'olya determinant.  Neither the dyadic level sequence nor its
 second trace is P-recursive.
 \item The gamma-free map \(\mathscr R\) has transcendental inverse images at
 every algebraic level in its range, giving the explicit family \(r_m\).
 Normalized finite \(\Ein\)-certificates are unique.  This turns all
 cross-index real Euler witnesses into one countable exclusion theorem and
 yields the relative five-class theorem.  The critical cancellation
 coefficients \(u_N(q)\) have algebraic defect at most one.
 \item For the 2026 Challenge recurrence, the terminating \({}_2F_2\)
 state is explicitly gauge-equivalent to the contragredient normalized
 official Meijer--\(G\) state.  The finite bilinear identity fixes the
 conserved pairing; a flat dual-number \(SL_3\), equivalently \(SL_6\),
 includes the numerator.  The Beta--log transform connects the Challenge to
 the cubic family and every fixed offset.  A single rational carrier yields
 the denominator, numerator, full zeta log-jets, and a monic contact
 polynomial.  The exact maxima mixture, Meijer stop-loss interpretation,
 Poisson--binomial law, zero-density limit, regularized Gamma determinant,
 strengthened divisor, factorial-six Casoratian, all-initial-value limit
 classification, and integer solution lattice complete the finite
 structure.
 \item The Euler--Pad\'e convergents have a finite residue criterion on the
 \(p\)-adic unit boundary and converge as a full sequence at
 \(t=-1\) for \(p=2,3,5,13\).  One lcm-indexed rational subsequence converges
 placewise at the real place and at every fixed prime.  The denominator
 interpolation \(q_p\) identifies unbounded valuations with a \(p\)-adic
 zero, without confusing that event with failure of convergence.
 \item The dyadic augmentation tail is the unique minimum in the
 cumulative-positive cone, while outside that cone three points already
 produce dense tail values.  For fixed \(q\), the exact denominator
 transition for consecutive augmentation-one triples occurs at
 \(e^{(q-1)q^N}/N\).  A moving-point lower bound covering the displayed
 fixed system cannot have sublinear point cost; at linear cost the example
 forces the stated height loss.  The two-base form
 \(Y_N(2)-Y_{M_N}(3)\) retains the same leading \(2^N\) logarithmic scale.
 The exact asymmetric Sondow formula
 replaces the invalid universal entropy-floor heuristic.
 \item In each nondegenerate fixed residual direction of the Pr\'evost grid,
 the real Hilbert quotient is exactly saturated.  For the uncombined
 optimally saddled form, and absent additional exponential divisibility or
 cancellation, the baseline fixed-place ledger leaves the deficit
 \(2-\log(3+2\sqrt2)\).  The bounded-degree reduction transfers a surviving
 one-prime problem to quantitative approximation of \(\log2\), but does not
 close it.
 The fixed-prime Rodrigues deformation produces integral polynomials,
 exact logarithmic linear forms, positivity, and the endpoint factor
 \(p^{(1-\theta)n}\).  Even after the whole factor is favorably credited as
 fixed-place gain, the universal saddle-cancellation inequality shows that
 this ledger does not improve the Pr\'evost baseline.  Balanced factorial
 ratios themselves have only \(O(\log n)\) fixed-prime valuation.
 \item The displayed unframed lateral transition of the elementary Euler
 differential operator is gamma-free.  The missing information is the
 canonical framing.  In the mixed arithmetic-Gevrey route, the intersection
 conjecture \(\mathbf E\cap\mathbf D=\Qbar\) yields simultaneous linear
 independence of the dilation family and of all derivatives
 \(1,\Gamma'(1),\Gamma''(1),\ldots\).  These conclusions remain
 conditional.  In the Euler--Mascheroni motive the same extension
 coordinate is the off-diagonal entry of the period matrix.  The explicit
 matched-line lifting problem \cref{prob:euler-tate-fullness} would force
 the \(\Qbar\)-linear independence of \(1,\pi,\gamma\), not merely the
 transcendence of \(\gamma\).
\end{enumerate}

The principal quantitative unresolved problem can now be stated without
spectral language.

\begin{problem}[Common-coefficient separation]
\label{prob:common-coefficient}
Prove that for every \(s\in\Qbar\),
\[
 \limsup_{N\to\infty}
 \frac{-\log|2sY_N^2+Y_N-2|}{2^N}<1,
\]
or establish any equivalent separation that exploits the same coefficient
\(s\) across all \(N\).
\end{problem}

By \cref{thm:critical-slope}, this problem is equivalent to the
transcendence of \(S_2\), hence to the transcendence of \(\gamma\).  It is
therefore a precise quantitative target, not a smaller lemma already within
current fixed-point \(E\)-function theory.  Independently,
\cref{prob:euler-tate-fullness} is a categorical target for the canonical
rank-two Euler extension; its \(\Qbar\)-version would additionally prove
the transcendence of \(\gamma/\pi\).  The mixed-Gevrey intersection
conjecture gives a logically independent conditional route and the stronger
all-order linear-independence package, but it is not presently known.

We emphasize the corresponding negative conclusions.  The algebraic independence of the
finite stages does not determine the arithmetic nature of their limit.  The
Laguerre--P\'olya determinant is an exact spectral carrier but is not proved
arithmetic or \(D\)-finite; the regularized Gamma determinant does not prove
the irrationality of any odd zeta value.  Certificate uniqueness shows that producing more
linear witnesses of the same normalized type creates mutually exclusive
alternatives rather than independent chances to prove transcendence.  The
specific Pr\'evost and Rodrigues families analyzed here do not pay their
complete height ledger, even under the favorable fixed-place accounting
stated above.  The \(p\)-adic convergence theorems prove neither the
irrationality of \(E_p(-1)\), nor Kurepa's conjecture, nor full convergence
for every prime.  The external relation classifications do not close a
famous transcendence conjecture.  Finally, unframed Stokes data forget the
Euler extension coordinate, and the mixed-Gevrey theorems that recover it
are explicitly conditional.
Accordingly, no unconditional theorem in this paper proves the irrationality or
transcendence of \(\gamma\), \(\delta\), \(S_2\), or any odd zeta value; the
exact Euler apertures above state what is still missing.
\section{Consolidated theorem-status and dependency ledger}
\label{sec:master-status-ledger}

This ledger governs every statement in the monograph when a theorem heading
does not itself carry a badge.  Analytic identities and monodromy results
listed as \UPASS{} are independent of the companion arithmetic-independence
paper.  Every transcendence or algebraic-independence transfer whose input is
the algebraic-point theorem of [001] is marked \RPASS{} here, even when the
subsequent deduction is elementary.

\begin{longtable}{@{}p{.27\textwidth}p{.14\textwidth}p{.52\textwidth}@{}}
\toprule
\textbf{Result family}&\textbf{Status}&\textbf{Exact scope}\\
\midrule
\endfirsthead
\toprule
\textbf{Result family}&\textbf{Status}&\textbf{Exact scope}\\
\midrule
\endhead
Universal inverse traces&\UPASS&
\(T_m(c)=P_m(c^{-1})\), exact two-trace recovery, the fixed-level trace
ring, and the unique \((T_2,T_3)\) collision.\\

Moving-level arithmetic&\RPASS&
Finite-rank/cofinite independence, the dyadic algebraically independent
tower, inverse-level transcendence, algebraic matroid realizations, and
transcendental algebraic-level spectral divisors.\\

Polyexponential Faber structure&\UPASS&
Absolute indecomposability for all \(q\ge1,m\ge2\), and
\(G^{\rm geom}_{q,m}\in\{A_m,S_m\}\) for \(m\ge3\).\\

Polyexponential symmetric monodromy&\UPASS&
Every even degree; every prime degree; \(m=p+2\) for prime \(p\ge5\); and,
for each fixed \(m\), all sufficiently large integer orders \(q\).\\

Polyexponential finite monodromy&\XPASS&
Every \(q\ge1\) for \(m\le1001\), and the original \(q=1\) tower for
\(m\le10^6\).  These are group statements, not all-degree Morse statements.\\

Fixed-order polylogarithmic window&\UPASS/\XPASS&
For \(\tau=1,2,3,4\), every odd \(7\le m\le1001\) has geometric and
arithmetic group \(S_m\).  The row \((\tau,m)=(4,17)\) is covered by the
unconditional prime-degree theorem; all other rows use the exact archived
window audit.\\

Derivative separability&\UPASS/\XPASS&
At least half of the roots of \(P'_{q,2s}\) are simple and
\(P'_{q,p+1}\) is square-free; exact residual computations cover the
displayed finite ranges.  The global DS assertion remains open.\\

Polylogarithmic Faber tower&\UPASS&
For every integer \(\tau\ge1\), all degrees are absolutely indecomposable;
every prime degree is Morse with group \(S_m\); and the degrees \(m=p+2\)
have group \(S_m\).  The corresponding off-diagonal collision polynomials
are absolutely irreducible in those two symmetric families.\\

Factorially weighted Faber transfer&\UPASS&
For \(f_{\sigma,\tau}\), the single-spike hypotheses give an isolated
transposition.  Full \(S_m\) additionally requires a separate
indecomposability input, so only this local conclusion is asserted.\\

Log--Stirling/supernecklace&\UPASS&
All-degree indecomposability, the \(m=p+2\) symmetric family, the
supernecklace EGF, and off-diagonal absolute irreducibility in that family.\\

Inverse surfaces&\UPASS&
No finite asymptotic values, maximal finitary symmetric inverse monodromy,
algebraic freedom of finite distinct branches and their first jets, and the
local-freedom/global-trace-collapse synthesis.\\

Ramanujan--Gamma package&\UPASS&
Recurrence closed forms, flat fields, Meijer--\(G\)/hypergeometric
transmutation, Beta--log transform, probability carrier, contact
polynomials, and the regularized Gamma determinant.\\

Euler Pad\'e/continued fraction&\UPASS&
The finite-residue unit criterion; full \(t=-1\) convergence for
\(p=2,3,5,13\); the common real/all-fixed-prime \(\operatorname{lcm}\)
subsequence; and the exact \(\Z_p\)-zero reformulation.\\

Total positivity and Jacobi tower&\UPASS/\RPASS&
The STP kernel, moment determinacy, Jacobi operator, interlacing, spectral
flow and asymptotics are unconditional.  Arithmetic independence, generic
Grassmannian/Galois, eigenvalue, node, and weight conclusions are relative
to [001].\\

Double scaling&\UPASS/\RPASS&
The ramp equilibrium, free energy, and Marchenko--Pastur-to-arcsine
transition are unconditional; the arithmetic independent phase sequence is
relative to [001].\\

Gamma-tail and spectral sieve&\RPASS&
Fixed-cut half-rank, differential transcendence, wholly transcendental
divisor density, all-order spectral sieves, and the order-axis arithmetic
consequences use [001].\\

Jossen reductions&\UPASS/\RPASS&
The factorial affine E-surface statement is unconditional (with priority
caveat); polyexponential polynomiality, reductions, and noncyclic subclasses
use [001].  The full multigenerator problem is not claimed.\\

Density-one odd monodromy&\CPASS&
For \(q=1\), the exceptional count is
\(O_\varepsilon(X^{15/16+\varepsilon})\) under the displayed derivative-
separability hypothesis DS.\\
\bottomrule
\end{longtable}

\subsection{Statements deliberately left open}

The following assertions are not consequences of this volume and are never
used as hidden inputs:

\begin{itemize}
\item full geometric \(S_m\) in every odd composite degree, or a uniform
odd-square arithmetic theorem;
\item Morse behavior for every trace polynomial, global DS, or global
critical-value noncollision for every \(q\ge2\);
\item general indecomposability, full \(S_m\), or a fixed-\(m\) large-\(\tau\)
Morse theorem for the factorially weighted \(f_{\sigma,\tau}\) family;
\item higher inverse-jet freedom and a minimal differential-order theorem in
every polyexponential order;
\item irrationality or transcendence of any individual \(E_p(-1)\), full
continued-fraction convergence at every prime, Kurepa's conjecture, or
convergence in the restricted adelic-product topology;
\item irrationality or transcendence of \(\gamma\), \(\delta\), or an odd
zeta value; comparison-map injectivity for the proposed Euler motive;
\item the full multigenerator Jossen problem, local kink universality for the
double-scaling model, or unproved quantitative Hilbert/DLS assertions.
\end{itemize}

\subsection{Corrections incorporated in Version 3.0}

Version 3.0 incorporates the sharp threshold \(m\ge5\) for the inverse-level
corollary; the explicit row-space step in the linear-image rank bound; the
scoped moving-level statement; the cancellation check for
\(v_p(b)=q-1\); the endpoint/interior Newton-polygon wording; the exact
\texttt{candidate-count=10} search policy; the missing \(p=11\) Euler--Pad\'e
row; and a direct link to the mean-square prime-gap input.  The values (7)
and (15) in the \(m=27\) table are retained because they are discriminants;
the alternative values (22) and (14) arise from a different resultant
calculation.  The invalid general criterion formerly labelled barriers4
Criterion~7.1 and the hidden-C1 versions of its Gamma claims are excluded.

\subsection{Additions incorporated in Version 3.1}

Version 3.1 supplies the complete Maynard bibliography for the mean-square
prime-gap theorem, proves strict negativity of every positive-degree
coefficient of \(z/\operatorname{Li}_\tau(z)\), and uses the resulting Faber
obstructions to establish all-degree absolute indecomposability for every
integer polylogarithmic order.  It also adds the polylogarithmic prime-degree
Morse theorem, the unconditional \(m=p+2\) symmetric family, the associated
off-diagonal absolute irreducibility, and exact fixed-order window
certificates through odd degree \(1001\).  These conclusions do not extend
the unresolved factorially weighted \(f_{\sigma,\tau}\) family.

\clearpage
\appendix
\part{Proof Certificates and Reproducibility}
\section{An \texorpdfstring{\(\Ein\)}{Ein}-linear realization of
characteristic-zero algebraic matroids}
\label{app:Ein-matroids}

This appendix is the only place where the paper uses the linearization
theorem of Rosen--Sidman--Theran.  It is kept separate so that the moving
trace and Ramanujan results do not appear to depend on it.

\begin{theorem}[\(\Ein\)-realization of algebraic matroids]
\label{thm:matroid-Ein-realization}
Let \(M\) be a finite characteristic-zero algebraic matroid of rank \(r\)
on a ground set \(E\).  There are a matrix
\(B=(b_{je})\in\Qbar^{r\times E}\) representing \(M\), distinct
\(\beta_1,\ldots,\beta_r\in\Qbar^\times\), and numbers
\begin{equation}\label{eq:matroid-Ein-realization}
 \xi_j=\Ein(\beta_j),\qquad
 W_e=\sum_{j=1}^r b_{je}\xi_j,
\end{equation}
such that, for every \(S\subseteq E\),
\begin{equation}\label{eq:matroid-rank-realization}
 \boxed{\displaystyle
 \trdeg_KK(W_e:e\in S)=r_M(S).}
\end{equation}
Moreover the complete relation ideal is
\begin{equation}\label{eq:matroid-linear-ideal}
 \boxed{\displaystyle
 \ker\!\left(K[X_e:e\in E]\longrightarrow K[W_e:e\in E]\right)
 =\left\langle\sum_{e\in E}c_eX_e:
 c\in\Qbar^E,\ Bc=0\right\rangle.}
\end{equation}
\end{theorem}

\begin{proof}
Extend a characteristic-zero field of algebraic realization to its
algebraic closure; the algebraic-dependence matroid is unchanged.  By
\cite[Thm.~1.10]{RosenSidmanTheran2025}, the resulting matroid is linearly
representable over that algebraically closed field.  Linear
representability of a fixed finite
matroid is described by the vanishing of its nonbasis minors and the
nonvanishing of its basis minors.  Introducing one auxiliary inverse for
the product of the basis minors converts these conditions into polynomial
equations over \(\Q\).  Since they have a solution in an algebraically
closed characteristic-zero field, the Nullstellensatz gives a solution in
\(\Qbar\).  Thus a representing matrix \(B\) may be chosen over \(\Qbar\).

Choose distinct nonzero algebraic \(\beta_j\).  By
\cref{thm:input-independence}, the \(\xi_j\) are algebraically independent
over \(K\).  If the columns of \(B\) indexed by \(S\) have rank \(q\),
choose \(q\) independent columns.  The corresponding \(W_e\) are
algebraically independent by \cref{lem:linear-images}, giving the lower
bound \(q\); every remaining column is a \(\Qbar\)-linear combination of
those columns, giving the upper bound.  This proves
\eqref{eq:matroid-rank-realization}.  Finally the symmetric algebra of the
column span is the polynomial ring on the \(X_e\) modulo the linear forms
coming from \(\ker B\), which proves \eqref{eq:matroid-linear-ideal}.
\end{proof}

\begin{remark}[Exact scope]
The construction realizes the rank function of \(M\) and the linearized
relation ideal \eqref{eq:matroid-linear-ideal}.  It does not reproduce the
nonlinear circuit polynomials of a chosen algebraic presentation of
\(M\), and it is not a claim that such a presentation becomes linear as an
algebraic variety.
\end{remark}

\section{Creative-telescoping certificates for the Euler recurrence}
\label{app:telescopers}

For \(a=n+3\), \(b=n+4\), define meromorphically
\[
 T_{n,k}=b\frac{(a!b!)^2}
 {\Gamma(k+1)^3\Gamma(a-k+1)\Gamma(b-k+1)},
 \qquad
 \mathcal L_{n,k}=3H_k-H_{a-k}-H_{b-k},
\]
where \(H_z=\psi(z+1)+\gamma\).  The regularized endpoint values are
\[
 T_{n,b}\mathcal L_{n,b}=a!,
 \qquad
 T_{n,k}\mathcal L_{n,k}=0\quad(k>b).
\]
Indeed, at \(k=b\) the simple zero of
\(1/\Gamma(a-k+1)\) cancels the simple pole of
\(-H_{a-k}\), with limiting product \(a!\).  For \(k>b\), the two
reciprocal-gamma zeros dominate the at-most-simple digamma pole, so the
regularized product vanishes.
Thus \eqref{eq:qn-closed} and \eqref{eq:pn-closed} are respectively
\(\sum_{k\ge0}T_{n,k}\) and
\(\sum_{k\ge0}T_{n,k}\mathcal L_{n,k}\).

Let
\[
 c_0=-A_n,
 \qquad c_1=B_n,
 \qquad c_2=-C_n,
 \qquad c_3=D_n,
\]
and \(\Delta_ku(k)=u(k+1)-u(k)\).  The certificates are
\begin{align}
 \sum_{j=0}^3c_jT_{n-j,k}
 &=\Delta_k\{T_{n,k}h_{n,k}\},
 \label{eq:q-telescoper}\\
 \sum_{j=0}^3c_jT_{n-j,k}\mathcal L_{n-j,k}
 &=\Delta_k\{T_{n,k}(h_{n,k}\mathcal L_{n,k}+g_{n,k})\}.
 \label{eq:p-telescoper}
\end{align}
Here
\[
 h_{n,k}=-\frac{k^3}{(n+3)^3(n+4)^3}\mathcal H(n,k),
 \qquad
 g_{n,k}=\frac{k^2}{(n+3)^3(n+4)^3}\mathcal G(n,k),
\]
where
\begin{align*}
\mathcal H(n,k)={}&
 (8n^3+75n^2+231n+232)k^3\\
&-(48n^4+562n^3+2436n^2+4626n+3248)k^2\\
&+(96n^5+1388n^4+7924n^3+22278n^2+30772n+16668)k\\
&-(64n^6+1120n^5+8071n^4+30590n^3+64148n^2+70347n+31388),
\end{align*}
and
\begin{align*}
\mathcal G(n,k)={}&
 (48n^3+450n^2+1386n+1392)k^3\\
&-(240n^4+2810n^3+12180n^2+23130n+16240)k^2\\
&+(384n^5+5552n^4+31696n^3+89112n^2+123088n+66672)k\\
&-(192n^6+3360n^5+24213n^4+91770n^3+192444n^2+211041n+94164).
\end{align*}

To verify \eqref{eq:q-telescoper}, divide by \(T_{n,k}\), use
\[
 \frac{T_{n,k+1}}{T_{n,k}}
 =\frac{(n+3-k)(n+4-k)}{(k+1)^3},
\]
and, for \(0\le j\le3\),
\[
 \frac{T_{n-j,k}}{T_{n,k}}
 =\frac{n-j+4}{n+4}
 \frac{(n+3-k)^{\underline j}(n+4-k)^{\underline j}}
 {\{(n+3)^{\underline j}(n+4)^{\underline j}\}^{2}},
\]
where \(x^{\underline j}=x(x-1)\cdots(x-j+1)\), and clear denominators.
For \eqref{eq:p-telescoper}, also use
\[
 \mathcal L_{n,k+1}-\mathcal L_{n,k}
 =\frac3{k+1}+\frac1{n+3-k}+\frac1{n+4-k}.
\]
The remaining shifts are reduced by
\[
 \mathcal L_{n-j,k}-\mathcal L_{n,k}
 =\sum_{r=0}^{j-1}
 \left(\frac1{n+3-k-r}+\frac1{n+4-k-r}\right).
\]
Both become polynomial identities in \(n,k\).  The lower boundary vanishes
because of the factors \(k^3,k^2\), and the reciprocal gamma factors give
the upper boundary.  Summing proves \eqref{eq:challenge-recurrence} for both
closed forms for every \(n\ge0\).  The accompanying source package contains
an exact symbolic verifier for these identities and for
\eqref{eq:cmf-det}--\eqref{eq:cmf-flatness}.

\section{Finite certificate for the integer initial lattice}
\label{app:integer-lattice-certificate}

The official dual transition and cyclic gauge may be taken as
\begin{equation}\label{eq:Rn-certificate}
 \mathsf R(n)=
 \begin{pmatrix}
 6n^2+21n+18&-8n-11&3\\
 3n^3+12n^2+16n+7&-3n^2-7n-4&n+1\\
 (n+1)^4&0&0
 \end{pmatrix}
\end{equation}
and
\begin{equation}\label{eq:Un-certificate}
\mathsf U(n)=
\begin{pmatrix}
1&6n^2+21n+18&15n^4+147n^3+520n^2+792n+440\\
0&3n^3+12n^2+16n+7&10n^5+96n^4+353n^3+625n^2+537n+179\\
0&(n+1)^4&6n^6+57n^5+213n^4+402n^3+408n^2+213n+45
\end{pmatrix}.
\end{equation}
Let the row-companion transfer of \eqref{eq:challenge-recurrence} be
\begin{equation}\label{eq:row-companion-certificate}
 \mathsf C_n^{\mathrm{comp}}=
 \begin{pmatrix}
 0&0&D_n/A_n\\
 1&0&-C_n/A_n\\
 0&1&B_n/A_n
 \end{pmatrix}.
\end{equation}
Direct multiplication gives the exact coboundary
\begin{equation}\label{eq:RU-companion-coboundary}
 \mathsf R(n)\mathsf U(n+1)
 =\mathsf U(n)\mathsf C_n^{\mathrm{comp}}.
\end{equation}
For a raw row \(c\), the integer recurrence state at level \(N\) is
\begin{equation}\label{eq:raw-state-certificate}
 c\,\mathsf R(0)\cdots\mathsf R(N-1)\mathsf U(N).
\end{equation}
Indeed \eqref{eq:RU-companion-coboundary} proves by induction that
\eqref{eq:raw-state-certificate} equals
\((u_{N-3},u_{N-2},u_{N-1})\) when \(c\mathsf U(0)\) is the initial row;
the empty product is used at \(N=0\).
Let
\[
 d_1=(8,-22,10),\qquad d_2=(0,2,-6),\qquad d_3=(0,-1,7).
\]
Direct multiplication gives
\begin{align}
 \frac{d_1}{8}\mathsf U(0)&=(1,0,4),&
 \frac{d_2}{8}\mathsf U(0)&=(0,1,11),&
 \frac{d_3}{8}\mathsf U(0)&=(0,0,17),
 \label{eq:lattice-basis-certificate}
\end{align}
and the finite congruence certificate
\begin{align}
 \mathsf R(0)\cdots\mathsf R(7)&\equiv0\pmod8,
 \label{eq:R8-zero}\\
 d_i\mathsf R(0)\cdots\mathsf R(N-1)\mathsf U(N)&\equiv0\pmod8
 \quad(1\le i\le3,\ 0\le N<8).
 \label{eq:finite-mod8-certificate}
\end{align}
For \(N\ge8\), \eqref{eq:R8-zero} supplies the factor eight; for
\(N<8\), \eqref{eq:finite-mod8-certificate} does.  Hence the three rows in
\eqref{eq:lattice-basis-certificate} generate integral solutions at all
levels, completing the sufficiency argument in
\cref{thm:ram-integer-lattice}.  The source verifier checks every displayed
matrix identity over \(\Z\).

\section{A numerical ledger}
\label{app:numerics}

The following values provide reproducibility checks for the indexing in
\eqref{eq:dN-def}.  They are not used in any proof.

\begin{table}[ht]
\centering
\caption{Positive increments and positive zeros of \(\Delta_2\).}
\begin{tabular}{@{}rcc@{}}
\toprule
\(N\)&\(d_N\)&\(r_N=\sqrt{2/d_N}\)\\
\midrule
0&1.474569389580200&1.1646150303\\
1&\(5.633945512124626\times10^{-1}\)&1.8841205074\\
2&\(9.586804440833174\times10^{-2}\)&4.5674948216\\
3&\(1.339700176289053\times10^{-3}\)&38.637693279\\
4&\(2.954643417466919\times10^{-7}\)&2601.7314103\\
5&\(2.050814843583590\times10^{-14}\)&\(9.8753335\times10^6\)\\
\bottomrule
\end{tabular}
\end{table}

The indexing is tied to the convention \(Y_N\) for \(N\ge1\):
\[
 d_0=\tau_1,\qquad d_N=\tau_{N+1}-\tau_N\quad(N\ge1).
\]
Starting a level sequence at \(N=0\) would give a different, though equally
valid, spectral decomposition.
\section{Exact computational certificates}
\label{app:computational-certificates}

The computer-assisted statements in this volume use exact integer or
finite-field arithmetic.  Floating-point computations are never used as
proof of a Galois-group, separability, or valuation assertion.  The complete
scripts and captured outputs are shipped in the source archive under
\nolinkurl{computation/}.

\subsection{All-order rectangle}

The command
\begin{center}
\texttt{python3 computation/trace\_allq\_modp\_audit.py
  --bound 1001 --candidate-count 10}
\end{center}
examines, for every eligible odd degree, at most the ten largest
strict-window primes in descending order.  It checks every relevant order
residue class, records the prime-degree fallback rows, and emits the combined
certificates at \(m=27\) and \(m=989\).  The seven \(m=27\) discriminant
rows use the values displayed in the main text; the search does not replace
them by resultants.

\subsection{Fixed-order polylogarithmic window}

The command
\begin{center}
\texttt{python3 computation/polylog\_window\_prime\_audit.py
  --bound 1001}
\end{center}
uses exact arithmetic in $\F_p$ to test every strict-window prime
$(m+1)/2<p<m$ for each $\tau\in\{1,2,3,4\}$ and odd
$7\le m\le1001$.  It constructs the residual polynomial $R_r$ from
\eqref{eq:R-general} and checks $\gcd(R_r,R_r')=1$.  There are no window
exceptions for $\tau=1,2,3$; for $\tau=4$ the sole window exception is
$m=17$, where the unconditional prime-degree part of
\Cref{thm:polylog-global} proves $S_{17}$.  Thus the exception concerns only
the window certificate, not the conclusion of \Cref{cor:polylog-window}.

The frozen files have SHA-256 hashes
\begin{itemize}
\item \nolinkurl{polylog_window_prime_audit.py}:\\
\texttt{\small\seqsplit{dd42b574f3eb5144cd49316530911f4cee0e3722ed03dda7382bd10ec8a31cd2}};
\item \nolinkurl{polylog_window_prime_audit_1001.txt}:\\
\texttt{\small\seqsplit{8f952488d2011c1391ae15a4f607715c7fc7de73a00ce8cae00fa973f938cbdf}}.
\end{itemize}

\subsection{Original tower through one million}

The command
\begin{center}
\texttt{python3 computation/window\_prime\_selection\_audit.py
  --bound 1000000}
\end{center}
checks the original \(q=1\) tower.  For each odd composite degree it first
tests the largest strict-window prime and then the next candidate if needed.
The output records the two failures of the largest-prime choice and the
successful second-prime certificates.

\subsection{Morse and derivative certificates}

The lower-degree exact Morse certificates are generated by
\begin{center}
\texttt{python3 computation/trace\_morse\_certificates.py
  --max-m 55 --prime-bound 3000}.
\end{center}
The residual even-derivative computations are implemented independently in
the supplied C++ source and their captured outputs are included for bounds
\(1000\) and \(5000\).  These finite certificates support exactly the ranges
stated in the theorem ledger; they are not extrapolated to a universal DS
theorem.

\subsection{Reproducibility convention}

Each source-archive file is covered by the top-level SHA-256 manifest.  A
successful independent reproduction consists of:

\begin{enumerate}
\item verifying the manifest;
\item compiling the monograph with the documented pdfLaTeX command;
\item running the quick certificate suite;
\item optionally running the extended \(10^6\)-degree suite and comparing
its terminal summary with the captured output.
\end{enumerate}

The source archive records the interpreter/compiler versions used for the
frozen build.  No network access is required to compile the PDF or rerun the
certificates.

\clearpage
\begin{thebibliography}{99}
\small
\raggedright

\bibitem{AdamczewskiFaverjon2025}
B.~Adamczewski and C.~Faverjon,
\newblock Algebraic independence measures for values of \(E\)-functions and
\(M\)-functions,
\newblock arXiv:2502.09999 (2025).
\newblock \url{https://arxiv.org/abs/2502.09999}.

\bibitem{DLMF}
NIST Digital Library of Mathematical Functions,
\newblock \S6.12, asymptotic expansions of exponential and logarithmic
integrals,
\newblock release 1.2.7 (2026).
\newblock \url{https://dlmf.nist.gov/6.12}.

\bibitem{Delaygue2018}
E.~Delaygue,
\newblock Arithmetic properties of Ap\'ery-like numbers,
\newblock \emph{Compos. Math.} \textbf{154} (2018), 249--274.
\newblock \url{https://doi.org/10.1112/S0010437X17007552}.

\bibitem{FischlerRivoal2016}
S.~Fischler and T.~Rivoal,
\newblock Arithmetic theory of \(E\)-operators,
\newblock \emph{J. \'Ec. polytech. Math.} \textbf{3} (2016), 31--65.
\newblock \url{https://doi.org/10.5802/jep.28}.

\bibitem{FischlerRivoal2024}
S.~Fischler and T.~Rivoal,
\newblock Values of \(E\)-functions are not Liouville numbers,
\newblock \emph{J. \'Ec. polytech. Math.} \textbf{11} (2024), 1--18.
\newblock \url{https://doi.org/10.5802/jep.249}.

\bibitem{FischlerRivoalMixed2024}
S.~Fischler and T.~Rivoal,
\newblock Relations between values of arithmetic Gevrey series, and
applications to values of the Gamma function,
\newblock \emph{J. Number Theory} \textbf{261} (2024), 36--54.
\newblock \url{https://arxiv.org/abs/2301.13518}.

\bibitem{FresanJossen2024}
J.~Fres\'an and P.~Jossen,
\newblock \emph{Exponential Motives},
\newblock monograph manuscript, version of April 2024.
\newblock \url{http://javier.fresan.perso.math.cnrs.fr/expmot.pdf}.

\bibitem{KEIO2026I}
KEIO,
\newblock Exponential-integral periods across arithmetic, Pad\'e, Hankel,
reciprocal geometry, and Frobenius framing,
\newblock Version 6.3, research manuscript, 108 pages (2026).

\bibitem{Kaluza1928}
Th.~Kaluza,
\newblock \emph{\"Uber die Koeffizienten reziproker Potenzreihen},
\newblock \emph{Math. Z.} \textbf{28} (1928), 161--170.
\newblock \url{https://doi.org/10.1007/BF01181155}.

\bibitem{Lagarias2013}
J.~C.~Lagarias,
\newblock Euler's constant: Euler's work and modern developments,
\newblock \emph{Bull. Amer. Math. Soc.} \textbf{50} (2013), 527--628.
\newblock \url{https://doi.org/10.1090/S0273-0979-2013-01423-X}.

\bibitem{Levin1996}
B.~Ya.~Levin,
\newblock \emph{Lectures on Entire Functions},
\newblock Translations of Mathematical Monographs 150,
American Mathematical Society, Providence, RI, 1996.

\bibitem{Marcovecchio2009}
R.~Marcovecchio,
\newblock The Rhin--Viola method for \(\log 2\),
\newblock \emph{Acta Arith.} \textbf{139} (2009), 147--184.
\newblock \url{https://doi.org/10.4064/aa139-2-5}.

\bibitem{Marcovecchio2025}
R.~Marcovecchio,
\newblock Multiple Legendre polynomials in Diophantine approximation,
\newblock arXiv:2512.12890 (2025).
\newblock \url{https://arxiv.org/abs/2512.12890}.

\bibitem{Maynard2012}
J.~Maynard,
\newblock \emph{On the difference between consecutive primes},
\newblock arXiv:1201.1787v3 [math.NT] (2012), Theorem~3.1.
\newblock \url{https://arxiv.org/abs/1201.1787}.

\bibitem{Prevost2008}
M.~Pr\'evost,
\newblock A family of criteria for irrationality of Euler's constant,
\newblock \emph{J. Comput. Appl. Math.} \textbf{219} (2008), no.~2,
484--492.
\newblock \url{https://arxiv.org/abs/math/0507231}.

\bibitem{RamanujanChallenge2026}
M.~Shalyt, R.~Kalisch, C.~Schneider, H.~Barkan, E.~Leibtag,
J.~Campbell, S.~Weinbaum, T.~Monderer, A.~Narayanan, and I.~Kaminer,
\newblock The Ramanujan Challenge for AI,
\newblock arXiv:2607.09721 (2026).
\newblock \url{https://arxiv.org/abs/2607.09721}.

\bibitem{Ridout1958}
D.~Ridout,
\newblock The \(p\)-adic generalization of the Thue--Siegel--Roth theorem,
\newblock \emph{Mathematika} \textbf{5} (1958), 40--48.
\newblock \url{https://doi.org/10.1112/S0025579300001339}.

\bibitem{Snodgrass2026}
B.~Snodgrass,
\newblock Periods of \(E\)-operators,
\newblock arXiv:2608.06005v2 (2026).
\newblock \url{https://arxiv.org/abs/2608.06005}.

\bibitem{WolfsVanAssche2025}
W.~Van Assche and T.~Wolfs,
\newblock Rational approximation of Euler's constant using multiple
orthogonal polynomials,
\newblock \emph{Combinatorics and Number Theory} \textbf{14} (2025),
no.~2, 141--162.
\newblock \url{https://doi.org/10.2140/cnt.2025.14.141}.

\bibitem{Akhiezer1965}
N.~I.~Akhiezer,
\newblock \emph{The Classical Moment Problem and Some Related Questions in
Analysis},
\newblock Oliver \& Boyd, Edinburgh, 1965.

\bibitem{Delaygue2022}
\'E.~Delaygue,
\newblock A Lindemann--Weierstrass theorem for \(E\)-functions,
\newblock \emph{J. Reine Angew. Math.} \textbf{820} (2025), 75--85.
\newblock \url{https://doi.org/10.1515/crelle-2024-0090}.

\bibitem{ErnvallMatalaSeppala2022}
A.-M.~Ernvall-Hyt\"onen, T.~Matala-Aho, and L.~Sepp\"al\"a,
\newblock Euler's factorial series, Hardy integral, and continued fractions,
\newblock \emph{J. Number Theory} \textbf{244} (2023), 224--250.
\newblock \url{https://doi.org/10.1016/j.jnt.2022.09.007}.

\bibitem{GrahovacEtAl2024}
D.~Grahovac, A.~Kovtun, N.~N.~Leonenko, and A.~Pepelyshev,
\newblock Dickman type stochastic processes with short- and long-range
dependence,
\newblock \emph{Stochastics} \textbf{97} (2025), no.~6, 722--743.
\newblock \url{https://doi.org/10.1080/17442508.2025.2522789}.

\bibitem{GrossCS}
B.~H.~Gross,
\newblock On the periods of abelian integrals and a formula of Chowla and
Selberg,
\newblock \emph{Invent. Math.} \textbf{45} (1978), no.~2, 193--211.
\newblock \url{https://doi.org/10.1007/BF01390273}.

\bibitem{Petrov1975}
V.~V.~Petrov,
\newblock \emph{Sums of Independent Random Variables},
\newblock Ergebnisse der Mathematik und ihrer Grenzgebiete 82,
Springer, Berlin, 1975.

\bibitem{MurtySaha2015}
M.~R.~Murty and E.~Saha,
\newblock Transcendental values of the incomplete gamma function and
related questions,
\newblock \emph{Arch. Math.} \textbf{105} (2015), 271--283.
\newblock \url{https://doi.org/10.1007/s00013-015-0800-3}.

\bibitem{Pilehrood2013}
Kh.~Hessami Pilehrood and T.~Hessami Pilehrood,
\newblock On a continued fraction expansion for Euler's constant,
\newblock \emph{J. Number Theory} \textbf{133} (2013), no.~2, 769--786.
\newblock \url{https://doi.org/10.1016/j.jnt.2012.08.016}.

\bibitem{PilehroodBell2012}
Kh.~Hessami Pilehrood and T.~Hessami Pilehrood,
\newblock Rational approximations to values of Bell polynomials at points
involving Euler's constant and zeta values,
\newblock \emph{J. Aust. Math. Soc.} \textbf{92} (2012), no.~1, 71--98.
\newblock \url{https://doi.org/10.1017/S1446788712000134}.

\bibitem{PolyaSchur1914}
G.~P\'olya and I.~Schur,
\newblock \"Uber zwei Arten von Faktorenfolgen in der Theorie der
algebraischen Gleichungen,
\newblock \emph{J. Reine Angew. Math.} \textbf{144} (1914), 89--113.

\bibitem{Powers2023}
M.~R.~Powers,
\newblock A sequence of nested exponential random variables with connections
to two constants of Euler,
\newblock arXiv:2302.04605v4 (2023).
\newblock \url{https://arxiv.org/abs/2302.04605}.

\bibitem{Powers2025}
M.~R.~Powers,
\newblock Transcendence results for
\(\Gamma^{(n)}(1)\) and related sequences of generalized constants,
\newblock arXiv:2511.01849v7 (2026).
\newblock \url{https://arxiv.org/abs/2511.01849}.

\bibitem{RamanujanOfficial2026}
Ramanujan Machine Group,
\newblock \emph{Ramanujan Challenge Proofs},
\newblock official solution manuscript, 45 pages (2026).
\newblock
\url{https://ramanujanmachine.com/wp-content/uploads/2026/08/Ramanujan_Challenge_Proofs_02_08.pdf}.
\newblock Supplementary code:
\url{https://github.com/RamanujanMachine/RamanujanChallengeSolutions}.

\bibitem{RosenSidmanTheran2025}
Z.~Rosen, J.~Sidman, and L.~Theran,
\newblock Linearizing algebraic matroids,
\newblock arXiv:2507.07220v3 (2026).
\newblock \url{https://arxiv.org/abs/2507.07220}.

\bibitem{Sondow2003}
J.~Sondow,
\newblock Criteria for irrationality of Euler's constant,
\newblock \emph{Proc. Amer. Math. Soc.} \textbf{131} (2003), no.~11,
3335--3344.
\newblock \url{https://doi.org/10.1090/S0002-9939-03-07081-3}.

\bibitem{VanAsscheWolfs2023}
W.~Van Assche and T.~Wolfs,
\newblock Multiple orthogonal polynomials associated with the exponential
integral,
\newblock \emph{Stud. Appl. Math.} \textbf{151} (2023), no.~2, 411--449.
\newblock \url{https://doi.org/10.1111/sapm.12608}.

\end{thebibliography}

\end{document}
